\documentclass[11pt]{article}

\usepackage[margin=1.1in]{geometry}
\usepackage{amsmath,amssymb,amsthm,mathtools}
\usepackage{enumitem}
\usepackage{hyperref}
\usepackage{microtype}
\usepackage{booktabs}
\usepackage{mathrsfs}

\hypersetup{
  colorlinks=true,
  linkcolor=blue,
  citecolor=blue,
  urlcolor=blue,
  hypertexnames=false,
  pdftitle={Boundary, Terminal, Diagonal, and Cycle Limit Laws for Luce Permutations},
  pdfauthor={Christopher D. Long}
}

\newtheorem{theorem}{Theorem}[section]
\newtheorem{proposition}[theorem]{Proposition}
\newtheorem{lemma}[theorem]{Lemma}
\newtheorem{corollary}[theorem]{Corollary}
\newtheorem{conjecture}[theorem]{Conjecture}
\theoremstyle{remark}
\newtheorem{remark}[theorem]{Remark}
\newtheorem{warning}[theorem]{Warning}
\newtheorem{definition}[theorem]{Definition}

\newcommand{\Pp}{\mathbb P}
\newcommand{\Ee}{\mathbb E}
\newcommand{\Rr}{\mathbb R}
\newcommand{\Nn}{\mathbb N}
\newcommand{\Zz}{\mathbb Z}
\newcommand{\one}{\mathbf 1}
\newcommand{\TV}{\mathrm{TV}}
\newcommand{\Pois}{\operatorname{Pois}}
\newcommand{\Exp}{\operatorname{Exp}}
\newcommand{\PPP}{\operatorname{PPP}}
\newcommand{\Var}{\operatorname{Var}}
\newcommand{\Cov}{\operatorname{Cov}}
\newcommand{\diag}{\operatorname{diag}}
\newcommand{\dd}{\,d}
\newcommand{\eps}{\varepsilon}
\newcommand{\calD}{\mathcal D}
\newcommand{\calC}{\mathcal C}
\newcommand{\calE}{\mathcal E}
\newcommand{\calG}{\mathcal G}

\title{Boundary, Terminal, Diagonal, and Cycle Limit Laws for Luce Permutations}
\author{Christopher D. Long}
\date{}

\begin{document}
\maketitle

\begin{abstract}
The stationary distribution of the Tsetlin library is the Luce distribution,
equivalently the order induced by independent exponential clocks with prescribed
rates.  Chatterjee, Diaconis, and Kim initiated an enumerative theory for this
non-uniform distribution, and Borga, Chatterjee, and Diaconis developed its
permuton and local-limit theory.  This paper proves boundary, terminal,
diagonal, and cycle limit laws for Luce permutations.

The boundary, diagonal, and endpoint results are proved by local-entry analysis
in exponential-clock and shifted-Gumbel coordinates.  For top windows we prove
exact total-variation and one-sided relative $L^\infty$ identities; for bottom
windows we identify the precise terminal leakage when the limiting integral is
defective.  Writing $C_r=C_r(\sigma_n)$ for the number of $r$-cycles, we prove,
for uniformly positive continuous profiles, joint Poisson limits for every
fixed vector $(C_1,\ldots,C_R)$, together with sparse deterministic-entry and
overlap laws.  The limiting cycle means $\lambda_r$ satisfy
\[
        \lambda_r=\frac1r+O_f(q^r/r)
\]
for some $q=q(f)\in(0,1)$.  Under the weaker assumption that the largest and
smallest weights are uniformly comparable, a complementary clock-slab mixing
and regeneration argument proves that the ranked macroscopic cycle lengths
converge to $\operatorname{PD}(1)$.

For the singular Sukhatme weights and, more generally, endpoint-regular profiles
with regular-variation index $\gamma>0$, fixed points form a local Poisson point
process on logarithmic terminal scale with intensity
\[
        c_\gamma\dd x,
        \qquad
        c_\gamma=\gamma K_\gamma e^{-K_\gamma},
        \qquad
        K_\gamma=\Gamma(1+1/\gamma)^\gamma.
\]
With $F_n=C_1(\sigma_n)$, we prove the global approximation
\[
        d_{\mathrm{TV}}\bigl(\mathcal L(F_n),
        \operatorname{Pois}(\Ee F_n)\bigr)
        =O\!\left(\frac{\log\log n}{\sqrt{\log n}}\right).
\]
For every fixed collection of endpoint cycle lengths we determine all mixed
factorial-moment asymptotics; in particular,
\[
        \Ee C_r\sim \frac{k_\gamma^{*r}(0)}r\log n,
        \qquad
        k_\gamma(x)=\gamma K_\gamma e^{-\gamma x}
        e^{-K_\gamma e^{-\gamma x}},
\]
and the normalized fixed-length cycle counts satisfy a joint law of large
numbers.  Finally, on every mesoscopic scale $b_n\to\infty$, $b_n=o(n)$, the
normalized endpoint graph converges to the deterministic bistochastic tangent
kernel
\[
        \mathcal K_\gamma(x,y)
        =y^{-1}k_\gamma(\log(y/x)).
\]
\end{abstract}

\medskip
\tableofcontents

\section{Introduction and scope}

Let $\theta_1,\ldots,\theta_n$ be positive weights.  The Luce distribution,
originating in Luce's choice model \cite{Luce1959}, is generated on $S_n$ by
sampling labels without replacement, choosing each remaining label with
probability proportional to its weight.  Equivalently, if
\[
        X_i\sim\Exp(\theta_i),\qquad 1\le i\le n,
\]
are independent, then the labels in increasing order of the $X_i$'s have the
Luce law.  This is the stationary distribution of the Tsetlin library; see
Chatterjee, Diaconis, and Kim \cite{CDK2024}.  The Tsetlin library also belongs
to the hyperplane-walk framework of Bidigare, Hanlon, and Rockmore
\cite{BHR1999} and Brown and Diaconis \cite{BD1998}.  Borga, Chatterjee, and
Diaconis developed the permuton and local-limit theory of Luce permutations
\cite{BCD2025}; for general permutation-limit background, see Hoppen,
Kohayakawa, Moreira, R\'ath, and Sampaio \cite{HKMRS2013}.

Mukherjee \cite[Theorems~1.4, 1.6, and~1.9]{Mukherjee2016} derives fixed-point,
overlap, and fixed short-cycle Poisson limits from a global uniform
approximation of finite point probabilities for convergent random permutations.
In the regular Luce setting, Proposition \ref{prop:bulk-local-final} proves the
corresponding multipoint approximation only on separated bulk configurations,
with microscopic entry density equal to the Luce permuton density identified
in \cite[Theorem~2.1 and Proposition~2.2]{BCD2025}.  Because the estimate is not
uniform through endpoint and collision regimes, we do not invoke Mukherjee's
theorems directly.  Instead, endpoint and collision envelopes are combined with
the separated estimate in a self-contained factorial-moment argument.  The
Sukhatme and endpoint-regular profiles lie outside the uniformly positive
regime and require a distinct terminal analysis.

We use \emph{terminal-condition flux} as an organizing description for the
local-entry mechanism in most of the proofs.  A statistic is controlled by the
asymptotic mass entering a relevant rank, boundary, diagonal, or terminal set.
The exponential-clock representation makes these entries explicit: top-window
statistics are governed by sampling without replacement, bottom-window laws by
leakage through a terminal endpoint, regular fixed points and short cycles by
microscopic point-entry estimates, and singular profiles by logarithmic
terminal saddles.  The equivalent shifted-Gumbel representation
\[
        Y_i=-\log X_i=\log\theta_i+G_i
\]
turns the same events into threshold crossings.  For local rank events, the
density of the specified priority and the derivative of the Bernoulli threshold
mean contain a common factor that cancels.  The terminology introduces no
additional assumption; all proofs are stated in standard probabilistic terms.

The results fall into four regimes.  First, at the boundary, we prove exact
top-window distance identities and identify the precise leakage of candidate
infinite bottom-window laws.  The shifted-Gumbel threshold formulas and the
vanishing-block construction in the preliminaries show, respectively, how local
rank probabilities are calculated and why sparse cycle statistics cannot be
recovered from permuton convergence alone.

Second, for uniformly positive continuous profiles, we prove joint Poisson
limits for every fixed vector of cycle counts, sparse-entry and
deterministic-overlap laws, and exponentially accurate long-trace universality.
Third, under the weaker assumption of uniformly comparable weights, the
clock-slab and permanental-matching argument of Section
\ref{sec:macroscopic-cycles} yields $\operatorname{PD}(1)$ macroscopic-cycle
universality.  A common minorization gives a spectral gap depending only on the
weight-comparability constant; quantitative prefix balance, pointwise rooted
smoothing, and regular-history deletion provide the regeneration step.

Fourth, for the singular Sukhatme profile, we establish the logarithmic
fixed-point mean, a local terminal Poisson point process, and all fixed
factorial moments.  The endpoint-regular theory extends these conclusions under
regular variation, proves a global total-variation Poisson approximation by an
approximate-Palm rank transposition, identifies finite-dimensional endpoint
displacements, and constructs microscopic and mesoscopic endpoint limits.  It
also determines the logarithmic means and all fixed mixed factorial moments of
the endpoint cycle counts of every fixed length, yielding a joint law of large
numbers.

These results answer, under explicit microscopic regularity hypotheses, the
fixed-point question and the joint-law question for every fixed vector of short
cycle counts raised by Borga, Chatterjee, and Diaconis
\cite[Section~5.2]{BCD2025}.  The asymptotic law of the total number of cycles
remains open; see Conjecture \ref{conj:total-cycles}.  Proposition
\ref{prop:implantation} explains why permuton convergence alone cannot imply
such microscopic conclusions.

Most results use local-entry analysis: identify the relevant event, write an
exact integral or sequential estimate, and evaluate it using a lattice local
central limit theorem, saddle analysis, or product bounds.  The macroscopic
$\operatorname{PD}(1)$ theorem instead uses a complementary clock-slab mixing
and regenerative-cycle argument.

\subsection*{Conventions}
The notation $\Exp(\theta)$ denotes the exponential distribution with rate
$\theta$.  Weight vectors may form triangular arrays and therefore depend on
$n$.  The phrase \emph{lattice local central limit theorem} below refers to
point probabilities for lattice-valued sums and is distinct from permutation
local-limit notions.  All exponential clocks are almost surely distinct.  If \(\pi\) is the
one-line ordering of labels, then \(\pi(r)\) is the label drawn at position
\(r\).  Let \(\sigma=\pi^{-1}\), so \(\sigma(i)\) is the position of label
\(i\).  Fixed points and cycle counts are invariant under inversion.  In
fixed-point and cycle arguments we use \(\sigma\), since
\[
        \sigma(i)=i \quad\Longleftrightarrow\quad
        \#\{j:X_j<X_i\}=i-1.
\]

\section{Preliminaries}

\subsection{The exponential race}
Let \(\theta_1,\ldots,\theta_n>0\), and let \(X_i\sim\Exp(\theta_i)\) be independent.  For distinct labels \(i_1,\ldots,i_n\),
\[
\Pp(X_{i_1}<\cdots<X_{i_n})
= \prod_{r=1}^n
 \frac{\theta_{i_r}}{\theta_{i_r}+\theta_{i_{r+1}}+\cdots+\theta_{i_n}}.
\]
This is the standard exponential race identity and follows from the memoryless property.


\subsection{Gumbel priorities}\label{subsec:gumbel-priorities}
A standard Gumbel random variable has distribution
\[
        \Pp(G\le y)=\exp(-e^{-y}),
        \qquad y\in\Rr .
\]
If \(E_i\sim\Exp(1)\) are independent, then \(-\log E_i\) are independent standard Gumbels.  Since
\[
        X_i\stackrel{d}=E_i/\theta_i,
\]
the exponential race may equivalently be generated by the shifted priorities
\[
        Y_i=-\log X_i=\log\theta_i+G_i,
        \qquad 1\le i\le n,
\]
ordered in decreasing order.  Thus
\[
        \sigma(i)=1+\#\{j:Y_j>Y_i\}.
\]
For later reference we record the exact threshold formulas.  The density and upper-tail function of \(Y_i\) are
\[
        f_i(y)=\theta_i e^{-y}\exp(-\theta_i e^{-y}),
        \qquad
        \Pp(Y_i>y)=1-\exp(-\theta_i e^{-y}).
\]

\begin{lemma}[Gumbel threshold identities]\label{lem:gumbel-threshold-identities}
For a single specified label \(p\),
\[
        \Pp(\sigma(p)=q)
        =\int_{-\infty}^{\infty}
        f_p(y)\,
        \Pp\left(\sum_{j\ne p}\one_{\{Y_j>y\}}=q-1\right)\dd y .
\]
More generally, let \(p_1,\ldots,p_\ell\) be distinct labels and let
\(q_1<\cdots<q_\ell\) be distinct target positions.  Put
\(P=\{p_1,\ldots,p_\ell\}\) and
\[
        S_a(y)=\sum_{j\notin P}\one_{\{Y_j>y_a\}},
        \qquad 1\le a\le\ell .
\]
Then
\[
\begin{aligned}
&\Pp(\sigma(p_a)=q_a,\ 1\le a\le\ell) \\
&\quad=\int_{y_1>\cdots>y_\ell}
        \prod_{a=1}^\ell f_{p_a}(y_a)\,
        \Pp(S_a(y)=q_a-a,\ 1\le a\le\ell)
        \dd y_1\cdots\dd y_\ell .
\end{aligned}
\]
\end{lemma}

\begin{proof}
Condition on the specified priorities.  For one label, the event
\(\sigma(p)=q\) is exactly the event that precisely \(q-1\) nonspecified
priorities exceed the observed value of \(Y_p\).  This gives the first
identity.

For the multi-point identity, the ordering of the target positions forces the
specified priorities to satisfy \(Y_{p_1}>\cdots>Y_{p_\ell}\).  Conditional on
\(Y_{p_a}=y_a\) in this chamber, exactly \(a-1\) specified priorities are above
\(y_a\).  Hence \(\sigma(p_a)=q_a\) is equivalent to
\(S_a(y)=q_a-a\), simultaneously for \(1\le a\le\ell\).  Integrating against the
independent specified-priority densities gives the formula.
\end{proof}

\begin{remark}[The derivative cancellation]\label{rem:gumbel-derivative-cancellation}
The one-point identity explains the local-entry density used later.  Let
\[
        M_{n,P}(y)=\sum_{j\notin P}\Pp(Y_j>y)
        =\sum_{j\notin P}(1-e^{-\theta_j e^{-y}}).
\]
At a rank saddle \(y_*\), where \(M_{n,P}(y_*)\) equals the target rank up to
\(O(1)\), a one-dimensional local CLT gives formally
\[
        \Pp(\sigma(p)=q)
        \approx
        \frac{f_p(y_*)}{|M_{n,P}'(y_*)|}.
\]
Here
\[
        |M_{n,P}'(y_*)|
        =\sum_{j\notin P}\theta_j e^{-y_*}e^{-\theta_j e^{-y_*}},
\]
while
\[
        f_p(y_*)=\theta_p e^{-y_*}e^{-\theta_p e^{-y_*}}.
\]
Thus the common factor \(e^{-y_*}\) cancels.  In the regular-profile notation of
Section \ref{sec:regular}, the exact finite-$n$ saddle uses
$A_{n,P}(t)=(q-1)/n$.  Replacing $(q-1)/n$ by $q/n$ changes only a lower-order
term, and putting \(t=e^{-y_*}\) gives
\[
        \frac{f_p(y_*)}{|M_{n,P}'(y_*)|}
        =\frac1n\frac{\theta_p e^{-\theta_p t}}{B_{n,P}(t)}
        \longrightarrow
        \frac1n\rho(p/n,q/n).
\]
This is the Gumbel form of the exponential-clock saddle calculation in
Proposition \ref{prop:bulk-local-final}.
\end{remark}

\begin{remark}[Hyperplane-walk priorities]\label{rem:hyperplane-gumbel}
The same representation applies to the face-sampling step in the Brown--Diaconis
stationary construction for a finite hyperplane walk whose face measure is
separating.  If the faces in the support have positive weights \(w_F\), then
sampling them without replacement with probabilities proportional to the
remaining weights is equivalent to ordering independent priorities
\[
        \log w_F+G_F
\]
in decreasing order and taking their ordered Tits product.  The resulting
chamber is therefore a deterministic function of a Gumbel-priority order.  This
is an exact probabilistic representation, although it does not by itself give a
closed formula for arbitrary arrangement statistics; see
\cite{BHR1999,BD1998} for the face-semigroup construction.
\end{remark}

\subsection{Analytic input}
We use standard estimates for sums of independent Bernoulli variables.  They are stated explicitly to make clear what is invoked.

\begin{lemma}[Point probability and Bernstein bounds]\label{lem:bernstein-point}
Let \(Y=\sum_{j=1}^N \xi_j\), where the \(\xi_j\) are independent, not necessarily identically distributed, Bernoulli variables, and set \(\mu=\Ee Y\), \(v=\Var Y\).  There is an absolute constant \(C\) such that
\[
        \sup_m \Pp(Y=m)\le \frac{C}{\sqrt{1+v}}.
\]
Moreover, there is an absolute constant \(c>0\) such that, for every \(u\ge0\),
\[
        \Pp(|Y-\mu|\ge u)\le
        2\exp\left[-c\min\left(\frac{u^2}{v+1},u\right)\right].
\]
\end{lemma}

\begin{proof}
The point bound is the classical concentration-function estimate for sums of Bernoulli variables.  It follows from Fourier inversion and the standard inequality
\[
        |1-p+pe^{it}|
        \le \exp\left[-\frac{2}{\pi^2}p(1-p)\operatorname{dist}(t,2\pi\mathbb Z)^2\right],
\]
where \(\operatorname{dist}(t,2\pi\mathbb Z)\) denotes distance from \(t\) to the nearest multiple of \(2\pi\).  Bernstein's inequality gives the second bound for bounded independent summands.  See Petrov \cite[Chapter VII]{Petrov1975} for both estimates.
\end{proof}

\begin{lemma}[Uniform lattice local CLT with a tilted local bound]
\label{lem:llt}
Fix a dimension $d$ and a support bound $B<\infty$.  Let
\[
        Y_{n,m}=\sum_{j=1}^{N_n}\zeta_{n,m,j},
        \qquad m\in\mathcal M_n,
\]
be sums of independent $\mathbb Z^d$-valued random vectors satisfying
$\|\zeta_{n,m,j}\|\le B$.  Write
\[
        \mu_{n,m}=\Ee Y_{n,m},
        \qquad \Sigma_{n,m}=\Var(Y_{n,m}),
\]
and suppose that, uniformly in $m$,
\[
        c s_n I_d\le\Sigma_{n,m}\le C s_n I_d,
        \qquad s_n\longrightarrow\infty.
\]
For $\lambda\in\mathbb R^d$, let $\Pp_\lambda$ denote the product exponential
tilt
\[
 \Pp_\lambda(\zeta_{n,m,j}=x)
 =\frac{e^{\lambda\cdot x}\Pp(\zeta_{n,m,j}=x)}
 {\Ee e^{\lambda\cdot\zeta_{n,m,j}}}.
\]
Assume that there is $\lambda_0>0$ for which the following two properties hold
uniformly for $\|\lambda\|\le\lambda_0$:
\begin{enumerate}[label=(\roman*)]
\item the tilted covariance matrices satisfy
\[
        c_0s_nI_d\le\Var_\lambda(Y_{n,m})\le C_0s_nI_d;
\]
\item for every $\varepsilon>0$ there are $\kappa_\varepsilon>0$ and
$n_\varepsilon$ such that
\[
 \left|\prod_{j=1}^{N_n}
 \Ee_\lambda e^{iu\cdot\zeta_{n,m,j}}\right|
 \le e^{-\kappa_\varepsilon s_n}
\]
whenever $n\ge n_\varepsilon$ and
$u\in[-\pi,\pi]^d$ has distance at least $\varepsilon$ from
$2\pi\mathbb Z^d$.
\end{enumerate}
Then, uniformly in $m$ and $z\in\mathbb Z^d$,
\begin{align*}
 \Pp(Y_{n,m}=z)
 &=(2\pi)^{-d/2}(\det\Sigma_{n,m})^{-1/2}\\
 &\quad\times
 \exp\left[-\frac12(z-\mu_{n,m})^T
 \Sigma_{n,m}^{-1}(z-\mu_{n,m})\right]
 +o(s_n^{-d/2}).
\end{align*}
Moreover,
\[
        \sup_z\Pp(Y_{n,m}=z)\le C_1s_n^{-d/2}.
\]
For every fixed $K<\infty$ and $\gamma<1/2$, there are $c_1,C_2>0$ such
that, uniformly when
\[
        \|z-\mu_{n,m}\|\le K s_n^{1/2+\gamma},
\]
one has
\[
 \Pp(Y_{n,m}=z)
 \le C_2s_n^{-d/2}
 \exp\left[-c_1\frac{\|z-\mu_{n,m}\|^2}{s_n}\right].
\]
All constants depend only on the fixed dimension, the support bound, and the
uniform constants in the hypotheses.
\end{lemma}

\begin{proof}
We suppress $n,m$ and write $Y,\mu,\Sigma,s$.  Let
$X_j=\zeta_j-\Ee\zeta_j$ and
\[
        \Psi(u)=\Ee e^{iu\cdot(Y-\mu)}.
\]
Since $\|X_j\|\le2B$ and
$\sum_j\Ee\|X_j\|^2=\operatorname{tr}\Sigma=O(s)$,
\[
        \sum_j\Ee\|X_j\|^3=O(s).
\]
Consequently, uniformly for $\|u\|$ sufficiently small,
\[
        \log\Psi(u)=-\frac12u^T\Sigma u+O(s\|u\|^3).
        \tag{2.1}
\]
A symmetrization argument also gives, on a fixed neighborhood of the origin,
\[
        |\Psi(u)|\le e^{-cs\|u\|^2}.
        \tag{2.2}
\]
Indeed, if $X_j'$ is an independent copy of $X_j$, then
$1-|\Ee e^{iu\cdot X_j}|^2=
\Ee[1-\cos(u\cdot(X_j-X_j'))]$; for small $u$ the right-hand side is
bounded below by a constant multiple of $\Var(u\cdot X_j)$, and multiplication
over $j$ yields (2.2).

Fourier inversion on $[-\pi,\pi]^d$ now proves the local expansion.  Choose
$L_s\to\infty$ with $L_s^3/\sqrt{s}\to0$ and split the integral into
\[
 \|u\|\le L_s/\sqrt{s},\qquad
 L_s/\sqrt{s}<\|u\|\le\varepsilon,
 \qquad \|u\|>\varepsilon.
\]
On the first region, (2.1), followed by the change of variables
$v=\Sigma^{1/2}u$, gives the Gaussian Fourier integral with an
$o(s^{-d/2})$ error.  The estimate (2.2) makes the second region
$o(s^{-d/2})$, and tilted aperiodicity with $\lambda=0$ makes the third
region exponentially small.  The oscillatory factor
$e^{-iu\cdot(z-\mu)}$ has modulus one, so every error is uniform in $z$.
Taking absolute values in the same inversion formula proves the concentration
bound.  This is the usual fixed-dimensional Fourier proof of the lattice local
central limit theorem; compare Petrov \cite[Chapter~VII]{Petrov1975}.

It remains to prove the Gaussian local upper bound.  Let
$h=z-\mu$ and define the centered cumulant generating function
\[
        K(\lambda)=\log\Ee e^{\lambda\cdot(Y-\mu)}.
\]
For $\|\lambda\|\le\lambda_0$, its Hessian is
$\nabla^2K(\lambda)=\Var_\lambda(Y)$ and is therefore between
$c_0sI_d$ and $C_0sI_d$.  If $\|h\|=o(s)$, strong convexity shows that the
minimum of $K(\lambda)-h\cdot\lambda$ over the ball
$\|\lambda\|\le\lambda_0/2$ is attained in its interior for all large $n$.
Thus there is a unique $\lambda_h$ in that ball satisfying
\[
        \nabla K(\lambda_h)=h,
        \qquad \|\lambda_h\|\le C\|h\|/s.
\]
The associated Legendre transform
$I(h)=\lambda_h\cdot h-K(\lambda_h)$ satisfies
\[
        I(h)\ge c\|h\|^2/s.
\]
For example, the last inequality follows by testing the variational formula
$I(h)=\sup_\lambda\{\lambda\cdot h-K(\lambda)\}$ at a sufficiently small
multiple of $h/s$ and using the tilted covariance upper bound.

Under $\Pp_{\lambda_h}$, the mean of $Y$ is $z$.  Repeating the absolute-value
Fourier estimate above, now using the two tilted hypotheses, gives
\[
        \sup_w\Pp_{\lambda_h}(Y=w)\le Cs^{-d/2}.
\]
The change-of-measure identity therefore yields
\[
 \Pp(Y=z)=e^{-I(h)}\Pp_{\lambda_h}(Y=z)
 \le Cs^{-d/2}e^{-c\|h\|^2/s}.
\]
The stated displacement range is $o(s)$ because $\gamma<1/2$, completing the
proof.
\end{proof}



\subsection{A permuton-level obstruction}\label{subsec:implantation}

The regularity assumptions in the cycle theorems below cannot be replaced by
permuton convergence alone.  A vanishing exceptional set of labels can carry an
arbitrary sublinear permutation while leaving the limiting profile unchanged
almost everywhere.

\begin{proposition}[Vanishing-block implantation]\label{prop:implantation}
Let $m_n\to\infty$ with $m_n=o(n)$, and let $\pi_n\in S_{m_n}$ be arbitrary.
There are Luce weights $\theta_{n,1},\ldots,\theta_{n,n}$ such that
\[
        \Pp\bigl(\sigma_n(i)=\pi_n(i)
        \text{ for every }1\le i\le m_n\bigr)\longrightarrow1,
\]
while the step profiles
\[
        f_n(x)=\theta_{n,\lceil nx\rceil},\qquad 0<x\le1,
\]
converge almost everywhere to the constant profile $f\equiv1$.  Consequently,
$\sigma_n$ has the uniform permuton limit,
but its cycle structure on $m_n=o(n)$ prescribed labels can be chosen
arbitrarily.
\end{proposition}

\begin{proof}
Write $\tau_r=\pi_n^{-1}(r)$, $1\le r\le m_n$, and choose $R_n$ so that
$R_n/n\to\infty$.  Set
\[
        \theta_{n,\tau_r}=R_n^{m_n-r+1},
        \qquad 1\le r\le m_n,
        \qquad
        \theta_{n,i}=1,
        \quad i>m_n.
\]
Let $X_{n,i}\sim\Exp(\theta_{n,i})$ be independent, and let
$Y_n=\min_{i>m_n}X_{n,i}\sim\Exp(n-m_n)$.  For $1\le r<m_n$,
\[
        \Pp(X_{n,\tau_r}\ge X_{n,\tau_{r+1}})
        =\frac1{R_n+1},
\]
while
\[
        \Pp(X_{n,\tau_{m_n}}\ge Y_n)
        =\frac{n-m_n}{R_n+n-m_n}.
\]
A union bound therefore gives
\[
\begin{aligned}
&\Pp\bigl(X_{n,\tau_1}<\cdots<X_{n,\tau_{m_n}}<Y_n\bigr)\\
&\qquad\ge1-\frac{m_n-1}{R_n+1}
        -\frac{n-m_n}{R_n+n-m_n}\longrightarrow1,
\end{aligned}
\]
because $n/R_n\to0$ and
$m_n/R_n=(m_n/n)(n/R_n)\to0$.
On this event, the first $m_n$ labels in the Luce draw order are
$\tau_1,\ldots,\tau_{m_n}$, so
$\sigma_n(\tau_r)=r=\pi_n(\tau_r)$ for every $r$.

For every fixed $x>0$, the index $\lceil nx\rceil$ exceeds $m_n$ for all
large $n$, and hence $f_n(x)=1$.  Thus $f_n\to1$ almost everywhere.  On the
high-probability event above, the remaining $n-m_n$ labels all have equal
weight and therefore appear in a uniformly random order after the implanted
block.  The implanted points carry permuton mass $m_n/n=o(1)$, while a uniform
permutation of the remaining labels has the uniform permuton limit.  Hence the
whole permutation has the uniform permuton limit.  The final assertion follows
by prescribing the cycle decomposition of $\pi_n$.
\end{proof}

\begin{remark}\label{rem:implantation-consequences}
Taking $\pi_n$ to be the identity produces at least $m_n$ fixed points with
probability tending to one, despite the uniform permuton limit.  More generally,
one may implant $\lfloor m_n/r\rfloor$ cycles of any fixed length $r$, or one
cycle of length $m_n$.  Thus fixed points, fixed short-cycle counts, and the
total number of cycles on sublinear scales are not continuous statistics of the
permuton topology.  The uniformly positive continuous-profile assumptions in
Section \ref{sec:regular} exclude this type of microscopic exceptional block.
\end{remark}

\section{Top windows: exact boundary identities}

Normalize \(\sum_i\theta_i=1\).  Let \(P_k\) be the law of the first
\(k\) labels under Luce sampling without replacement, and let \(Q_k\) be the
product law of \(k\) independent draws from \([n]\) with probabilities
\(\theta_i\).  Define
\[
 \calD_k=\{(i_1,\ldots,i_k)\in[n]^k:
 i_1,\ldots,i_k\text{ are pairwise distinct}\}
\]
and
\[
 e_k(\theta_1,\ldots,\theta_n)
 =\sum_{1\le j_1<\cdots<j_k\le n}
 \theta_{j_1}\cdots\theta_{j_k}.
\]

\begin{theorem}[Exact top-window total variation]\label{thm:top-tv}
For every \(k\le n\),
\[
        \|P_k-Q_k\|_{\TV}
        =Q_k(\calD_k^c)
        =1-k!e_k(\theta_1,\ldots,\theta_n).
\]
\end{theorem}

\begin{proof}
The measure \(P_k\) is supported on \(\calD_k\).  If \(i=(i_1,\ldots,i_k)\in\calD_k\), then
\[
        P_k(i)=\prod_{r=1}^k
        \frac{\theta_{i_r}}{1-\theta_{i_1}-\cdots-\theta_{i_{r-1}}},
\]
with the empty sum equal to zero.  Hence \(P_k(i)\ge Q_k(i)\) on \(\calD_k\), while \(P_k(i)=0\) on
\(\calD_k^c\).  The positive and negative parts therefore have equal mass, and
\[
\begin{aligned}
2\|P_k-Q_k\|_{\TV}
 &=\sum_{i\in\calD_k}(P_k(i)-Q_k(i))
   +\sum_{i\notin\calD_k}Q_k(i)\\
 &=2Q_k(\calD_k^c).
\end{aligned}
\]
Thus $\|P_k-Q_k\|_{\TV}=1-Q_k(\calD_k)$.
Finally,
\[
        Q_k(\calD_k)=\sum_{\{j_1,\ldots,j_k\}}
        k!\theta_{j_1}\cdots\theta_{j_k}=k!e_k(\theta_1,\ldots,\theta_n).
\]
\end{proof}

Define the one-sided relative defect
\[
 d_\infty^+(P_k,Q_k)
 =\left\|1-\frac{\dd(Q_k|_{\calD_k})}{\dd P_k}
 \right\|_{L^\infty(P_k)}
 =\max_{i\in\calD_k}
 \left(1-\frac{Q_k(i)}{P_k(i)}\right).
\]
This is not a metric.  Let
\(\theta_{(1)}\ge\cdots\ge\theta_{(n)}\) be the decreasing rearrangement.

\begin{theorem}[Exact one-sided relative $L^\infty$ defect]
\label{thm:dinfty}
For \(k\le n\),
\[
        d_\infty^+(P_k,Q_k)
        =1-\prod_{r=1}^{k-1}
        \left(1-\sum_{j=1}^r\theta_{(j)}\right),
\]
where the empty product is one when $k=1$.
\end{theorem}

\begin{proof}
For \(i=(i_1,\ldots,i_k)\in\calD_k\),
\[
        \frac{Q_k(i)}{P_k(i)}
        =\prod_{r=1}^{k-1}
        \left(1-\sum_{j=1}^r\theta_{i_j}\right).
\]
For every \(r\), \(\sum_{j=1}^r\theta_{i_j}\le \sum_{j=1}^r\theta_{(j)}\).  Since each factor \(1-x\) is decreasing in \(x\), the product is minimized by taking the first \(k-1\) selected weights to be \(\theta_{(1)},\ldots,\theta_{(k-1)}\).  The final selected label does not enter the product.  This proves the formula.
\end{proof}

\begin{corollary}[Birthday transfer]\label{cor:birthday}
Let \(\xi_1,\ldots,\xi_k\) be independent with law \(\theta\), and let
\[
        W_k=\sum_{1\le r<s\le k}\one_{\{\xi_r=\xi_s\}}.
\]
Then
\[
        \|P_k-Q_k\|_{\TV}=\Pp(W_k\ge1).
\]
Consequently, every asymptotic estimate for the no-collision probability
$\Pp(W_k=0)$, equivalently for the collision probability $\Pp(W_k\ge1)$,
transfers immediately to the top-window total-variation distance.
\end{corollary}

\begin{proof}
Under \(Q_k\), the event \(\{W_k=0\}\) is precisely \(\calD_k\).  Hence
\[
        Q_k(W_k=0)=Q_k(\calD_k)=k!e_k(\theta_1,\ldots,\theta_n).
\]
Theorem \ref{thm:top-tv} therefore gives
\[
        \|P_k-Q_k\|_{\TV}=1-Q_k(\calD_k)=Q_k(W_k\ge1),
\]
which is the claimed identity.
\end{proof}

\section{Bottom windows and terminal leakage}

Let \((\theta_i)_{i\ge1}\) be positive rates and let \(X_i\sim\Exp(\theta_i)\) be independent.  Define
\[
        f(x)=\sum_{i=1}^{\infty}e^{-\theta_i x},
        \qquad
        x_0=\inf\{x:f(x)<\infty\}.
\]
The function $f$ is nonincreasing.  For a finite set $A\subset\mathbb N$ and
$x>0$, define
\[
 \prod_{i\notin A}(1-e^{-\theta_i x})
 :=\lim_{N\to\infty}
 \prod_{\substack{i\le N\\i\notin A}}(1-e^{-\theta_i x}).
\]
The limit always exists.  It is positive exactly when
$\sum_{i\notin A}e^{-\theta_i x}<\infty$, and otherwise it equals zero.  Indeed,
for $u_i=e^{-\theta_i x}\in(0,1)$, the partial products are decreasing, and
$\log(1-u_i)\sim-u_i$ once $u_i\to0$; if $u_i$ does not tend to zero, the
product plainly vanishes.

For distinct \(a_1,\ldots,a_k\), define
\[
\begin{aligned}
\mu_k(a_1,\ldots,a_k)
  :=& \int_{x_1>\cdots>x_k>0}
     \prod_{j=1}^k \theta_{a_j}e^{-\theta_{a_j}x_j}
     \prod_{i\notin\{a_1,\ldots,a_k\}}
       (1-e^{-\theta_i x_k})
     \dd x_1\cdots\dd x_k .
\end{aligned}
\]
The following finite-volume statement identifies $\mu_k$ as an actual limiting
subprobability law rather than only a formal candidate.

\begin{proposition}[Finite bottom-window convergence]
\label{prop:finite-bottom-convergence}
For $N\ge\max_j a_j$, let
$\mu_k^{(N)}(a_1,\ldots,a_k)$ be the probability that the $k$ largest values
among $X_1,\ldots,X_N$, listed in decreasing order, have labels
$a_1,\ldots,a_k$.  Then, for every fixed ordered distinct tuple,
\[
        \mu_k^{(N)}(a_1,\ldots,a_k)
        \longrightarrow\mu_k(a_1,\ldots,a_k).
\]
Thus $\mu_k$ is the coordinatewise limit of the finite bottom-window laws.
\end{proposition}

\begin{proof}
For $N\ge\max_j a_j$,
\begin{align*}
\mu_k^{(N)}(a_1,\ldots,a_k)
=&\int_{x_1>\cdots>x_k>0}
 \prod_{j=1}^k\theta_{a_j}e^{-\theta_{a_j}x_j}\\
&\quad\times
 \prod_{\substack{1\le i\le N\\i\notin\{a_1,\ldots,a_k\}}}
 (1-e^{-\theta_i x_k})
 \dd x_1\cdots\dd x_k.
\end{align*}
The finite products decrease pointwise to the defining infinite product and are
bounded by one.  The remaining factor is the joint density of the specified
independent clocks and is integrable over the chamber.  Dominated convergence
proves the assertion.
\end{proof}

\begin{theorem}[Terminal leakage]\label{thm:leakage}
Assume \(x_0<\infty\).  Let
\[
        N_0=\sum_{i=1}^{\infty}\one_{\{X_i>x_0\}}.
\]
Then
\[
        \sum_{a_1,\ldots,a_k}^{*}\mu_k(a_1,\ldots,a_k)
        =\Pp(N_0\ge k),
\]
where \(\sum^*\) denotes summation over distinct ordered \(k\)-tuples.  In
particular, if \(f(x_0)=\infty\), then the mass is one.  If
\(f(x_0)<\infty\), the escaped mass is \(\Pp(N_0<k)>0\), so the limiting law
is strictly defective.
\end{theorem}

\begin{proof}
For every \(x>x_0\), \(f(x)<\infty\).  Hence
\[
        \sum_i \Pp(X_i>x)=\sum_i e^{-\theta_i x}<\infty,
\]
so Borel--Cantelli implies that only finitely many of the events
\(\{X_i>x\}\) occur almost surely.  Applying this with
\(x=x_0+1/m\), \(m\ge1\), shows that the point set
\(\{X_i:X_i>x_0\}\) is locally finite in \((x_0,\infty)\).  Thus its
largest, second largest, and subsequent points are well-defined whenever
\(N_0\ge k\); possible accumulation can occur only at the endpoint \(x_0\),
not above it.  If \(x_0=0\), then \(N_0=\infty\) almost surely, since
\(X_i>0\) almost surely for every \(i\), and the same argument applies with
boundary \(x_k=0\).

For a fixed ordered distinct \(k\)-tuple \(a=(a_1,\ldots,a_k)\), the event that
the largest \(k\) clock values above \(x_0\), equivalently the ordered
bottom-window labels, are \(a_1,\ldots,a_k\) and have heights in
\(\dd x_1\cdots\dd x_k\), with \(x_1>\cdots>x_k>x_0\), has density
\[
     \prod_{j=1}^k \theta_{a_j}e^{-\theta_{a_j}x_j}
     \prod_{i\notin\{a_1,\ldots,a_k\}}(1-e^{-\theta_i x_k}).
\]
The first factor fixes the specified heights and the product says that every other variable is below \(x_k\).  Conversely, if this density contributes with \(x_k>x_0\), then these are exactly the largest \(k\) clock values above \(x_0\).  If \(x_k<x_0\), the infinite product is zero because \(f(x_k)=\infty\).  The hyperplane \(x_k=x_0\) has Lebesgue measure zero.  Therefore integrating and summing over all ordered distinct \(a\)'s partitions the event \(\{N_0\ge k\}\), giving the identity.

If \(f(x_0)=\infty\), then \(\sum_i\Pp(X_i>x_0)=\infty\), so
independence and the second Borel--Cantelli lemma imply \(N_0=\infty\) almost
surely.  If \(f(x_0)<\infty\), then necessarily $x_0>0$, and $N_0$ is a
finite Poisson-binomial variable with parameters $e^{-\theta_i x_0}$.  Since
these parameters are all strictly less than one and have finite sum,
\[
        \Pp(N_0=0)=\prod_{i=1}^\infty
        (1-e^{-\theta_i x_0})>0.
\]
Hence $\Pp(N_0<k)>0$.
\end{proof}

\begin{remark}[Complete leakage when $x_0=\infty$]
If $x_0=\infty$, then $f(x)=\infty$ for every finite $x$.  The infinite
product in the definition of $\mu_k$ therefore vanishes for every point of the
integration chamber, so $\mu_k\equiv0$.  In this case all finite-label
coordinate mass escapes.
\end{remark}

\begin{corollary}[A log-log phase diagram]\label{cor:loglog}
Suppose, as $i\to\infty$,
\[
        \theta_i=\beta\log i+\alpha\log\log i+O(1),
        \qquad \beta>0.
\]
Then \(x_0=1/\beta\), and the bottom fixed-window law is non-defective if and only if \(\alpha\le\beta\).
\end{corollary}

\begin{proof}
For large \(i\),
\[
        e^{-\theta_i x}\asymp i^{-\beta x}(\log i)^{-\alpha x}.
\]
Thus \(\sum_i e^{-\theta_i x}\) converges for \(x>1/\beta\) and diverges for \(x<1/\beta\), so \(x_0=1/\beta\).  At \(x_0\),
\[
        f(x_0)\asymp \sum_i\frac{1}{i(\log i)^{\alpha/\beta}},
\]
which diverges exactly when \(\alpha/\beta\le1\).  Apply Theorem \ref{thm:leakage}.
\end{proof}

\section{Regular positive profiles: fixed points and short cycles}\label{sec:regular}

Let \(\sigma_n\) be the label-to-position Luce permutation generated by independent exponentials
\[
        X_{n,i}\sim \Exp(\theta_{n,i}),\qquad 1\le i\le n,
\]
so that
\[
        \sigma_n(i)=1+\#\{j:X_{n,j}<X_{n,i}\}.
\]
Throughout this section assume
\[
        \theta_{n,i}=a_n(f(i/n)+\eps_{n,i}),
        \qquad \max_{1\le i\le n}|\eps_{n,i}|\to0,
\]
where
\[
        f\in C[0,1],\qquad 0<m_f\le f(x)\le M_f<\infty.
\]
Multiplying all rates by \(a_n\) only rescales the exponential-clock time axis
and does not change their order, so take \(a_n=1\).  Fix constants
\[
        0<\underline\theta<\overline\theta<\infty
\]
such that, for every sufficiently large $n$,
$\underline\theta\le\theta_{n,i}\le\overline\theta$ for all $i$; for
example, one may take $\underline\theta=m_f/2$ and
$\overline\theta=2M_f$.

Define
\[
        A(t)=\int_0^1(1-e^{-f(u)t})\dd u,
        \qquad
        B(t)=A'(t)=\int_0^1 f(u)e^{-f(u)t}\dd u.
\]
Then
\[
        A(0)=0,\qquad A'(t)=B(t)>0,\qquad A(t)\uparrow1
        \quad(t\to\infty),
\]
so $\tau=A^{-1}:(0,1)\to(0,\infty)$ is well defined.  The density
\[
        \rho(x,y)=\frac{f(x)e^{-f(x)\tau(y)}}{B(\tau(y))},
        \qquad 0\le x\le1,\quad 0<y<1,
\]
is the Luce permuton density identified by Borga, Chatterjee, and Diaconis
\cite[Proposition~2.2]{BCD2025}.  Here it will also arise as the microscopic
point-entry density.  It is strictly positive and continuous, and is bounded
on every strip $[0,1]\times[\delta,1-\delta]$.

For \(r\ge1\) define
\[
        \lambda_r
        =\lim_{\delta\downarrow0}\frac1r
        \int_{[\delta,1-\delta]^r}
        \rho(x_1,x_2)\rho(x_2,x_3)\cdots\rho(x_r,x_1)
        \dd x_1\cdots\dd x_r,
\]
where \(x_{r+1}=x_1\).  The proof below shows that these improper limits exist and are finite.

\begin{theorem}[Regular-profile fixed points and short cycles]\label{thm:regular-cycles}
For every fixed \(R\),
\[
        (C_1(\sigma_n),\ldots,C_R(\sigma_n))
        \Rightarrow (Z_1,\ldots,Z_R),
\]
where \(Z_1,\ldots,Z_R\) are independent Poisson variables with means
\(\lambda_1,\ldots,\lambda_R\).  In particular,
\[
        C_1(\sigma_n)\Rightarrow\Pois(\lambda_1),
        \qquad
        \lambda_1=\lim_{\delta\downarrow0}\int_\delta^{1-\delta}\rho(x,x)\dd x .
\]
\end{theorem}

The proof has three parts.  First, we prove a uniform local estimate for finitely many point constraints \(\sigma_n(p_a)=q_a\) when the target positions lie in a separated bulk region.  Second, we prove endpoint and collision envelopes showing that the complement of this separated bulk region has negligible total contribution in fixed-order factorial-moment sums.  Third, we sum the separated estimate over fixed points and over disjoint directed cycles.  Thus the argument uses only fixed-order factorial moments; no global total-variation approximation is needed.

\subsection{Separated bulk local estimate}

For a finite set of specified labels \(P\subset[n]\), define empirical functions
\[
        A_{n,P}(t)=\frac1n\sum_{j\notin P}(1-e^{-\theta_{n,j}t}),
        \qquad
        B_{n,P}(t)=A_{n,P}'(t)=\frac1n\sum_{j\notin P}\theta_{n,j}e^{-\theta_{n,j}t}.
\]
If \(|P|\) is fixed, then \(A_{n,P}\to A\) and \(B_{n,P}\to B\) uniformly on compact subsets of \([0,\infty)\).

\begin{lemma}[Uniform saddle geometry]\label{lem:regular-saddle-geometry}
Fix \(\ell\ge1\), \(\delta>0\), and \(\eta>0\).  Uniformly over sets \(P\subset[n]\) of cardinality at most \(\ell\) and over integers
\[
        q_1<\cdots<q_\ell,
        \qquad q_a/n\in[\delta,1-\delta],
        \qquad q_{a+1}-q_a\ge \eta n,
\]
put \(y_a=(q_a-a)/n\).  For all large \(n\), the equations
\[
        A_{n,P}(t_{n,a}^*)=y_a,
        \qquad 1\le a\le\ell,
\]
have unique solutions.  These solutions lie in a compact interval
\([c_\delta,C_\delta]\subset(0,\infty)\), satisfy
\[
        t_{n,a}^*=\tau(q_a/n)+o(1),
        \qquad
        B_{n,P}(t_{n,a}^*)=B(\tau(q_a/n))+o(1),
\]
uniformly, and are separated:
\[
        t_{n,a+1}^*-t_{n,a}^*\ge c_{\delta,\eta}>0,
        \qquad 1\le a<\ell .
\]
\end{lemma}

\begin{proof}
Because \(A\) is continuous, strictly increasing, \(A(0)=0\), and \(A(t)\uparrow1\), choose \(0<c_\delta<C_\delta<\infty\) so that
\[
        A(c_\delta)<\delta/3,
        \qquad
        A(C_\delta)>1-\delta/3 .
\]
For large \(n\), uniformly over bounded-size \(P\), we have
\[
        \sup_{0\le t\le C_\delta}|A_{n,P}(t)-A(t)|\to0.
\]
Also \(y_a\in[\delta/2,1-\delta/2]\) for all large \(n\), uniformly in the allowed positions.  Hence \(A_{n,P}(c_\delta)<y_a<A_{n,P}(C_\delta)\).  Since \(|P|\le\ell\) and the rates are uniformly bounded in \([\underline\theta,\overline\theta]\), for all large \(n\),
\[
        B_{n,P}(t)
        =\frac1n\sum_{j\notin P}\theta_{n,j}e^{-\theta_{n,j}t}
        \ge \frac{\underline\theta}{2}e^{-\overline\theta t}>0
\]
on bounded intervals.  Thus \(A_{n,P}\) is strictly increasing, and the solution exists uniquely in \([c_\delta,C_\delta]\).

Uniform convergence of \(A_{n,P}\) to \(A\), together with the uniform continuity of \(\tau=A^{-1}\) on \([\delta/2,1-\delta/2]\), gives
\[
        t_{n,a}^*=\tau(y_a)+o(1)=\tau(q_a/n)+o(1),
\]
uniformly.  Uniform convergence of \(B_{n,P}\) to \(B\) on \([c_\delta,C_\delta]\) gives the displayed convergence of derivatives.

Finally, \(\tau'(y)=1/A'(\tau(y))\).  Since \(A'\) is bounded above on \([c_\delta,C_\delta]\), \(\tau'\) is bounded below on \([\delta/2,1-\delta/2]\).  Moreover,
\[
        y_{a+1}-y_a=\frac{q_{a+1}-q_a-1}{n}\ge \eta/2
\]
for all large \(n\).  Hence the limiting saddle points \(\tau(y_a)\) are separated by a positive constant depending only on \(\delta\) and \(\eta\).  The uniform convergence just proved transfers this lower bound to \(t_{n,a}^*\) for all large \(n\).
\end{proof}

\begin{lemma}[Uniform LLT for Luce cumulative counts]\label{lem:luce-cumulative-llt}
Fix \(\ell\ge1\), \(0<c<C<\infty\), \(\eta>0\), and \(K<\infty\).  Let \(P\subset[n]\) satisfy \(|P|\le K\), and let
\[
        c\le t_1<\cdots<t_\ell\le C,
        \qquad t_{a+1}-t_a\ge \eta .
\]
For independent \(X_{n,j}\sim\Exp(\theta_{n,j})\), define
\[
        S_a(t)=\sum_{j\notin P}\one_{\{X_{n,j}\le t_a\}},
        \qquad 1\le a\le\ell,
\]
and \(S(t)=(S_1(t),\ldots,S_\ell(t))\).  Then \(S(t)\) satisfies the fixed-dimensional local CLT of Lemma \ref{lem:llt} uniformly over the above parameters with scale \(s_n=n\).  In particular, if \(\Sigma_{n,P}(t)=n^{-1}\Var S(t)\), then
\[
        c_0 I_\ell\le \Sigma_{n,P}(t)\le C_0 I_\ell
\]
for constants depending only on \(\ell,c,C,\eta,K,\underline\theta,\overline\theta\), and uniformly for
\(\|z-\Ee S(t)\|=O(n^{1/2+\gamma})\), \(\gamma<1/6\),
\[
\begin{aligned}
        \Pp(S(t)=z)
        &=(2\pi n)^{-\ell/2}(\det\Sigma_{n,P}(t))^{-1/2} \\
        &\quad\times
        \exp\left[-\frac1{2n}(z-\Ee S(t))^T\Sigma_{n,P}(t)^{-1}(z-\Ee S(t))\right]
        +o(n^{-\ell/2}).
\end{aligned}
\]
Moreover, in the same range,
\[
        \Pp(S(t)=z)
        \le C n^{-\ell/2}
        \exp\left[-c\frac{\|z-\Ee S(t)\|^2}{n}\right].
\]
\end{lemma}

\begin{proof}
For a single nonspecified label $j$, the increment vector
\[
 V_j(t)=(\one_{\{X_{n,j}\le t_1\}},\ldots,
 \one_{\{X_{n,j}\le t_\ell\}})
\]
takes the $\ell+1$ adjacent cumulative values
\[
 (1,\ldots,1),\ (0,1,\ldots,1),\ldots,
 (0,\ldots,0,1),\ (0,\ldots,0).
\]
The corresponding categories are
\[
 (0,t_1],\ (t_1,t_2],\ldots,(t_{\ell-1},t_\ell],\ (t_\ell,\infty).
\]
Because the rates lie in
$[\underline\theta,\overline\theta]$, the times stay in $[c,C]$, and
successive times are separated by at least $\eta$, every category probability
is bounded below by a constant $c_*>0$, uniformly in all parameters.

The adjacent support differences are the coordinate vectors.  For any
$u\in\mathbb R^\ell$, the variance identity
\[
 \Var(u\cdot V_j)
 =\frac12\sum_{v,w}p_vp_w\{u\cdot(v-w)\}^2
\]
therefore gives
\[
        \Var(u\cdot V_j)\ge c_*^2\sum_{a=1}^\ell u_a^2.
\]
Summing over $n-|P|$ labels gives the covariance lower bound, and bounded
increments give the upper bound.

Fix $\lambda_0<\infty$.  Under any exponential tilt
$\|\lambda\|\le\lambda_0$, the probabilities of the finitely many atoms are
multiplied by factors lying between positive constants depending only on
$\ell$ and $\lambda_0$.  Thus every tilted category probability is still
bounded below uniformly, and the same covariance comparison holds under all
such tilts.  The same observation gives tilted strong aperiodicity.  Indeed,
for $u\in[-\pi,\pi]^\ell$ away from $2\pi\mathbb Z^\ell$, at least one
adjacent phase difference $u_a$ is bounded away from zero.  The two
corresponding adjacent atoms have uniformly positive tilted mass, so compactness
gives
\[
        |\Ee_\lambda e^{iu\cdot V_j(t)}|\le1-\kappa
\]
for some $\kappa>0$.  Multiplication over $n-|P|$ labels yields the exponential
Fourier bound required in Lemma \ref{lem:llt}.  All hypotheses of that lemma
are therefore satisfied with scale $s_n=n$, proving the stated local expansion
and Gaussian local upper bound.
\end{proof}

\begin{lemma}[Off-saddle exclusion]\label{lem:regular-off-saddle}
Fix \(\ell\ge1\), \(\delta>0\), \(\eta>0\), and \(\gamma\in(0,1/6)\).  In the setting of Lemma \ref{lem:regular-saddle-geometry}, let
\[
        K_\delta=[c_\delta,C_\delta]
\]
be a compact interval containing all saddles.  In the integral representation for finitely many point constraints below, the contribution from
\[
        \max_a |t_a-t_{n,a}^*|>n^{-1/2+\gamma}
\]
is \(o(n^{-\ell})\), uniformly over the allowed labels and positions.
\end{lemma}

\begin{proof}
Enlarge \(K_\delta\) slightly, if necessary, so that outside \(K_\delta\) the limiting mean \(A(t)\) is separated from \([\delta/2,1-\delta/2]\).  Uniform convergence of \(A_{n,P}\) to \(A\) then implies that if some \(t_a\notin K_\delta\), the corresponding Bernoulli sum has mean separated from its target by order \(n\).  Bernstein's inequality gives an exponentially small bound for the marginal event \(\{S_a(t)=q_a-a\}\).  The full joint point-constraint event is contained in this marginal event, so the same bound applies to the integrand.  Since the specified-clock density integrates to at most one over the full positive chamber, this part is \(O(e^{-c n})\).

It remains to consider \(t\in K_\delta^\ell\) with \(\max_a |t_a-t_{n,a}^*|>n^{-1/2+\gamma}\).  On \(K_\delta\), the derivatives \(B_{n,P}\) are uniformly bounded below.  Hence, for at least one coordinate \(a\),
\[
        |nA_{n,P}(t_a)-nA_{n,P}(t_{n,a}^*)|
        \ge c n^{1/2+\gamma}.
\]
Since \(nA_{n,P}(t_{n,a}^*)=q_a-a\), Bernstein's inequality applied to the one-dimensional count \(S_a(t)\) gives
\[
        \Pp(S_a(t)=q_a-a)\le \exp(-c n^{2\gamma}).
\]
The full joint point-constraint event is contained in this one-dimensional marginal event, so the same bound applies to the integrand.  The specified-clock density is uniformly bounded on \(K_\delta^\ell\), and the volume is bounded.  Thus the contribution of this region is exponentially small, hence \(o(n^{-\ell})\).
\end{proof}


\begin{lemma}[Gaussian-scale saddle tail]\label{lem:regular-gaussian-tail}
Fix \(\ell\ge1\), \(\delta>0\), \(\eta>0\), and \(\gamma\in(0,1/6)\).  Let \(p_1,\ldots,p_\ell\) be distinct labels and let
\(q_1<\cdots<q_\ell\) be target positions satisfying
\[
        q_a/n\in[\delta,1-\delta],\qquad q_{a+1}-q_a\ge \eta n .
\]
Let \(t_n^*=(t_{n,1}^*,\ldots,t_{n,\ell}^*)\) be the saddle vector from Lemma \ref{lem:regular-saddle-geometry}.  In the corresponding point-constraint integral, put
\[
        S(t)=\bigl(S_1(t),\ldots,S_\ell(t)\bigr),
        \qquad z_n=(q_1-1,\ldots,q_\ell-\ell).
\]
For
\[
        t=t_n^*+h/\sqrt n,\qquad \|h\|\le n^\gamma,
\]
the integrand
\[
        \prod_{a=1}^\ell\theta_{n,p_a}e^{-\theta_{n,p_a}t_a}\,
        \Pp(S(t)=z_n)
\]
is bounded by
\[
        C n^{-\ell/2}e^{-c\|h\|^2},
\]
with constants independent of the allowed labels and positions.  Consequently, uniformly in those labels and positions, the part of the saddle window with \(\|h\|>R\) contributes at most
\[
        C n^{-\ell}\int_{\|h\|>R}e^{-c\|h\|^2}\dd h .
\]
\end{lemma}

\begin{proof}
For \(\|h\|\le n^\gamma\), the saddle separation from Lemma \ref{lem:regular-saddle-geometry} implies that the coordinates of \(t\) remain ordered and separated for all large \(n\).  Moreover, \(t\) remains in a fixed compact subinterval of \((0,\infty)\).  Hence Lemma \ref{lem:luce-cumulative-llt}, including the Gaussian local upper bound inherited from Lemma \ref{lem:llt}, applies uniformly.

Let
\[
        A_n(t)=\bigl(A_{n,P}(t_1),\ldots,A_{n,P}(t_\ell)\bigr),
        \qquad z_n=(q_1-1,\ldots,q_\ell-\ell).
\]
Since \(A_n(t_n^*)=z_n/n\), Taylor's theorem and the uniform \(C^2\) bounds for \(A_{n,P}\) on the compact saddle region give
\[
        nA_n(t)-z_n
        =\sqrt n\,DA_n(t_n^*)h+O(\|h\|^2).
\]
The diagonal matrix \(DA_n(t_n^*)\) has entries \(B_{n,P}(t_{n,a}^*)\), uniformly bounded above and below.  Since \(\gamma<1/2\), the quadratic remainder is \(o(\sqrt n\|h\|)\) uniformly when \(\|h\|\le n^\gamma\) and \(\|h\|\to\infty\).  For bounded \(\|h\|\), the desired Gaussian upper bound follows from the uniform local bound \(C n^{-\ell/2}\), after decreasing \(c\).  Thus, in the range relevant for Gaussian decay,
\[
        \|nA_n(t)-z_n\|\ge c\sqrt n\|h\|.
\]
The Gaussian local upper bound gives
\[
        \Pp(S(t)=z_n)\le C n^{-\ell/2}e^{-c\|h\|^2}.
\]
The specified-clock density \(\prod_a\theta_{n,p_a}e^{-\theta_{n,p_a}t_a}\) is uniformly bounded on the compact region, which proves the integrand bound.  Since \(\dd t=n^{-\ell/2}\dd h\), the displayed tail estimate follows.
\end{proof}

\begin{lemma}[Uniform integrated lattice saddle]\label{lem:integrated-lattice-saddle}
Fix \(\ell\ge1\), \(\gamma\in(0,1/6)\), and an open set \(U\subset\Rr^\ell\).  Let \(K\Subset U\) be compact.  Suppose \(t_n^*\in K\), \(z_n\in\mathbb Z^\ell\), and \(S_n(t)\) is an integer-valued random vector such that
\[
        \Ee S_n(t)=nA_n(t),\qquad A_n(t_n^*)=z_n/n.
\]
Assume, uniformly for \(t_n^*\in K\), that:
\begin{enumerate}
\item \(A_n\) is \(C^2\) in a fixed neighborhood of \(K\), with uniformly bounded second derivatives;
\item \(D_n=DA_n(t_n^*)\) is nonsingular and \(\|D_n^{-1}\|\) is uniformly bounded;
\item the covariance matrices \(\Sigma_n(t)=n^{-1}\Var S_n(t)\) are uniformly positive definite and uniformly continuous on a fixed neighborhood of \(K\);
\item a uniform local CLT holds for \(S_n(t)\) when \(t=t_n^*+O(n^{-1/2+\gamma})\);
\item the contribution outside \(\|t-t_n^*\|\le n^{-1/2+\gamma}\) is \(o(n^{-\ell})\), and inside this window the Gaussian tail bound
\[
        |g_n(t)|\Pp(S_n(t)=z_n)\le C n^{-\ell/2}e^{-c\|h\|^2},
        \qquad t=t_n^*+h/\sqrt n,
\]
holds for some \(C,c>0\).
\end{enumerate}
If \(g_n\) is uniformly bounded and equicontinuous near \(K\), then
\[
        \int_U g_n(t)\Pp(S_n(t)=z_n)\dd t
        =n^{-\ell}g_n(t_n^*)|\det D_n|^{-1}+o(n^{-\ell}),
\]
uniformly over the family.
\end{lemma}

\begin{proof}
The off-saddle hypothesis reduces the integral to \(t=t_n^*+h/\sqrt n\) with \(\|h\|\le n^\gamma\).  By the Gaussian tail hypothesis, the part with \(\|h\|>R\) is bounded by
\[
        C n^{-\ell}\int_{\|h\|>R}e^{-c\|h\|^2}\dd h,
\]
which is \(o_R(n^{-\ell})\) uniformly as \(R\to\infty\).  It therefore suffices to evaluate the integral on each fixed ball \(\|h\|\le R\).

On such a ball, Taylor's theorem gives
\[
        z_n-nA_n(t_n^*+h/\sqrt n)
        =-\sqrt n\,D_nh+O(1),
\]
uniformly.  The covariance matrices satisfy
\[
        \Sigma_n(t_n^*+h/\sqrt n)=\Sigma_n(t_n^*)+o(1)
\]
uniformly on \(\|h\|\le R\), and \(g_n(t_n^*+h/\sqrt n)=g_n(t_n^*)+o(1)\) uniformly.  The local CLT yields
\[
\begin{aligned}
        \Pp(S_n(t_n^*+h/\sqrt n)=z_n)
        &=(2\pi n)^{-\ell/2}(\det\Sigma_n(t_n^*))^{-1/2}  \\
        &\quad\times
        \exp\left[-\frac12 h^T D_n^T\Sigma_n(t_n^*)^{-1}D_nh\right]
        +o(n^{-\ell/2}),
\end{aligned}
\]
uniformly on the fixed ball.  All \(o(1)\) terms in this fixed-ball argument are uniform over the family and over \(\|h\|\le R\).  Since \(\dd t=n^{-\ell/2}\dd h\), the contribution from \(\|h\|\le R\) equals
\[
\begin{aligned}
        &n^{-\ell}g_n(t_n^*)(2\pi)^{-\ell/2}(\det\Sigma_n(t_n^*))^{-1/2} \\
        &\qquad\times
        \int_{\|h\|\le R}
        \exp\left[-\frac12 h^T D_n^T\Sigma_n(t_n^*)^{-1}D_nh\right]\dd h
        +o_R(n^{-\ell}).
\end{aligned}
\]
Letting \(R\to\infty\), uniformly in the family, the Gaussian integral is
\[
        (2\pi)^{\ell/2}(\det\Sigma_n(t_n^*))^{1/2}|\det D_n|^{-1}.
\]
Indeed,
\[
        \det(D_n^T\Sigma_n(t_n^*)^{-1}D_n)
        =(\det D_n)^2(\det\Sigma_n(t_n^*))^{-1},
\]
so the covariance determinant cancels and the Gaussian integral contributes \(|\det D_n|^{-1}\).  Combining this with the preceding tail estimate proves the claim.
\end{proof}

\begin{proposition}[Separated bulk local estimate]\label{prop:bulk-local-final}
Fix \(\ell\ge1\), \(\delta>0\), and \(\eta>0\).  Uniformly over distinct labels \(p_1,\ldots,p_\ell\) and distinct positions \(q_1,\ldots,q_\ell\) satisfying
\[
        q_a/n\in[\delta,1-\delta],
        \qquad |q_a-q_b|\ge \eta n\quad(a\ne b),
\]
we have
\[
        \Pp(\sigma_n(p_a)=q_a,1\le a\le\ell)
        =n^{-\ell}\prod_{a=1}^\ell
        \rho(p_a/n,q_a/n)+o(n^{-\ell}).
\]
\end{proposition}

\begin{proof}
Reorder the pairs so that \(q_1<\cdots<q_\ell\), relabeling the corresponding \(p_a\)'s.  Conditional on \(X_{n,p_a}=t_a\) with \(0<t_1<\cdots<t_\ell\), define
\[
        S_a(t)=\sum_{j\notin\{p_1,\ldots,p_\ell\}}\one_{\{X_{n,j}\le t_a\}}.
\]
The event \(\sigma_n(p_a)=q_a\) for all \(a\) is equivalent to
\[
        S_a(t)=q_a-a,
        \qquad 1\le a\le\ell,
\]
because exactly \(a-1\) specified clocks occur before \(t_a\).  Thus
\[
\begin{aligned}
\Pp(\sigma_n(p)=q)
  &=\int_{0<t_1<\cdots<t_\ell}
     \prod_{a=1}^\ell\theta_{n,p_a}e^{-\theta_{n,p_a}t_a}  \\
  &\quad\times
     \Pp(S_a(t)=q_a-a,1\le a\le\ell)\dd t_1\cdots\dd t_\ell .
\end{aligned}
\]
Put \(P=\{p_1,\ldots,p_\ell\}\), \(y_a=(q_a-a)/n\), and let \(t_{n,a}^*\) be the saddle point from Lemma \ref{lem:regular-saddle-geometry}.  The ordered integration domain contains a fixed neighborhood of \(t_n^*=(t_{n,1}^*,\ldots,t_{n,\ell}^*)\), because the saddle coordinates are uniformly separated.

For fixed \(t\), the vector \(S(t)=(S_1(t),\ldots,S_\ell(t))\) has mean
\[
        \Ee S_a(t)=nA_{n,P}(t_a),
\]
and covariance \(n\Sigma_{n,P}(t)\), where
\[
        \Sigma_{n,P,ab}(t)=\frac1n\sum_{j\notin P}
        (1-e^{-\theta_{n,j}t_{a\wedge b}})e^{-\theta_{n,j}t_{a\vee b}}.
\]
In the saddle neighborhood, Lemma \ref{lem:luce-cumulative-llt} gives the required uniform local CLT.  Lemma \ref{lem:regular-gaussian-tail} supplies the Gaussian-scale tail bound, and Lemma \ref{lem:regular-off-saddle} removes the complementary region.  Apply Lemma \ref{lem:integrated-lattice-saddle} with
\[
        A_n(t)=\bigl(A_{n,P}(t_1),\ldots,A_{n,P}(t_\ell)\bigr),
        \qquad
        z_n=(q_1-1,\ldots,q_\ell-\ell),
\]
and
\[
        g_n(t)=\prod_{a=1}^\ell\theta_{n,p_a}e^{-\theta_{n,p_a}t_a}.
\]
The functions \(g_n\) are uniformly bounded and equicontinuous in the saddle neighborhood, since the specified rates are uniformly bounded and the neighborhood is compact.  Thus the hypotheses of Lemma \ref{lem:integrated-lattice-saddle} are satisfied uniformly over the separated bulk configurations.  Here
\[
        DA_n(t_n^*)=
        \diag(B_{n,P}(t_{n,1}^*),\ldots,B_{n,P}(t_{n,\ell}^*)).
\]
Therefore the integral equals
\[
        n^{-\ell}\prod_{a=1}^\ell
        \frac{\theta_{n,p_a}e^{-\theta_{n,p_a}t_{n,a}^*}}{B_{n,P}(t_{n,a}^*)}
        +o(n^{-\ell}).
\]
Finally, Lemma \ref{lem:regular-saddle-geometry} and the uniform convergence of the rates give
\[
        \theta_{n,p_a}=f(p_a/n)+o(1),
        \qquad
        t_{n,a}^*=\tau(q_a/n)+o(1),
        \qquad
        B_{n,P}(t_{n,a}^*)=B(\tau(q_a/n))+o(1),
\]
uniformly.  This is exactly
\[
        n^{-\ell}\prod_{a=1}^\ell
        \rho(p_a/n,q_a/n)+o(n^{-\ell}).
\]
\end{proof}

\subsection{Endpoint and collision domination}

We next prove global domination estimates sufficient for the factorial-moment sums.  Let a Luce permutation of \(N\) labels have all rates in \([\underline\theta,\overline\theta]\), and let \(R_i\) be the rank of label \(i\).  Put
\[
        \alpha=\frac{\underline\theta}{\overline\theta}\in(0,1].
\]
For \(1\le q\le N\), define
\[
        w_N(q)=
        \begin{cases}
        N^{-1}, & 1\le q\le N/2,\\[4pt]
        N^{-\alpha}(N-q+1)^{\alpha-1}, & N/2<q\le N.
        \end{cases}
\]

\begin{lemma}[One-point rank envelope]\label{lem:one-rank-envelope}
There is \(C=C(\underline\theta,\overline\theta)\) such that, for all \(N\), all rate vectors in \([\underline\theta,\overline\theta]^N\), all labels \(i\), and all \(q\),
\[
        \Pp(R_i=q)\le Cw_N(q).
\]
\end{lemma}

\begin{proof}
First suppose \(q\le N/2\).  Reveal the Luce order sequentially.  Conditional on label \(i\) not having appeared before position \(q\), at position \(q\) there are at least \(N-q+1\ge N/2\) labels remaining.  The total remaining weight is at least \(m(N-q+1)\), and the weight of label \(i\) is at most \(M\).  Hence
\[
        \Pp(R_i=q)
        \le \frac{\overline\theta}{\underline\theta(N-q+1)}
        \le \frac{2\overline\theta}{\underline\theta N}.
\]
Thus \(\Pp(R_i=q)\le C N^{-1}\) for \(q\le N/2\).

Now suppose \(q>N/2\).  Put \(s=N-q+1\).  The event \(R_i=q\) means that label \(i\) survives until there are \(s\) labels remaining and is then selected.  When there are \(r\) labels remaining and \(i\) is among them, the conditional probability that \(i\) is selected at the next draw is at least
\[
        \frac{\underline\theta}{\overline\theta r}=\frac{\alpha}{r}.
\]
Therefore the conditional probability that \(i\) is not selected at that draw is at most \(1-\alpha/r\).  Hence the probability that \(i\) survives from \(N\) labels down to \(s\) labels is at most
\[
        \prod_{r=s+1}^{N}\left(1-\frac{\alpha}{r}\right).
\]
Indeed,
\[
        \sum_{r=s+1}^{N}\frac1r\ge \log(N/s)-O(1),
\]
and therefore
\[
        \prod_{r=s+1}^{N}\left(1-\frac{\alpha}{r}\right)
        \le \exp\left(-\alpha\sum_{r=s+1}^{N}\frac1r\right)
        \le C_\alpha\left(\frac{s}{N}\right)^\alpha.
\]
Given survival to \(s\) labels, the conditional probability that \(i\) is selected next is at most
\[
        \min\left(1,\frac{\overline\theta}{\underline\theta s}\right)\le \frac{1}{\alpha s}.
\]
Therefore
\[
        \Pp(R_i=q)
        \le C\left(\frac{s}{N}\right)^\alpha\frac1s
        =C N^{-\alpha}s^{\alpha-1}.
\]
Since \(s=N-q+1\), this is the desired bound.
\end{proof}

\begin{lemma}[Gap envelope]\label{lem:gap-envelope}
Fix \(\ell\ge1\).  There is \(C_\ell\) such that the following holds.  Let \(p_1,\ldots,p_\ell\) be distinct labels and \(q_1,\ldots,q_\ell\) distinct positions in a Luce permutation of \([n]\) with rates in \([\underline\theta,\overline\theta]\).  Sort the target positions:
\[
        q_{(1)}<q_{(2)}<\cdots<q_{(\ell)}.
\]
Set \(q_{(0)}=0\), \(r_a=q_{(a)}-q_{(a-1)}\), and \(N_a=n-q_{(a-1)}\).  Then
\[
        \Pp(\sigma_n(p_a)=q_a,1\le a\le\ell)
        \le C_\ell\prod_{a=1}^\ell w_{N_a}(r_a).
\]
\end{lemma}


\begin{proof}
Let \(p_{(a)}\) denote the specified label paired with the sorted target position \(q_{(a)}\).  Reveal the Luce order from left to right, stopping at the specified target positions.  At the first specified position \(q_{(1)}\), the label \(p_{(1)}\) must appear at relative rank \(r_1=q_{(1)}\) in a Luce order of \(N_1=n\) labels.  By Lemma \ref{lem:one-rank-envelope}, this has probability at most \(Cw_{N_1}(r_1)\).

At step \(a\), condition on the full history revealed up to \(q_{(a-1)}\).  If \(p_{(a)}\) has already appeared, the conditional probability of the joint event is zero.  If \(p_{(a)}\) is still alive, the Plackett--Luce deletion property says that the unexposed tail is again a Luce order on the remaining labels, with the remaining rates.  This property follows immediately from the exponential-clock representation and memorylessness: after the revealed clocks are removed, the residual order of the unrevealed labels is generated by independent exponential clocks with their original rates.  In this tail there are \(N_a=n-q_{(a-1)}\) labels, and \(p_{(a)}\) must appear at relative rank
\[
        r_a=q_{(a)}-q_{(a-1)}.
\]
Lemma \ref{lem:one-rank-envelope} gives a conditional probability at most \(Cw_{N_a}(r_a)\).  Iterating over \(a=1,\ldots,\ell\) gives the stated bound, with \(C_\ell=C^\ell\).
\end{proof}

\begin{lemma}[Summability of the gap envelope]\label{lem:gap-sum}
There is $C<\infty$ such that
\[
        \sum_{r=1}^{N}w_N(r)\le C
\]
for every $N$.  Also, for $1\le h\le N/2$,
\[
        \sum_{r>N-h}w_N(r)\le C\left(\frac hN\right)^\alpha.
\]
For $\ell\ge1$, define
\[
H_\ell(N,h)=
\sum_{\substack{r_1,\ldots,r_\ell\ge1\\
 r_1+\cdots+r_\ell\le N\\
 N-r_1-\cdots-r_\ell\le h}}
\prod_{a=1}^\ell
w_{N-r_1-\cdots-r_{a-1}}(r_a),
\]
where the empty partial sum is zero.  Then, for every $0<\beta<\alpha$,
\[
        H_\ell(N,h)\le C_{\ell,\beta}
        \min\left\{1,\left(\frac hN\right)^\beta\right\},
        \qquad N\ge1,\ h\ge1.
\]
\end{lemma}

\begin{proof}
The first estimate follows from
\[
 \sum_{r\le N/2}N^{-1}\le\frac12
\]
and, after writing $s=N-r+1$,
\[
 \sum_{r>N/2}N^{-\alpha}(N-r+1)^{\alpha-1}
 =N^{-\alpha}\sum_{s\le N/2}s^{\alpha-1}\le C.
\]
The same substitution gives the terminal-tail estimate.

We prove the final assertion by induction on $\ell$.  If $h\ge N/2$, repeated
use of the first estimate gives $H_\ell(N,h)\le C^\ell$, which is bounded by
the asserted right-hand side after increasing the constant.  It therefore
suffices to assume $h<N/2$.

For $\ell=1$, the defining condition implies $r>N-h-1$, so the terminal-tail
bound, with a harmless endpoint term, gives
$H_1(N,h)\le C(h/N)^\alpha\le C(h/N)^\beta$.
Assume the result for $\ell-1$.  After the first gap $r$, put $s=N-r$.
The induction hypothesis gives
\[
 H_\ell(N,h)
 \le C\sum_{r=1}^{N-1}w_N(r)
 \min\left\{1,\left(\frac h{N-r}\right)^\beta\right\}.
\]
For $s\le h$, the contribution is at most $C(h/N)^\alpha$.  For
$h<s<N/2$, one has $r>N/2$ and hence
\[
\begin{aligned}
 \sum_{h<s<N/2}w_N(N-s)\left(\frac hs\right)^\beta
 &\le C N^{-\alpha}h^\beta
 \sum_{h<s<N/2}s^{\alpha-1-\beta}\\
 &\le C(h/N)^\beta.
\end{aligned}
\]
Finally, for $s\ge N/2$ one has $r\le N/2$, so $w_N(r)=N^{-1}$ and each
summand is at most $CN^{-1}(h/N)^\beta$.  Summing over at most $N$ values
completes the induction.
\end{proof}

Let
\[
        E_{\delta,n}
        =\{1,\ldots,\lfloor\delta n\rfloor\}
        \cup
        \{\lceil(1-\delta)n\rceil,\ldots,n\}.
\]

\begin{lemma}[Endpoint domination]\label{lem:endpoint-dom}
Fix \(\ell\ge1\), \(0<\beta<\alpha\), and \(0<\delta<1/2\).  Then, for all sufficiently large \(n\),
\[
\sum_{\substack{q_1,
\ldots,q_\ell\ {\rm distinct}\\
q_a\in E_{\delta,n}\ {\rm for\ some}\ a}}
\sup_{p_1,
\ldots,p_\ell}
\Pp(\sigma_n(p_a)=q_a,1\le a\le\ell)
\le C_{\ell,\beta}(\delta+\delta^\beta),
\]
where \(C_{\ell,\beta}\) is independent of \(n\) and \(\delta\).
\end{lemma}

\begin{proof}
Sort the target positions.  There are at most \(\ell!\) possible orderings, so it suffices to sum over ordered target positions and use the gap envelope.

If some target position lies in the first endpoint block, then \(q_{(1)}\le\delta n\), equivalently \(r_1\le\delta n\).  For \(\delta<1/2\),
\[
        \sum_{r_1\le\delta n}w_n(r_1)
        \le \sum_{r_1\le\delta n}\frac1n\le\delta.
\]
The remaining \(\ell-1\) gaps have total envelope mass bounded by a constant by Lemma \ref{lem:gap-sum}.  Thus this endpoint block contributes at most \(C_\ell\delta\).

If some target position lies in the last endpoint block, then \(q_{(\ell)}\ge(1-\delta)n\).  Equivalently, after the last specified target position, at most \(\delta n\) labels remain.  For fixed \(\delta>0\) and all sufficiently large \(n\), we have \(1\le\delta n\le n/2\), so Lemma \ref{lem:gap-sum} applies with \(h=\delta n\).  The final estimate in that lemma gives total envelope mass at most \(C_{\ell,\beta}\delta^\beta\).  Multiplying by the finite number of orderings proves the lemma.
\end{proof}



\begin{lemma}[Bulk collision domination]\label{lem:collision-dom}
Fix \(\ell\ge1\) and \(\delta>0\).  Let \(B_{\delta,\eta,n}\) be the set of distinct \(\ell\)-tuples \(q_1,\ldots,q_\ell\) such that
\[
        q_a/n\in[\delta,1-\delta]
\]
for all \(a\), but \(|q_a-q_b|\le\eta n\) for some \(a\ne b\).  Then
\[
\sum_{q\in B_{\delta,\eta,n}}
\sup_{p_1,\ldots,p_\ell}
\Pp(\sigma_n(p_a)=q_a,1\le a\le\ell)
\le C_{\ell,\delta}\eta.
\]
\end{lemma}

\begin{proof}
Sort the positions.  If two target positions are within \(\eta n\), then some adjacent ordered pair has gap \(r_j=q_{(j)}-q_{(j-1)}\le\eta n\) with \(j\ge2\).  Since all target positions are in \([\delta n,(1-\delta)n]\), before that gap the remaining deck size satisfies \(N_j\ge\delta n\).  If \(\eta<\delta/2\), then \(r_j\le N_j/2\) and
\[
        \sum_{r_j\le\eta n}w_{N_j}(r_j)
        \le \frac{\eta n}{N_j}
        \le \frac{\eta}{\delta}.
\]
For all other gaps, we sum the gap envelope iteratively and use the uniform bound
\[
        \sum_{r=1}^N w_N(r)\le C
\]
from Lemma \ref{lem:gap-sum}.  This gives a factor depending only on \(\ell\), uniformly in the intermediate deck sizes.  Summing over the possible adjacent pair and orderings gives the result.  The case \(\eta\ge\delta/2\) is absorbed by increasing the constant.
\end{proof}

\begin{corollary}[Endpoint tightness and improper integrability]
\label{cor:endpoint-tightness}
Fix $\ell\ge1$.  For every $0<\beta<\alpha$,
\[
\lim_{\delta\downarrow0}\limsup_{n\to\infty}
\sum_{\substack{q_1,\ldots,q_\ell\ {\rm distinct}\\
q_a\in E_{\delta,n}\ {\rm for\ some}\ a}}
\sup_{p_1,\ldots,p_\ell}
\Pp(\sigma_n(p_a)=q_a,1\le a\le\ell)=0.
\]
Moreover, all limiting integrals obtained from fixed-order fixed-point and
cycle factorial moments by Proposition \ref{prop:bulk-local-final} exist as
finite improper integrals.
\end{corollary}

\begin{proof}
The displayed endpoint-tightness assertion follows immediately from Lemma
\ref{lem:endpoint-dom} by first taking $\limsup_{n\to\infty}$ and then
letting $\delta\downarrow0$.

For the integrability assertion, fix $\delta,\eta>0$.  On the truncated region
$[\delta,1-\delta]^\ell$ with pairwise target separation greater than $\eta$,
Proposition \ref{prop:bulk-local-final} identifies the relevant finite sums
with Riemann sums for the corresponding product of $\rho$-factors.  Lemma
\ref{lem:collision-dom} bounds the omitted finite-$n$ collision contribution by
$O_{\ell,\delta}(\eta)$.  Since $\rho$ is bounded and continuous on the
truncated strip, the limiting integral over the collision neighborhood also
vanishes as $\eta\downarrow0$.  Thus the full truncated integral is obtained by
first taking $n\to\infty$ and then $\eta\downarrow0$.

The gap envelope and Lemma \ref{lem:gap-sum} bound the complete fixed-order
finite sums uniformly in $n$.  Endpoint tightness shows that the portion lost
outside $[\delta,1-\delta]^\ell$ tends to zero as $\delta\downarrow0$.
Consequently the nonnegative truncated integrals are uniformly bounded and
increase to finite limits.  This applies in particular to
\[
 \int_\delta^{1-\delta}\rho(x,x)\dd x
\]
and to every cyclic trace integral
\[
 \int_{[\delta,1-\delta]^r}
 \rho(x_1,x_2)\cdots\rho(x_r,x_1)
 \dd x_1\cdots\dd x_r,
\]
as well as to their finite products in mixed factorial moments.
\end{proof}

\subsection{Factorial moments}

\begin{proposition}[Fixed-point factorial moments]\label{prop:regular-fixed-factorial}
For every fixed \(k\ge1\),
\[
        \Ee(C_1(\sigma_n))_k\to \lambda_1^k.
\]
\end{proposition}

\begin{proof}
Write
\[
        C_1(\sigma_n)=\sum_{i=1}^n\one_{\{\sigma_n(i)=i\}}.
\]
For fixed \(k\),
\[
        \Ee(C_1(\sigma_n))_k
        =\sum_{i_1,\ldots,i_k}^{*}
        \Pp(\sigma_n(i_a)=i_a,1\le a\le k),
\]
where \(*\) means that the indices are distinct.  Fix \(0<\delta<1/2\) and \(\eta>0\).  Split the sum into endpoint terms, non-separated bulk terms, and separated bulk terms.  Lemma \ref{lem:endpoint-dom} bounds the endpoint terms by \(O_k(\delta+\delta^\beta)\), for any fixed \(0<\beta<\alpha\).  Lemma \ref{lem:collision-dom} bounds the non-separated bulk terms by \(O_{k,\delta}(\eta)\).  On the separated bulk part, Proposition \ref{prop:bulk-local-final} gives
\[
        \Pp(\sigma_n(i_a)=i_a,1\le a\le k)
        =n^{-k}\prod_{a=1}^k\rho(i_a/n,i_a/n)+o(n^{-k}),
\]
uniformly.  The error is uniform over the separated bulk configurations, and there are \(O(n^k)\) such configurations, so its total contribution is \(o(1)\).  Hence, as \(n\to\infty\), the separated bulk contribution tends to
\[
        \left(\int_\delta^{1-\delta}\rho(x,x)\dd x\right)^k,
\]
up to an error that vanishes as \(\eta\downarrow0\).  Indeed, on \([\delta,1-\delta]\) the function \(x\mapsto\rho(x,x)\) is bounded and continuous, so the contribution of the \(\eta\)-neighborhood of the diagonals in \([\delta,1-\delta]^k\) is \(O_{k,\delta}(\eta)\).  Letting \(\eta\downarrow0\) and then \(\delta\downarrow0\), and using Corollary \ref{cor:endpoint-tightness}, gives
\[
        \Ee(C_1(\sigma_n))_k\to \lambda_1^k.
\]
\end{proof}

\begin{proposition}[Mixed short-cycle factorial moments]\label{prop:regular-cycle-factorial}
For every fixed \(R\ge1\) and fixed nonnegative integers \(k_1,\ldots,k_R\),
\[
        \Ee\prod_{r=1}^R(C_r(\sigma_n))_{k_r}
        \to \prod_{r=1}^R\lambda_r^{k_r}.
\]
\end{proposition}

\begin{proof}
Let \(\calC_{n,r}\) be the set of directed \(r\)-cycles on \([n]\), counted up to cyclic rotation.  For \(\gamma=(i_1,\ldots,i_r)\in\calC_{n,r}\), let
\[
        E_\gamma=\{\sigma_n(i_1)=i_2,
        \sigma_n(i_2)=i_3,
        \ldots,
        \sigma_n(i_r)=i_1\}.
\]
Then \(C_r(\sigma_n)=\sum_{\gamma\in\calC_{n,r}}\one_{E_\gamma}\).  Put \(L=\sum_{r=1}^R r k_r\).  Expanding the mixed factorial moment gives a sum over ordered collections of distinct cycles.  For \(k_r\) cycles of the same length \(r\), the factorial moment \((C_r(\sigma_n))_{k_r}\) keeps those cycles ordered, so no additional \(k_r!\) denominator appears; the only normalization inside a single directed \(r\)-cycle is the factor \(1/r\) for cyclic rotations.  If two distinct cycles share a label, their events are incompatible; hence only disjoint collections contribute.

For a disjoint collection, the event imposes \(L\) distinct point constraints \(\sigma_n(p_a)=q_a\).  The labels \(p_a\) are distinct, and the target positions \(q_a\) are the same labels cyclically shifted within each cycle, hence are also distinct.  The separated local estimate requires separation of the target positions; no separate metric separation of the labels is needed.

Fix \(0<\delta<1/2\) and \(\eta>0\).  Decompose the sum according to whether some target position lies in \(E_{\delta,n}\), whether all target positions lie in the bulk but two are within \(\eta n\), or whether the target positions are separated in the bulk.  For fixed \(R,k_1,\ldots,k_R\), the number of ways to arrange a given ordered \(L\)-tuple of distinct point constraints into the prescribed ordered cycle collection is bounded by a constant depending only on \(R\) and the \(k_r\)'s.  Therefore Lemmas \ref{lem:endpoint-dom} and \ref{lem:collision-dom} bound the endpoint terms by \(O_{R,k}(\delta+\delta^\beta)\) and the non-separated bulk terms by \(O_{R,k,\delta}(\eta)\).  On separated bulk configurations, Proposition \ref{prop:bulk-local-final} applies.

Consider first a single \(r\)-cycle.  On the separated bulk, the local estimate gives the Riemann sum
\[
        \frac1r
        \int_{[\delta,1-\delta]^r}
        \rho(x_1,x_2)\rho(x_2,x_3)\cdots\rho(x_r,x_1)
        \dd x_1\cdots\dd x_r,
\]
where the factor \(1/r\) appears because an ordered \(r\)-tuple represents the same directed cycle under \(r\) cyclic rotations.  For an ordered collection of disjoint cycles, the local estimate gives
\[
        n^{-L}\prod_{\gamma}\prod_{a=1}^{|\gamma|}
        \rho(x_{\gamma,a},x_{\gamma,a+1})+o(n^{-L}),
        \qquad x_{\gamma,|\gamma|+1}=x_{\gamma,1},
\]
uniformly on the separated bulk.  The disjointness restriction only removes configurations lying on equality diagonals among the \(L\) labels.  There are \(O(n^{L-1})\) such grid configurations, and on the truncated bulk the limiting product density is bounded, so their Riemann-sum weight is \(O(n^{L-1})\cdot O(n^{-L})=O(n^{-1})\).  Thus the corresponding limiting integral factors over the cycle components.  After \(n\to\infty\) and \(\eta\downarrow0\), the contribution from separated bulk collections tends to the product of the truncated cycle means.  Finally \(\delta\downarrow0\) is justified by Corollary \ref{cor:endpoint-tightness}.  We obtain
\[
        \Ee\prod_{r=1}^R(C_r(\sigma_n))_{k_r}
        \to \prod_{r=1}^R\lambda_r^{k_r}.
\]
\end{proof}

\begin{proof}[Proof of Theorem \ref{thm:regular-cycles}]
Propositions \ref{prop:regular-fixed-factorial} and
\ref{prop:regular-cycle-factorial} give convergence of all fixed mixed
factorial moments:
\[
        \Ee\prod_{r=1}^R(C_r(\sigma_n))_{k_r}
        \to
        \prod_{r=1}^R\lambda_r^{k_r}.
\]
These are exactly the mixed factorial moments of independent Poisson variables
\(Z_r\sim\Pois(\lambda_r)\).  Since ordinary mixed moments are finite linear
combinations of mixed factorial moments, all fixed mixed ordinary moments also
converge to those of \((Z_1,\ldots,Z_R)\).  The limiting vector has finite
exponential moments and is therefore moment-determinate, so the multivariate
method of moments gives
\[
        (C_1(\sigma_n),\ldots,C_R(\sigma_n))
        \Rightarrow (Z_1,\ldots,Z_R),
\]
with independent \(Z_r\sim\Pois(\lambda_r)\).
\end{proof}

\begin{corollary}[Total-variation convergence for fixed cycle vectors]\label{cor:regular-cycles-tv}
For every fixed $R$,
\[
 d_{\TV}\!\left(
 \mathcal L(C_1(\sigma_n),\ldots,C_R(\sigma_n)),
 \bigotimes_{r=1}^R\Pois(\lambda_r)
 \right)\longrightarrow0.
\]
\end{corollary}

\begin{proof}
Let $p_n$ and $p$ be the probability mass functions of the two laws on the
countable discrete space $\Zz_{\ge0}^R$.  Theorem \ref{thm:regular-cycles} implies
$p_n(z)\to p(z)$ for every $z\in\Zz_{\ge0}^R$.  Given $\varepsilon>0$, choose a
finite set $A\subset\Zz_{\ge0}^R$ with $p(A)>1-\varepsilon$.  Then
$\sum_{z\in A}|p_n(z)-p(z)|\to0$ and $p_n(A^c)\to p(A^c)$.  Hence
\[
 \sum_{z\in\Zz_{\ge0}^R}|p_n(z)-p(z)|
 \le \sum_{z\in A}|p_n(z)-p(z)|+p_n(A^c)+p(A^c)
 \le 3\varepsilon+o(1).
\]
Letting $\varepsilon\downarrow0$ proves the claim.
\end{proof}

\subsection{Bistochasticity and long trace universality}\label{subsec:long-traces}

The fixed-length cycle means admit an operator interpretation that also determines
their large-length behavior.  Define the integral operator $K$ on
$L^2([0,1])$ by
\[
        (Kh)(x)=\int_0^1\rho(x,y)h(y)\dd y,
\]
and let $\Pi$ be the rank-one projection onto the constants,
\[
        (\Pi h)(x)=\int_0^1h(y)\dd y.
\]

\begin{lemma}[Bistochastic Hilbert--Schmidt kernel]\label{lem:bistochastic-hs}
The kernel $\rho$ is bistochastic:
\[
        \int_0^1\rho(x,y)\dd y=1
        \quad\text{for every }x\in[0,1],
        \qquad
        \int_0^1\rho(x,y)\dd x=1
        \quad\text{for every }y\in(0,1).
\]
Moreover, $K$ is a positivity-preserving Hilbert--Schmidt Markov contraction on $L^2([0,1])$ and
\[
        \|K\|_{\mathrm{HS}}^2
        =\int_{[0,1]^2}\rho(x,y)^2\dd x\dd y
        \le \frac{M_f}{m_f}.
\]
In particular, $K$ is compact, $K\one=\one$, $K^*\one=\one$, and
$K\Pi=\Pi K=\Pi$.
\end{lemma}

\begin{proof}
Put $y=A(t)$, so that $\dd y=B(t)\dd t$.  Then
\[
\begin{aligned}
        \int_0^1\rho(x,y)\dd y
        &=\int_0^\infty f(x)e^{-f(x)t}\dd t=1,\\
        \int_0^1\rho(x,y)\dd x
        &=\frac{\int_0^1f(x)e^{-f(x)\tau(y)}\dd x}
        {B(\tau(y))}=1.
\end{aligned}
\]
Thus $K\one=K^*\one=\one$.  Jensen's inequality and the two marginal
identities give
\[
        |Kh(x)|^2\le\int_0^1\rho(x,y)|h(y)|^2\dd y,
        \qquad
        \|Kh\|_2^2\le\|h\|_2^2,
\]
so $K$ is a contraction.

For the Hilbert--Schmidt norm, again use $y=A(t)$:
\[
\begin{aligned}
        \|K\|_{\mathrm{HS}}^2
        &=\int_0^\infty
        \frac{\int_0^1f(x)^2e^{-2f(x)t}\dd x}{B(t)}\dd t\\
        &\le M_f\int_0^\infty\frac{B(2t)}{B(t)}\dd t.
\end{aligned}
\]
Since $f\ge m_f$,
\[
        B(2t)=\int_0^1f(x)e^{-f(x)t}e^{-f(x)t}\dd x
        \le e^{-m_ft}B(t),
\]
and hence $\|K\|_{\mathrm{HS}}^2\le M_f/m_f$.  Hilbert--Schmidt operators are
compact.  The identities involving $\Pi$ follow from bistochasticity.
\end{proof}

\begin{lemma}[Cyclic trace formula]\label{lem:cyclic-trace-formula}
Let $K$ be a Hilbert--Schmidt integral operator on $L^2([0,1])$ with
nonnegative kernel $\rho\in L^2([0,1]^2)$.  For every $r\ge2$, $K^r$ is trace
class and
\[
 \operatorname{Tr}(K^r)
 =\int_{[0,1]^r}\rho(x_1,x_2)\rho(x_2,x_3)\cdots
 \rho(x_r,x_1)\dd x_1\cdots\dd x_r.
\]
\end{lemma}

\begin{proof}
Approximate $\rho$ in $L^2$ by nonnegative step kernels $\rho_m$, and let
$K_m$ be the corresponding finite-rank operators.  The formula is the ordinary
matrix trace identity for $K_m$.  Since $K_m\to K$ in Hilbert--Schmidt norm and
the Hilbert--Schmidt norms remain bounded, the Schatten ideal inequality
\[
        \|AB\|_1\le\|A\|_2\|B\|_2
\]
and a telescoping expansion give $K_m^r\to K^r$ in trace norm.  Hence
$\operatorname{Tr}(K_m^r)\to\operatorname{Tr}(K^r)$.  Applying the same
Cauchy--Schwarz inequalities to the cyclic products shows that their integrals
converge to the displayed integral.  This proves both its finiteness and the
identity.
\end{proof}

\begin{theorem}[Long trace universality]\label{thm:long-trace}
There are constants $C<\infty$ and $q\in(0,1)$, depending only on the limiting
profile $f$, such that, for every $r\ge2$,
\[
        \boxed{\displaystyle
        \lambda_r=\frac1r+O\!\left(\frac{q^r}{r}\right).}
\]
Equivalently,
\[
        r\lambda_r=1+O(q^r).
\]
Consequently,
\[
        \lim_{r\to\infty}\lim_{n\to\infty}
        r\,\Ee C_r(\sigma_n)=1,
        \qquad
        \sum_{r=2}^\infty\left|\lambda_r-\frac1r\right|<\infty.
\]
\end{theorem}

\begin{proof}
We first identify the peripheral spectrum of $K$.  Suppose $Kh=\zeta h$ with
$|\zeta|=1$.  Since $K$ is a contraction,
\[
        \|h\|_2=\|Kh\|_2.
\]
Equality must therefore hold in the pointwise Jensen inequality used in the
proof of Lemma \ref{lem:bistochastic-hs} for almost every $x$.  Because
$\rho(x,y)>0$ for every $x\in[0,1]$ and every $y\in(0,1)$, equality in Jensen's
inequality forces $h(y)$ to be almost everywhere constant.  It follows that
$\zeta=1$.  Thus the only eigenvectors with eigenvalue of modulus one are the
constants.  The eigenvalue $1$ is algebraically simple as well.  Indeed, if
$(K-I)g=c\one$ with $c\ne0$, then $K^ng=g+nc\one$, contradicting
$\|K^n\|\le1$.

Since $K$ is compact, its nonzero spectrum consists of isolated eigenvalues of
finite algebraic multiplicity, with possible accumulation only at zero.  Hence
there is $q\in(0,1)$ such that every eigenvalue of $K$ other than $1$ has
modulus at most $q$.  Put $T=K-\Pi$.  Then $T$ is Hilbert--Schmidt,
$T\Pi=\Pi T=0$, and
\[
        K^r=\Pi+T^r,
        \qquad r\ge1.
\]

For $r\ge2$, Lemma \ref{lem:cyclic-trace-formula} gives
\[
        \operatorname{Tr}(K^r)
        =\int_{[0,1]^r}
        \rho(x_1,x_2)\rho(x_2,x_3)\cdots\rho(x_r,x_1)
        \dd x_1\cdots\dd x_r
        =r\lambda_r.
\]
The equality with the improper integral defining $\lambda_r$ follows by
monotone convergence and Corollary \ref{cor:endpoint-tightness}.

Let $(\zeta_j)_{j\ge2}$ be the nontrivial eigenvalues of $K$, repeated according
to algebraic multiplicity.  Lidskii's theorem \cite[Theorem~3.7]{Simon2005} gives
\[
        r\lambda_r=\operatorname{Tr}(K^r)
        =1+\operatorname{Tr}(T^r)
        =1+\sum_{j\ge2}\zeta_j^r.
\]
Weyl's eigenvalue inequality for Schatten classes
\cite[Theorem~2.3]{Simon2005} gives
\[
        \sum_{j\ge2}|\zeta_j|^2\le\|T\|_{\mathrm{HS}}^2.
\]
Therefore
\[
        \left|r\lambda_r-1\right|
        \le q^{r-2}\sum_{j\ge2}|\zeta_j|^2
        \le q^{r-2}\|T\|_{\mathrm{HS}}^2.
\]
After changing the constant, this is $O(q^r)$.  The two consequences follow
immediately.
\end{proof}


\subsection{Diagonal bands and transverse curves}\label{subsec:regular-diagonal-consequences}

The separated local estimate has several consequences for sparse deterministic
entry sets.  We first record the fixed-width diagonal-band law, and then prove a
master sparse-entry theorem.  The latter removes the transversality and
smoothness assumptions from the thin-curve statement and yields deterministic
overlap laws as a further corollary.

\begin{theorem}[Fixed-width diagonal bands]\label{thm:diagonal-band}
Under the hypotheses of Theorem \ref{thm:regular-cycles}, fix an integer
\(h\ge0\), and define
\[
        N_h=\#\{i:1\le i\le n,\ |\sigma_n(i)-i|\le h\}.
\]
Then
\[
        N_h\Rightarrow
        \Pois\left((2h+1)\lambda_1\right),
        \qquad
        \lambda_1=\lim_{\delta\downarrow0}
        \int_\delta^{1-\delta}\rho(x,x)\dd x .
\]
\end{theorem}

\begin{proof}
Write
\[
        N_h=\sum_{d=-h}^h\sum_{\substack{1\le i\le n\\1\le i+d\le n}}
        \one_{\{\sigma_n(i)=i+d\}} .
\]
Fix \(k\ge1\).  Expanding \(\Ee(N_h)_k\) gives a sum over ordered \(k\)-tuples
\((i_a,d_a)\).  Terms with repeated labels or repeated target positions are
zero unless they are identical summands, which are excluded by the factorial
expansion.  Thus the contributing terms impose \(k\) distinct point constraints
\(\sigma_n(i_a)=i_a+d_a\).

Fix \(0<\delta<1/2\) and \(\eta>0\).  Endpoint target positions contribute
\(O_{k,h}(\delta+\delta^\beta)\), for every \(0<\beta<\alpha\), by Lemma
\ref{lem:endpoint-dom}; the finite number of offsets only changes the constant.
Among the remaining bulk terms, the contribution from target positions within
\(\eta n\) of one another is \(O_{k,h,\delta}(\eta)\) by Lemma
\ref{lem:collision-dom}.  On the separated bulk part, Proposition
\ref{prop:bulk-local-final} gives, uniformly,
\[
        \Pp(\sigma_n(i_a)=i_a+d_a,\ 1\le a\le k)
        =n^{-k}\prod_{a=1}^k
        \rho(i_a/n,(i_a+d_a)/n)+o(n^{-k}).
\]
For each fixed offset \(d\), uniformly on the truncated bulk,
\[
        \rho(i/n,(i+d)/n)=\rho(i/n,i/n)+o(1),
\]
because \(\rho\) is continuous on compact subintervals of
\([0,1]\times(0,1)\).  Hence, after summing over the \((2h+1)^k\) offset
choices and passing to Riemann sums, the separated bulk contribution tends to
\[
        \left((2h+1)\int_\delta^{1-\delta}\rho(x,x)\dd x\right)^k,
\]
up to an error vanishing as \(\eta\downarrow0\).  Letting
\(\eta\downarrow0\) and then \(\delta\downarrow0\), using Corollary
\ref{cor:endpoint-tightness}, gives
\[
        \Ee(N_h)_k\to ((2h+1)\lambda_1)^k .
\]
These are the factorial moments of \(\Pois((2h+1)\lambda_1)\), which proves the
claim by the method of moments.
\end{proof}

\subsection{Sparse entry sets, thin graphs, and overlaps}\label{subsec:sparse-entry}

The preceding diagonal arguments are instances of a single sparse-entry theorem.
For a deterministic set $E_n\subset[n]^2$, put
\[
        N(E_n)=\sum_{(i,j)\in E_n}\one_{\{\sigma_n(i)=j\}},
        \qquad
        \nu_n=\frac1n\sum_{(i,j)\in E_n}
        \delta_{(i/n,j/n)}.
\]
The measures $\nu_n$ need not be probability measures.

\begin{theorem}[Sparse bulk entry sets]\label{thm:sparse-entry}
Assume the hypotheses of Theorem \ref{thm:regular-cycles}.  Fix
$\delta\in(0,1/2)$, and suppose that $E_n\subset[n]\times
[\delta n,(1-\delta)n]$ satisfies
\[
        |E_n|=O(n),
        \qquad
        \max_{1\le i\le n}\#\{j:(i,j)\in E_n\}=o(n),
\]
and
\[
        \lim_{\eta\downarrow0}\limsup_{n\to\infty}
        \frac1n\sup_{q\in[n]}
        \#\{(i,j)\in E_n:|j-q|\le\eta n\}=0.
\]
If $\nu_n\Rightarrow\nu$ weakly as finite measures on $[0,1]^2$, then
\[
        N(E_n)\Rightarrow
        \Pois\!\left(\int_{[0,1]^2}\rho(x,y)\,\nu(\dd x,\dd y)\right).
\]
\end{theorem}

\begin{proof}
Fix $k\ge1$.  Expanding the factorial moment gives
\[
        \Ee(N(E_n))_k
        =\sum_{e_1,\ldots,e_k\in E_n}^{*}
        \Pp\bigl(\sigma_n(i_a)=j_a,
        1\le a\le k\bigr),
        \qquad e_a=(i_a,j_a),
\]
where the star excludes repeated edges.  If two distinct selected edges have the
same row or the same target column, the corresponding event is impossible.

We first record a uniform bulk bound.  If $j_1,\ldots,j_k$ are distinct and all
belong to $[\delta n,(1-\delta)n]$, then Lemma
\ref{lem:gap-envelope} gives
\[
        \sup_{i_1,\ldots,i_k}
        \Pp(\sigma_n(i_a)=j_a,1\le a\le k)
        \le C_{k,\delta}n^{-k}.
\]
Indeed, after sorting the target positions, every remaining deck size is at
least $\delta n$, and every gap in Lemma \ref{lem:gap-envelope} is separated
from the terminal end of its remaining deck by at least $\delta n$; hence each
gap factor is $O_\delta(n^{-1})$.

Let
\[
        D_n=\max_i\#\{j:(i,j)\in E_n\},
        \qquad
        M_n(\eta)=\sup_q
        \#\{(i,j)\in E_n:|j-q|\le\eta n\}.
\]
The number of ordered $k$-tuples containing two distinct edges in the same row
is at most
\[
        C_k D_n|E_n|^{k-1}=o(n^k).
\]
The number containing a pair of target columns at distance at most $\eta n$ is
at most
\[
        C_k |E_n|^{k-1}M_n(\eta).
\]
Consequently, by the preceding $O(n^{-k})$ probability bound, the total
factorial-moment contribution from repeated rows or $\eta n$-close target
columns is
\[
        o(1)+O_k(M_n(\eta)/n).
\]
After $n\to\infty$ and then $\eta\downarrow0$, this contribution vanishes.

On the remaining tuples the rows are distinct and the target columns are
pairwise separated by more than $\eta n$.  Proposition
\ref{prop:bulk-local-final} gives, uniformly,
\[
\begin{aligned}
        &\Pp(\sigma_n(i_a)=j_a,1\le a\le k)\\
        &\qquad=n^{-k}\prod_{a=1}^k
        \rho(i_a/n,j_a/n)+o(n^{-k}).
\end{aligned}
\]
There are $O(n^k)$ tuples, so the accumulated error is $o(1)$.  Since $\rho$ is
bounded and continuous on $[0,1]\times[\delta,1-\delta]$, the removal of the
repeated-row and close-target tuples also changes the deterministic product sum
by $o(1)$ after $n\to\infty$ and $\eta\downarrow0$.  Hence
\[
\begin{aligned}
        \Ee(N(E_n))_k
        &=\left(\frac1n\sum_{(i,j)\in E_n}
        \rho(i/n,j/n)\right)^k+o(1)\\
        &\longrightarrow
        \left(\int\rho(x,y)\,\nu(\dd x,\dd y)\right)^k.
\end{aligned}
\]
These are the factorial moments of the stated Poisson law.  The method of
moments proves the theorem.
\end{proof}

\begin{corollary}[General thin-graph Poisson law]\label{cor:thin-graph-general}
Assume the hypotheses of Theorem \ref{thm:regular-cycles}.  Let
$q_{n,i}\in[\delta n,(1-\delta)n]$ be deterministic and suppose
\[
        \frac1n\sum_{i=1}^n
        \delta_{(i/n,q_{n,i}/n)}\Rightarrow\nu,
\]
where $\nu$ is a probability measure on $[0,1]\times[\delta,1-\delta]$.
If
\[
        \lim_{\eta\downarrow0}\limsup_{n\to\infty}
        \frac1n\sup_q
        \#\{i:|q_{n,i}-q|\le\eta n\}=0,
\]
then
\[
        \#\{i:\sigma_n(i)=q_{n,i}\}
        \Rightarrow
        \Pois\!\left(\int\rho(x,y)\,\nu(\dd x,\dd y)\right).
\]
If, in addition to the preceding no-crowding assumption,
$q_{n,i}/n\to g(i/n)$ uniformly for a continuous function $g$ with values in
$[\delta,1-\delta]$, then the limiting mean is
\[
        \int_0^1\rho(x,g(x))\dd x.
\]
\end{corollary}

\begin{proof}
Apply Theorem \ref{thm:sparse-entry} to
$E_n=\{(i,q_{n,i}):1\le i\le n\}$.  Every row degree is one, and the remaining
hypotheses are exactly those in the statement.  The final assertion follows
from convergence of the empirical graph measures to the pushforward of
Lebesgue measure under $x\mapsto(x,g(x))$.
\end{proof}

\begin{corollary}[Overlaps with deterministic permutations]\label{cor:overlap-permutation}
Assume the hypotheses of Theorem \ref{thm:regular-cycles}.  Let
$\tau_n\in S_n$ be deterministic and suppose its graph permutons
\[
        \mu_{\tau_n}=\frac1n\sum_{i=1}^n
        \delta_{(i/n,\tau_n(i)/n)}
\]
converge weakly to a permuton $\mu$.  Then
\[
        \#\{i:\sigma_n(i)=\tau_n(i)\}
        \Rightarrow\Pois(\Lambda_\mu),
\]
where
\[
        \Lambda_\mu=
        \lim_{\delta\downarrow0}
        \int_{[0,1]\times[\delta,1-\delta]}
        \rho(x,y)\,\mu(\dd x,\dd y)<\infty.
\]
\end{corollary}

\begin{proof}
Fix $\delta\in(0,1/2)$ and retain only the edges
$(i,\tau_n(i))$ whose target belongs to
$[\delta n,(1-\delta)n]$.  The row degrees and column degrees are one, and
\[
        \sup_q\#\{i:|\tau_n(i)-q|\le\eta n\}
        \le 2\eta n+2.
\]
Since the second marginal of every permuton is Lebesgue measure, the boundary
lines $y=\delta$ and $y=1-\delta$ have $\mu$-measure zero.  Theorem
\ref{thm:sparse-entry} therefore gives a Poisson limit with mean
\[
        \Lambda_\mu(\delta)=
        \int_{[0,1]\times[\delta,1-\delta]}
        \rho(x,y)\,\mu(\dd x,\dd y).
\]

The expected number of omitted overlaps is bounded by
\[
        \sum_{q\in E_{\delta,n}}
        \sup_i\Pp(\sigma_n(i)=q),
\]
which tends to zero as $\delta\downarrow0$, uniformly in $n$ along the limit,
by Corollary \ref{cor:endpoint-tightness} with $\ell=1$.  Hence the omitted
count tends to zero in probability after first taking $n\to\infty$ and then
$\delta\downarrow0$.  To see that the limiting mean is finite, fix $0<\delta'<\delta$.  Applying
Theorem \ref{thm:sparse-entry} to the graph edges with target in
$[\delta',\delta]\cup[1-\delta,1-\delta']$ identifies the corresponding
limiting integral with the limit of their expected overlap count.  The latter
is bounded by the endpoint estimate above and is therefore
$O(\delta+\delta^\beta)$.  Thus $\Lambda_\mu(\delta)$ is Cauchy as
$\delta\downarrow0$ and has a finite monotone limit.  The
converging-together lemma now yields the stated Poisson limit.
\end{proof}

\begin{theorem}[Transverse thin-curve Poisson law]\label{thm:thin-curve-transverse}
Assume the hypotheses of Theorem \ref{thm:regular-cycles}.  Let
$g\in C^1[0,1]$ satisfy
\[
        0<\delta_g\le g(x)\le1-\delta_g<1,
        \qquad
        \inf_{0\le x\le1}|g'(x)|>0.
\]
For all large $n$, put
\[
        q_{n,i}=\lfloor ng(i/n)\rfloor,
        \qquad
        N_g=\#\{i:\sigma_n(i)=q_{n,i}\}.
\]
Then
\[
        N_g\Rightarrow
        \Pois\left(\int_0^1\rho(x,g(x))\dd x\right).
\]
\end{theorem}

\begin{proof}
The derivative condition makes $g$ monotone and bi-Lipschitz on its image.
Hence
\[
        \sup_q\#\{i:|q_{n,i}-q|\le\eta n\}
        \le C_g(\eta n+1),
\]
while $q_{n,i}/n\to g(i/n)$ uniformly.  Corollary
\ref{cor:thin-graph-general} applies.
\end{proof}

\section{Macroscopic cycles under uniformly comparable weights}
\label{sec:macroscopic-cycles}

The preceding results concern cycles whose lengths remain fixed as $n\to\infty$.
We now consider cycles of order $n$.  Under a uniform bound on the ratio of the
largest to the smallest weight, the macroscopic cycle partition is universal.
We write $\operatorname{GEM}(1)$ for the residual-allocation sequence
\[
 U_1,\quad (1-U_1)U_2,\quad
 (1-U_1)(1-U_2)U_3,\ldots,
\]
where the $U_j$ are independent uniform random variables, and
$\operatorname{PD}(1)$ for its decreasing rearrangement; see
Pitman \cite[Chapter~3]{Pitman2006}.  Let
\[
 \mathcal S^\downarrow
 =\left\{x_1\ge x_2\ge\cdots\ge0:
 \sum_{j\ge1}x_j=1\right\},
\]
equipped with the $\ell^1$ metric.

\begin{theorem}[Macroscopic cycle universality]
\label{thm:macroscopic-pd}
Let $\sigma_n$ be a Luce permutation with positive weights
$\theta_{n,1},\ldots,\theta_{n,n}$ and assume
\[
 \sup_n\frac{\max_i\theta_{n,i}}{\min_i\theta_{n,i}}<\infty.
\]
If $L_{n,1}\ge L_{n,2}\ge\cdots$ are the cycle lengths of $\sigma_n$, then
\[
 \left(\frac{L_{n,1}}n,\frac{L_{n,2}}n,\ldots\right)
 \Rightarrow\operatorname{PD}(1)
\]
in $\mathcal S^\downarrow$.  Moreover, if
$\widetilde L_{n,1},\widetilde L_{n,2},\ldots$ are obtained by choosing a
uniform vertex, listing the length of its cycle, deleting that cycle, and
repeating on the residual vertex set, then
\[
 \left(\frac{\widetilde L_{n,1}}n,
 \frac{\widetilde L_{n,2}}n,\ldots\right)
 \Rightarrow\operatorname{GEM}(1)
\]
in $\ell^1$.  In particular, if $I_n$ is a uniformly chosen label and
$\mathscr L_n(I_n)$ is the length of its cycle, then
\[
        \frac{\mathscr L_n(I_n)}n\Rightarrow\operatorname{Unif}(0,1).
\]
\end{theorem}

The proof uses one fixed finite partition of clock time.  Conditional on the
clock slabs, each slab row is a permanental matching whose entries are nearly
constant.  A common minorization gives a spectral gap depending only on the
weight-comparability constant.  A relative-prefix theorem then controls every
target category under arbitrary adaptive exploration.  We formulate all
concentration statements uniformly over an explicit class of regular residual
histories, and prove that deleting an interior macroscopic cycle preserves that
class.  This supplies the regeneration needed for the $\operatorname{GEM}(1)$
induction.

Multiplying all weights by a common positive factor does not change the Luce
order.  We therefore normalize throughout the proof so that, for a fixed
$\kappa<\infty$,
\begin{equation}
        1\le\theta_{n,i}\le\kappa,
        \qquad 1\le i\le n.
        \tag{6.0}\label{eq:pd-normalized-rates}
\end{equation}

\subsection{Subsequential marked kernels and a uniform contraction}

Fix an arbitrary subsequence.  Passing to a further subsequence, compactness
gives
\[
 \eta_n:=\frac1n\sum_{i=1}^n\delta_{(i/n,\theta_{n,i})}
 \Longrightarrow\eta
\]
on $[0,1]\times[1,\kappa]$.  The first marginal of $\eta$ is Lebesgue
measure.  Choose a disintegration
\[
        \eta(\dd x,\dd\vartheta)=\dd x\,\eta_x(\dd\vartheta),
\]
defined for almost every $x$.  Put
\[
 A(t)=\int(1-e^{-\vartheta t})\,\eta(\dd x,\dd\vartheta),
 \qquad
 B(t)=A'(t)=\int\vartheta e^{-\vartheta t}\,
 \eta(\dd x,\dd\vartheta),
\]
let $\tau=A^{-1}$, and define, for almost every $x$,
\[
 g_x(t)=\int\vartheta e^{-\vartheta t}\,\eta_x(\dd\vartheta),
 \qquad
 \rho(x,y)=\frac{g_x(\tau(y))}{B(\tau(y))}.
\]
Let $K$ be the associated Markov operator on $L^2[0,1]$.

\begin{lemma}[Bistochasticity and common minorization]
\label{lem:pd-kernel-contraction}
The kernel $\rho$ is bistochastic up to null sets and
\[
        \iint_{[0,1]^2}\rho(x,y)^2\dd x\dd y\le\kappa.
\]
On the mean-zero subspace,
\[
        \|K\|_{L^2_0\to L^2_0}
        \le q_\kappa:=\sqrt{1-\kappa^{-1}}<1,
\]
and the same bound holds for $K^*$.
\end{lemma}

\begin{proof}
Using $y=A(t)$ and $\dd y=B(t)\dd t$,
\[
 \int_0^1\rho(x,y)\dd y=\int_0^\infty g_x(t)\dd t=1
\]
for almost every $x$, while
\[
 \int_0^1\rho(x,y)\dd x
 =\frac{\int_0^1g_x(\tau(y))\dd x}{B(\tau(y))}=1.
\]
Jensen's inequality and $\vartheta\le\kappa$ give
\[
 \int_0^1g_x(t)^2\dd x
 \le\int\vartheta^2e^{-2\vartheta t}\,
 \eta(\dd x,\dd\vartheta)
 \le\kappa B(2t).
\]
Since $\vartheta\ge1$, $B(2t)\le e^{-t}B(t)$, and hence
\[
 \iint\rho^2
 =\int_0^\infty\frac{\int g_x(t)^2\dd x}{B(t)}\dd t
 \le\kappa.
\]

For the contraction, regard $g_x(t)\dd t$ as the conditional law of a clock
time $T$ given $x$.  From \eqref{eq:pd-normalized-rates},
$g_x(t)\ge e^{-\kappa t}$.  Let
\[
        \alpha=\kappa^{-1},
        \qquad \nu(\dd t)=\kappa e^{-\kappa t}\dd t.
\]
There are probability kernels $R_x$ such that
\[
        g_x(t)\dd t=\alpha\nu(\dd t)+(1-\alpha)R_x(\dd t).
\]
The marginal law of $T$ is
$\mu=\alpha\nu+(1-\alpha)\bar R$, where
$\bar R=\int_0^1R_x\dd x$.  If $f\in L^2(\mu)$ and $\mu(f)=0$, then
\[
 \int f(t)g_x(t)\dd t
 =(1-\alpha)\{R_x(f)-\bar R(f)\}.
\]
Jensen's inequality therefore yields
\[
\begin{aligned}
 \int_0^1\left|\int f(t)g_x(t)\dd t\right|^2\dd x
 &\le(1-\alpha)^2\bar R(f^2)\\
 &\le(1-\alpha)\mu(f^2).
\end{aligned}
\]
The change of variables $y=A(t)$ is an isometry from $L^2(\dd y)$ to
$L^2(\mu)$, and under this isometry the displayed conditional-expectation
operator is $K$.  This proves the bound; the adjoint has the same norm.
\end{proof}

\subsection{Clock slabs and residual permanental matchings}

Fix $T>0$ and a partition
\[
        0=t_0<t_1<\cdots<t_d=T,
        \qquad \max_a(t_a-t_{a-1})\le h,
\]
and append the tail interval $I_\star=(T,\infty)$.  Put
\[
 I_a=(t_{a-1},t_a]\quad(a\le d),
 \qquad H_{n,a}=\{i:X_{n,i}\in I_a\},
 \qquad m_{n,a}=|H_{n,a}|.
\]
If $M_{n,a}=\sum_{c\le a}m_{n,c}$, then the ranks occupied by $H_{n,a}$ are
\[
        Q_{n,a}=\{M_{n,a-1}+1,\ldots,M_{n,a}\};
\]
the tail row is defined analogously.  Set
\[
 N_{n,ab}=|Q_{n,a}\cap H_{n,b}|,
 \qquad p_{n,a}=m_{n,a}/n.
\]
When $m_{n,a}>0$, put $P_{n,ab}=N_{n,ab}/m_{n,a}$; if
$m_{n,a}=0$, define that row of $P_n$ arbitrarily.  On the good events used
below every row is nonempty.  Thus $P_n$ is stochastic and $p_nP_n=p_n$.

Let
\[
 J_a=(A(t_{a-1}),A(t_a)]\quad(a\le d),
 \qquad J_\star=(A(T),1],
 \qquad p_a=|J_a|,
\]
and define
\[
 q_{ab}=\int_{J_a\times J_b}\rho(x,y)\dd x\dd y,
 \qquad P_{ab}=q_{ab}/p_a.
\]
Every $p_a$ and $q_{ab}$ is strictly positive.

\begin{lemma}[Slab law and deterministic good events]
\label{lem:pd-slab-law}
For every fixed partition,
\[
        p_{n,a}\xrightarrow{\Pp}p_a,
        \qquad N_{n,ab}/n\xrightarrow{\Pp}q_{ab}
\]
simultaneously in $a,b$.  There are deterministic $\varepsilon_n\downarrow0$
and events $\mathcal G_n$ with $\Pp(\mathcal G_n)\to1$ such that, on
$\mathcal G_n$:
\begin{enumerate}[label=(\roman*)]
\item all of the preceding convergences hold with error at most
$\varepsilon_n^2$, and, for all sufficiently large $n$,
\[
        p_{n,a}\ge\tfrac12p_a,
        \qquad N_{n,ab}\ge\tfrac12q_{ab}n
\]
for every $a,b$;
\item for a constant $q_0\in(q_\kappa,1)$,
\[
        \|P_n\|_{L^2_0(p_n)\to L^2_0(p_n)}\le q_0;
\]
\item if $S$ is a partition endpoint and $J_b\subset[0,A(S)]$, then
\[
        P_{n,ab}\le 2M_Sp_{n,b}\qquad\text{for every }a,
        \qquad M_S=\kappa e^{(\kappa-1)S};
        \tag{6.1}
\]
\item for every union $D$ of slab states,
\[
 \sum_{a\in D}\sum_b p_{n,a}\frac{P_{n,ab}^2}{p_{n,b}}
 \le2\kappa;
        \tag{6.2}
\]
\item if $U$ is a partition endpoint and $D(U)$ is the union of rows whose
clock intervals lie in $(U,\infty)$, then
\[
        \sum_{a\in D(U)}p_{n,a}\le2e^{-U}+\varepsilon_n^2.
        \tag{6.3}
\]
\end{enumerate}
The constants may depend on the fixed partition and on $S$, but not on $n$.
\end{lemma}

\begin{proof}
For a fixed clock cell $I_b=(u,v]$, the variables
$\one_{\{X_{n,i}\in I_b\}}$ are independent and have means
$e^{-\theta_{n,i}u}-e^{-\theta_{n,i}v}$.  Kolmogorov's maximal inequality
gives
\[
 \max_{k\le n}\left|\sum_{i\le k}
 \left(\one_{\{X_{n,i}\in I_b\}}
 -\Ee\one_{\{X_{n,i}\in I_b\}}\right)\right|=o_{\Pp}(n).
\]
Marked empirical convergence and the Lebesgue first marginal imply, uniformly
in $x\in[0,1]$,
\[
 \frac1n\sum_{i\le nx}
 \Ee\one_{\{X_{n,i}\in I_b\}}
 \longrightarrow
 \int_0^x\int
 (e^{-\vartheta u}-e^{-\vartheta v})
 \eta_z(\dd\vartheta)\dd z.
\]
This proves the slab-size convergence and convergence of the endpoints of
$Q_{n,a}/n$ to those of $J_a$.  Applying the same uniform partial-sum estimate
at those random endpoints gives
\[
 \frac{N_{n,ab}}n\longrightarrow
 \int_{J_a}\int
 (e^{-\vartheta t_{b-1}}-e^{-\vartheta t_b})
 \eta_x(\dd\vartheta)\dd x=q_{ab}.
\]
The last equality follows by integrating $\rho(x,y)$ over $J_b$ and using
$y=A(t)$.

The matrix $P$ is the compression of $K$ to functions that are constant on the
sets $J_a$: explicitly, it is $\Pi_{\mathcal J}K\Pi_{\mathcal J}$.  Conditional
expectation is a contraction, so Lemma \ref{lem:pd-kernel-contraction} gives
$\|P\|_{L^2_0(p)}\le q_\kappa$.  Entrywise convergence of the conjugated
matrices
$D_{p_n}^{1/2}P_nD_{p_n}^{-1/2}$ proves (ii).

If $y\le A(S)$, then $\tau(y)\le S$ and
\[
 \rho(x,y)
 =\frac{g_x(\tau(y))}{B(\tau(y))}
 \le\frac{\kappa e^{-\tau(y)}}{e^{-\kappa\tau(y)}}
 \le\kappa e^{(\kappa-1)S}=:M_S.
\]
Hence $P_{ab}\le M_Sp_b$, which yields (6.1) by entrywise convergence.  On
every rectangle, Cauchy's inequality gives
\[
 p_a\frac{P_{ab}^2}{p_b}
 =\frac{(\int_{J_a\times J_b}\rho)^2}{p_ap_b}
 \le\int_{J_a\times J_b}\rho^2.
\]
Summing and using Lemma \ref{lem:pd-kernel-contraction} proves the limiting
version of (6.2), and entrywise convergence gives its finite-$n$ form.  Finally,
the expected proportion of clocks exceeding $U$ is at most $e^{-U}$ because
all rates are at least one; the slab law gives (6.3).  Since the partition is finite, choose deterministic $\delta_n\downarrow0$
so that all convergence assertions hold with error at most $\delta_n$ on a
single event of probability tending to one, and set
$\varepsilon_n=\sqrt{\delta_n}$.  This gives the stated event
$\mathcal G_n$ and the errors $\varepsilon_n^2$ in (i) and (6.3).
\end{proof}

Condition on the complete slab data: the sets $H_{n,a}$, the rank intervals
$Q_{n,a}$, and the ordered clock values
\[
        s_{a,1}<\cdots<s_{a,m_{n,a}}
\]
in every slab.  Write
$Q_{n,a}=\{q_{a,1}<\cdots<q_{a,m_{n,a}}\}$, and let
$\phi_a:H_{n,a}\to[m_{n,a}]$ be defined by
$\sigma_n(i)=q_{a,\phi_a(i)}$.  Then
\begin{equation}
 \Pp(\phi_a=\phi\mid\text{slab data})
 =\frac{\prod_{i\in H_{n,a}}e^{-\theta_{n,i}s_{a,\phi(i)}}}
 {\operatorname{per}(e^{-\theta_{n,i}s_{a,j}})_{i,j}}.
 \tag{6.3a}\label{eq:pd-permanental-row}
\end{equation}
The row matchings are conditionally independent.  The factors
$\prod_i\theta_{n,i}$ cancel from the conditional density.

\begin{lemma}[Residual permanental factorization]
\label{lem:pd-residual-factorization}
Condition on the slab data and on any compatible partial collection of exposed
matching edges.  In every row, delete the exposed source and target vertices.
The remaining row matchings are conditionally independent, and each has the
permanental law \eqref{eq:pd-permanental-row} on its remaining submatrix.
\end{lemma}

\begin{proof}
The conditional density of the complete matching is a product over rows and,
within each row, a product over its edges.  Conditioning on fixed exposed edges
removes their constant factors and restricts the sum to bijections of the
remaining rows and columns.  The normalizing sum is the permanent of the
remaining submatrix, and the factorization over rows is unchanged.
\end{proof}

\begin{lemma}[Hereditary flatness]
\label{lem:pd-hereditary-flatness}
Let $A=(a_{ij})$ be a positive $m\times m$ matrix satisfying
\[
        e^{-\epsilon}\le a_{ij}/a_{ik}\le e^\epsilon
\]
for every row $i$ and every pair of columns $j,k$.  Under the permanental
matching law proportional to $\prod_i a_{i,\phi(i)}$, after conditioning on an
arbitrary partial matching, every remaining edge probability satisfies
\[
 \frac{e^{-2\epsilon}}s
 \le\Pp(\phi(i)=j\mid\text{partial matching})
 \le\frac{e^{2\epsilon}}s,
\]
where $s$ rows and columns remain.  For a Luce slab of width at most $h$, this
holds with $\epsilon=\kappa h$; write
\[
        \beta=2\kappa h.
        \tag{6.4}
\]
\end{lemma}

\begin{proof}
It suffices to work with the current remaining submatrix.  For two remaining
columns $j,k$,
\[
 \frac{\Pp(\phi(i)=j)}{\Pp(\phi(i)=k)}
 =\frac{a_{ij}\operatorname{per}A_{-i,-j}}
 {a_{ik}\operatorname{per}A_{-i,-k}}.
\]
Map a term of the first minor to a term of the second by locating the row sent
to $k$ and sending that row to $j$ instead.  Corresponding products differ by
a factor in $[e^{-\epsilon},e^\epsilon]$.  Including the factor
$a_{ij}/a_{ik}$ shows that every pair of edge probabilities has ratio in
$[e^{-2\epsilon},e^{2\epsilon}]$.  Since the $s$ probabilities sum to one,
the individual bounds follow.  In a clock slab,
\[
        e^{-\kappa h}\le
        e^{-\theta_i(s_j-s_k)}\le e^{\kappa h}.
\]
\end{proof}
\begin{lemma}[Predictable row exposure]
\label{lem:pd-predictable-exposure}
Condition on the slab data and on a compatible partial collection of exposed
matching edges.  Let $(\mathscr F_r)$ be an exploration filtration such that,
before the $(r+1)$st edge is exposed, its source is an unexposed source that is
$\mathscr F_r$-measurable.  Then, conditional on $\mathscr F_r$, the target of
that source has the one-edge marginal of the residual permanental matching in
its row.  In particular, if that row has $s$ targets remaining and has width at
most $h$, every individual target has conditional probability in
\[
        [e^{-\beta}/s,e^\beta/s],
        \qquad \beta=2\kappa h.
\]
If a distinguished target $v$ in that row is forbidden and one conditions the
next target not to equal $v$, then every one of the other $s-1$ targets has
conditional probability in
\[
        [e^{-\beta}/(s-1),e^\beta/(s-1)].
\]
The conclusions remain true when the next source is selected adaptively using
all previously exposed edges, including edges from the same row.
\end{lemma}

\begin{proof}
By Lemma \ref{lem:pd-residual-factorization}, conditional on
$\mathscr F_r$ the unexposed part of the relevant row has the permanental law
on its current submatrix.  Lemma \ref{lem:pd-hereditary-flatness} therefore
gives the first interval of probabilities.  If $v$ is forbidden, conditioning
away from $v$ merely renormalizes the probabilities of the other targets.
Their pairwise ratios are unchanged and lie in $[e^{-\beta},e^\beta]$; since
their conditional probabilities sum to one, each lies in the second displayed
interval.  The measurability assumption is precisely what is needed to apply
the residual-factorization statement before the new edge is revealed.  No
independence between the chosen source and the previously exposed edges is
required.
\end{proof}


\subsection{A quantitative relative-prefix theorem}

\begin{lemma}[Relative prefix law]
\label{lem:pd-relative-prefix}
Fix $d\ge1$ and $c_*>0$.  There are absolute constants
$\beta_0>0$ and $C<\infty$, and deterministic sequences
$\chi_m\downarrow0$ and $\pi_m\downarrow0$, depending only on $d$ and $c_*$,
with the following property.  Let $m$ targets be
partitioned into $d$ categories of sizes $m_1,\ldots,m_d$ satisfying
$m_b\ge c_*m$.  Targets are selected sequentially without replacement in an
arbitrary predictable, adaptive source order.  Suppose that, whenever $s$
targets remain, every individual remaining target has conditional selection
probability in
\[
        [e^{-\beta}/s,e^\beta/s],
        \qquad 0\le\beta\le\beta_0.
\]
If $X_b(r)$ is the number of category-$b$ targets among the first $r$
selections, then
\[
 \Pp\left(
 \max_{b\le d}\max_{0\le r\le m}
 \frac{|X_b(r)-rm_b/m|}{m_b}>C\beta+\chi_m
 \right)\le\pi_m.
 \tag{6.5}\label{eq:pd-relative-prefix}
\]
The estimate is uniform over all adaptive source rules satisfying the displayed
one-target bounds.  It remains valid after deleting any fixed number of
distinguished targets before the procedure, or after forbidding those targets
throughout the procedure, after changing $c_*$ and the two deterministic
sequences harmlessly.
\end{lemma}

\begin{proof}
Fix one category and write
\[
        N_r=m_b-X_b(r),\qquad s_r=m-r.
\]
Let $Y_{r+1}$ be the indicator that the next selected target belongs to this
category.  By summing the individual-target bounds over the $N_r$ remaining
targets in the category,
\[
 \pi_r:=\Ee(Y_{r+1}\mid\mathscr F_r)
 \in\left[e^{-\beta}\frac{N_r}{s_r},
          e^\beta\frac{N_r}{s_r}\right].
 \tag{6.6}\label{eq:pd-prefix-pi}
\]
Set $a_m=m_b^{2/3}$ and stop at
\[
        \tau=\inf\{r:N_r<a_m\text{ or }s_r<a_m\}.
\]
For $r<\tau$, put $Z_r=\log(N_r/s_r)$.  Since
\[
 Z_{r+1}-Z_r
 =-\log(1-s_r^{-1})
  +Y_{r+1}\log(1-N_r^{-1}),
\]
Taylor's formula and \eqref{eq:pd-prefix-pi} give
\[
 \left|\Ee(Z_{r+1}-Z_r\mid\mathscr F_r)\right|
 \le C_1\frac\beta{s_r}
 +C_1\left(\frac1{s_r^2}+\frac1{s_rN_r}\right),
 \tag{6.7}\label{eq:pd-prefix-drift}
\]
and the conditional variance of the centered increment is at most
$C_1/(s_rN_r)$.  Because $m_b\ge c_*m$,
\[
 \sum_{r<\tau}\frac1{s_rN_r}
 \le a_m^{-1}\sum_{s\ge a_m}\frac1s
 \le C_{c_*}\frac{1+\log m}{m^{2/3}}=:v_m,
 \tag{6.8}\label{eq:pd-prefix-qv}
\]
and the sum of the two second-order drift terms in
\eqref{eq:pd-prefix-drift} is $O(v_m)$.

Let $M_r$ be the stopped martingale obtained by summing the centered
increments.  Doob's $L^2$ inequality and \eqref{eq:pd-prefix-qv} imply
\[
 \Pp\left(\max_{r\le\tau}|M_r|>v_m^{1/4}\right)
 \le C v_m^{1/2}.
 \tag{6.9}\label{eq:pd-prefix-doob}
\]
On the complementary event, for every $r<\tau$,
\[
 \left|\log\frac{N_r/m_b}{s_r/m}\right|
 \le C_2\beta\log\frac m{s_r}+C_2(v_m^{1/4}+v_m).
 \tag{6.10}\label{eq:pd-prefix-log}
\]
Choose $\beta_0$ so that $C_2\beta_0\le1/2$.  Writing $u=s_r/m$ and
$\eta_m=C_2(v_m^{1/4}+v_m)$, \eqref{eq:pd-prefix-log} yields
\[
 e^{-\eta_m}u^{1+C_2\beta}
 \le\frac{N_r}{m_b}
 \le e^{\eta_m}u^{1-C_2\beta}.
 \tag{6.11}\label{eq:pd-prefix-power}
\]
For $|a|\le1/2$,
\[
        \sup_{0\le u\le1}|u^{1+a}-u|\le C|a|,
\]
which follows from the mean-value theorem in the exponent and
$\sup_{0<u\le1}u|\log u|<\infty$.  Hence, before $\tau$,
\[
        \left|\frac{N_r}{m_b}-\frac{s_r}{m}\right|
        \le C_3\beta+C_3\eta_m.
 \tag{6.12}\label{eq:pd-prefix-before-stop}
\]

At the stopping time, if $s_\tau<a_m$, then
$s_\tau/m=O(m^{-1/3})$ and the upper bound in
\eqref{eq:pd-prefix-power}, applied one step earlier, gives
$N_\tau/m_b=o(1)$.  If $N_\tau<a_m$, then the lower bound in
\eqref{eq:pd-prefix-power} gives $s_\tau/m=o(1)$.  In either case, after
$\tau$ both normalized remaining counts are nonincreasing and their difference
is bounded by $C_3\beta+o(1)$.  Thus
\eqref{eq:pd-prefix-before-stop} extends to all $0\le r\le m$.

For example, one may take
\[
 \chi_m=C_{d,c_*}\left(v_m^{1/4}+m^{-1/12}\right),
 \qquad
 \pi_m=C_{d,c_*}v_m^{1/2}.
\]
A union bound over the fixed number of categories proves
\eqref{eq:pd-relative-prefix}.  Deleting or forbidding finitely many targets
changes every relevant category size and the total size by $O(1)$; for all
large $m$ the positive-proportion hypothesis and the same one-target ratio
bounds remain valid after changing $c_*$ by a fixed factor.
\end{proof}

In every application below the mesh is chosen sufficiently fine that
$\beta=2\kappa h\le\beta_0$.

\subsection{Regular residual histories and scale-sensitive orbit balance}

A residual history $\mathcal H$ consists of the complete slab data, all roots
selected so far, and the matching edges in a collection of fully exposed and
deleted cycles; it contains no additional information about unexposed row
matchings.  After such a history, let $R\subset[n]$ be the remaining vertex set
and $N=|R|$.  Put
\[
 m_a^R=|H_{n,a}\cap R|,
 \qquad N_{ab}^R=|Q_{n,a}\cap H_{n,b}\cap R|,
 \qquad r=N/n.
\]
For every deleted cycle $C$ and every slab $a$,
\begin{equation}
        |C\cap H_{n,a}|=|C\cap Q_{n,a}|,
        \label{eq:pd-cycle-block-equality}
\end{equation}
because $\sigma_n$ maps $C\cap H_{n,a}$ bijectively onto
$C\cap Q_{n,a}$.  Consequently $m_a^R=|Q_{n,a}\cap R|$.

Let $U$ be a master-partition endpoint.  A residual history is called
$(U,\delta,\xi)$-regular if $N\ge\xi n$, $0\le\delta<1/4$, and, for every row
whose clock interval is contained in $[0,U]$ and every target category $b$,
\begin{equation}
 \left|\frac{m_a^R}{r m_{n,a}}-1\right|\le\delta,
 \qquad
 \left|\frac{N_{ab}^R}{r N_{n,ab}}-1\right|\le\delta.
 \label{eq:pd-regular-history}
\end{equation}
The initial history is $(U,0,1)$-regular.  By Lemma
\ref{lem:pd-residual-factorization}, conditional on any residual history, the
unexposed row matchings are independent permanental matchings on their
remaining submatrices.

An \emph{admissible exploration} starts from a residual history and reveals one
outgoing matching edge at a time.  Before an edge is revealed, its source must
be measurable with respect to the filtration generated by the slab data, the
residual history, and all previously exposed edges.  We also permit a fixed
number of distinguished targets to be forbidden, with the next edge in a row
containing such a target sampled conditionally on avoiding it.  Lemma
\ref{lem:pd-predictable-exposure} shows that every actual cycle exploration and
every root-avoiding exploration used below is admissible.

Consider an injective orbit segment $v_0,v_1,\ldots,v_\ell$, where
$v_{j+1}=\sigma_n(v_j)$ for the actual exploration, and allow
$v_\ell=v_0$ only at a terminal closing step.  For an auxiliary exploration,
$v_{j+1}$ denotes the target selected under its conditional avoid-target law.
Put
\[
 R_a(\ell)=\#\{0\le j<\ell:v_j\in H_{n,a}\},
\]
and
\[
 R_{ab}(\ell)=\#\{0\le j<\ell:v_j\in H_{n,a},\ v_{j+1}\in H_{n,b}\}.
\]
Whenever no ambiguity is possible we suppress the argument $\ell$.  The
deterministic flow identity is
\begin{equation}
 R_b-\sum_aR_{ab}
 =\one_{\{v_0\in H_{n,b}\}}-
  \one_{\{v_\ell\in H_{n,b}\}}.
 \label{eq:pd-flow}
\end{equation}

\begin{lemma}[Scale-sensitive two-scale orbit balance]
\label{lem:pd-two-scale-balance}
Fix master-partition endpoints $0<S<U\le T$, $\xi>0$, and
$0\le\delta<1/4$.  There are deterministic sequences
$\chi_n\downarrow0$ and $\pi_n\downarrow0$, depending on the fixed partition,
$U$, and $\xi$ but not on the residual history, with the following property.
On $\mathcal G_n$, uniformly over every $(U,\delta,\xi)$-regular history and
every admissible exploration, with conditional probability at least
$1-\pi_n$ the following assertions hold simultaneously at every defined prefix
length $\ell\ge1$.

For every row $a$ whose clock interval is contained in $[0,U]$ and every target
category $b$,
\begin{equation}
 \left|R_{ab}(\ell)-\frac{R_a(\ell)}{m_a^R}N_{ab}^R\right|
 \le (C\beta+\chi_n)N_{ab}^R,
 \label{eq:pd-row-prefix}
\end{equation}
where $C$ is absolute and $\beta=2\kappa h$.  Define
\[
        \Delta_n(U,\delta,h)
        =\delta+\beta+\chi_n+\varepsilon_n+e^{-U/2}.
\]
Then, for every inner state $b$ with $J_b\subset[0,A(S)]$,
\begin{equation}
 \left|\frac{R_b(\ell)}{\ell p_{n,b}}-1\right|
 \le C_{\kappa,S}\frac n\ell\,
        \Delta_n(U,\delta,h)
      +\frac{C_{\mathcal J}}\ell.
 \label{eq:pd-scale-sensitive-balance}
\end{equation}
Here $C_{\kappa,S}$ is independent of the mesh and of $\ell$; the finite
constant $C_{\mathcal J}$ may depend on the fixed master partition.  In
particular, for every fixed $\lambda>0$, all prefixes with $\ell\ge\lambda n$
satisfy
\[
 \max_{b:J_b\subset[0,A(S)]}
 \left|\frac{R_b(\ell)}{\ell p_{n,b}}-1\right|
 \le C_{\kappa,S,\lambda}\Delta_n(U,\delta,h)+o(1).
\]
\end{lemma}

\begin{proof}
On $\mathcal G_n$, every $m_{n,a}$ and every $N_{n,ab}$ is bounded below by a
positive constant times $n$.  If the history is regular, the same is true for
$m_a^R$ and $N_{ab}^R$, uniformly over the stated histories.  Consequently,
for each row below $U$, the residual target categories have proportions bounded
below by a positive constant depending only on the fixed partition.  Apply
Lemmas \ref{lem:pd-predictable-exposure} and
\ref{lem:pd-relative-prefix} to the adaptive sequence of queried sources in
that row, extending it predictably after the orbit segment ends until all
ordinary targets in that row have been selected.  In a row with a forbidden
target, this means a complete order of the allowable targets.  A union bound over the finitely many rows and categories gives
deterministic sequences $\chi_n,\pi_n\downarrow0$ for which
\eqref{eq:pd-row-prefix} holds simultaneously for every prefix.  If finitely
many distinguished targets are forbidden, remove them from the corresponding
category counts; the resulting $O(1)$ changes are absorbed into $\chi_n$.

Write $p=p_n$, $P=P_n$, $q_a=R_a/\ell$, and
\[
        E_{ab}=R_{ab}-R_aP_{ab}.
\]
For a row below $U$, regularity gives
\[
 \frac{N_{ab}^R}{m_a^R}=P_{ab}\{1+O(\delta)\},
\]
uniformly, and $R_a\le m_a^R$.  Hence
\begin{equation}
        |E_{ab}|
        \le C(\delta+\beta+\chi_n)N_{n,ab}
        \qquad(a\notin D(U)),
 \label{eq:pd-good-row-error}
\end{equation}
where $D(U)$ is the union of rows whose clock intervals lie in $(U,\infty)$.
For an outer row we use only injectivity of the explored path:
$R_{ab}\le N_{ab}^R\le N_{n,ab}$ and
$R_aP_{ab}\le m_a^RP_{ab}\le N_{n,ab}$.  Thus
\begin{equation}
        |E_{ab}|\le2N_{n,ab}
        \qquad(a\in D(U)).
 \label{eq:pd-bad-row-error}
\end{equation}

Let
\[
        G_b=\sum_{a\notin D(U)}E_{ab},
        \qquad H_b=\sum_{a\in D(U)}E_{ab}.
\]
Since $\sum_aN_{n,ab}=m_{n,b}=np_b$, equation
\eqref{eq:pd-good-row-error} gives
\begin{equation}
 \left\|\frac{G}{\ell p}\right\|_{L^2(p)}
 \le C\frac n\ell(\delta+\beta+\chi_n).
 \label{eq:pd-good-L2}
\end{equation}
For the outer rows, Cauchy's inequality and the Hilbert--Schmidt block estimate
on $\mathcal G_n$ give
\begin{align}
 \sum_b\frac{H_b^2}{m_{n,b}}
 &\le4\sum_b\frac{(\sum_{a\in D(U)}N_{n,ab})^2}{m_{n,b}}\notag\\
 &\le4m_{n,D(U)}
 \sum_{a\in D(U)}\sum_b
 \frac{m_{n,a}P_{n,ab}^2}{m_{n,b}}
 \le C m_{n,D(U)}.
 \label{eq:pd-bad-L2-raw}
\end{align}
Therefore, using the tail-mass estimate on $\mathcal G_n$,
\begin{equation}
 \left\|\frac{H}{\ell p}\right\|_{L^2(p)}
 \le C\frac n\ell
       \left(e^{-U/2}+\varepsilon_n\right).
 \label{eq:pd-bad-L2}
\end{equation}

Let $d_b$ denote the right-hand side of \eqref{eq:pd-flow}.  It has at most
two nonzero coordinates, each of absolute value at most one.  Dividing
\eqref{eq:pd-flow} by $\ell$ gives
\[
        q-qP=e,
        \qquad e_b=\frac{d_b+G_b+H_b}{\ell}.
\]
Since the partition is fixed,
\begin{equation}
 \left\|\frac e p\right\|_{L^2(p)}
 \le C\frac n\ell\Delta_n(U,\delta,h)
      +\frac{C_{\mathcal J}}\ell.
 \label{eq:pd-g-L2-scale}
\end{equation}
Put $u_a=q_a/p_a-1$.  Because $pP=p$ and $\sum_ap_au_a=0$,
\[
        u-P^*u=e/p.
\]
The spectral-gap event in $\mathcal G_n$ yields
\begin{equation}
 \|u\|_{L^2(p)}
 \le(1-q_0)^{-1}\left\|\frac e p\right\|_{L^2(p)}.
 \label{eq:pd-u-L2-scale}
\end{equation}

It remains to recover a coordinatewise estimate for an inner state $b$.
The inner-column bound and stationarity imply
\[
 \sum_a p_a\left(\frac{P_{ab}}{p_b}\right)^2
 \le \left(\max_a\frac{P_{ab}}{p_b}\right)
      \sum_a\frac{p_aP_{ab}}{p_b}
 \le C_{\kappa,S}.
\]
Hence
\[
        |(P^*u)_b|\le C_{\kappa,S}\|u\|_{L^2(p)}.
\]
For the coordinate $(e/p)_b$, equations
\eqref{eq:pd-good-row-error} and \eqref{eq:pd-bad-row-error}, together with
$P_{ab}\le C_{\kappa,S}p_b$ for an inner column, give
\[
 \left|\frac{e_b}{p_b}\right|
 \le C_{\kappa,S}\frac n\ell
      \left(\delta+\beta+\chi_n+e^{-U}+\varepsilon_n^2\right)
      +\frac{C_{\mathcal J}}\ell.
\]
Combining this estimate with
$u_b=(e/p)_b+(P^*u)_b$, \eqref{eq:pd-g-L2-scale}, and
\eqref{eq:pd-u-L2-scale} proves
\eqref{eq:pd-scale-sensitive-balance}.  Every probabilistic input was the
conditional relative-prefix event, so the failure probability $\pi_n$ is
uniform over the entire stated class of histories and explorations.
\end{proof}

\subsection{The cycle of a uniform residual root}

For a vertex $v$, let $c(v)$ be the unique row index for which
$v\in Q_{n,c(v)}$.  An orbit can hit $v$ only when its current source lies in
$H_{n,c(v)}$.

\begin{lemma}[Root-avoiding change of measure]
\label{lem:pd-root-avoiding}
Fix a residual history and a distinguished remaining target $v$.  Explore the
path from $v$.  At the $j$th opportunity to hit $v$, let $\lambda_j$ be the
conditional probability, under the original residual matching law and given
that $v$ has not yet been selected, that the next target is $v$.  Let
$\Pp^\circ$ be the auxiliary law obtained by conditioning each such choice not
to select $v$.  If $S(k)$ is the number of closure opportunities among the
first $k$ completed edges, then
\begin{equation}
        \Pp(|\mathcal C(v)|>k\mid\mathcal H,v)
        =\Ee^\circ\prod_{j=0}^{S(k)-1}(1-\lambda_j),
 \label{eq:pd-root-rn}
\end{equation}
where the product is interpreted as zero if the auxiliary exploration is
forced to stop before completing $k$ edges.  If the closure row initially has
$m$ available targets, including $v$, then at its $j$th closure opportunity,
\begin{equation}
        \frac{e^{-\beta}}{m-j}\le\lambda_j\le
        \frac{e^\beta}{m-j},
        \qquad \beta=2\kappa h.
 \label{eq:pd-root-hazard}
\end{equation}
\end{lemma}

\begin{proof}
At a closure opportunity, the original transition law is the mixture of an
atom of mass $\lambda_j$ at $v$ and the conditional avoid-$v$ law of mass
$1-\lambda_j$.  A non-root previously visited vertex cannot be selected,
because its target column has already been used; thus the only possible first
repeat is the root.  Multiplying the sequential Radon--Nikodym factors gives
\eqref{eq:pd-root-rn}.  Before the $j$th opportunity, exactly $j$ non-$v$
targets have been removed from the closure row, so $m-j$ targets remain.
Equation \eqref{eq:pd-root-hazard} follows from hereditary flatness.  If only
$v$ remains, the next closure probability is one, and the survival product is
zero, which agrees with the stated convention.
\end{proof}

\begin{lemma}[Harmonic-product approximation]
\label{lem:pd-harmonic-product}
Fix $\zeta\in(0,1)$.  Uniformly for $c\in[1/2,2]$, integers $m\ge1$, and
$0\le s\le(1-\zeta)m$ for which all factors are nonnegative,
\[
 \prod_{j=0}^{s-1}\left(1-\frac c{m-j}\right)
 =(1-s/m)^c\left\{1+O_\zeta(m^{-1})\right\}.
\]
Moreover, uniformly for $x\in[\zeta,1-\zeta]$ and $|a|\le1/2$,
\[
        |(1-x)^{1+a}-(1-x)|\le C_\zeta|a|.
\]
\end{lemma}

\begin{proof}
Since $m-j\ge\zeta m$,
\[
 \sum_{j<s}\log\left(1-\frac c{m-j}\right)
 =-c\sum_{j<s}\frac1{m-j}+O_\zeta(m^{-1}).
\]
The harmonic sum equals
$\log(m/(m-s))+O_\zeta(m^{-1})$, which proves the first assertion after
exponentiation.  The second follows from the mean-value theorem in the exponent
and the boundedness of $|(1-x)^b\log(1-x)|$ for
$x\in[\zeta,1-\zeta]$ and $b\in[1/2,3/2]$.
\end{proof}

\begin{proposition}[Uniform conditional residual fraction]
\label{prop:pd-uniform-residual-fraction}
Fix master-partition endpoints $0<S<U\le T$, $\xi>0$,
$\zeta\in(0,1/2)$, and $0\le\delta<1/4$.  On $\mathcal G_n$, let
$\mathcal H$ be a $(U,\delta,\xi)$-regular history.  Choose a root uniformly
from its $N$ remaining vertices, independently of the unexposed matchings, and
let $L$ be the root-cycle length.  With $\chi_n,\pi_n$ as in Lemma
\ref{lem:pd-two-scale-balance}, uniformly over all such histories,
\begin{align}
 &\sup_{x\in[\zeta,1-\zeta]}
 \left|\Pp\left(\left.\frac LN\le x\right|\mathcal H\right)-x\right|
 \notag\\
 &\qquad\le
 C_{\kappa,S,\xi,\zeta}
 \bigl(\delta+\beta+\chi_n+\varepsilon_n+e^{-U/2}\bigr)
 +C_\xi(e^{-S}+\varepsilon_n^2)
 +\pi_n+\frac{C_{\mathcal J,\xi,\zeta}}n.
 \label{eq:pd-root-cdf}
\end{align}
The constants multiplying the principal error are independent of the mesh;
$C_{\mathcal J,\xi,\zeta}$ may depend on the fixed master partition.
\end{proposition}

\begin{proof}
Let $v$ be the uniform root and put $C=c(v)$.  Complete-cycle deletion gives
$|Q_{n,a}\cap R|=m_a^R$, so, conditional on the history,
\[
        \Pp(C=a\mid\mathcal H)=m_a^R/N.
\]
Consequently, using $N\ge\xi n$ and the tail bound on $\mathcal G_n$,
\begin{equation}
 \Pp(C\in D(S)\mid\mathcal H)
 \le C_\xi(e^{-S}+\varepsilon_n^2).
 \label{eq:pd-closure-tail}
\end{equation}
We henceforth work on the event $C\notin D(S)$.  Put
$m=m_C^R$ and $p_C=p_{n,C}$.  Regularity gives
\begin{equation}
        \frac{m}{Np_C}=1+O(\delta).
 \label{eq:pd-closure-size}
\end{equation}
For the fixed master partition, $p_C$ is bounded below over the finitely many
rows below $S$; thus $m\ge c_{\mathcal J,S}\xi n$ for all large $n$.

We first rule out premature termination of the root-avoiding exploration.  Let
$r$ be the number of completed edges if that exploration stops because only
$v$ remains in the closure row.  Then exactly $m-1$ closure-row targets have
been used, so
\[
        R_C(r)=m-1,
        \qquad r\ge m-1.
\]
On the simultaneous prefix event of Lemma
\ref{lem:pd-two-scale-balance}, its scale-sensitive estimate at this terminal
prefix gives
\[
 \left|\frac{m-1}{r p_C}-1\right|
 \le C_{\kappa,S}\frac nr\Delta_n(U,\delta,h)
      +\frac{C_{\mathcal J}}r.
\]
Writing $\theta=r/N$ and multiplying by $\theta$ yields
\[
 \left|\frac{m-1}{Np_C}-\theta\right|
 \le C_{\kappa,S,\xi}\Delta_n(U,\delta,h)
      +\frac{C_{\mathcal J}}N.
\]
Together with \eqref{eq:pd-closure-size},
\begin{equation}
        |1-r/N|
        \le C_{\kappa,S,\xi}
        \Delta_n(U,\delta,h)+\frac{C_{\mathcal J,S,\xi}}n.
 \label{eq:pd-no-premature}
\end{equation}
If the right-hand side is smaller than $\zeta/2$, the auxiliary exploration
cannot stop before $(1-\zeta)N$ completed edges.  If it is not smaller than
$\zeta/2$, the right-hand side of \eqref{eq:pd-root-cdf}, after increasing its
constant, is at least one and the proposition is trivial.  Thus we may assume
that the auxiliary exploration is defined through every
$k\le(1-\zeta)N$.

Fix $x\in[\zeta,1-\zeta]$ and put $k=\lfloor xN\rfloor$.  The number $S(k)$
of closure opportunities is exactly $R_C(k)$ under the avoid-$v$ law.  Applying
\eqref{eq:pd-scale-sensitive-balance} at $\ell=k$, and using
$k\ge\zeta\xi n$, gives, on the simultaneous prefix event,
\[
 \left|\frac{S(k)}{k p_C}-1\right|
 \le C_{\kappa,S,\xi,\zeta}
       \Delta_n(U,\delta,h)+\frac{C_{\mathcal J,\xi,\zeta}}n.
\]
Combining this with \eqref{eq:pd-closure-size} and
$k/N=x+O(n^{-1})$ gives
\begin{equation}
 \sup_{x\in[\zeta,1-\zeta]}
 \left|\frac{S(\lfloor xN\rfloor)}m-x\right|
 \le C_{\kappa,S,\xi,\zeta}
       \Delta_n(U,\delta,h)+\frac{C_{\mathcal J,\xi,\zeta}}n.
 \label{eq:pd-closure-opportunities}
\end{equation}
The exceptional conditional probability is at most $\pi_n$.

On the event in \eqref{eq:pd-closure-opportunities}, the ratio
$S(k)/m$ lies in $[\zeta/2,1-\zeta/2]$ once the displayed error is small; in the
complementary parameter regime the final estimate is again trivial.  Equations
\eqref{eq:pd-root-rn} and \eqref{eq:pd-root-hazard} squeeze the survival
product between
\[
 \prod_{j=0}^{S(k)-1}\left(1-\frac{e^\beta}{m-j}\right)
 \quad\text{and}\quad
 \prod_{j=0}^{S(k)-1}\left(1-\frac{e^{-\beta}}{m-j}\right).
\]
Lemma \ref{lem:pd-harmonic-product}, followed by
\eqref{eq:pd-closure-opportunities}, shows that both products differ from
$1-x$ by at most
\[
 C_{\kappa,S,\xi,\zeta}
 \Delta_n(U,\delta,h)+\frac{C_{\mathcal J,\xi,\zeta}}n.
\]
On the exceptional prefix event we use the trivial bound by one.  Averaging
under $\Pp^\circ$, then adding the closure-tail probability
\eqref{eq:pd-closure-tail}, gives
\[
 \sup_{x\in[\zeta,1-\zeta]}
 \left|\Pp(L>\lfloor xN\rfloor\mid\mathcal H)-(1-x)\right|
\]
bounded by the right-hand side of \eqref{eq:pd-root-cdf}.  Since
$\{L/N\le x\}=\{L\le\lfloor xN\rfloor\}$, this is the claimed CDF estimate.
\end{proof}

\subsection{Deletion stability and regeneration}

\begin{lemma}[Deletion of an interior macroscopic cycle]
\label{lem:pd-deletion-stability}
Fix master-partition endpoints $0<V<U\le T$, $\xi>0$,
$\zeta\in(0,1/2)$, and $0\le\delta<1/4$.  On $\mathcal G_n$, start from a
$(U,\delta,\xi)$-regular history and expose the cycle of a uniform residual
root.  Let $N$ and $L$ be the residual size and the exposed cycle length.  With
$\chi_n,\pi_n$ as in Lemma \ref{lem:pd-two-scale-balance}, there is a constant
$C=C(\kappa,V,\xi,\zeta)$ such that
\begin{align}
&\Pp\Bigl(
 \zeta N\le L\le(1-\zeta)N,\ 
 \text{the post-deletion history is not }
 (V,\delta',\zeta\xi)\text{-regular}
 \ \Bigm|\ \mathcal H\Bigr)
 \le\pi_n,
 \label{eq:pd-deletion-joint}
\end{align}
where
\begin{equation}
 \delta'
 =C\bigl(\delta+\beta+\chi_n+\varepsilon_n+e^{-U/2}\bigr)
  +\frac{C_{\mathcal J,\xi,\zeta}}n.
 \label{eq:pd-deletion-error}
\end{equation}
The conclusion is asserted whenever the displayed $\delta'$ is smaller than
$1/4$, which is the only regime used below.  This is a joint-event estimate; it
does not condition on the exact random value of $L$.
\end{lemma}

\begin{proof}
Let $d_a$ be the number of vertices of the exposed cycle in $H_{n,a}$, and let
$d_{ab}$ be the number of its directed edges from $H_{n,a}$ to $H_{n,b}$.
The conditional prefix event in Lemma \ref{lem:pd-two-scale-balance} is
simultaneous over all prefix lengths.  It may therefore be evaluated at the
random closing time without conditioning on $L$.

On the event $\zeta N\le L\le(1-\zeta)N$, we have
$L\ge\zeta\xi n$.  Apply the scale-sensitive estimate
\eqref{eq:pd-scale-sensitive-balance} with inner cutoff $V$ and $\ell=L$.
For every row below $V$,
\[
 \left|\frac{d_a}{Lp_{n,a}}-1\right|
 \le C_{\kappa,V,\xi,\zeta}
       \Delta_n(U,\delta,h)+\frac{C_{\mathcal J,\xi,\zeta}}n.
\]
Regularity gives $m_a^R=Np_{n,a}(1+O(\delta))$, and hence
\begin{equation}
 \frac{d_a}{m_a^R}
 =\frac LN+O\left(
 \delta+\beta+\chi_n+\varepsilon_n+e^{-U/2}
 +\frac{C_{\mathcal J,\xi,\zeta}}n\right)
 \label{eq:pd-deleted-row-fraction}
\end{equation}
uniformly over rows below $V$.

The row-prefix estimate \eqref{eq:pd-row-prefix}, evaluated at the same random
closing time, gives
\[
 d_{ab}
 =\frac{d_a}{m_a^R}N_{ab}^R
  +O((\beta+\chi_n)N_{ab}^R).
\]
Combining this with \eqref{eq:pd-deleted-row-fraction},
\begin{equation}
 d_{ab}
 =\frac LN N_{ab}^R
  +O\left(
 \delta+\beta+\chi_n+\varepsilon_n+e^{-U/2}
 +\frac{C_{\mathcal J,\xi,\zeta}}n\right)N_{ab}^R.
 \label{eq:pd-deleted-cell-fraction}
\end{equation}
The quantity $d_{ab}$ is exactly the number of vertices removed from
$Q_{n,a}\cap H_{n,b}$, because every cycle edge from row $a$ to category $b$
uses one such target.

Write $x=L/N$ and let $R'$ be the post-deletion vertex set.  Equations
\eqref{eq:pd-deleted-row-fraction}--\eqref{eq:pd-deleted-cell-fraction} give
\[
 m_a^{R'}=(1-x)m_a^R+O(\Delta_n^*m_a^R),
 \qquad
 N_{ab}^{R'}=(1-x)N_{ab}^R+O(\Delta_n^*N_{ab}^R),
\]
where $\Delta_n^*$ denotes the parenthesized error in
\eqref{eq:pd-deleted-row-fraction}.  Since $1-x\ge\zeta$ and
$|R'|=(1-x)N\ge\zeta\xi n$, division by the new total fraction preserves the
regularity inequalities with error at most the quantity in
\eqref{eq:pd-deletion-error}.  Residual permanental factorization and
hereditary flatness remain valid after deleting the exposed edges.  The only
conditional failure event is the simultaneous prefix failure, whose
probability is at most $\pi_n$.  This proves
\eqref{eq:pd-deletion-joint}.
\end{proof}

\subsection{Independent residual fractions}

Order the cycles size-biased by choosing a uniform root from the current
residual vertex set, deleting its cycle, and repeating.  Let
\[
 V_{n,j}=\frac{\widetilde L_{n,j}}
 {n-\widetilde L_{n,1}-\cdots-\widetilde L_{n,j-1}}
\]
whenever the denominator is positive, and set $V_{n,j}=0$ after exhaustion.

\begin{lemma}[Conditional test functions from an interior CDF bound]
\label{lem:pd-cdf-test}
Let $W\in[0,1]$ and let $\mathscr F$ be a sigma-field.  Suppose that, almost
surely,
\[
        \sup_{x\in[a,b]}
        |\Pp(W\le x\mid\mathscr F)-x|\le\eta,
        \qquad 0<a<b<1.
\]
Then
\[
        \Pp(W\notin[a,b]\mid\mathscr F)
        \le a+(1-b)+2\eta.
\]
If $f$ is absolutely continuous, vanishes outside $[a,b]$, and satisfies
$f(a)=f(b)=0$, then
\[
 \left|\Ee(f(W)\mid\mathscr F)-\int_0^1f(u)\dd u\right|
 \le\eta\int_a^b|f'(u)|\dd u.
\]
\end{lemma}

\begin{proof}
The endpoint estimate follows by evaluating the CDF bound at $a$ and $b$.
For the second assertion, conditional Stieltjes integration by parts and the
vanishing boundary values give
\[
        \Ee(f(W)\mid\mathscr F)
        =-\int_a^b\Pp(W\le u\mid\mathscr F)f'(u)\dd u,
\]
whereas
\[
        \int_a^bf(u)\dd u=-\int_a^bu f'(u)\dd u.
\]
Subtracting proves the claim.
\end{proof}

\begin{proposition}[Finite-dimensional $\operatorname{GEM}(1)$ limits]
\label{prop:pd-residual-uniforms}
For every fixed $k$,
\[
        (V_{n,1},\ldots,V_{n,k})
        \Rightarrow(U_1,\ldots,U_k),
\]
where $U_1,\ldots,U_k$ are independent $\operatorname{Unif}(0,1)$ variables.
\end{proposition}

\begin{proof}
Fix $k$ and $\zeta\in(0,1/4)$, and put
$\xi_*=\zeta^{k-1}$.  We first record a noncircular parameter hierarchy.
Let $\eta>0$ be a target error.

Choose $T_0$ so large that the closure-tail term
$C_{\xi_*}e^{-T_0}$ in Proposition
\ref{prop:pd-uniform-residual-fraction} is below a prescribed small fraction of
$\eta$.  Suppose $T_0,\ldots,T_j$ have been chosen.  Let $R_j$ and $D_j$ be,
respectively, upper bounds for the principal-error coefficients in the root law
and deletion lemma with inner cutoff $T_j$, residual fraction at least
$\xi_*$, and truncation parameter $\zeta$.  These constants are now fixed.  An
error created at this level can subsequently be amplified only by
$D_{j-1},\ldots,D_0$, all of which have already been fixed.  Put
\[
        A_j=(1+R_j+D_j)\prod_{i=0}^{j-1}(1+D_i),
\]
with the empty product equal to one, and choose $T_{j+1}>T_j$ so large that
\[
        A_je^{-T_{j+1}/2}\le 2^{-j-3}\eta.
\]
Continue until $T_k$ has been chosen.

Use one master partition of $[0,T_k]$ containing
$T_0,\ldots,T_k$, append the tail interval, and let its mesh be at most $h$.
After the cutoffs are fixed, all coefficients in the finitely many root and
deletion estimates are finite.  Choose $h$ so small that the errors involving
$\beta=2\kappa h$ remain below their assigned fractions of $\eta$ after the
finite deletion recursion.  More explicitly, if $D_j$ denotes the coefficient
in \eqref{eq:pd-deletion-error} with inner cutoff $T_j$, the deterministic
recursion
\[
        d_0=0,\qquad
        d_{r+1}=D_{k-r-1}(d_r+\beta+e^{-T_{k-r}/2})
\]
has $\max_{r\le k}d_r\to0$ as all displayed input errors tend to zero.
The preceding choices make this maximum smaller than a prescribed fraction of
$\eta$, and in particular smaller than $1/8$.  Finally, for this fixed master
partition take the maxima of the finitely many sequences $\chi_n$ and $\pi_n$
that occur at the different cutoff levels, and choose $n$ so large that
\[
        \varepsilon_n+\chi_n+\pi_n+C_{\mathcal J}/n
\]
is below its assigned fraction of $\eta$ and
$\Pp(\mathcal G_n^c)\le\eta$.  This order of choices is important: the inner
cutoff is fixed before its outer cutoff, and the mesh is chosen only after all
cutoff-dependent constants are known.

On the event that
\[
        \zeta\le V_{n,j}\le1-\zeta,
        \qquad 1\le j<r,
 \tag{6.20}\label{eq:pd-interior-history}
\]
the residual size before stage $r$ is at least
$\zeta^{r-1}n\ge\xi_*n$.  At stage $r$, use outer cutoff
$T_{k-r+1}$ and inner cutoff $T_{k-r}$.  Inductively, the current history is
regular at the outer cutoff with a regularity error bounded by the deterministic
recursion above, except on a conditional event of probability $O_k(\eta)$.
Proposition \ref{prop:pd-uniform-residual-fraction} therefore gives, uniformly
over every such current history,
\begin{equation}
 \sup_{x\in[\zeta,1-\zeta]}
 \left|\Pp(V_{n,r}\le x\mid\mathscr F_{r-1})-x\right|
 \le C_k\eta.
 \label{eq:pd-stage-cdf}
\end{equation}
On the event $V_{n,r}\in[\zeta,1-\zeta]$, Lemma
\ref{lem:pd-deletion-stability} makes the next history regular at cutoff
$T_{k-r}$, except on a conditional event of probability at most $C_k\eta$.
All these estimates are uniform over the preceding regular history.

We now turn the conditional estimates into joint convergence.  Let
$f_1,\ldots,f_k$ be bounded Lipschitz functions on $[0,1]$.  Choose a piecewise
linear cutoff $\varphi_\zeta$ which is zero outside
$[\zeta,1-\zeta]$ and equal to one on
$[2\zeta,1-2\zeta]$, and put $g_j=f_j\varphi_\zeta$.  Lemma
\ref{lem:pd-cdf-test}, applied to \eqref{eq:pd-stage-cdf}, gives
\[
 \left|\Ee(g_r(V_{n,r})\mid\mathscr F_{r-1})
       -\int_0^1g_r(u)\dd u\right|
 \le C_{f_r,\zeta}\eta
\]
on every regular history.  Let $\mathcal R_r$ be the event that, after the
first $r$ deletions, the history is regular at cutoff $T_{k-r}$ with the
prescribed error; put $\mathcal R_0=\mathcal G_n$.  Because $g_r$ vanishes
outside $[\zeta,1-\zeta]$, the joint-event deletion lemma also gives, on
$\mathcal R_{r-1}$,
\[
 \left|\Ee\bigl(g_r(V_{n,r})\one_{\mathcal R_r}
       \mid\mathscr F_{r-1}\bigr)
       -\int_0^1g_r(u)\dd u\right|
 \le C_{f_r,\zeta}\eta.
\]
Multiplying by
$\prod_{j<r}g_j(V_{n,j})\one_{\mathcal R_{r-1}}$, taking expectations, and
inducting on $r$ yields
\[
 \left|\Ee\left(\prod_{j=1}^kg_j(V_{n,j})
        \one_{\mathcal R_k}\right)
       -\prod_{j=1}^k\int_0^1g_j(u)\dd u\right|
 \le C_{f,k,\zeta}\eta+o_n(1).
\]
On the event that every $g_j(V_{n,j})$ is nonzero, every residual fraction is
in the interior range, so a union bound over the conditional deletion failures
shows that removing $\one_{\mathcal R_k}$ costs another
$C_{f,k}\eta+o_n(1)$.  Consequently
\begin{equation}
 \left|\Ee\prod_{j=1}^kg_j(V_{n,j})
       -\prod_{j=1}^k\int_0^1g_j(u)\dd u\right|
 \le C_{f,k,\zeta}\eta+o_n(1).
 \label{eq:pd-truncated-product}
\end{equation}
Here the $o_n(1)$ includes $\Pp(\mathcal G_n^c)$.

Evaluating \eqref{eq:pd-stage-cdf} at $2\zeta$ and $1-2\zeta$, and adding the
regularity-failure probabilities, gives
\begin{equation}
 \limsup_{n\to\infty}\Pp\left(
 V_{n,j}\notin[2\zeta,1-2\zeta]
 \text{ for some }j\le k\right)
 \le4k\zeta+C_k\eta.
 \label{eq:pd-actual-endpoints}
\end{equation}
Consequently, replacing $g_j$ by $f_j$ in
\eqref{eq:pd-truncated-product} costs at most
$C_{f,k}(\zeta+\eta)$, and the same is true for replacing
$\int g_j$ by $\int f_j$.  Let first $n\to\infty$, then
$\eta\downarrow0$, and finally $\zeta\downarrow0$.  We obtain
\[
        \Ee\prod_{j=1}^kf_j(V_{n,j})
        \longrightarrow
        \prod_{j=1}^k\int_0^1f_j(u)\dd u.
\]
Products of bounded Lipschitz functions form a convergence-determining class on
$[0,1]^k$.  This proves convergence to $k$ independent uniform variables.
The convention $V_{n,j}=0$ after exhaustion causes no difficulty, because
exhaustion before stage $k$ is contained in the endpoint event in
\eqref{eq:pd-actual-endpoints}.
\end{proof}

\begin{proof}[Proof of Theorem \ref{thm:macroscopic-pd}]
Put
\[
 X_{n,j}=\frac{\widetilde L_{n,j}}n
 =V_{n,j}\prod_{i<j}(1-V_{n,i}),
 \qquad
 R_{n,k}=1-\sum_{j=1}^kX_{n,j}
        =\prod_{i=1}^k(1-V_{n,i}).
\]
Proposition \ref{prop:pd-residual-uniforms} and the continuous mapping theorem
give, for every fixed $k$,
\begin{align}
 &(X_{n,1},\ldots,X_{n,k},R_{n,k})\notag\\
 &\qquad\Rightarrow
 \left(U_1,(1-U_1)U_2,\ldots,
 U_k\prod_{i<k}(1-U_i),\prod_{i=1}^k(1-U_i)\right).
 \label{eq:pd-finite-gem}
\end{align}
Because $0\le R_{n,k}\le1$,
\[
        \Ee R_{n,k}\longrightarrow2^{-k}.
\]
Thus, for every $\epsilon>0$,
\[
 \limsup_{n\to\infty}\Pp(R_{n,k}>\epsilon)
 \le\epsilon^{-1}2^{-k}.
 \label{eq:pd-tail-mass}
\]
For a bounded Lipschitz functional on $\ell^1$, truncate after $k$ coordinates,
use \eqref{eq:pd-finite-gem}, and bound the truncation error by the residual
tail mass.  Equation \eqref{eq:pd-tail-mass}, followed by $k\to\infty$, proves
\[
        (X_{n,1},X_{n,2},\ldots)
        \Rightarrow\operatorname{GEM}(1)
\]
in $\ell^1$.

For finite nonnegative sequences, the usual two-point uncrossing inequality
shows that matching coordinates in decreasing order minimizes the sum of
absolute differences.  Passing to truncations gives
\[
        \|x^\downarrow-y^\downarrow\|_1\le\|x-y\|_1
\]
for nonnegative $\ell^1$ sequences.  The continuous mapping theorem therefore
gives the $\operatorname{PD}(1)$ limit.  The first residual fraction is the
cycle length of a uniform root divided by $n$, which proves the final assertion
of the theorem.

The argument was carried out along an arbitrary subsequence after extracting a
marked-profile limit.  The resulting law does not depend on that profile.
Every subsequence therefore has a further subsequence with the same limit, so
the full sequence converges.
\end{proof}

\begin{corollary}[Power sums of macroscopic cycles]
\label{cor:pd-power-sums}
Under the hypotheses of Theorem \ref{thm:macroscopic-pd}, for every $q>1$,
\[
 \sum_{j\ge1}\left(\frac{L_{n,j}}n\right)^q
 \Rightarrow\sum_{j\ge1}P_j^q,
 \qquad (P_j)\sim\operatorname{PD}(1),
\]
and
\[
 \Ee\sum_{j\ge1}\left(\frac{L_{n,j}}n\right)^q
 \longrightarrow\frac1q.
\]
\end{corollary}

\begin{proof}
The map $x\mapsto\sum_jx_j^q$ is continuous and bounded by one on
$\mathcal S^\downarrow$.  If $\widetilde P$ is a size-biased pick from a
partition $(P_j)$, then, conditionally on the partition,
\[
        \Ee(\widetilde P^{\,q-1}\mid(P_j))=\sum_jP_j^q.
\]
For a $\operatorname{PD}(1)$ partition, its size-biased ordering has the
$\operatorname{GEM}(1)$ law, so $\widetilde P$ is uniform on $(0,1)$.  Hence
\[
        \Ee\sum_jP_j^q
        =\Ee\widetilde P^{\,q-1}
        =\int_0^1u^{q-1}\dd u=1/q.
\]
Bounded convergence transfers the expectation limit.
\end{proof}

\section{Sukhatme fixed points}

We now prove the singular Sukhatme fixed-point mean, a local terminal Poisson point-process limit, and fixed-order factorial moment asymptotics.  Throughout this section
\[
        \theta_i=n+1-i.
\]
We call this the Sukhatme profile because the same decreasing rates \(n,n-1,\ldots,1\) occur in the classical exponential-spacing representation associated with Sukhatme and R\'enyi \cite{Sukhatme1937,Renyi1953}.  No external result from that theory is used below; the citations record only the historical source of the rate pattern and its standard order-statistic interpretation.
It is convenient to use the terminal rate coordinate
\[
        a=n+1-i.
\]
Thus rate coordinate \(a\) is the actual exponential rate of the corresponding clock, and it corresponds to label \(i=n+1-a\).  We write
\[
        T_a\sim\Exp(a),\qquad 1\le a\le n,
\]
for these independent clocks.  This is a change of notation from the regular-profile section, where the rates were denoted \(\theta_{n,i}\).  The fixed-point indicator in rate coordinates is
\[
        I_{n,a}=\one\left\{
        \sum_{b\ne a}\one_{\{T_b>T_a\}}=a-1
        \right\},
\]
and the total number of fixed points is
\[
        F_n=\sum_{a=1}^n I_{n,a}.
\]
Thus \(F_n\) is the fixed-point count \(C_1(\sigma_n)\) specialized to the singular Sukhatme profile, written in terminal rate coordinates.

The scale in this section is terminal rather than bulk.  The relevant rates satisfy
\(a\to\infty\), \(a=o(n)\), and their clock times are of order \(1/a\).  A single rate \(a\) contributes asymptotically \(e^{-1}/a\).  Summing over a logarithmic range gives the mean \(e^{-1}\log n\).  Joint probabilities factor for rates separated on logarithmic scale, which gives a local Poisson point process on the coordinate \(\log a\).  The proof parallels the regular-profile argument, but with terminal rather than bulk geometry: a saddle computation and local central limit estimate again produce the local entry probability, while the integrated saddle is anisotropic and the residual domination comes from the terminal death process rather than the bulk endpoint and collision envelopes.  The proof has three ingredients: finite-tail saddle geometry, an anisotropic terminal local central limit estimate, and a death-process product bound for non-separated tuples.

\begin{lemma}[Death-process product bound]\label{lem:death-bound}
For every finite set \(A\subset[n]\),
\[
        \Pp(I_{n,a}=1\text{ for all }a\in A)
        \le \prod_{a\in A}\frac{2}{a+1}
        \le 2^{|A|}\prod_{a\in A}\frac1a .
\]
\end{lemma}

\begin{proof}
Expose the exponential race as a pure death process.  When the alive set is \(S\), the next deleted rate is \(b\in S\) with conditional probability
\[
        \frac{b}{\sum_{c\in S}c}.
\]
The event \(I_{n,a}=1\) says that rate \(a\) is deleted at the unique transition at which the number of alive particles drops from \(a\) to \(a-1\).  Equivalently, immediately before \(T_a\) there are \(a\) alive rates, including \(a\), and immediately after \(T_a\) there are \(a-1\).  Thus \(I_{n,a}=1\) is exactly the event that rate \(a\) is deleted at this death transition.  Condition on the entire history before this transition.  If rate \(a\) is already dead, the conditional probability is zero.  If it is alive, the alive set has size \(a\) and contains \(a\).  The smallest possible total rate of such a set is
\[
        1+2+\cdots+a=\frac{a(a+1)}2,
\]
so the conditional probability that rate \(a\) is next deleted is at most
\[
        \frac{a}{a(a+1)/2}=\frac{2}{a+1}.
\]
Order the elements of \(A\) decreasingly and condition successively on the histories before the corresponding death levels.  Multiplying the conditional bounds gives the result.
\end{proof}

\subsection{The one-point terminal saddle}

We first record the elementary finite-tail estimates used throughout the terminal analysis.  They are stated in relative form because this is the form needed below: when \(x=o(n)\), the omitted infinite tail is negligible relative to \(x\), but it need not be \(O(1)\).

\begin{lemma}[Finite-tail geometric estimates]\label{lem:terminal-geometric}
Let \(x=x_n\to\infty\) and \(x=o(n)\).  Uniformly for \(s\) in compact subsets of \((0,\infty)\),
\[
        \sum_{b=1}^n e^{-bs/x}=\frac{x}{s}(1+o(1)),
\]
\[
        \sum_{b=1}^n \frac{b}{x}e^{-bs/x}=\frac{x}{s^2}(1+o(1)),
        \qquad
        \sum_{b=1}^n e^{-2bs/x}=\frac{x}{2s}(1+o(1)).
\]
Also
\[
        \sum_{b=1}^n b^2 e^{-bs/x}=O(x^3),
\]
uniformly for the same range of \(s\).  The same estimates remain true if a bounded number of summands is omitted.
\end{lemma}

\begin{proof}
For example,
\[
        \sum_{b=1}^n e^{-bs/x}
        =\frac{e^{-s/x}(1-e^{-sn/x})}{1-e^{-s/x}}.
\]
Since \(s\) is bounded away from zero and infinity, \(e^{-sn/x}=o(1)\) uniformly and \((1-e^{-s/x})^{-1}=x/s+O(1)\).  This gives the first estimate in relative form.

For the weighted sum, put \(q=e^{-s/x}\).  Then
\[
        \sum_{b=1}^n bq^b
        =
        \frac{q\bigl(1-(n+1)q^n+nq^{n+1}\bigr)}{(1-q)^2}.
\]
Since \(n/x\to\infty\), the truncation terms are exponentially small on the \(n/x\) scale, while
\[
        \frac{q}{(1-q)^2}=\frac{x^2}{s^2}(1+o(1))
\]
uniformly on compact \(s\)-intervals.  Dividing by \(x\) gives the second estimate.  The third estimate is the first estimate with \(s\) replaced by \(2s\).  For the final bound, compare with the infinite tail:
\[
        \sum_{b=1}^n b^2e^{-bs/x}\le \sum_{b=1}^{\infty}b^2e^{-bs/x}=O(x^3),
\]
uniformly on compact \(s\)-intervals bounded away from zero.  Omitting a bounded number of summands changes the first, second, and third displayed sums by \(O(1)\), \(O(1)\), and \(O(1)\), respectively, and changes the final sum by at most \(O(x^2)\).  Indeed,
\[
        \sup_{b\ge1}\frac{b}{x}e^{-bs/x}=O(1),
        \qquad
        \sup_{b\ge1}b^2e^{-bs/x}=O(x^2),
\]
uniformly for \(s\) in compact subsets of \((0,\infty)\).  These errors are negligible relative to the corresponding main scales, or harmless for the displayed \(O(x^3)\) bound.
\end{proof}

\begin{lemma}[One-point terminal saddle geometry]\label{lem:suk-one-saddle}
Assume \(a\to\infty\) and \(a=o(n)\).  Define
\[
        \mu_{n,a}(s)=\sum_{b\ne a}e^{-bs/a},
        \qquad
        v_{n,a}(s)=\sum_{b\ne a}e^{-bs/a}(1-e^{-bs/a}).
\]
Let \(s_{n,a}^*\) be the unique solution of
\[
        \mu_{n,a}(s_{n,a}^*)=a-1.
\]
Then, uniformly in this range,
\[
        s_{n,a}^*\to1,
        \qquad
        -\mu_{n,a}'(s_{n,a}^*)\sim a,
        \qquad
        v_{n,a}(s_{n,a}^*)\sim \frac a2 .
\]
Moreover, for all sufficiently large \(n\), uniformly for \(s\in[1/2,2]\),
\[
        |\mu_{n,a}(s)-(a-1)|\ge c a |s-s_{n,a}^*|
\]
whenever \(|s-s_{n,a}^*|\le 1/4\), with an absolute constant \(c>0\).
\end{lemma}

\begin{proof}
The function \(\mu_{n,a}\) is continuous and strictly decreasing, tends to \(n-1\) as \(s\downarrow0\), and tends to zero as \(s\to\infty\).  Hence the saddle exists and is unique.  By Lemma \ref{lem:terminal-geometric}, uniformly for \(s\) in compact subsets of \((0,\infty)\),
\[
        \mu_{n,a}(s)=\frac{a}{s}(1+o(1)),
        \qquad
        -\mu_{n,a}'(s)=\sum_{b\ne a}\frac{b}{a}e^{-bs/a}=\frac{a}{s^2}(1+o(1)),
\]
and
\[
        v_{n,a}(s)
        =\sum_{b\ne a}e^{-bs/a}-\sum_{b\ne a}e^{-2bs/a}
        =\frac{a}{2s}(1+o(1)).
\]
These estimates also locate the saddle.  Fix \(\varepsilon\in(0,1)\).  Uniformly,
\[
        \mu_{n,a}(1-\varepsilon)
        =\frac{a}{1-\varepsilon}(1+o(1))>a-1,
\]
whereas
\[
        \mu_{n,a}(1+\varepsilon)
        =\frac{a}{1+\varepsilon}(1+o(1))<a-1
\]
for all sufficiently large \(n\).  Since \(\mu_{n,a}\) is strictly decreasing, this implies
\(1-\varepsilon<s_{n,a}^*<1+\varepsilon\), and hence \(s_{n,a}^*\to1\).  The displayed derivative and variance asymptotics follow by substitution.

Finally, on \([1/2,2]\) the derivative satisfies \(-\mu_{n,a}'(s)\ge c a\) for all large \(n\).  The mean-value theorem gives the stated linear separation from the saddle.
\end{proof}

\begin{lemma}[Sukhatme one-point estimate]\label{lem:suk-one}
Let \(p_{n,a}=\Pp(I_{n,a}=1)\).  Uniformly for \(a\to\infty\) and \(a=o(n)\),
\[
        p_{n,a}\sim \frac1{ea}.
\]
Moreover, for all \(1\le a\le n\),
\[
        p_{n,a}\le \frac{2}{a+1}.
\]
\end{lemma}

\begin{proof}
The upper bound is Lemma \ref{lem:death-bound} with \(A=\{a\}\).  It remains to prove the asymptotic.

Conditioning on \(T_a=t\) and then putting \(t=s/a\) gives
\[
        p_{n,a}=\int_0^\infty e^{-s}
        \Pp(Y_{n,a}(s)=a-1)\dd s,
\]
where
\[
        Y_{n,a}(s)=\sum_{b\ne a}\xi_b(s),
        \qquad \xi_b(s)\sim\operatorname{Bernoulli}(e^{-bs/a})
\]
are independent.  Let \(s_{n,a}^*\) be the finite saddle from Lemma \ref{lem:suk-one-saddle}.

In the saddle window write
\[
        s=s_{n,a}^*+\frac{u}{\sqrt a},
        \qquad |u|\le a^{1/10}.
\]
For $s\in[1/2,2]$, a positive proportion of the indices $b\asymp a$ have
Bernoulli parameters $e^{-bs/a}$ in a fixed compact subinterval of $(0,1)$.
Under every bounded scalar exponential tilt these parameters remain uniformly
nondegenerate.  The total tilted variance is comparable to $a$, and the product
characteristic function has exponential decay away from $2\pi\mathbb Z$.
Thus all tilted hypotheses of Lemma \ref{lem:llt} hold uniformly with scale
$s_n=a$.  We use its local estimate in two forms.  First, on each fixed
\(u\)-ball, Lemma \ref{lem:llt}, together with Lemma
\ref{lem:suk-one-saddle}, gives
\[
        \Pp(Y_{n,a}(s)=a-1)
        =\frac{1}{\sqrt{\pi a}}e^{-u^2}(1+o(1)).
\]
The exponent is obtained from
\[
        a-1-\mu_{n,a}(s)
        =-\mu_{n,a}'(s_{n,a}^*)(s-s_{n,a}^*)+o(\sqrt a)
        =u\sqrt a(1+o(1)),
\]
and \(v_{n,a}(s_{n,a}^*)\sim a/2\).  Second, the Gaussian upper bound in the same one-dimensional form of Lemma \ref{lem:llt} gives
\[
        \Pp(Y_{n,a}(s)=a-1)
        \le C a^{-1/2}e^{-c u^2},
        \qquad |u|\le a^{1/10}.
\]
This bound justifies extending the fixed \(u\)-ball to the whole growing saddle window.

The contribution away from the saddle is negligible.  If \(s\notin[1/2,2]\), the mean of \(Y_{n,a}(s)\) is separated from \(a-1\) by a positive multiple of \(a\), so Bernstein's inequality gives an exponentially small contribution.  If \(s\in[1/2,2]\) but \(|s-s_{n,a}^*|>a^{-2/5}\), Lemma \ref{lem:suk-one-saddle} gives
\[
        |\mu_{n,a}(s)-(a-1)|
        \ge c a |s-s_{n,a}^*|.
\]
Since \(v_{n,a}(s)=O(a)\) uniformly on \([1/2,2]\), Bernstein's inequality gives
\[
        \Pp(Y_{n,a}(s)=a-1)
        \le C\exp[-c a(s-s_{n,a}^*)^2].
\]
The integral of this bound over \(|s-s_{n,a}^*|>a^{-2/5}\) is \(o(a^{-1})\).  Also \(e^{-s}=e^{-s_{n,a}^*}(1+o(1))\) uniformly for \(|u|\le a^{1/10}\).  Therefore
\[
\begin{aligned}
        p_{n,a}
        &\sim \int_{-\infty}^{\infty}
        e^{-s_{n,a}^*}
        \frac{1}{\sqrt{\pi a}}e^{-u^2}\frac{\dd u}{\sqrt a}  \\
        &\sim \frac{e^{-1}}{a\sqrt\pi}\int_{-\infty}^{\infty}e^{-u^2}\dd u
        =\frac1{ea}.
\end{aligned}
\]
\end{proof}

\begin{theorem}[Sukhatme fixed-point mean]\label{thm:suk-mean}
For Sukhatme weights \(\theta_i=n+1-i\),
\[
        \Ee F_n\sim e^{-1}\log n.
\]
\end{theorem}

\begin{proof}
Choose integer cutoffs
\[
        L_n=\lceil\log n\rceil,
        \qquad
        U_n=\lfloor n/\log n\rfloor .
\]
On this range \(L_n\to\infty\) and \(U_n=o(n)\).  Hence the uniform estimate in Lemma \ref{lem:suk-one} gives
\[
        \sup_{L_n\le a\le U_n}
        \left|ea\,p_{n,a}-1\right|\to0.
\]
Therefore
\[
        \sum_{L_n\le a\le U_n}p_{n,a}
        =e^{-1}(1+o(1))
        \sum_{L_n\le a\le U_n}\frac1a
        \sim e^{-1}\log n.
\]
The lower range contributes
\[
        \sum_{a<L_n}p_{n,a}
        \le 2\sum_{a<L_n}\frac1{a+1}
        =O(\log\log n)=o(\log n),
\]
and the upper range contributes
\[
        \sum_{a>U_n}p_{n,a}
        \le 2\sum_{a>U_n}\frac1{a+1}
        =O(\log\log n)=o(\log n).
\]
Thus \(\Ee F_n=\sum_a p_{n,a}\sim e^{-1}\log n\).
\end{proof}

\subsection{Separated terminal joint estimates}

The joint estimate is anisotropic: the covariance scale in the \(j\)-th terminal coordinate is \(a_j\), not a common scalar scale.  We therefore use a terminal local CLT adapted to logarithmically separated rates.  The one-point saddle above differentiated with respect to the scaled variable \(s=at\); in the joint saddle below we differentiate with respect to the clock time \(t\), which accounts for the change from derivative scale \(a\) to derivative scale \(a^2\).

\begin{lemma}[Joint terminal saddle geometry]\label{lem:suk-joint-saddle}
Fix \(r\ge1\) and \(\eta>0\).  Uniformly over integers
\[
        1\ll a_1<\cdots<a_r\ll n,
        \qquad a_{j+1}/a_j\ge1+\eta,
\]
where the ratio-separation condition is empty when \(r=1\).  Let \(A=\{a_1,\ldots,a_r\}\) and define
\[
        \mu_j(t)=\sum_{b\notin A}e^{-bt},
        \qquad 1\le j\le r.
\]
There is a unique \(t_j^*>0\) satisfying
\[
        \mu_j(t_j^*)=a_j-j.
\]
Moreover
\[
        a_jt_j^*\to1,
        \qquad
        -\mu_j'(t_j^*)\sim a_j^2,
        \qquad
        \mu_j''(t_j^*)=O(a_j^3),
\]
uniformly in \(j\).  Finally, for all large \(n\),
\[
        t_1^*>t_2^*>\cdots>t_r^*,
\]
and the ratios \(t_j^*/t_{j+1}^*\) are bounded below by a constant greater than one depending only on \(\eta\).
\end{lemma}

\begin{proof}
Fix \(j\).  For \(t=s/a_j\), Lemma \ref{lem:terminal-geometric} gives, uniformly for \(s\) in compact subsets of \((0,\infty)\),
\[
        \mu_j(s/a_j)=\frac{a_j}{s}(1+o(1)),
        \qquad
        -\mu_j'(s/a_j)=\frac{a_j^2}{s^2}(1+o(1)).
\]
The omitted set \(A\) has fixed size \(r\).  It changes the first sum by \(O_r(1)\).  For the derivative sum, each omitted term satisfies
\[
        b e^{-bs/a_j}\le C a_j
\]
uniformly for \(s\) in compact subsets of \((0,\infty)\), and hence the omitted terms change the derivative sum by \(O_r(a_j)=o(a_j^2)\).  Thus both displayed estimates remain valid after removing \(A\).

The equation \(\mu_j(t)=a_j-j\) has a unique solution, since \(\mu_j\) is continuous and strictly decreasing from \(n-r\) to \(0\), while \(0<a_j-j<n-r\) for all large \(n\).  To locate it, fix \(\varepsilon\in(0,1)\).  The first displayed estimate gives
\[
        \mu_j((1-\varepsilon)/a_j)>a_j-j,
        \qquad
        \mu_j((1+\varepsilon)/a_j)<a_j-j
\]
for all sufficiently large \(n\), uniformly in \(j\).  Therefore
\[
        1-\varepsilon<a_jt_j^*<1+\varepsilon,
\]
and \(a_jt_j^*\to1\).  Substituting \(t=t_j^*\) into the preceding derivative estimate, equivalently taking \(s=a_jt_j^*\to1\), gives \(-\mu_j'(t_j^*)\sim a_j^2\) uniformly in \(j\).  The second-derivative bound follows from
\[
        \mu_j''(t)=\sum_{b\notin A}b^2e^{-bt}=O(a_j^3)
        \qquad(t\asymp a_j^{-1}),
\]
again by the same finite-geometric estimate.  Since \(t_j^*\sim1/a_j\) and \(a_{j+1}/a_j\ge1+\eta\), the saddle times are eventually strictly ordered and uniformly separated multiplicatively.
\end{proof}

\paragraph{Terminal interval notation.}
For the remainder of this subsection, fix an admissible tuple
\[
        1\ll a_1<\cdots<a_r\ll n,
        \qquad a_{j+1}/a_j\ge1+\eta,
\]
put \(A=\{a_1,\ldots,a_r\}\), and consider ordered times
\(t_1>\cdots>t_r>0\).  Define
\[
        Z_j(t)=\sum_{b\notin A}\one_{\{T_b>t_j\}},
        \qquad z_j=a_j-j,
        \qquad \mu_j(t_j)=\Ee Z_j(t).
\]
Let
\[
        I_1=(t_1,\infty),\qquad
        I_j=(t_j,t_{j-1}]\quad(2\le j\le r),\qquad
        I_{r+1}=(0,t_r],
\]
and write
\[
        Y_1(t)=Z_1(t),\qquad
        Y_j(t)=Z_j(t)-Z_{j-1}(t)\quad(2\le j\le r).
\]
The corresponding interval-count target is
\[
        y_1=a_1-1,
        \qquad
        y_j=a_j-a_{j-1}-1\quad(2\le j\le r).
\]
Finally, put
\[
\begin{aligned}
        m_j(t)&=\Ee Y_j(t),\\
        H_j(t)&=z_j-\mu_j(t_j),\qquad H_0(t)=0,\\
        d_j(t)&=y_j-m_j(t)=H_j(t)-H_{j-1}(t),\\
        s_j&=a_jt_j,\qquad s_j^*=a_jt_j^*.
\end{aligned}
\]
The transformation between \(Z\) and \(Y\) is triangular with determinant
one, so \(\Pp(Z(t)=z)=\Pp(Y(t)=y)\).

\begin{lemma}[Coordinatewise Fourier concentration]\label{lem:coord-fourier}
Let
\[
        W=\sum_{\nu\in\mathcal I}W_\nu
\]
be a sum of independent \(\mathbb Z^r\)-valued random vectors, each supported in a fixed cube \([-M_0,M_0]^r\).  Let
\(A_1,\ldots,A_r\ge1\).  Suppose that, for every coordinate \(j\), there is a set
\(\mathcal I_j\subset\mathcal I\) with
\(|\mathcal I_j|\ge c_0A_j\), and that the sets
\(\mathcal I_1,\ldots,\mathcal I_r\) are disjoint.  Suppose also that, for every
\(\nu\in\mathcal I_j\), the law of \(W_\nu\) gives mass at least \(c_0\) to two atoms whose difference is \(e_j\), the \(j\)-th coordinate vector.  Then there are constants
\(c,C>0\), depending only on \(r,M_0,c_0\), such that the characteristic function
\(\phi_W(u)=\Ee e^{iu\cdot W}\) satisfies
\[
        |\phi_W(u)|
        \le \exp\left[-c\sum_{j=1}^r
        A_j\operatorname{dist}(u_j,2\pi\mathbb Z)^2\right],
        \qquad u\in[-\pi,\pi]^r.
\]
Consequently,
\[
        \sup_z\Pp(W=z)\le C\prod_{j=1}^r A_j^{-1/2}.
\]
The same conclusions hold uniformly after any exponential tilt
\[
        \Pp_\lambda(W_\nu=x)
        =\frac{e^{\lambda\cdot x}\Pp(W_\nu=x)}
        {\Ee e^{\lambda\cdot W_\nu}},
        \qquad \|\lambda\|_\infty\le \Lambda,
\]
with constants depending additionally on \(\Lambda\).
\end{lemma}

\begin{proof}
For a single summand \(Y\), write \(p_x=\Pp(Y=x)\).  If two atoms
\(x\) and \(x+e_j\) both have mass at least \(c_0\), then
\[
\begin{aligned}
        1-|\Ee e^{iu\cdot Y}|^2
        &=\frac12\sum_{y,z}p_yp_z
        |e^{iu\cdot y}-e^{iu\cdot z}|^2  \\
        &\ge c_0^2 |1-e^{iu_j}|^2
        \ge c\,\operatorname{dist}(u_j,2\pi\mathbb Z)^2 .
\end{aligned}
\]
After decreasing \(c\), this gives
\[
        |\Ee e^{iu\cdot Y}|
        \le \exp\{-c\operatorname{dist}(u_j,2\pi\mathbb Z)^2\}.
\]
Multiplying this estimate over the disjoint blocks \(\mathcal I_j\) gives the displayed characteristic-function bound.  Fourier inversion then gives
\[
\begin{aligned}
        \Pp(W=z)
        &\le (2\pi)^{-r}\int_{[-\pi,\pi]^r}
        \exp\left[-c\sum_{j=1}^r
        A_j\operatorname{dist}(u_j,2\pi\mathbb Z)^2\right]\dd u \\
        &\le C\prod_{j=1}^r A_j^{-1/2}.
\end{aligned}
\]
Under a tilt with \(\|\lambda\|_\infty\le\Lambda\), all probabilities of atoms in the fixed support are multiplied by factors bounded between positive constants depending only on \(M_0,r,\Lambda\).  Hence the two distinguished atom probabilities remain bounded below, with a possibly smaller constant, and the same proof applies.
\end{proof}

\begin{lemma}[Weighted chain comparison]\label{lem:weighted-chain-comparison}
Fix $r\ge1$ and $\eta>0$.  Let
$0<A_1<\cdots<A_r$ satisfy $A_{j+1}/A_j\ge1+\eta$, and put
$u_{r+1}=0$.  Then
\[
        \sum_{j=1}^r A_j u_j^2
        \le C_{r,\eta}\sum_{j=1}^r A_j(u_j-u_{j+1})^2 .
\]
Equivalently, if $d_j=H_j-H_{j-1}$ with $H_0=0$, then
\[
        \sum_{j=1}^r \frac{d_j^2}{A_j}
        \asymp_{r,\eta}
        \sum_{j=1}^r \frac{H_j^2}{A_j} .
\]
\end{lemma}

\begin{proof}
Write $\Delta_j=u_j-u_{j+1}$.  Then $u_j=\sum_{k=j}^r\Delta_k$.  By Cauchy's inequality with the geometrically decaying weights $A_j/A_k\le(1+\eta)^{-(k-j)}$,
\[
        A_j u_j^2
        \le C_{r,\eta}\sum_{k=j}^r A_k\Delta_k^2 .
\]
Summing over $j$ gives the first display.  The second follows by applying the first inequality to the inverse triangular map and its reverse; since $r$ is fixed and the ratios are bounded below, both triangular maps have uniformly bounded operator norms in the weighted norms.
\end{proof}

\begin{lemma}[Truncated triangular comparison]\label{lem:truncated-triangular-comparison}
Fix $r\ge1$, $\eta>0$, and $M<\infty$.  Let
$0<A_1<\cdots<A_r$ satisfy $A_{j+1}/A_j\ge1+\eta$.  If
$d_j=H_j-H_{j-1}$, $H_0=0$, and $|d_j|\le M A_j$ for every $j$, then
\[
        \sum_{j=1}^r
        \min\left\{\frac{d_j^2}{A_j},A_j\right\}
        \ge c_{r,\eta,M}
        \sum_{j=1}^r
        \min\left\{\frac{H_j^2}{A_j},A_j\right\}.
\]
\end{lemma}

\begin{proof}
Write $x_j=d_j/A_j$.  Then $|x_j|\le M$ and
\[
        \frac{H_j}{A_j}=\sum_{k\le j}\frac{A_k}{A_j}x_k,
        \qquad
        \frac{A_k}{A_j}\le (1+\eta)^{-(j-k)} .
\]
Since $\min\{y^2,1\}\le y^2$, Cauchy's inequality and the geometric
bound on the ratios give, for each fixed $j$,
\[
\begin{aligned}
        A_j\min\left\{\frac{H_j^2}{A_j^2},1\right\}
        &\le A_j\left(\sum_{k\le j}\frac{A_k}{A_j}x_k\right)^2  \\
        &\le C_{r,\eta} A_j\sum_{k\le j}\frac{A_k}{A_j}x_k^2
         = C_{r,\eta}\sum_{k\le j} A_k x_k^2 .
\end{aligned}
\]
Because $|x_k|\le M$, we have
$x_k^2\le C_M\min\{x_k^2,1\}$.  Summing the preceding inequality over
$j$ and using that $r$ is fixed gives
\[
        \sum_{j=1}^r A_j
        \min\left\{\frac{H_j^2}{A_j^2},1\right\}
        \le C_{r,\eta,M}
        \sum_{k=1}^r A_k\min\{x_k^2,1\}.
\]
This is exactly the claimed inequality after rewriting
$A_j\min\{H_j^2/A_j^2,1\}=\min\{H_j^2/A_j,A_j\}$ and
$A_k\min\{x_k^2,1\}=\min\{d_k^2/A_k,A_k\}$.
\end{proof}

\begin{lemma}[Boundary-aware truncated chain comparison]
\label{lem:boundary-truncated-chain}
Fix \(r\ge1\) and \(M<\infty\).  Let
\(0<A_1<\cdots<A_r\), let \(B_-\) and \(B_+\) be disjoint subsets of
\(\{1,\ldots,r\}\), and put
\[
        G=\{1,\ldots,r\}\setminus(B_-\cup B_+).
\]
Let \(H_0=0\), \(d_j=H_j-H_{j-1}\), and
\[
        \phi_j(x)=\min\left\{\frac{x^2}{A_j},A_j\right\}.
\]
Assume
\[
        |d_j|\le M A_j\quad(j\in G),
        \qquad
        |H_j|\le M A_j\quad(j\in B_+).
\]
Then
\[
\begin{aligned}
&\sum_{j\in G}\phi_j(d_j)
 +\sum_{j\in B_-}|H_j|
 +\sum_{j\in B_+}A_j\\
&\qquad\ge c_{r,M}
 \sum_{j\in G}\phi_j(H_j).
\end{aligned}
\]
\end{lemma}

\begin{proof}
Decompose \(G\) into maximal consecutive runs.  Let
\(\{\ell,\ell+1,\ldots,u\}\) be one such run and put
\(b=\ell-1\).  For every \(j\) in the run,
\[
        H_j=H_b+\sum_{k=\ell}^j d_k,
\]
where \(H_b=H_0=0\) if \(\ell=1\).  For a fixed positive number \(A\) and
at most \(r+1\) real numbers \(x_i\),
\[
        \min\left\{\frac{(\sum_i x_i)^2}{A},A\right\}
        \le C_r\sum_i
        \min\left\{\frac{x_i^2}{A},A\right\}.
        \tag{7.1}
\]
Indeed, if one term on the right equals \(A\), the claim is immediate; otherwise
it follows from Cauchy's inequality.

If \(k\le j\) belongs to the run, then \(A_k\le A_j\) and
\(|d_k|\le M A_k\).  Hence
\[
        \phi_j(d_k)\le C_M\phi_k(d_k).
        \tag{7.2}
\]
To see this, the assertion is immediate when \(|d_k|\le A_k\); when
\(|d_k|>A_k\), one has
\(\phi_k(d_k)=A_k\) and
\(\phi_j(d_k)\le d_k^2/A_j\le M^2A_k\).

It remains to control the left boundary.  If \(b\in B_-\), then
\[
        \phi_j(H_b)\le |H_b|.
\]
If \(b\in B_+\), then \(|H_b|\le M A_b\) and \(A_b\le A_j\), so
\[
        \phi_j(H_b)\le M^2 A_b.
\]
Combining these estimates with (7.1) and (7.2), and then summing first over
\(j\) in the run and finally over the finitely many runs, proves the lemma.
\end{proof}

\begin{lemma}[Edge Fourier concentration]\label{lem:edge-fourier}
Let $e_1,\ldots,e_r$ be the coordinate vectors in $\mathbb Z^r$ and put $e_{r+1}=0$.  Let
\[
        W=\sum_{\nu\in\mathcal I}W_\nu
\]
be a sum of independent random vectors supported on
$\{e_1,\ldots,e_r,0\}$.  Let $A_1,\ldots,A_r\ge1$ satisfy
$A_{j+1}/A_j\ge1+\eta$ for $j<r$.  Suppose that, for every edge $j$, there is a set
$\mathcal I_j\subset\mathcal I$ with $|\mathcal I_j|\ge c_0A_j$, the sets $\mathcal I_j$ are disjoint, and for every $\nu\in\mathcal I_j$ the two atoms $e_j$ and $e_{j+1}$ each have probability at least $c_0$.  Then
\[
        |\Ee e^{iu\cdot W}|
        \le \exp\left[-c\sum_{j=1}^r
        A_j\operatorname{dist}(u_j-u_{j+1},2\pi\mathbb Z)^2\right],
        \qquad u_{r+1}=0,
\]
and consequently
\[
        \sup_z\Pp(W=z)\le C\prod_{j=1}^r A_j^{-1/2}.
\]
The same conclusions hold uniformly after any exponential tilt
$\|\lambda\|_\infty\le\Lambda$, with constants depending additionally on
$\Lambda$.
\end{lemma}

\begin{proof}
If a summand $Y$ gives mass at least $c_0$ to $e_j$ and to $e_{j+1}$, then
\[
        1-|\Ee e^{iu\cdot Y}|^2
        \ge c_0^2|e^{iu_j}-e^{iu_{j+1}}|^2
        \ge c\operatorname{dist}(u_j-u_{j+1},2\pi\mathbb Z)^2 .
\]
As in Lemma \ref{lem:coord-fourier}, this implies
\[
        |\Ee e^{iu\cdot Y}|
        \le \exp\{-c\operatorname{dist}(u_j-u_{j+1},2\pi\mathbb Z)^2\}.
\]
Multiplying over the disjoint edge blocks gives the displayed Fourier bound.
For the concentration estimate, Fourier inversion and Lemma
\ref{lem:weighted-chain-comparison} give
\[
\begin{aligned}
        \Pp(W=z)
        &\le (2\pi)^{-r}\int_{[-\pi,\pi]^r}
        \exp\left[-c\sum_{j=1}^r A_j
        \operatorname{dist}(u_j-u_{j+1},2\pi\mathbb Z)^2\right]\dd u \\
        &\le C\prod_{j=1}^r A_j^{-1/2}.
\end{aligned}
\]
Bounded tilts only multiply the finitely many atom probabilities by bounded factors, so the edge-block lower bounds persist after changing constants.
\end{proof}

\begin{lemma}[Forward-edge projected concentration]\label{lem:projected-forward-edge-concentration}
Let $G\ne\emptyset$ be a subset of $\{1,\ldots,r\}$ and let
$0<A_1<\cdots<A_r$ satisfy $A_{j+1}/A_j\ge1+\eta$.
Let $W$ be a sum of independent random vectors supported on
$\{e_j:j\in G\}\cup\{0\}$ after projection onto the coordinates in $G$.
Suppose that for each $j\in G$ there is an index
$\kappa(j)\in (G\cap\{j+1,\ldots,r\})\cup\{0\}$ and a set of at
least $c_0A_j$ summands, disjoint as $j$ varies, for which the two
projected atoms $e_j$ and $e_{\kappa(j)}$ have probabilities at least
$c_0$; here $e_0=0$.  Then, uniformly in $z_G$,
\[
        \Pp(W_G=z_G)\le C\prod_{j\in G} A_j^{-1/2}.
\]
The same estimate holds after any exponential tilt whose coordinates are
bounded in absolute value by a fixed constant.
\end{lemma}

\begin{proof}
For a summand in the block assigned to $j$, the same two-atom argument used
in Lemma \ref{lem:edge-fourier} gives the factor
\[
        \exp\{-c\operatorname{dist}(u_j-u_{\kappa(j)},2\pi\mathbb Z)^2\},
        \qquad u_0=0 .
\]
Multiplying over the disjoint blocks gives
\[
        |\phi(u)|\le
        \exp\left[-c\sum_{j\in G}A_j
        \operatorname{dist}(u_j-u_{\kappa(j)},2\pi\mathbb Z)^2\right].
\]
Because every directed edge goes from $j$ to a larger index, or to the
anchor $0$, each path reaches $0$ after at most $r$ steps.  Along such a
path the weights only increase.  Therefore, for $u_0=0$,
\[
        \sum_{j\in G}A_j u_j^2
        \le C_{r,\eta}\sum_{j\in G}A_j(u_j-u_{\kappa(j)})^2 .
\]
The torus version follows from the same path argument.  For each $j$, choose
integer multiples of $2\pi$ successively along the directed path $P(j)$ from
$j$ to the anchor $0$.  The triangle inequality on the circle gives
\[
        \operatorname{dist}(u_j,2\pi\mathbb Z)
        \le \sum_{\ell\in P(j)}
        \operatorname{dist}(u_\ell-u_{\kappa(\ell)},2\pi\mathbb Z).
\]
Cauchy's inequality, together with monotonicity of the weights along this path,
gives the displayed comparison with distances to $2\pi\mathbb Z$.
Fourier inversion then gives
\[
        \Pp(W_G=z_G)
        \le C\int_{[-\pi,\pi]^G}
        \exp\left[-c\sum_{j\in G}A_ju_j^2\right]\dd u_G
        \le C\prod_{j\in G}A_j^{-1/2}.
\]
Bounded tilts distort all atoms in the finite support by bounded factors, so
the two-atom lower bounds persist after changing constants.
\end{proof}

\begin{lemma}[Terminal interval edge blocks]\label{lem:terminal-interval-edge-blocks}
Fix \(r\ge1\), \(0<c_*<C_*<\infty\), and \(c_0>0\).  In the terminal
interval notation above, fix \(j\in\{1,\ldots,r\}\) and suppose
\[
        c_*\le s_j=a_jt_j\le C_*,
        \qquad
        m_j(t)=\sum_{b\notin A}\Pp(T_b\in I_j)\ge c_0a_j.
\]
Then there are constants \(0<\alpha<\beta<\infty\) and \(c>0\), depending
only on \(r,c_*,C_*\), and \(c_0\), and a set of at least \(ca_j\) rates
\(b\in[\alpha a_j,\beta a_j]\setminus A\) such that
\[
        \Pp(T_b\in I_j)\ge c,
        \qquad
        \Pp(T_b\le t_j)\ge c.
\]
Consequently, after projecting any chosen subset of the interval coordinates
larger than \(j\) either to their own coordinates or to the anchor \(0\), a
further subblock of size at least \(ca_j\) gives two projected atoms
\(e_j\) and \(e_{\kappa(j)}\), where \(\kappa(j)>j\) is a retained
coordinate or \(\kappa(j)=0\), each with probability at least \(c\).
The conclusion remains valid under every interval-coordinate exponential tilt
whose coordinates are bounded by a fixed constant, after changing \(c\).
\end{lemma}

\begin{proof}
Choose \(\alpha>0\) small and \(\beta<\infty\) large.  The contribution to
\(m_j(t)\) from \(b<\alpha a_j\) is at most \(\alpha a_j+1\).  Since
\(t_j\ge c_*/a_j\), the contribution from \(b>\beta a_j\) is at most
\[
        \sum_{b>\beta a_j}e^{-bt_j}
        \le C a_j e^{-c_*\beta}.
\]
After first decreasing \(\alpha\) and then increasing \(\beta\), at least a
fixed positive multiple of \(a_j\) of the total interval mass comes from
\([\alpha a_j,\beta a_j]\).  This block contains \(O(a_j)\) rates, so at
least \(ca_j\) of them satisfy \(\Pp(T_b\in I_j)\ge c\).  For every such
rate,
\[
        bt_j\ge\alpha c_*,
        \qquad
        \Pp(T_b\le t_j)=1-e^{-bt_j}\ge1-e^{-\alpha c_*}>0.
\]
The event \(\{T_b\le t_j\}\) is the disjoint union of the later categories
\(I_{j+1},\ldots,I_{r+1}\).  For each selected rate, one of these finitely
many categories has probability bounded below.  A second pigeonhole argument
produces a common later category on a subblock of size \(ca_j\).  Projection
maps that category either to a retained atom \(e_{\kappa(j)}\) or to the anchor
\(0\).  Removing the fixed set \(A\) changes the block size by only \(O_r(1)\).
A bounded exponential tilt changes every atom probability in the finite
categorical support by a bounded multiplicative factor, which proves the final
assertion.
\end{proof}

\begin{lemma}[Projected good--bad interval bound]\label{lem:projected-good-bad-interval-bound}
Fix $r\ge1$, $\eta>0$, and a compact scaled chamber
$K\subset(0,\infty)^r$.  In the terminal interval notation above, suppose
$s=(a_1t_1,\ldots,a_rt_r)\in K$ and $t_1>\cdots>t_r$.  Suppose the deterministic targets satisfy
\[
        c_0a_j\le y_j\le C_0a_j,
        \qquad 1\le j\le r .
\]
Then, uniformly for such $t$ and $y$,
\[
        \Pp(Y(t)=y)
        \le C\prod_{j=1}^r a_j^{-1/2}
        \exp\left[-c\sum_{j=1}^r
        \min\left\{\frac{(y_j-m_j(t))^2}{a_j},a_j\right\}\right].
\]
The constants depend only on $r,\eta,K,c_0$, and $C_0$.
\end{lemma}

\begin{proof}
Write $d_j=y_j-m_j(t)$.  Since $s\in K$, finite-geometric estimates give
$m_j(t)=O_K(a_j)$ and hence $|d_j|\le M_Ka_j$.  Fix a small
$\varepsilon>0$.  Let
\[
        G=\{j: |d_j|\le\varepsilon a_j\},
        \qquad B=\{1,\ldots,r\}\setminus G .
\]
For $j\in G$ the target lower bound gives
$m_j(t)\ge(c_0-\varepsilon)a_j$; choose
$\varepsilon<c_0/2$, so Lemma \ref{lem:terminal-interval-edge-blocks}
applies to every good coordinate.

Use one combined exponential tilt.  Put
\[
        \lambda_j=
        \begin{cases}
        \delta d_j/a_j, & j\in G,\\
        \delta\,\operatorname{sgn}(d_j), & j\in B,
        \end{cases}
\]
where $\delta>0$ is a sufficiently small fixed constant.  Let
$K_t(\lambda)=\log\Ee\exp\{\lambda\cdot(Y(t)-m(t))\}$.
For every bounded tilt $\xi$, the categorical support and the finite-geometric
bound $\Var_\xi Y_j(t)\le C a_j$ give
\[
        \Var_\xi(v\cdot Y(t))\le C\sum_{j=1}^r a_jv_j^2.
\]
Therefore
\[
        K_t(\lambda)\le C\sum_{j=1}^r a_j\lambda_j^2 .
\]
On the other hand,
\[
        \lambda\cdot d
        \ge \delta\sum_{j\in G}\frac{d_j^2}{a_j}
            +\delta\varepsilon\sum_{j\in B}a_j .
\]
Choosing $\delta$ small enough yields
\[
        -\lambda\cdot d+K_t(\lambda)
        \le -c\sum_{j\in G}\frac{d_j^2}{a_j}
             -c\sum_{j\in B}a_j .
\]
Under this bounded tilt all categorical probabilities are changed by bounded
multiplicative factors.  Project the tilted array onto the good coordinates
and ignore the bad-coordinate constraints.  For each $j\in G$, Lemma \ref{lem:terminal-interval-edge-blocks} gives a candidate block $B_j$ of at least $c_0a_j$
rates with two projected atoms $e_j$ and $e_{\kappa(j)}$, where
$\kappa(j)>j$ is another good coordinate or $\kappa(j)=0$.  We now extract
disjoint subblocks rigorously.  Process the indices of $G$ in increasing order.
Since $a_{i+1}/a_i\ge1+\eta$,
\[
        \sum_{i<j}a_i\le C_{r,\eta}a_j.
\]
Choose $c'>0$ so small that
$c'(1+C_{r,\eta})\le c_0$.  After subblocks of sizes $c'a_i$ have been chosen
for all earlier good indices $i<j$, at most
$c'\sum_{i<j}a_i\le c'C_{r,\eta}a_j$ elements of $B_j$ have been used.
At least $c'a_j$ unused elements of $B_j$ therefore remain.  Continuing
greedily produces disjoint subblocks of size at least $c'a_j$ for every
$j\in G$.
Lemma \ref{lem:projected-forward-edge-concentration} gives
\[
        \Pp_\lambda(Y_G(t)=y_G)\le C\prod_{j\in G}a_j^{-1/2}.
\]
Consequently
\begin{align*}
        \Pp(Y(t)=y)
        &\le \exp\{-\lambda\cdot d+K_t(\lambda)\}
        \Pp_\lambda(Y_G(t)=y_G)  \\
        &\le C\prod_{j\in G}a_j^{-1/2}
        \exp\left[-c\sum_{j\in G}\frac{d_j^2}{a_j}
                  -c\sum_{j\in B}a_j\right].
\end{align*}
For $j\in B$, the factor $e^{-ca_j}$ absorbs the missing local factor:
$e^{-ca_j}\le C a_j^{-1/2}e^{-c'a_j}$.  Since for $j\in G$ we have
$|d_j|\le\varepsilon a_j$, this is the desired bound, after changing
constants.

\end{proof}


\begin{lemma}[Actual-mean cumulative Chernoff bounds]\label{lem:actual-mean-lower-cumulative-chernoff}
Fix \(r\ge1\).  Let \(B\subset\{1,\ldots,r\}\), and put
\[
        \mu_j=\Ee Z_j(t),\qquad j\in B.
\]
There are constants \(\delta_0,C>0\), depending only on \(r\), such that,
for every \(|\theta|\le\delta_0\),
\[
\log \Ee\exp\left\{\theta\sum_{j\in B}(Z_j(t)-\mu_j)\right\}
        \le C\theta^2\sum_{j\in B}\mu_j .
\]
After decreasing \(\delta_0\), if necessary, fix
\(0<\delta\le\delta_0\) and \(C_0<\infty\).  There are
\(c=c(\delta,r,C_0)>0\) and \(L_0=L_0(\delta,r,C_0)\) such that, whenever
\(L\ge L_0\), both of the following assertions hold.

If \(z_j\le C_0L^{-1}\mu_j\) for every \(j\in B\), then
\[
        -\sum_{j\in B}(-\delta)(z_j-\mu_j)
        +\log \Ee\exp\left\{-\delta\sum_{j\in B}(Z_j(t)-\mu_j)\right\}
        \le -c\sum_{j\in B}\mu_j .
\]
If \(\mu_j\le C_0L^{-1}z_j\) for every \(j\in B\), then
\[
        -\sum_{j\in B}\delta(z_j-\mu_j)
        +\log \Ee\exp\left\{\delta\sum_{j\in B}(Z_j(t)-\mu_j)\right\}
        \le -c\sum_{j\in B}z_j .
\]
In particular, the first exponent is at most
\(-cL\sum_{j\in B}a_j\) if \(\mu_j\ge c_0La_j\), while the second is at most
\(-c\sum_{j\in B}a_j\) if \(z_j\ge c_0a_j\).
\end{lemma}

\begin{proof}
For a single rate \(b\notin A\), define
\[
        R_b=\sum_{j\in B}\one_{\{T_b>t_j\}},
        \qquad m_b=\Ee R_b .
\]
Then \(0\le R_b\le r\), and
\[
        \sum_{b\notin A}m_b=\sum_{j\in B}\mu_j .
\]
For \(|\theta|\le1\), let
\[
        K_b(\theta)=\log \Ee e^{\theta(R_b-m_b)} .
\]
Then \(K_b(0)=K_b'(0)=0\), and
\[
        K_b''(\theta)=\Var_\theta(R_b)
        \le \Ee_\theta R_b^2
        \le r\,\Ee_\theta R_b .
\]
Since \(0\le R_b\le r\), a bounded tilt changes the expectation of the
nonnegative variable \(R_b\) by at most a factor depending only on \(r\):
\[
        \Ee_\theta R_b
        =\frac{\Ee R_be^{\theta R_b}}{\Ee e^{\theta R_b}}
        \le e^{2r}m_b,
        \qquad |\theta|\le1 .
\]
Therefore \(K_b''(\theta)\le C_rm_b\) for \(|\theta|\le1\), and Taylor's
formula gives \(K_b(\theta)\le C_r\theta^2m_b\).  Summing over
\(b\notin A\) proves the first displayed inequality for both signs of
\(\theta\).

For the lower-tail assertion, use the preceding bound and
\(z_j\le C_0L^{-1}\mu_j\):
\[
\begin{aligned}
&-\sum_{j\in B}(-\delta)(z_j-\mu_j)
 +\log \Ee\exp\left\{-\delta\sum_{j\in B}(Z_j(t)-\mu_j)\right\} \\
&\qquad\le
        -\delta\left(1-\frac{C_0}{L}\right)\sum_{j\in B}\mu_j
        +C\delta^2\sum_{j\in B}\mu_j .
\end{aligned}
\]
Decrease \(\delta_0\), if necessary, so that \(C\delta_0\le1/4\), and then
choose \(L_0\) so that \(C_0/L\le1/4\).  The last display is at most
\(-\delta\sum_{j\in B}\mu_j/2\).

For the upper-tail assertion, the same cumulant bound and
\(\mu_j\le C_0L^{-1}z_j\) give
\[
\begin{aligned}
&-\sum_{j\in B}\delta(z_j-\mu_j)
 +\log \Ee\exp\left\{\delta\sum_{j\in B}(Z_j(t)-\mu_j)\right\} \\
&\qquad\le
        -\delta\sum_{j\in B}z_j
        +(\delta+C\delta^2)\sum_{j\in B}\mu_j\\
&\qquad\le
        -\delta\left(1-\frac{2C_0}{L}\right)
        \sum_{j\in B}z_j,
\end{aligned}
\]
after decreasing \(\delta_0\) once more.  Enlarging \(L_0\) makes this at
most \(-\delta\sum_{j\in B}z_j/2\).  The two final scale bounds follow
immediately from the corresponding assumptions.
\end{proof}

\begin{lemma}[Mixed cumulative--interval localization]
\label{lem:projected-cumulative-tail-localization}
Fix $r\ge1$ and $\eta>0$.  There is $L_0=L_0(r,\eta)$ such that the
following holds for every $L\ge L_0$.  In the terminal interval notation, let
$s_j=a_jt_j$ and suppose $s\notin[L^{-1},L]^r$.  Put
\[
        B_-(s)=\{j:s_j<L^{-1}\},\qquad
        B_+(s)=\{j:s_j>L\},\qquad
        G(s)=\{1,\ldots,r\}\setminus(B_-(s)\cup B_+(s)).
\]
On every region on which these three sets are fixed,
\[
\begin{aligned}
        \Pp(Y(t)=y)
        &\le C_L\prod_{j=1}^r a_j^{-1/2}
        \exp\left[-c\sum_{j\in B_-}|H_j(t)|
                   -c(\log L)\sum_{j\in B_+}a_j\right]\\
        &\quad\times
        \exp\left[-c_L\sum_{j\in G(s)}
        \min\left\{\frac{H_j(t)^2}{a_j},a_j\right\}\right].
\end{aligned}\tag{7.3}
\]
Here $c>0$ is independent of $L$, while $C_L<\infty$ and $c_L>0$
depend only on $L,r$, and $\eta$.  Consequently,
\[
\begin{aligned}
        \Pp(Y(t)=y)
        &\le C_L\prod_{j=1}^r a_j^{-1/2}
        \exp\left[-cL\sum_{j\in B_-}a_j
                   -c(\log L)\sum_{j\in B_+}a_j\right]\\
        &\quad\times
        \exp\left[-c_L\sum_{j\in G(s)}
        a_j\min\{(s_j-s_j^*)^2,1\}\right].
\end{aligned}\tag{7.4}
\]
\end{lemma}

\begin{proof}
Throughout the proof the offender sets are fixed.  Write $B_-=B_-(s)$,
$B_+=B_+(s)$, and $G=G(s)$.  Finite-geometric estimates give, uniformly in
the admissible terminal range,
\[
\begin{array}{lll}
 j\in B_-:& \mu_j(t_j)\ge cLa_j,& H_j(t)=z_j-\mu_j(t_j)<0,\\[2mm]
 j\in B_+:& \mu_j(t_j)\le Ca_j/L,& 0<H_j(t)\le z_j\le a_j,
\end{array}
\tag{7.5}
\]
for $L$ sufficiently large.  In particular,
$|H_j(t)|\ge cLa_j$ on $B_-$ and $H_j(t)\ge ca_j$ on $B_+$.
For $j\in G$, one has $s_j\in[L^{-1},L]$, and therefore
\[
        m_j(t)\le\mu_j(t_j)\le C_La_j,
        \qquad |d_j(t)|\le M_La_j.
        \tag{7.6}
\]

Let
\[
        K_t(\theta)=\log\Ee\exp\{\theta\cdot(Y(t)-m(t))\}
\]
be the centered interval-coordinate cumulant.  If
$\gamma=(\gamma_1,\ldots,\gamma_r)$ is a cumulative-coordinate vector, put
\[
        (\mathcal T\gamma)_k=\sum_{j\ge k}\gamma_j.
\]
Since $Z_j=Y_1+\cdots+Y_j$ and $d_j=H_j-H_{j-1}$,
\[
        (\mathcal T\gamma)\cdot d=\gamma\cdot H.
        \tag{7.7}
\]
Choose $\delta>0$ so small that $3\delta$ is admissible in Lemma
\ref{lem:actual-mean-lower-cumulative-chernoff} and
$3\delta r<1/2$.  Define
\[
        \gamma_j^-=-\delta\one_{\{j\in B_-\}},
        \qquad
        \gamma_j^+=\delta(\log L)\one_{\{j\in B_+\}}.
\]

We also introduce a local interval tilt on $G$.  Ratio separation implies that
there is $c_\eta>0$ such that $y_j\ge c_\eta a_j$ for all $j$ and all
sufficiently large $n$.  Fix $0<\varepsilon<c_\eta/4$ and put
\[
        G_0=\{j\in G:|d_j|\le\varepsilon a_j\},
        \qquad G_1=G\setminus G_0.
\]
For a constant $\delta_L>0$ to be chosen below, set
\[
        \lambda_j=
        \begin{cases}
        \delta_L d_j/a_j,&j\in G_0,\\
        \delta_L\operatorname{sgn}(d_j),&j\in G_1,\\
        0,&j\notin G.
        \end{cases}
\]
Finally, let
\[
        \theta=\mathcal T\gamma^-+\mathcal T\gamma^++\lambda.
\]
Generalized H\"older inequality gives the precise combined-cumulant estimate
\[
\begin{aligned}
        K_t(\theta)
        &\le\frac13K_t(3\mathcal T\gamma^-)
             +\frac13K_t(3\mathcal T\gamma^+)
             +\frac13K_t(3\lambda).
\end{aligned}\tag{7.8}
\]
Thus, by (7.7),
\[
\begin{aligned}
-\theta\cdot d+K_t(\theta)
&\le\frac13\{-3\gamma^-\cdot H+K_t(3\mathcal T\gamma^-)\}\\
&\quad+\frac13\{-3\gamma^+\cdot H+K_t(3\mathcal T\gamma^+)\}\\
&\quad+\frac13\{-3\lambda\cdot d+K_t(3\lambda)\}.
\end{aligned}\tag{7.9}
\]
We estimate the three braces separately.

For $B_-$, (7.5) implies $z_j\le C L^{-1}\mu_j(t_j)$.  Lemma
\ref{lem:actual-mean-lower-cumulative-chernoff}, applied with parameter
$3\delta$, gives
\[
        -3\gamma^-\cdot H+K_t(3\mathcal T\gamma^-)
        \le-c\sum_{j\in B_-}\mu_j(t_j)
        \le-c\sum_{j\in B_-}|H_j(t)|.
        \tag{7.10}
\]
The same bound also retains the weaker loss
$-cL\sum_{j\in B_-}a_j$.

For $B_+$, put
\[
        R_b^+=\sum_{j\in B_+}\one_{\{T_b>t_j\}},
        \qquad 0\le R_b^+\le r.
\]
Since the summands indexed by $b\notin A$ are independent,
\[
\begin{aligned}
K_t(3\mathcal T\gamma^+)
&=\sum_{b\notin A}
 \log\Ee\exp\{3\delta(\log L)(R_b^+-\Ee R_b^+)\}\\
&\le \sum_{b\notin A}
 \log\Ee\exp\{3\delta(\log L)R_b^+\}\\
&\le L^{3\delta r}\sum_{b\notin A}\Pp(R_b^+>0)\\
&\le L^{3\delta r}\sum_{j\in B_+}\mu_j(t_j)
 \le C L^{-1/2}\sum_{j\in B_+}a_j.
\end{aligned}
\tag{7.11}
\]
On the other hand, (7.5) gives
\[
        -3\gamma^+\cdot H
        \le-c(\log L)\sum_{j\in B_+}a_j.
\]
After increasing $L_0$, (7.11) is absorbed by this linear gain, and hence
\[
        -3\gamma^+\cdot H+K_t(3\mathcal T\gamma^+)
        \le-c(\log L)\sum_{j\in B_+}a_j.
        \tag{7.12}
\]

For the local tilt, let $v$ vanish outside $G$.  Under every bounded
interval-coordinate tilt $\xi$, a categorical summand contributes at most one
coordinate, and therefore
\[
        \Var_\xi(v\cdot Y(t))
        \le\sum_{j\in G}v_j^2\Ee_\xi Y_j(t)
        \le C_L\sum_{j\in G}a_jv_j^2,
        \tag{7.13}
\]
where the last inequality follows from (7.6) and bounded distortion of atom
probabilities.  Taylor's formula for the cumulant gives
\[
        K_t(3\lambda)\le C_L\sum_{j\in G}a_j\lambda_j^2.
\]
Choose $\delta_L>0$ sufficiently small.  Then
\[
        -3\lambda\cdot d+K_t(3\lambda)
        \le-c_L\sum_{j\in G}
        \min\left\{\frac{d_j^2}{a_j},a_j\right\}.
        \tag{7.14}
\]
Indeed, on $G_0$ the linear term is a fixed multiple of $d_j^2/a_j$,
whereas on $G_1$ it is at least a fixed multiple of $a_j$.

We next recover the product local factor under the single total tilt $\theta$.
For fixed $L$, its interval coordinates are bounded by a constant depending
only on $L$ and $r$.  If $j\in G_0$, then
\[
        m_j(t)\ge y_j-\varepsilon a_j\ge(c_\eta/2)a_j.
\]
Lemma \ref{lem:terminal-interval-edge-blocks}, with
$s_j\in[L^{-1},L]$, supplies a candidate forward-edge block of order $a_j$.
Process the indices of $G_0$ increasingly and choose small subblocks greedily.
Since $a_{i+1}/a_i\ge1+\eta$ and $r$ is fixed,
\[
        \sum_{i<j}a_i\le C_{r,\eta}a_j,
\]
so disjoint subblocks of size $c_{L,r,\eta}a_j$ can be retained for every
$j\in G_0$.  Boundedness of the total tilt preserves the two-atom lower
bounds.  Lemma \ref{lem:projected-forward-edge-concentration}, applied after
projection onto $G_0$, yields
\[
        \Pp_\theta(Y_{G_0}(t)=y_{G_0})
        \le C_L\prod_{j\in G_0}a_j^{-1/2}.
        \tag{7.15}
\]
The exact change-of-measure identity and (7.9)--(7.15) now give
\[
\begin{aligned}
\Pp(Y(t)=y)
&\le C_L\prod_{j\in G_0}a_j^{-1/2}
 \exp\left[-c\sum_{j\in B_-}|H_j|
            -c(\log L)\sum_{j\in B_+}a_j\right]\\
&\quad\times
 \exp\left[-c_L\sum_{j\in G}
 \min\left\{\frac{d_j^2}{a_j},a_j\right\}\right].
\end{aligned}\tag{7.16}
\]
For $j\in G_1$, the last exponent contains $-c_La_j$; for $j\in B_-$ and
$j\in B_+$, (7.5) and the preceding tail losses contain the same fixed-order
exponential cost.  These factors absorb all missing $a_j^{-1/2}$ terms, after
changing constants.

It remains to pass from interval residuals to cumulative residuals without
losing boundary terms.  The assumptions of Lemma
\ref{lem:boundary-truncated-chain} hold with $A_j=a_j$ and $M=M_L$: (7.6)
controls $d_j$ on $G$, while (7.5) gives $|H_j|\le a_j$ on $B_+$.  Split a
fixed fraction of each offender loss in (7.16) for this comparison.  Since
$\log L$ is large, Lemma \ref{lem:boundary-truncated-chain} yields
\[
\begin{aligned}
&c_L\sum_{j\in G}\min\left\{\frac{d_j^2}{a_j},a_j\right\}
 +c\sum_{j\in B_-}|H_j|
 +c(\log L)\sum_{j\in B_+}a_j\\
&\qquad\ge
 c_L'\sum_{j\in G}\min\left\{\frac{H_j^2}{a_j},a_j\right\}
 +c'\sum_{j\in B_-}|H_j|
 +c'(\log L)\sum_{j\in B_+}a_j.
\end{aligned}
\tag{7.17}
\]
Combining (7.16) and (7.17) proves (7.3).

Finally, on $G$ the scaled coordinate $s_j$ belongs to $[L^{-1},L]$.
Uniform finite-geometric derivative estimates and strict monotonicity give
\[
        |H_j(t)|=|\mu_j(t_j^*)-\mu_j(t_j)|
        \ge c_La_j\min\{|s_j-s_j^*|,1\}.
        \tag{7.18}
\]
Together with (7.5), this converts (7.3) into (7.4).
\end{proof}


\begin{lemma}[Anisotropic tilted local upper bound]\label{lem:anisotropic-tilted-upper}
Let \(W=\sum_{\nu\in\mathcal I}W_\nu\) be as in Lemma \ref{lem:coord-fourier}.  Assume that the hypotheses of Lemma \ref{lem:coord-fourier} hold uniformly under all tilts \(\|\lambda\|_\infty\le\Lambda\).  Assume also that, for those tilted laws,
\[
        \Var_\lambda(v\cdot W)
        \le C_0\sum_{j=1}^r A_jv_j^2,
        \qquad v\in\Rr^r .
\]
Then there are constants \(C,c>0\), depending only on the constants in the assumptions, such that for every lattice point \(z\), with
\(h=z-\Ee W\),
\[
        \Pp(W=z)
        \le C\prod_{j=1}^r A_j^{-1/2}
        \exp\left[-c\sum_{j=1}^r
        \min\left\{\frac{h_j^2}{A_j}, |h_j|\right\}\right].
\]
In particular, if \(|h_j|\le A_j\) for every \(j\), then the exponent may be written as
\(-c\sum_j h_j^2/A_j\).
\end{lemma}

\begin{proof}
Let
\[
        K(\lambda)=\log\Ee e^{\lambda\cdot(W-\Ee W)}.
\]
For \(\|\lambda\|_\infty\le\Lambda\), the covariance upper bound gives
\[
        K(\lambda)
        =\int_0^1(1-s)\lambda^T
        \Var_{s\lambda}(W)\lambda\dd s
        \le C\sum_{j=1}^r A_j\lambda_j^2 .
\]
Choose a sufficiently small \(\delta>0\), depending only on the constants in the assumptions and with \(\delta\le\Lambda\), and put
\[
        \lambda_j=\delta\,\mathrm{sgn}(h_j)
        \min\{|h_j|/A_j,1\}.
\]
By tilting and then applying Lemma \ref{lem:coord-fourier} under the tilted law,
\[
\begin{aligned}
        \Pp(W=z)
        &=\exp\{-\lambda\cdot h+K(\lambda)\}
          \Pp_\lambda(W=z) \\
        &\le C\prod_{j=1}^r A_j^{-1/2}
        \exp\{-\lambda\cdot h+K(\lambda)\}.
\end{aligned}
\]
For the chosen \(\lambda\),
\[
        \lambda\cdot h
        =\delta\sum_{j=1}^r |h_j|\min\{|h_j|/A_j,1\},
\]
while
\[
        K(\lambda)
        \le C\delta^2\sum_{j=1}^r
        A_j\min\{h_j^2/A_j^2,1\}.
\]
Taking \(\delta\) small enough makes the second display at most one half of the first, after changing constants.  This proves the claim.
\end{proof}

\begin{lemma}[Anisotropic terminal local CLT]\label{lem:terminal-llt}
Fix \(r\ge1\), \(\eta>0\), \(\kappa>0\), and \(0<c_0<C_0<\infty\).  Let
\[
        1\ll a_1<\cdots<a_r\ll n,
        \qquad a_{j+1}/a_j\ge1+\eta,
\]
where the ratio-separation condition is empty when \(r=1\).  Let
\(t_1>\cdots>t_r\) satisfy
\[
        \frac{c_0}{a_j}\le t_j\le \frac{C_0}{a_j},
        \qquad
        \frac{t_j}{t_{j+1}}\ge1+\kappa
        \quad (1\le j<r).
\]
Let \(A=\{a_1,\ldots,a_r\}\), and put
\[
        Z_j(t)=\sum_{b\notin A}\one_{\{T_b>t_j\}},
        \qquad Z(t)=(Z_1(t),\ldots,Z_r(t)).
\]
If \(\Sigma(t)=\Var Z(t)\), then there are constants \(c,C>0\), depending only on \(r,\eta,\kappa,c_0,C_0\), such that
\[
        c\,\diag(a_1,\ldots,a_r)
        \le \Sigma(t)\le
        C\,\diag(a_1,\ldots,a_r)
\]
in quadratic-form order.  Moreover,
\[
\begin{aligned}
        \Pp(Z(t)=z)
        &=(2\pi)^{-r/2}(\det\Sigma(t))^{-1/2} \\
        &\quad\times
        \exp\left[-\frac12(z-\Ee Z(t))^T\Sigma(t)^{-1}(z-\Ee Z(t))\right]
        +o\left((a_1\cdots a_r)^{-1/2}\right)
\end{aligned}
\]
uniformly whenever the Mahalanobis norm of \(z-\Ee Z(t)\) is bounded.  In addition, for every fixed \(\gamma\in(0,1/6)\), the coordinatewise moderate-deviation upper bound
\[
        \Pp(Z(t)=z)
        \le C(\det\Sigma(t))^{-1/2}
        \exp\left[-c\sum_{j=1}^r
        \frac{(z_j-\Ee Z_j(t))^2}{a_j}\right]
\]
holds uniformly whenever
\[
        |z_j-\Ee Z_j(t)|\le a_j^{1/2+\gamma},
        \qquad 1\le j\le r .
\]
\end{lemma}

\begin{proof}
For a single rate \(b\notin A\), write
\[
        V_b(t)=(\one_{\{T_b>t_1\}},\ldots,\one_{\{T_b>t_r\}}).
\]
The vector \(V_b(t)\) takes the cumulative values
\[
        (1,\ldots,1),\quad (0,1,\ldots,1),\quad\ldots,
        \quad (0,\ldots,0,1),\quad (0,\ldots,0),
\]
according as \(T_b\) falls in
\[
        (t_1,\infty),\quad (t_2,t_1],\quad\ldots,
        \quad(t_r,t_{r-1}],\quad (0,t_r].
\]
Adjacent support points differ by the coordinate vectors, so the lattice span is \(\mathbb Z^r\).

Choose constants \(0<d_1<d_2<\infty\), depending only on
\(r,\eta,\kappa,c_0,C_0\), with \(d_2/d_1<1+\eta\).  Then the rate intervals
\([d_1a_j,d_2a_j]\) are disjoint for different \(j\)'s.  For every integer rate \(b\in[d_1a_j,d_2a_j]\) and all large \(n\), the two adjacent time categories separated by the threshold \(t_j\) have probabilities bounded below.  For instance,
\[
        e^{-bt_j}\ge e^{-d_2C_0},
        \qquad
        b(t_{j-1}-t_j)\ge d_1c_0\frac{\kappa}{1+\kappa}
\]
for the upper adjacent interval, with the evident endpoint modifications when \(j=1\) or \(j=r\).  The fixed omitted set \(A\) removes only \(r\) rates, so each block still contains order \(a_j\) available rates.

These disjoint adjacent blocks imply the covariance lower bound: for every \(u\in\Rr^r\), the \(j\)-th block contributes at least \(c a_j u_j^2\) to \(\Var(u\cdot Z(t))\).  Thus
\[
        \Var(u\cdot Z(t))\ge c\sum_{j=1}^r a_j u_j^2 .
\]
The upper bound follows from \(\Var Z_j(t)=O(a_j)\), which is a direct finite-geometric estimate, and Cauchy's inequality:
\[
        \Var(u\cdot Z(t))
        \le r\sum_{j=1}^r u_j^2\Var Z_j(t)
        \le C\sum_{j=1}^r a_j u_j^2 .
\]
This proves the covariance comparison.

The same blocks give the coordinatewise Fourier bound of Lemma \ref{lem:coord-fourier}, with \(A_j=a_j\).  This is the point at which anisotropy is used: rather than bounding the whole off-origin region by a scalar factor \(e^{-ca_1}\), we retain the product-scale estimate
\[
        |\Ee e^{iu\cdot Z(t)}|
        \le \exp\left[-c\sum_{j=1}^r
        a_j\operatorname{dist}(u_j,2\pi\mathbb Z)^2\right].
\]
This bound is stable under bounded exponential tilts, because the support is finite and the two distinguished adjacent atom probabilities remain bounded below.

We next prove the local CLT.  Fourier inversion gives
\[
        \Pp(Z(t)=z)
        =(2\pi)^{-r}\int_{[-\pi,\pi]^r}
        e^{-iu\cdot(z-\Ee Z(t))}\phi_t(u)\dd u,
\]
where \(\phi_t(u)=\Ee e^{iu\cdot(Z(t)-\Ee Z(t))}\).  Put
\(D_a=\diag(a_1,\ldots,a_r)\).  On the anisotropic small box
\(|u_j|\le R a_j^{-1/2}\), the finite-tail estimates and boundedness of the increments give the uniform expansion
\[
        \log\phi_t(u)
        =-\frac12u^T\Sigma(t)u
        +O\left((\max_j |u_j|)\sum_{j=1}^r a_j u_j^2\right).
\]
Indeed, the third cumulant in the direction \(u\) is bounded by a constant times \(\max_j|u_j|\) times the corresponding variance contribution, and the variance upper bound just proved controls the latter by \(C\sum_j a_j u_j^2\).  After the change of variables \(u=D_a^{-1/2}v\), the error term is
\(O_R(a_1^{-1/2})=o(1)\).  Since the Mahalanobis norm of
\(z-\Ee Z(t)\) is bounded, this gives the Gaussian integral on every fixed \(v\)-ball.  The contribution from the complement of such a ball is controlled uniformly by the coordinatewise Fourier bound above:
\[
        \int_{\max_j |u_j|>R a_j^{-1/2}}
        \exp\left[-c\sum_j a_j\operatorname{dist}(u_j,2\pi\mathbb Z)^2\right]\dd u
        \le C\left(\prod_j a_j^{-1/2}\right)e^{-cR^2}.
\]
Letting \(R\to\infty\) yields the stated local CLT and the error
\(o((a_1\cdots a_r)^{-1/2})\), uniformly over the allowed parameters.

For the moderate-deviation upper bound, apply Lemma \ref{lem:anisotropic-tilted-upper} to the terminal array.  The covariance upper bound required there has already been proved and remains valid under the small tilts used below.  If
\(|z_j-\Ee Z_j(t)|\le a_j^{1/2+\gamma}\), then for large \(n\) this displacement is at most \(a_j\), and Lemma \ref{lem:anisotropic-tilted-upper} gives
\[
        \Pp(Z(t)=z)
        \le C\prod_{j=1}^r a_j^{-1/2}
        \exp\left[-c\sum_{j=1}^r
        \frac{(z_j-\Ee Z_j(t))^2}{a_j}\right].
\]
The covariance comparison implies
\(\prod_j a_j^{-1/2}\asymp(\det\Sigma(t))^{-1/2}\), proving the displayed form.
\end{proof}


\begin{lemma}[Scaled interval-count localization]
\label{lem:terminal-interval-tilt}
Fix $r\ge1$, $\eta>0$, and a compact set $K\subset(0,\infty)^r$.  In the
terminal interval notation, define
\[
        \Phi(t)=\prod_{j=1}^r a_j e^{-a_jt_j}.
\]
Uniformly for $s=(s_1,\ldots,s_r)\in K$ with $t_1>\cdots>t_r$,
\[
        \Pp(Y(t)=y)
        \le C\left(\prod_{j=1}^r a_j^{-1/2}\right)
        \exp\left[-c\sum_{j=1}^r
        a_j\min\{(s_j-s_j^*)^2,1\}\right],
        \tag{7.19}
\]
where $C,c>0$ depend only on $r,\eta$, and $K$.

For a precise uniform tail statement, let $\mathfrak A_N(\eta)$ be the set of
integer tuples $(n,a_1,\ldots,a_r)$ satisfying
\[
        n\ge N,\qquad
        N\le a_1<\cdots<a_r\le n/N,
        \qquad a_{j+1}/a_j\ge1+\eta.
\]
Then
\[
\begin{aligned}
\lim_{L\to\infty}\limsup_{N\to\infty}
\sup_{(n,a)\in\mathfrak A_N(\eta)}
&\frac{1}{\prod_{j=1}^r a_j^{-1}}
\int_{\substack{t_1>\cdots>t_r\\
        (a_1t_1,\ldots,a_rt_r)\notin[L^{-1},L]^r}}
        \Phi(t)\Pp(Z(t)=z)\dd t
=0.
\end{aligned}\tag{7.20}
\]
In particular, for every $\varepsilon>0$ there are $L<\infty$ and $N_0$ such
that the integral in (7.20) is at most
$\varepsilon\prod_j a_j^{-1}$ for every
$(n,a)\in\mathfrak A_{N_0}(\eta)$.
\end{lemma}

\begin{proof}
Ratio separation gives
\[
        c_\eta a_j\le y_j\le C_\eta a_j,
        \qquad 1\le j\le r,
\]
for all sufficiently large $a_1$.  On the fixed compact scaled chamber $K$,
finite-geometric estimates give $m_j(t)=O_K(a_j)$ and hence
$|d_j(t)|\le M_Ka_j$.  Lemma
\ref{lem:projected-good-bad-interval-bound} therefore yields
\[
        \Pp(Y(t)=y)
        \le C\prod_{j=1}^r a_j^{-1/2}
        \exp\left[-c\sum_{j=1}^r
        \min\left\{\frac{d_j(t)^2}{a_j},a_j\right\}\right].
\]
Since $d_j=H_j-H_{j-1}$ and $H_0=0$, Lemma
\ref{lem:truncated-triangular-comparison} gives
\[
        \sum_{j=1}^r
        \min\left\{\frac{d_j(t)^2}{a_j},a_j\right\}
        \ge c_K\sum_{j=1}^r
        \min\left\{\frac{H_j(t)^2}{a_j},a_j\right\}.
\]
On $K$, the derivative of
$s\mapsto\mu_j(s/a_j)$ has magnitude comparable to $a_j$, uniformly in $j$.
The mean-value theorem near the saddle, followed by compactness and strict
monotonicity away from it, gives
\[
        |H_j(t)|
        \ge c_Ka_j\min\{|s_j-s_j^*|,1\}.
\]
This proves (7.19).

We turn to (7.20).  Fix $L\ge L_0(r,\eta)$ and partition the complement of
$[L^{-1},L]^r$ according to the finitely many possible pairs of offender sets
$(B_-,B_+)$.  On each such region, Lemma
\ref{lem:projected-cumulative-tail-localization} gives (7.4).  Since
\[
        \Phi(t)\dd t_1\cdots\dd t_r
        =\prod_{j=1}^r e^{-s_j}\dd s_j,
\]
and discarding the chamber constraint only enlarges the integration domain,
the resulting bound factors coordinatewise.  Every nonoffending coordinate
contributes at most
\[
        C_L\int_{L^{-1}}^L
        e^{-c_La_j\min\{(s-s_j^*)^2,1\}}\dd s
        \le C_La_j^{-1/2}.
        \tag{7.21}
\]
For a lower offender, the integration interval has length at most $L^{-1}$,
and
\[
        a_j^{1/2}e^{-cLa_j}
        \le \alpha_{L,N}^-
        :=\sup_{a\ge N}a^{1/2}e^{-cLa}
        \longrightarrow0
        \qquad(N\to\infty).
        \tag{7.22}
\]
For an upper offender,
\[
        \int_L^\infty e^{-s}\dd s=e^{-L},
\]
and
\[
        a_j^{1/2}L^{-ca_j}
        \le \alpha_{L,N}^+
        :=\sup_{a\ge N}a^{1/2}L^{-ca}
        \longrightarrow0
        \qquad(N\to\infty).
        \tag{7.23}
\]
The original factor $\prod_j a_j^{-1/2}$, together with (7.21)--(7.23),
therefore supplies a factor $a_j^{-1}$ for every coordinate and an additional
quantity tending to zero for every nonempty offender set.  The number of
possible offender sets is finite.  Hence, for each fixed $L\ge L_0$,
\[
\limsup_{N\to\infty}
\sup_{(n,a)\in\mathfrak A_N(\eta)}
\frac{
\displaystyle
\int_{\substack{t_1>\cdots>t_r\\s\notin[L^{-1},L]^r}}
\Phi(t)\Pp(Z(t)=z)\dd t}
{\prod_j a_j^{-1}}=0.
\]
This implies (7.20) and the final assertion.
\end{proof}


\begin{lemma}[Product-scale terminal localization]\label{lem:terminal-product-localization}
Fix \(r\ge1\) and \(\eta>0\).  In the setting of Lemma \ref{lem:suk-joint-saddle}, put
\[
        z=(a_1-1,a_2-2,\ldots,a_r-r),
        \qquad
        \Phi(t)=\prod_{j=1}^r a_j e^{-a_jt_j},
\]
and let \(Z(t)\) have the same definition as in Lemma \ref{lem:terminal-llt}.  There is \(\delta_0>0\), depending only on \(r\) and \(\eta\), such that for every fixed \(0<\delta<\delta_0\) there are constants \(C,c>0\), depending only on \(r,\eta,\delta\), for which the following holds.  If
\[
        \mathcal N_\delta
        =\left\{t_1>\cdots>t_r:
        |t_j/t_j^*-1|\le\delta\text{ for all }j\right\},
\]
and
\[
        u_j=a_j^{3/2}(t_j-t_j^*),
        \qquad 1\le j\le r,
\]
then, uniformly for \(t\in\mathcal N_\delta\),
\[
        \Phi(t)\Pp(Z(t)=z)\dd t_1\cdots\dd t_r
        \le C\left(\prod_{j=1}^r a_j^{-1}\right)
        \exp\left[-c\sum_{j=1}^r\min(u_j^2,a_j)\right]
        \dd u_1\cdots\dd u_r .
\]
Moreover,
\[
\begin{aligned}
&\int_{\{t_1>\cdots>t_r\}\setminus\mathcal N_\delta}
        \Phi(t)\Pp(Z(t)=z)\dd t_1\cdots\dd t_r  \\
&\qquad =o\left(\prod_{j=1}^r a_j^{-1}\right),
\end{aligned}
\]
uniformly over the log-separated terminal range.
\end{lemma}

\begin{proof}
Choose \(\delta_0\) small enough that, whenever
\(|t_j/t_j^*-1|\le\delta_0\) for all \(j\), the ratios
\(t_j/t_{j+1}\) are bounded below by \(1+\kappa\), for some
\(\kappa>0\) depending only on \(\eta\).  On this neighborhood, Lemma \ref{lem:suk-joint-saddle} gives
\[
        t_j\asymp a_j^{-1},\qquad
        -\mu_j'(t_j)\asymp a_j^2,
        \qquad
        \mu_j''(t_j)=O(a_j^3),
\]
uniformly.

For \(t\in\mathcal N_\delta\), put \(s_j=a_jt_j\).  Since
\[
        s_j-s_j^*=a_j(t_j-t_j^*)=u_j/\sqrt{a_j},
\]
Lemma \ref{lem:terminal-interval-tilt}, applied on the compact scaled neighborhood determined by \(\delta\), gives
\[
        \Pp(Z(t)=z)
        \le C\left(\prod_{j=1}^r a_j^{-1/2}\right)
        \exp\left[-c\sum_{j=1}^r\min(u_j^2,a_j)\right].
\]
Also
\[
        \Phi(t)\dd t_1\cdots\dd t_r
        =\left(\prod_{j=1}^r a_j e^{-s_j}\right)
        \left(\prod_{j=1}^r a_j^{-3/2}\right)
        \dd u_1\cdots\dd u_r
        \le C\prod_{j=1}^r a_j^{-1/2}\dd u_1\cdots\dd u_r,
\]
because \(s_j\) remains in a compact subset of \((0,\infty)\).  Multiplying the two displays proves the product-scale bound.

It remains to remove the complement of \(\mathcal N_\delta\).  In the scaled variables, Lemma \ref{lem:terminal-interval-tilt} first removes the complement of a fixed compact set at cost
\(o(\prod_j a_j^{-1})\).  On the remaining compact set, the same lemma gives
\[
        \Pp(Z(t)=z)
        \le C\left(\prod_{j=1}^r a_j^{-1/2}\right)
        \exp\left[-c\sum_{j=1}^r
        a_j\min\{(s_j-s_j^*)^2,1\}\right].
\]
If \(t\notin\mathcal N_\delta\), then, on this compact set and for all large \(n\), at least one scaled coordinate satisfies
\(|s_j-s_j^*|\ge c_\delta\).  Since
\[
        \Phi(t)\dd t_1\cdots\dd t_r
        =\prod_{j=1}^r e^{-s_j}\dd s_j,
\]
the compact contribution outside \(\mathcal N_\delta\) is bounded by
\[
        C\left(\prod_{j=1}^r a_j^{-1/2}\right)
        \int_{\max_j|s_j-s_j^*|\ge c_\delta}
        \exp\left[-c\sum_{j=1}^r
        a_j\min\{(s_j-s_j^*)^2,1\}\right]\dd s .
\]
Changing variables \(v_j=\sqrt{a_j}(s_j-s_j^*)\) gives
\[
        O\left(e^{-c_\delta a_1}\prod_{j=1}^r a_j^{-1}\right)
        =o\left(\prod_{j=1}^r a_j^{-1}\right).
\]
Together with the compact-tail estimate, this proves the second assertion.
\end{proof}

\begin{lemma}[Covariance stability on the terminal saddle scale]
\label{lem:terminal-covariance-stability}
Fix \(r\ge1\), \(\eta>0\), and \(R<\infty\).  In the setting of Lemma
\ref{lem:suk-joint-saddle}, let
\[
        \mathsf D_a=\diag(a_1,\ldots,a_r),
        \qquad
        \mathsf H_a=\diag(a_1^{-3/2},\ldots,a_r^{-3/2}),
\]
and let \(\Sigma(t)=\Var Z(t)\).  Then, uniformly over the log-separated
terminal range,
\[
\sup_{\|u\|\le R}
\left\|
\mathsf D_a^{-1/2}
\bigl(\Sigma(t^*+\mathsf H_au)-\Sigma(t^*)\bigr)
\mathsf D_a^{-1/2}
\right\|_{\mathrm{op}}
\longrightarrow0.
\tag{7.24}
\]
Consequently, uniformly on the same fixed \(u\)-ball,
\[
        \frac{\det\Sigma(t^*+\mathsf H_au)}{\det\Sigma(t^*)}
        \longrightarrow1
\tag{7.25}
\]
and
\[
\left\|
\mathsf D_a^{1/2}
\bigl(\Sigma(t^*+\mathsf H_au)^{-1}-\Sigma(t^*)^{-1}\bigr)
\mathsf D_a^{1/2}
\right\|_{\mathrm{op}}
\longrightarrow0.
\tag{7.26}
\]
\end{lemma}

\begin{proof}
For \(j\le k\), independence of the nonspecified clocks gives
\[
        \Sigma_{jk}(t)
        =\sum_{b\notin A}e^{-bt_j}(1-e^{-bt_k}).
        \tag{7.27}
\]
The same formula with \(j=k\) is the Bernoulli variance.  On a fixed saddle
ball \(t=t^*+\mathsf H_au\), Lemma \ref{lem:suk-joint-saddle} gives
\(t_j\asymp a_j^{-1}\), uniformly.  Differentiating (7.27) and using the
finite-geometric estimates yields
\[
        |\partial_{t_j}\Sigma_{jk}(t)|
        \le\sum_{b\ge1}b e^{-bt_j}\le C_Ra_j^2,
\]
while
\[
        |\partial_{t_k}\Sigma_{jk}(t)|
        \le\sum_{b\ge1}b e^{-b(t_j+t_k)}\le C_Ra_j^2.
\]
For \(j=k\), the derivative is bounded by the same geometric sum.  Hence
\[
\begin{aligned}
|\Sigma_{jk}(t^*+\mathsf H_au)-\Sigma_{jk}(t^*)|
&\le C_Ra_j^2(a_j^{-3/2}+a_k^{-3/2}).
\end{aligned}
\]
After division by \(\sqrt{a_ja_k}\), the right-hand side is at most
\[
        C_R\left(a_k^{-1/2}+\frac{a_j^{3/2}}{a_k^2}\right)
        \le C_Ra_1^{-1/2}.
\]
Since \(r\) is fixed, entrywise convergence implies (7.24).

Lemma \ref{lem:terminal-llt} gives uniform positive lower and upper bounds for
the eigenvalues of the normalized matrices
\[
        \mathsf D_a^{-1/2}\Sigma(t)\mathsf D_a^{-1/2}
\]
throughout the saddle neighborhood.  The determinant and inverse are
continuous, uniformly Lipschitz functions on this compact set of positive
definite matrices.  Applying these facts to (7.24) proves (7.25) and (7.26).
\end{proof}

\begin{lemma}[Terminal integrated saddle]\label{lem:terminal-integrated-saddle}
Fix \(r\ge1\) and \(\eta>0\).  In the setting of Lemma \ref{lem:suk-joint-saddle}, put
\[
        z=(a_1-1,a_2-2,\ldots,a_r-r),
        \qquad
        \Phi(t)=\prod_{j=1}^r a_j e^{-a_jt_j}.
\]
Let \(A=\{a_1,\ldots,a_r\}\), and put
\[
        Z_j(t)=\sum_{b\notin A}\one_{\{T_b>t_j\}},
        \qquad Z(t)=(Z_1(t),\ldots,Z_r(t)).
\]
Then
\[
\begin{aligned}
&\int_{t_1>\cdots>t_r}
        \Phi(t)\Pp(Z(t)=z)\dd t_1\cdots\dd t_r  \\
&\qquad =
        (1+o(1))\Phi(t^*)
        \prod_{j=1}^r\bigl(-\mu_j'(t_j^*)\bigr)^{-1},
\end{aligned}
\]
uniformly over the log-separated terminal range.
\end{lemma}

\begin{proof}
Let \(\mathcal N_\delta\) be the product saddle neighborhood from Lemma \ref{lem:terminal-product-localization}, with \(\delta>0\) fixed and sufficiently small.  That lemma gives
\[
        \int_{\{t_1>\cdots>t_r\}\setminus\mathcal N_\delta}
        \Phi(t)\Pp(Z(t)=z)\dd t_1\cdots\dd t_r
        =o\left(\prod_{j=1}^r a_j^{-1}\right).
\]
Since Lemma \ref{lem:suk-joint-saddle} gives
\[
        \Phi(t^*)\prod_{j=1}^r(-\mu_j'(t_j^*))^{-1}
        \asymp \prod_{j=1}^r a_j^{-1},
\]
the contribution outside \(\mathcal N_\delta\) is negligible on the required relative scale.

Inside \(\mathcal N_\delta\), put
\[
        H=\diag(a_1^{-3/2},\ldots,a_r^{-3/2}),
        \qquad
        t=t^*+Hu .
\]
This scale is forced by the balance
\(-\mu_j'(t_j^*)\asymp a_j^2\) and \(\Var Z_j(t^*)\asymp a_j\): changing \(t_j\) by \(a_j^{-3/2}\) moves the mean by order \(a_j^{1/2}\), the natural standard-deviation scale.  Let
\[
        D=\diag\bigl(-\mu_1'(t_1^*),\ldots,-\mu_r'(t_r^*)\bigr).
\]
For bounded \(u\), Lemma \ref{lem:suk-joint-saddle} and Taylor's theorem give the normalized mean expansion
\[
        \left\|
        \mathsf D_a^{-1/2}
        \bigl(\Ee Z(t)-z+DHu\bigr)
        \right\|\longrightarrow0,
        \qquad \mathsf D_a=\diag(a_1,\ldots,a_r),
        \tag{7.28}
\]
uniformly on every fixed \(u\)-ball.  Indeed, the \(j\)-th Taylor remainder is
\[
        O(a_j^3|t_j-t_j^*|^2)=O(u_j^2)=o(a_j^{1/2}),
\]
whereas
\[
        \mu_j'(t_j^*)(t_j-t_j^*)=-u_j a_j^{1/2}(1+o(1)).
\]
The neighborhood \(\mathcal N_\delta\) supplies constants \(\kappa,c_0,C_0\) for the hypotheses of Lemma \ref{lem:terminal-llt}; hence the anisotropic local CLT applies uniformly for bounded \(u\).  Lemma
\ref{lem:terminal-covariance-stability}, together with (7.28), gives
\[
\begin{aligned}
        (\det\Sigma(t))^{-1/2}
        &=(\det\Sigma(t^*))^{-1/2}(1+o(1)),\\
        (z-\Ee Z(t))^T\Sigma(t)^{-1}(z-\Ee Z(t))
        &=u^TH^TD^T\Sigma(t^*)^{-1}DHu+o(1),
\end{aligned}
        \tag{7.29}
\]
uniformly on every fixed ball.  Also,
\[
        \Phi(t^*+Hu)=\Phi(t^*)(1+o(1))
\]
uniformly there, because \(a_j(t_j-t_j^*)=u_j/\sqrt{a_j}=o(1)\).
Consequently, on \(\|u\|\le R\), the integrand has the uniform asymptotic
Gaussian form
\[
        \Phi(t^*)(2\pi)^{-r/2}(\det\Sigma(t^*))^{-1/2}
        \exp\left[-\frac12 u^T H^TD^T\Sigma(t^*)^{-1}DHu\right]
        |\det H|\,\dd u .
\]
The matrix \(H^TD^T\Sigma(t^*)^{-1}DH\) is uniformly positive definite and uniformly bounded in quadratic-form order, because \(D_{jj}\asymp a_j^2\), \(H_{jj}=a_j^{-3/2}\), and \(\Sigma(t^*)\asymp\diag(a_1,\ldots,a_r)\).  Hence, as \(n\to\infty\) with \(R\) fixed, the contribution from \(\|u\|\le R\) is
\[
\begin{aligned}
        &\Phi(t^*)(2\pi)^{-r/2}(\det\Sigma(t^*))^{-1/2}|\det H| \\
        &\quad\times
        \int_{\|u\|\le R}
        \exp\left[-\frac12 u^TH^TD^T\Sigma(t^*)^{-1}DHu\right]
        \dd u
        +o_R\left(\prod_{j=1}^r a_j^{-1}\right).
\end{aligned}
\]
Letting \(R\to\infty\), uniformly over the allowed rate tuples, the Gaussian integral gives the determinant cancellation
\[
\begin{aligned}
&(2\pi)^{-r/2}(\det\Sigma(t^*))^{-1/2}|\det H|
 \int_{\mathbb R^r}
        \exp\left[-\frac12 u^TH^TD^T\Sigma(t^*)^{-1}DHu\right]\dd u  \\
&\qquad = |\det D|^{-1}.
\end{aligned}
\]
Indeed,
\[
        \det(H^TD^T\Sigma(t^*)^{-1}DH)
        =|\det H|^2|\det D|^2(\det\Sigma(t^*))^{-1}.
\]
Thus the fixed-ball saddle contribution tends to
\[
        \Phi(t^*)\prod_{j=1}^r(-\mu_j'(t_j^*))^{-1}.
\]

It remains only to justify the passage from fixed \(R\) to the whole neighborhood.  This is precisely the product-scale bound from Lemma \ref{lem:terminal-product-localization}: on \(\mathcal N_\delta\),
\[
        \Phi(t)\Pp(Z(t)=z)\dd t
        \le C\left(\prod_{j=1}^r a_j^{-1}\right)
        \exp\left[-c\sum_{j=1}^r\min(u_j^2,a_j)\right]\dd u .
\]
The functions on the right have uniformly integrable tails, since
\[
        \sup_{a\ge1}\int_{\mathbb R}e^{-c\min(u^2,a)}
        \one_{\{|u|>R\}}\one_{\{|u|\le C\sqrt a\}}\dd u\to0
        \qquad (R\to\infty).
\]
Consequently the contribution from \(\mathcal N_\delta\cap\{\|u\|>R\}\) is
\[
        o_R\left(\prod_{j=1}^r a_j^{-1}\right)
\]
uniformly as \(R\to\infty\).  Combining this tail estimate with the negligible contribution outside \(\mathcal N_\delta\) proves the lemma.
\end{proof}

\begin{proposition}[Separated terminal joint estimate]\label{prop:suk-joint}
Fix \(r\ge1\) and \(\eta>0\).  Uniformly over integers
\[
        1\ll a_1<\cdots<a_r\ll n
\]
satisfying
\[
        a_{j+1}/a_j\ge 1+\eta \quad (1\le j<r),
\]
one has
\[
        \Pp(I_{n,a_1}=\cdots=I_{n,a_r}=1)
        \sim \frac{e^{-r}}{a_1\cdots a_r}.
\]
\end{proposition}

\begin{proof}
Let \(A=\{a_1,\ldots,a_r\}\).  On the event that all \(I_{n,a_j}=1\) occur, the time ordering is forced.  Rate \(a\) is deleted at the transition at which the number of alive particles drops from \(a\) to \(a-1\).  Since \(a_1<\cdots<a_r\), the deletion at level \(a_r\) occurs before the deletion at level \(a_{r-1}\), and so on.  Thus necessarily
\[
        T_{a_1}>T_{a_2}>\cdots>T_{a_r}.
\]
Condition on \(T_{a_j}=t_j\) in this chamber.  For \(b\notin A\), define
\[
        Z_j(t)=\sum_{b\notin A}\one_{\{T_b>t_j\}},
        \qquad 1\le j\le r.
\]
At time \(t_j\), precisely \(j-1\) of the specified clocks have larger times, namely \(a_1,\ldots,a_{j-1}\).  Therefore the joint event is equivalent to
\[
        Z_j(t)=a_j-j,
        \qquad 1\le j\le r.
\]
Conversely, in the chamber \(t_1>\cdots>t_r\), these equations imply that each specified rate \(a_j\) has exactly \(a_j-1\) clocks with larger time, and hence all the events \(I_{n,a_j}=1\) occur.  Consequently the joint probability is exactly
\[
\begin{aligned}
\int_{t_1>\cdots>t_r}
        \left(\prod_{j=1}^r a_j e^{-a_j t_j}\right)
        \Pp(Z_j(t)=a_j-j,
        1\le j\le r)
        \dd t_1\cdots\dd t_r .
\end{aligned}
\]
By Lemma \ref{lem:terminal-integrated-saddle}, this integral equals
\[
        (1+o(1))
        \prod_{j=1}^r
        \frac{a_j e^{-a_jt_j^*}}{-\mu_j'(t_j^*)}.
\]
Lemma \ref{lem:suk-joint-saddle} gives \(a_jt_j^*\to1\) and \(-\mu_j'(t_j^*)\sim a_j^2\).  Hence
\[
        \prod_{j=1}^r
        \frac{a_j e^{-a_jt_j^*}}{-\mu_j'(t_j^*)}
        \sim
        \prod_{j=1}^r\frac{a_je^{-1}}{a_j^2}
        =\frac{e^{-r}}{a_1\cdots a_r}.
\]
\end{proof}

\subsection{Terminal point process and full factorial moments}

\begin{theorem}[Local terminal Poisson point process]\label{thm:suk-ppp}
Let \(b_n\to\infty\) and \(b_n=o(n)\), and define
\[
        \Xi_n=\sum_{a=1}^n I_{n,a}\delta_{\log(a/b_n)}.
\]
For every bounded interval \(J\subset\mathbb R\),
\[
        \Xi_n|_J\Rightarrow \PPP(e^{-1}\dd x)|_J
\]
in the space of locally finite counting measures on \(J\), equipped with the
vague topology.  Equivalently, counts in finitely many disjoint bounded
intervals converge jointly to independent Poisson random variables with means
\(e^{-1}\) times their lengths.
\end{theorem}

\begin{proof}
It suffices to prove convergence of mixed factorial moments.  If \(\log(a/b_n)\in J\), then \(a\to\infty\), \(a=o(n)\), and all such \(a\)'s are within a bounded multiplicative factor.  Also
\[
        \sum_{\log(a/b_n)\in J}\frac1a\to |J|.
\]
For fixed \(r\), expand \(\Ee(\Xi_n(J))_r\) over ordered distinct \(r\)-tuples.  First restrict to tuples satisfying \(|\log a_i-\log a_j|>\eta\) for all \(i\ne j\).  After sorting such a tuple increasingly, the separation condition implies
\[
        \frac{a_{j+1}}{a_j}\ge e^\eta=1+(e^\eta-1),
\]
so Proposition \ref{prop:suk-joint} applies uniformly, with separation parameter \(e^\eta-1\), and gives
\[
        e^{-r}\sum_{\eta\text{-sep}}^*\frac1{a_1\cdots a_r}+o(1).
\]
As \(n\to\infty\), this tends to \(e^{-r}\) times the Lebesgue measure of the \(\eta\)-separated part of \(J^r\).  Letting \(\eta\downarrow0\) gives \((e^{-1}|J|)^r\).

The excluded tuples are controlled by Lemma \ref{lem:death-bound}.  Indeed,
\[
\begin{aligned}
&\sum_{\substack{\log(a_i/b_n)\in J\,\,1\le i\le r\\
        |\log a_i-\log a_j|\le\eta\text{ for some }i\ne j}}
        \prod_{i=1}^r\frac1{a_i} \\
&\qquad\longrightarrow
        \operatorname{Leb}\{(x_1,\ldots,x_r)\in J^r:
        |x_i-x_j|\le\eta\text{ for some }i\ne j\}
        =O_{r,J}(\eta).
\end{aligned}
\]
Thus the excluded contribution vanishes as \(\eta\downarrow0\).  For finitely many disjoint intervals \(J_1,\ldots,J_m\), the mixed factorial moment expansion is the same expansion with each selected rate colored by the interval from which it comes.  The separated part factors into the corresponding products of lengths, and the same non-separated estimate removes the diagonals.  These limits are the mixed factorial moments of independent Poisson variables
with means \(e^{-1}|J_1|,\ldots,e^{-1}|J_m|\), so the corresponding count
vectors converge in distribution.

It remains to pass from interval counts to convergence of point measures.  The
death-process bound gives, for every fixed bounded interval \(J\),
\[
        \sup_n\Ee\Xi_n(J)
        \le C_J<\infty.
        \tag{7.30}
\]
Indeed, the relevant terminal coordinates lie in a bounded multiplicative
window and
\(\Pp(I_{n,a}=1)\le2/(a+1)\).  Let \(f\ge0\) be continuous with compact
support in \(J\).  Choose a step function
\(g=\sum_{\ell=1}^m c_\ell\one_{J_\ell}\), where the \(J_\ell\) are
disjoint half-open intervals, such that \(\|f-g\|_\infty\le\varepsilon\).
Since \(x\mapsto e^{-x}\) is one-Lipschitz on \([0,\infty)\),
\[
\left|
\Ee e^{-\langle f,\Xi_n\rangle}
-\Ee e^{-\langle g,\Xi_n\rangle}
\right|
\le\varepsilon\Ee\Xi_n(J)
\le C_J\varepsilon.
\tag{7.31}
\]
The already proved count-vector convergence gives convergence of the Laplace
functional for \(g\).  The same approximation applies to the limiting Poisson
process, whose expected mass on \(J\) is finite.  Letting first \(n\to\infty\)
and then \(\varepsilon\downarrow0\) proves convergence of Laplace functionals
for every nonnegative compactly supported continuous \(f\).  These Laplace
functionals determine the law in the vague topology, completing the proof.
\end{proof}

\begin{theorem}[Sukhatme fixed factorial moments]\label{thm:suk-fact}
For every fixed \(r\ge1\),
\[
        \Ee(F_n)_r\sim (e^{-1}\log n)^r.
\]
Consequently,
\[
        \frac{F_n}{e^{-1}\log n}\to1
\]
in probability.
\end{theorem}

\begin{proof}
Let
\[
        L_n=\lceil\log n\rceil,
        \qquad
        U_n=\lfloor n/\log n\rfloor,
\]
and put
\[
        G_n=\sum_{L_n\le a\le U_n}I_{n,a}.
\]
We first prove the factorial asymptotic for \(G_n\).  Fix \(\eta>0\).  In the expansion of \(\Ee(G_n)_r\), restrict to ordered distinct \(r\)-tuples whose logarithms are pairwise separated by at least \(\eta\).  For each such ordered tuple, sort the rates increasingly before applying Proposition \ref{prop:suk-joint}.  Pairwise logarithmic separation implies
\[
        \frac{a_{j+1}}{a_j}\ge e^\eta=1+(e^\eta-1),
\]
so Proposition \ref{prop:suk-joint} applies uniformly with separation parameter \(e^\eta-1\).  Hence the separated contribution is
\[
        e^{-r}
        \sum_{\substack{L_n\le a_1,\ldots,a_r\le U_n\\
                         a_i\ne a_j,\ |\log a_i-\log a_j|>\eta}}
        \frac1{a_1\cdots a_r}(1+o(1)),
\]
with the error uniform for fixed \(r\) and \(\eta\).  For fixed \(\eta\), the excluded logarithmic diagonal neighborhood has weighted lattice sum \(O_{r,\eta}((\log n)^{r-1})\).  Hence the separated weighted sum is
\[
        \left(\sum_{L_n\le a\le U_n}\frac1a\right)^r+o((\log n)^r).
\]
The non-separated part itself is negligible as well.  By Lemma \ref{lem:death-bound}, its contribution is bounded by a constant times the same logarithmic diagonal weighted sum, hence it is \(O_{r,\eta}((\log n)^{r-1})=o((\log n)^r)\).  Since
\[
        \sum_{L_n\le a\le U_n}\frac1a\sim \log n,
\]
we obtain
\[
        \Ee(G_n)_r\sim (e^{-1}\log n)^r.
\]

It remains to show that the omitted ranges do not affect the leading term.  By Lemma \ref{lem:death-bound}, any factorial-moment term containing at least one index outside \([L_n,U_n]\) is bounded by a constant times the corresponding product \(\prod 1/a_j\).  The logarithmic mass of the omitted set is
\[
        \sum_{a<L_n}\frac1a+
        \sum_{a>U_n}\frac1a
        =O(\log\log n)=o(\log n).
\]
Therefore the total contribution of tuples with at least one omitted index is
\[
        O((\log n)^{r-1}\log\log n)=o((\log n)^r).
\]
This proves \(\Ee(F_n)_r\sim(e^{-1}\log n)^r\).

For \(r=1,2\), the factorial moment asymptotics imply \(\Ee F_n\sim e^{-1}\log n\) and
\[
        \Ee F_n^2=\Ee(F_n)_2+\Ee F_n
        =(e^{-1}\log n)^2+o((\log n)^2).
\]
Consequently
\[
        \Var(F_n)=\Ee F_n^2-(\Ee F_n)^2=o((\log n)^2).
\]
Let \(\lambda_n=e^{-1}\log n\).  Since \(\Ee F_n/\lambda_n\to1\) and
\(\Var(F_n)=o(\lambda_n^2)\), Chebyshev's inequality gives
\(F_n/\lambda_n\to1\) in probability.
\end{proof}

Theorem \ref{thm:suk-fact} proves the fixed-order factorial moments and a law of large numbers for the full Sukhatme count.  The global centered fluctuation theory is completed in Section \ref{sec:global-endpoint-poisson}, where a terminal-rank Palm coupling yields total-variation approximation by a Poisson law with the exact mean.


\section{Endpoint-regular fixed-point universality}\label{sec:endpoint-fixed-universality}

We now promote the endpoint fixed-point prediction to a theorem.  The statement
is formulated in terminal rate coordinates.  Put
\[
        \lambda_{n,a}=\theta_{n,n+1-a},
        \qquad 1\le a\le n,
\]
so that terminal coordinate \(a\) is label \(n+1-a\).  Let
\[
        T_{n,a}\sim\Exp(\lambda_{n,a}),
        \qquad
        \tau_n(a)=1+\sum_{b\ne a}\one_{\{T_{n,b}>T_{n,a}\}}
        =n+1-\sigma_n(n+1-a).
\]
Thus $\tau_n$ is conjugate to $\sigma_n$ by the reversal permutation and has
the same cycle counts.  In Gumbel variables
\[
        Y_{n,a}=\log \lambda_{n,a}+G_a,
\]
the fixed-point indicator is
\[
        I_{n,a}=\one\left\{
        \sum_{b\ne a}\one_{\{Y_{n,b}<Y_{n,a}\}}=a-1
        \right\}.
\]
Indeed, label \(n+1-a\) is fixed exactly when there are \(a-1\) labels below
it in priority order.  We write
\[
        F_n=\sum_{a=1}^n I_{n,a}=C_1(\sigma_n)
\]
inside this section.

Multiplying all rates by a common deterministic factor adds the same constant to
all priorities and does not change the Luce order.  The following definition
therefore fixes a convenient normalization.

\begin{definition}[Uniform endpoint regular variation]\label{def:endpoint-rv}
The terminal rate array \((\lambda_{n,a})_{1\le a\le n}\) is
\emph{uniformly endpoint regularly varying with index \(\gamma>0\)} if there
exist positive constants \(s_n\), a constant \(C_0>0\), and a positive
measurable slowly varying function \(L\) at infinity such that the normalized
rates
\[
        \bar\lambda_{n,a}=s_n\lambda_{n,a}
\]
satisfy the following conditions.

\begin{enumerate}[label=(ER\arabic*)]
\item For every sequence \(A_n\to\infty\) with \(A_n=o(n)\),
\[
        \sup_{1\le a\le A_n}
        \left|
        \frac{\bar\lambda_{n,a}}
        {C_0(a/n)^\gamma L(n/a)}-1
        \right|\longrightarrow0.
\]
\item For every \(\delta>0\), there are constants
\(0<c_\delta<C_\delta<\infty\) such that
\[
        c_\delta\le \bar\lambda_{n,a}\le C_\delta,
        \qquad \delta n\le a\le n,
\]
for all sufficiently large \(n\).
\item The normalized terminal rates are nondecreasing:
\[
        \bar\lambda_{n,1}\le\bar\lambda_{n,2}
        \le\cdots\le\bar\lambda_{n,n}.
\]
\end{enumerate}
\end{definition}

A common multiplication of all rates does not change the Luce order.  We
therefore replace \(\lambda_{n,a}\) by \(\bar\lambda_{n,a}\) throughout this
section and suppress the bar.  We use repeatedly the uniform convergence
theorem and Potter bounds for slowly varying functions in the forms of
Bingham, Goldie, and Teugels \cite[Theorems~1.2.1 and~1.5.6]{BGT1987}.

For example, suppose \(\lambda_{n,a}=f(1-a/n)\), where
\[
        f(1-u)\sim C_0u^\gamma L(1/u),\qquad u\downarrow0,
\]
where \(f\) is bounded above and below away from zero on every compact
subinterval of \([0,1)\), and where the function
\(u\mapsto f(1-u)\) is nondecreasing on \((0,1]\).  Then the sampled array
satisfies Definition \ref{def:endpoint-rv}.  The global monotonicity in this
example is intentional: eventual endpoint monotonicity alone would not imply
condition \textup{(ER3)} for the full triangular array.

\begin{lemma}[The endpoint partial-sum bound is automatic]\label{lem:endpoint-partial-sum}
Under Definition \ref{def:endpoint-rv}, there is $c_\Sigma>0$ such that, for
all sufficiently large $n$,
\[
        \sum_{b=1}^a\lambda_{n,b}
        \ge c_\Sigma a\lambda_{n,a},
        \qquad 1\le a\le n.
\]
\end{lemma}

\begin{proof}
Suppose otherwise.  Then there are $n_k\to\infty$ and $a_k\in[n_k]$ such that
\[
        \frac{\sum_{b\le a_k}\lambda_{n_k,b}}
        {a_k\lambda_{n_k,a_k}}\longrightarrow0.
\]
Pass to a subsequence.  If $(a_k)$ is bounded, the ratio is at least $1/a_k$,
a contradiction.  Assume therefore that $a_k\to\infty$.

If $a_k/n_k\to0$, uniform endpoint regular variation and Potter bounds give
\[
        \frac1{a_k}\sum_{b=1}^{a_k}
        \frac{\lambda_{n_k,b}}{\lambda_{n_k,a_k}}
        \longrightarrow\int_0^1x^\gamma\dd x
        =\frac1{\gamma+1}.
\]
For completeness, on $\varepsilon a_k\le b\le a_k$ this is an ordinary
Riemann sum by the uniform convergence theorem for slowly varying functions.
On $b<\varepsilon a_k$, Potter's bound gives
$\lambda_{n_k,b}/\lambda_{n_k,a_k}\le C(b/a_k)^{\gamma/2}$, so the omitted
part is $O(\varepsilon^{1+\gamma/2})$.  Letting $\varepsilon\downarrow0$ proves
the display, again a contradiction.

The remaining case has $a_k\ge\delta n_k$ for some $\delta>0$.  By
monotonicity and (ER2),
\[
        \sum_{b=1}^{a_k}\lambda_{n_k,b}
        \ge\sum_{b=\lceil a_k/2\rceil}^{a_k}\lambda_{n_k,b}
        \ge \frac{a_k}{2}c_{\delta/2},
\]
whereas $\lambda_{n_k,a_k}\le C_\delta$.  The ratio is therefore bounded below
by $c_{\delta/2}/(2C_\delta)>0$, the final contradiction.
\end{proof}


Set
\[
        A_\gamma=\Gamma(1+1/\gamma),
        \qquad
        K_\gamma=A_\gamma^\gamma,
        \qquad
        c_\gamma=\gamma K_\gamma e^{-K_\gamma}.
\]

For fixed \(r\ge1\) and \(\eta>0\), define the admissible family
\[
\mathfrak E_N(r,\eta)=
\left\{
(n,a_1,\ldots,a_r):
N\le a_1<\cdots<a_r\le n/N,\quad
\frac{a_{j+1}}{a_j}\ge1+\eta
\right\}.
\]
Here and below, a statement asserted uniformly for
\(1\ll a_1<\cdots<a_r\ll n\) means that its supremum over
\(\mathfrak E_N(r,\eta)\) tends to zero as \(N\to\infty\).  This convention
makes all endpoint uniformity statements independent of a particular choice of
intermediate cutoff sequences.

\begin{lemma}[Weighted regularly varying terminal sums]
\label{lem:endpoint-rv-sums}
Assume Definition \ref{def:endpoint-rv}.  Let \(a=a_n\to\infty\) with
\(a=o(n)\).  For every fixed integer \(m\ge0\), put
\[
J_{\gamma,m}(\kappa)
 =\int_0^\infty x^{m\gamma}e^{-\kappa x^\gamma}\dd x
 =\frac1\gamma\Gamma\!\left(m+\frac1\gamma\right)
  \kappa^{-m-1/\gamma}.
\]
Then, uniformly for \(\kappa\) in compact subintervals of \((0,\infty)\),
\[
\sum_{b=1}^n
\left(\frac{\lambda_{n,b}}{\lambda_{n,a}}\right)^m
\exp\left(-\kappa\frac{\lambda_{n,b}}{\lambda_{n,a}}\right)
 =aJ_{\gamma,m}(\kappa)(1+o(1)).
\tag{8.1}\label{eq:endpoint-weighted-sums}
\]
The same assertion holds after deleting any fixed number of summands, uniformly
in the deleted indices.  In particular,
\[
\sum_{b=1}^n
\exp\left(-\kappa\frac{\lambda_{n,b}}{\lambda_{n,a}}\right)
 =aA_\gamma\kappa^{-1/\gamma}(1+o(1)),
\tag{8.2}\label{eq:endpoint-sum-zero}
\]
\[
\sum_{b=1}^n
\frac{\lambda_{n,b}}{\lambda_{n,a}}
\exp\left(-\kappa\frac{\lambda_{n,b}}{\lambda_{n,a}}\right)
 =\frac{aA_\gamma}{\gamma}
 \kappa^{-1-1/\gamma}(1+o(1)),
\tag{8.3}\label{eq:endpoint-sum-one}
\]
and, with
\(p_{n,b}(\kappa)=
\exp(-\kappa\lambda_{n,b}/\lambda_{n,a})\),
\[
\sum_{b=1}^n p_{n,b}(\kappa)(1-p_{n,b}(\kappa))
 =aA_\gamma\kappa^{-1/\gamma}
 (1-2^{-1/\gamma})(1+o(1)).
\tag{8.4}\label{eq:endpoint-variance-sum}
\]
\end{lemma}

\begin{proof}
Fix a compact interval
\([\kappa_-,\kappa_+]\Subset(0,\infty)\), and put
\[
        z_n=\frac na\longrightarrow\infty,
        \qquad
        R_n=\bigl(\log(e+z_n)\bigr)^{4/\gamma},
        \qquad
        B_n=\lceil aR_n\rceil.
\]
Since \(R_n=o(z_n)\), one has \(B_n=o(n)\).  Definition
\ref{def:endpoint-rv} therefore applies uniformly to all indices at most
\(B_n\).

Fix \(0<\varepsilon<1<R<\infty\).  On
\(\varepsilon a\le b\le Ra\), the uniform convergence theorem for slowly
varying functions gives
\[
        \frac{\lambda_{n,b}}{\lambda_{n,a}}
        =\left(\frac ba\right)^\gamma(1+o(1))
\]
uniformly.  Hence the contribution of this range to the left-hand side of
\eqref{eq:endpoint-weighted-sums}, divided by \(a\), converges uniformly in
\(\kappa\in[\kappa_-,\kappa_+]\) to
\[
        \int_\varepsilon^R x^{m\gamma}e^{-\kappa x^\gamma}\dd x.
\]
The range \(b<\varepsilon a\) contains at most \(\varepsilon a\) terms, and
\[
        \sup_{x\ge0,\,\kappa\in[\kappa_-,\kappa_+]}
        x^m e^{-\kappa x}<\infty.
\]
Its normalized contribution is therefore \(O_m(\varepsilon)\).

It remains to control the upper tail without applying regular variation at a
fixed positive fraction of \(n\).  Potter bounds, with exponent
\(\gamma/2\), imply that for all sufficiently large \(n\), uniformly for
\(Ra<b\le B_n\),
\[
 c\left(\frac ba\right)^{\gamma/2}
 \le \frac{\lambda_{n,b}}{\lambda_{n,a}}
 \le C\left(\frac ba\right)^{3\gamma/2}.
\]
Consequently the normalized contribution of this range is at most
\[
 C_m\int_R^\infty
 x^{3m\gamma/2}\exp(-c\kappa_-x^{\gamma/2})\dd x+o(1),
\]
which tends to zero as \(R\to\infty\).

For \(b>B_n\), monotonicity and the preceding Potter lower bound give
\[
 \frac{\lambda_{n,b}}{\lambda_{n,a}}
 \ge \frac{\lambda_{n,B_n}}{\lambda_{n,a}}
 \ge cR_n^{\gamma/2}
 =c\bigl(\log(e+z_n)\bigr)^2.
\]
Therefore
\[
\begin{aligned}
&\frac1a\sum_{b>B_n}
 \left(\frac{\lambda_{n,b}}{\lambda_{n,a}}\right)^m
 e^{-\kappa\lambda_{n,b}/\lambda_{n,a}}\\
&\qquad\le z_n
 \sup_{x\ge c(\log(e+z_n))^2}x^m e^{-\kappa_-x}
 \longrightarrow0.
\end{aligned}
\]
Letting first \(n\to\infty\), then \(\varepsilon\downarrow0\) and
\(R\uparrow\infty\), proves \eqref{eq:endpoint-weighted-sums}.  A fixed
deleted summand is bounded uniformly by
\(\sup_{x\ge0}x^m e^{-\kappa_-x}\), so finitely many deletions change the sum
by \(O(1)=o(a)\).  The special cases
\eqref{eq:endpoint-sum-zero}--\eqref{eq:endpoint-sum-one} follow from
\(J_{\gamma,0}(\kappa)=A_\gamma\kappa^{-1/\gamma}\) and
\(J_{\gamma,1}(\kappa)=A_\gamma\kappa^{-1-1/\gamma}/\gamma\).
Finally, \eqref{eq:endpoint-variance-sum} is the difference of
\eqref{eq:endpoint-sum-zero} at \(\kappa\) and at \(2\kappa\).
\end{proof}

\begin{lemma}[Endpoint product envelope]\label{lem:endpoint-rv-product-bound}
Assume Definition \ref{def:endpoint-rv}.  For every finite set
\(A\subset[n]\),
\[
        \Pp(I_{n,a}=1\text{ for all }a\in A)
        \le C^{|A|}\prod_{a\in A}\frac1a,
\]
where $C$ depends only on the constant $c_\Sigma$ from Lemma \ref{lem:endpoint-partial-sum}.
\end{lemma}

\begin{proof}
Reveal the Luce order as a death process.  When the alive set is \(S\), the next
label deleted from \(S\) is \(b\in S\) with probability
\(\lambda_{n,b}/\sum_{c\in S}\lambda_{n,c}\).  The event \(I_{n,a}=1\) says that
terminal coordinate \(a\) is deleted at the transition at which the number of
alive labels drops from \(a\) to \(a-1\).  Condition on the history before that
transition.  If coordinate \(a\) is already dead, the conditional probability is
zero.  If it is alive, the alive set has size \(a\) and contains \(a\).  Since
the rates are nondecreasing, the total rate of any alive set of size \(a\) is at
least $\sum_{b=1}^a\lambda_{n,b}$, and Lemma \ref{lem:endpoint-partial-sum} gives
\[
        \frac{\lambda_{n,a}}{\sum_{c\in S}\lambda_{n,c}}
        \le \frac{1}{c_\Sigma a}.
\]
Order the elements of \(A\) decreasingly and condition successively before the
corresponding death levels.  Multiplying the conditional bounds proves the
claim.
\end{proof}


\subsection{Separated endpoint fixed-point asymptotics}
\label{subsec:endpoint-separated-fixed}

Fix \(r\ge1\) and \(\eta>0\), and let
\((n,a_1,\ldots,a_r)\in\mathfrak E_N(r,\eta)\).  Put
\[
        B=\{a_1,\ldots,a_r\},
        \qquad z_j=a_j-j,
        \qquad 1\le j\le r.
\]
For \(g=(g_1,\ldots,g_r)\), set
\[
        y_j(g)=\log\lambda_{n,a_j}+g_j,
        \qquad
        Z_j(g)=\sum_{b\notin B}
        \one_{\{Y_{n,b}<y_j(g)\}},
\]
and write
\[
        M_j(g_j)=\Ee Z_j(g)
        =\sum_{b\notin B}
        \exp\left(-e^{-g_j}
        \frac{\lambda_{n,b}}{\lambda_{n,a_j}}\right).
\]
The relevant chamber is
\[
        \mathscr C_n=\{g:y_1(g)<\cdots<y_r(g)\}.
\]
When \(g\in\mathscr C_n\), put \(y_0=-\infty\),
\(y_{r+1}=\infty\), and define interval counts
\[
        W_1=Z_1,
        \qquad W_j=Z_j-Z_{j-1}\quad(2\le j\le r).
\]
Thus a nonspecified priority contributes \(e_j\) to \(W\) when it lies in
\([y_{j-1},y_j)\), and contributes \(0\) when it lies in
\([y_r,\infty)\).  The target interval vector is
\[
        w_1=a_1-1,
        \qquad
        w_j=a_j-a_{j-1}-1\quad(2\le j\le r),
\]
and ratio separation gives
\[
        c_\eta a_j\le w_j\le a_j
        \qquad(1\le j\le r)
\tag{8.5}\label{eq:endpoint-interval-target-scale}
\]
for all sufficiently large \(N\).

\begin{lemma}[Endpoint Gumbel saddle geometry]
\label{lem:endpoint-gumbel-saddle}
For each \(j\), the equation
\[
        M_j(g_j^*)=z_j
\]
has a unique solution.  Uniformly over
\(\mathfrak E_N(r,\eta)\), as \(N\to\infty\),
\[
        g_j^*\longrightarrow g_\gamma:=-\log K_\gamma,
        \qquad
        M_j'(g_j^*)\sim\frac{a_j}{\gamma}.
\tag{8.6}\label{eq:endpoint-saddle-asymptotics}
\]
Moreover, for every fixed \(L<\infty\),
\[
        \sup_{|g|\le L}|M_j''(g)|\le C_La_j,
        \qquad
        c_La_j\le M_j'(g)\le C_La_j
\tag{8.7}\label{eq:endpoint-mean-derivative-bounds}
\]
uniformly whenever \(|g-g_j^*|\le1\).  There are constants
\(d_\eta>0\), \(\delta_\eta>0\), and \(N_0\) such that, for
\(N\ge N_0\),
\[
        y_{j+1}(g_{j+1})-y_j(g_j)\ge d_\eta
\tag{8.8}\label{eq:endpoint-threshold-gap}
\]
for every \(g\) satisfying
\(|g_j-g_j^*|\le\delta_\eta\), \(1\le j\le r\).  In particular, this
entire product neighborhood lies in \(\mathscr C_n\).
\end{lemma}

\begin{proof}
The function \(M_j\) is continuous and strictly increasing from \(0\) to
\(n-r\), so the saddle exists uniquely.  With \(\kappa=e^{-g}\), Lemma
\ref{lem:endpoint-rv-sums}, including its finite-deletion assertion, gives
uniformly on compact \(g\)-sets
\[
        M_j(g)=a_jA_\gamma e^{g/\gamma}(1+o(1)).
\tag{8.9}\label{eq:endpoint-mean-asymptotic}
\]
Differentiating term by term and applying
\eqref{eq:endpoint-weighted-sums} with \(m=1\) gives
\[
\begin{aligned}
        M_j'(g)
        &=\sum_{b\notin B}e^{-g}
        \frac{\lambda_{n,b}}{\lambda_{n,a_j}}
        \exp\left(-e^{-g}
        \frac{\lambda_{n,b}}{\lambda_{n,a_j}}\right)\\
        &=\frac{a_jA_\gamma}{\gamma}e^{g/\gamma}(1+o(1)).
\end{aligned}
\tag{8.10}\label{eq:endpoint-mean-prime-asymptotic}
\]
The saddle equation and \(z_j/a_j\to1\) imply
\(A_\gamma e^{g_j^*/\gamma}\to1\), which is equivalent to
\(g_j^*\to-\gamma\log A_\gamma=-\log K_\gamma\).  Substitution in
\eqref{eq:endpoint-mean-prime-asymptotic} gives the derivative asymptotic.
Since
\[
        M_j''(g)=\sum_{b\notin B}
        \left[
        \left(e^{-g}\frac{\lambda_{n,b}}{\lambda_{n,a_j}}\right)^2
        -e^{-g}\frac{\lambda_{n,b}}{\lambda_{n,a_j}}
        \right]
        e^{-e^{-g}\lambda_{n,b}/\lambda_{n,a_j}},
\]
\eqref{eq:endpoint-weighted-sums} with \(m=1,2\) proves the second-derivative
bound.  The lower and upper derivative bounds then follow from
\eqref{eq:endpoint-mean-prime-asymptotic} on a compact neighborhood of the
saddles.

It remains to separate the thresholds.  Uniform endpoint regular variation and
Potter bounds imply that there is \(c_\eta>0\) such that
\[
        \log\frac{\lambda_{n,a_{j+1}}}{\lambda_{n,a_j}}
        \ge2c_\eta
\tag{8.11}\label{eq:endpoint-rate-gap}
\]
for all sufficiently large \(N\).  To verify this uniformly, first restrict
\(a_{j+1}/a_j\) to a fixed compact interval
\([1+\eta,R]\) and use the uniform convergence theorem; for ratios larger
than \(R\), use the Potter lower bound
\(\lambda_{n,a_{j+1}}/\lambda_{n,a_j}
\ge c(a_{j+1}/a_j)^{\gamma/2}\), and then choose \(R\) large.
Because all \(g_j^*\) converge to the same constant, the saddle threshold gaps
are at least \(c_\eta\) for large \(N\).  Taking
\(\delta_\eta<c_\eta/4\) proves \eqref{eq:endpoint-threshold-gap}.
\end{proof}

\begin{lemma}[Endpoint interval edge blocks]
\label{lem:endpoint-interval-edge-blocks}
Fix \(L_0>0\), and consider any vector \(g\in\mathscr C_n\).  Write
\[
        \mathcal I_j(g)=[y_{j-1}(g),y_j(g))\quad(1\le j\le r),
        \qquad
        \mathcal I_0(g)=[y_r(g),\infty).
\]
Let
\[
        m_j(g)=\Ee W_j(g).
\]
For every \(c_0>0\), there are constants
\(0<\alpha<\beta<\infty\) and \(c>0\), depending only on
\(r,\eta,L_0,c_0\), with the following property.  If
\(|g_j|\le L_0\) and \(m_j(g)\ge c_0a_j\), then at least \(ca_j\) indices
\(b\in[\alpha a_j,\beta a_j]\setminus B\) satisfy
\[
        \Pp(Y_{n,b}\in\mathcal I_j(g))\ge c
\]
and, for one fixed
\(\kappa(j)\in\{j+1,\ldots,r,0\}\),
\[
        \Pp(Y_{n,b}\in\mathcal I_{\kappa(j)}(g))\ge c.
\]
Here \(\kappa(j)=0\) denotes the anchor interval.
The selected subblocks may be chosen disjoint as \(j\) varies, after reducing
\(c\).

In the saddle neighborhood from Lemma
\ref{lem:endpoint-gumbel-saddle}, one may instead choose disjoint blocks of
size at least \(ca_j\) on which the two adjacent atoms \(e_j\) and
\(e_{j+1}\) each have probability at least \(c\), with
\(e_{r+1}=0\).  All these conclusions remain valid, with modified constants,
under every interval-coordinate exponential tilt whose coordinates are bounded
by a fixed constant.
\end{lemma}

\begin{proof}
The contribution to \(m_j(g)\) from \(b<\alpha a_j\) is at most
\(\alpha a_j\).  Since membership in the \(j\)-th interval implies
\(Y_{n,b}<y_j\), the contribution from \(b>\beta a_j\) is bounded by
\[
        \sum_{b>\beta a_j}
        \exp\left(-e^{-g_j}
        \frac{\lambda_{n,b}}{\lambda_{n,a_j}}\right).
\]
The moving-cutoff proof of Lemma \ref{lem:endpoint-rv-sums}, with a Potter
bound on the intermediate tail and monotonicity beyond the moving cutoff,
shows that this is \(o_\beta(a_j)\), uniformly whenever \(|g_j|\le L_0\).  Choose \(\alpha\)
small and \(\beta\) large.  A positive multiple of \(a_j\) of the expected
interval mass then comes from \(b\in[\alpha a_j,\beta a_j]\).  This range
contains \(O(a_j)\) indices, so at least \(ca_j\) of them have
\(j\)-th interval probability at least \(c\).

On the same middle range, regular variation on fixed multiplicative intervals
keeps \(\lambda_{n,b}/\lambda_{n,a_j}\) between two positive constants.
Since \(g_j\) is bounded, the probability
\(\Pp(Y_{n,b}\ge y_j)\) is also bounded below.  Split this upper event among
the finitely many later intervals and the anchor; the pigeonhole principle
produces a common \(\kappa(j)>j\), or the anchor \(0\), on a subblock of order
\(a_j\).  Disjoint subblocks are extracted greedily.  Indeed,
\(a_{i+1}/a_i\ge1+\eta\) and \(r\) is fixed, so
\(\sum_{i<j}a_i\le C_{r,\eta}a_j\).

In the saddle neighborhood, choose a smaller multiplicative block
\([(1-\delta)a_j,(1+\delta)a_j]\), with \(\delta\) small enough that these
blocks are disjoint.  On this block the centered threshold
\(y_j-\log\lambda_{n,b}\) remains in a compact interval, and
\eqref{eq:endpoint-threshold-gap} gives uniformly positive intervals on both
sides of \(y_j\).  The standard Gumbel density is positive and continuous, so
the immediate lower and upper interval probabilities are both bounded below.
This gives the adjacent atoms.  A bounded exponential tilt multiplies all
probabilities in the finite categorical support by factors bounded above and
below, so every two-atom lower bound persists.
\end{proof}

\begin{lemma}[Endpoint Gumbel anisotropic local CLT]
\label{lem:endpoint-gumbel-llt}
Let \(\mathcal U_n=\{g:|g_j-g_j^*|\le\delta_\eta,
1\le j\le r\}\), where \(\delta_\eta\) is from Lemma
\ref{lem:endpoint-gumbel-saddle}.  Let
\[
        \Sigma(g)=\Var Z(g),
        \qquad
        D_a=\diag(a_1,\ldots,a_r).
\]
There are constants \(0<c<C<\infty\) such that, uniformly over
\(\mathfrak E_N(r,\eta)\) and \(g\in\mathcal U_n\),
\[
        cD_a\le\Sigma(g)\le CD_a
\tag{8.12}\label{eq:endpoint-covariance-comparison}
\]
in quadratic-form order.  For every fixed \(R<\infty\), uniformly for lattice
points \(q\in\mathbb Z^r\) satisfying
\[
        \|D_a^{-1/2}(q-M(g))\|\le R,
\]
\[
\begin{aligned}
\Pp(Z(g)=q)
&=(2\pi)^{-r/2}(\det\Sigma(g))^{-1/2}\\
&\quad\times
\exp\left[-\frac12(q-M(g))^T\Sigma(g)^{-1}(q-M(g))\right]
+o\left(\prod_{j=1}^r a_j^{-1/2}\right).
\end{aligned}
\tag{8.13}\label{eq:endpoint-gumbel-llt}
\]
The error is uniform over the displayed family.  Moreover, for every fixed
\(\nu<1/6\),
\[
\Pp(Z(g)=q)
\le C\prod_{j=1}^r a_j^{-1/2}
\exp\left[-c\sum_{j=1}^r
\frac{(q_j-M_j(g_j))^2}{a_j}\right]
\tag{8.14}\label{eq:endpoint-gumbel-local-upper}
\]
uniformly whenever
\(|q_j-M_j(g_j)|\le a_j^{1/2+\nu}\) for every \(j\).
\end{lemma}

\begin{proof}
For \(b\notin B\), let
\[
        V_b(g)=
        (\one_{\{Y_{n,b}<y_1\}},\ldots,
         \one_{\{Y_{n,b}<y_r\}}).
\]
These independent vectors have the cumulative categorical support
\[
        (1,\ldots,1),(0,1,\ldots,1),\ldots,(0,\ldots,0,1),0.
\]
The adjacent blocks in Lemma \ref{lem:endpoint-interval-edge-blocks} imply
that, for every \(u\in\mathbb R^r\), the \(j\)-th block contributes at least
\(ca_ju_j^2\) to \(\Var(u\cdot Z(g))\).  Hence the covariance lower bound in
\eqref{eq:endpoint-covariance-comparison} holds.  The upper bound follows from
\(\Var Z_j(g)=O(a_j)\), which is
\eqref{eq:endpoint-variance-sum}, and Cauchy's inequality.

The same adjacent blocks give, uniformly under bounded tilts,
\[
        |\Ee e^{iu\cdot(Z(g)-M(g))}|
        \le\exp\left[-c\sum_{j=1}^r
        a_j\operatorname{dist}(u_j,2\pi\mathbb Z)^2\right].
\tag{8.15}\label{eq:endpoint-fourier-decay}
\]
Indeed, the two adjacent cumulative atoms differ by the coordinate vector
\(e_j\), and the blocks are disjoint.  Thus the lattice generated by the
support differences is the full lattice \(\mathbb Z^r\).

Fourier inversion now gives the local expansion.  On the anisotropic small box
\(|u_j|\le Ra_j^{-1/2}\), boundedness of the categorical increments and the
covariance upper bound give
\[
\log\Ee e^{iu\cdot(Z-M)}
=-\frac12u^T\Sigma(g)u
+O\left((\max_j|u_j|)\sum_{j=1}^ra_ju_j^2\right).
\]
After the change of variables \(u=D_a^{-1/2}v\), the error is
\(O_R(a_1^{-1/2})=o(1)\).  The complement of a fixed \(v\)-ball is uniformly
negligible on the product scale by \eqref{eq:endpoint-fourier-decay}.  This
proves \eqref{eq:endpoint-gumbel-llt}.

For \eqref{eq:endpoint-gumbel-local-upper}, apply Lemma
\ref{lem:anisotropic-tilted-upper} with \(W=Z(g)\) and \(A_j=a_j\).  The
coordinatewise Fourier hypothesis under bounded tilts follows from the
adjacent blocks, and the required tilted covariance upper bound follows from
bounded distortion of the categorical probabilities and the argument proving
\eqref{eq:endpoint-covariance-comparison}.  Since
\(|q_j-M_j(g_j)|\le a_j^{1/2+\nu}=o(a_j)\), the truncated exponent in that
lemma is exactly \((q_j-M_j(g_j))^2/a_j\).  This proves
\eqref{eq:endpoint-gumbel-local-upper}.
\end{proof}

\begin{lemma}[Endpoint mixed cumulative--interval localization]
\label{lem:endpoint-mixed-localization}
There is \(L_0=L_0(r,\eta,\gamma)\) such that the following holds for
\(L\ge L_0\).  For \(g\in\mathscr C_n\), put
\[
\begin{aligned}
        B_R(g)&=\{j:g_j>L\},\\
        B_L(g)&=\{j:g_j<-L\},\\
        G(g)&=\{1,\ldots,r\}\setminus(B_R(g)\cup B_L(g)),
\end{aligned}
\]
and let
\[
        H_j=z_j-M_j(g_j),\qquad H_0=0,
        \qquad d_j=H_j-H_{j-1}.
\]
On every region on which the three index sets are fixed,
\[
\begin{aligned}
\Pp(Z(g)=z)
&\le C_L\prod_{j=1}^ra_j^{-1/2}
\exp\left[-c\sum_{j\in B_R}|H_j|
           -c\sum_{j\in B_L}a_j\right]\\
&\quad\times
\exp\left[-c_L\sum_{j\in G}
\min\left\{\frac{H_j^2}{a_j},a_j\right\}\right].
\end{aligned}
\tag{8.16}\label{eq:endpoint-mixed-cumulative-bound}
\]
Consequently,
\[
\begin{aligned}
\Pp(Z(g)=z)
&\le C_L\prod_{j=1}^ra_j^{-1/2}
\exp\left[-ce^{L/(2\gamma)}\sum_{j\in B_R}a_j
           -c\sum_{j\in B_L}a_j\right]\\
&\quad\times
\exp\left[-c_L\sum_{j\in G}
 a_j\min\{(g_j-g_j^*)^2,1\}\right].
\end{aligned}
\tag{8.17}\label{eq:endpoint-mixed-scaled-bound}
\]
Here \(c>0\) is independent of \(L\), while \(C_L,c_L\) may depend on
\(L,r,\eta,\gamma\).
\end{lemma}

\begin{proof}
Fix the offender sets and abbreviate them by \(B_R,B_L,G\).  From
\eqref{eq:endpoint-mean-asymptotic}, uniformly for large \(L\),
\[
\begin{array}{lll}
 j\in B_R:&M_j(g_j)\ge ce^{L/\gamma}a_j,&H_j<0,\\[1mm]
 j\in B_L:&M_j(g_j)\le Ce^{-L/\gamma}a_j,&0<H_j\le a_j.
\end{array}
\tag{8.18}\label{eq:endpoint-offender-geometry}
\]
For \(j\in G\), \(g_j\in[-L,L]\), so
\(m_j(g)=O_L(a_j)\) and \(|d_j|\le M_La_j\).

Let
\[
        K_g(\theta)=\log\Ee
        \exp\{\theta\cdot(W(g)-m(g))\}
\]
be the centered interval-coordinate cumulant.  For a cumulative-coordinate
vector \(\alpha\), define
\[
        (\mathcal T\alpha)_k=\sum_{j\ge k}\alpha_j.
\]
Because \(Z_j=W_1+\cdots+W_j\),
\[
        (\mathcal T\alpha)\cdot d=\alpha\cdot H.
\tag{8.19}\label{eq:endpoint-cumulative-interval-duality}
\]
Choose \(\delta>0\) so small that \(3\delta\) is admissible in Lemma
\ref{lem:actual-mean-lower-cumulative-chernoff}.  Set
\[
        \alpha_j^-=-\delta\one_{\{j\in B_R\}},
        \qquad
        \alpha_j^+=\delta\one_{\{j\in B_L\}}.
\]
For the nonoffending interval coordinates, fix
\(0<\varepsilon<c_\eta/4\), where \(c_\eta\) is from
\eqref{eq:endpoint-interval-target-scale}, and put
\[
        G_0=\{j\in G:|d_j|\le\varepsilon a_j\},
        \qquad G_1=G\setminus G_0.
\]
For a sufficiently small constant \(\delta_L>0\), define
\[
        \lambda_j=
        \begin{cases}
        \delta_Ld_j/a_j,&j\in G_0,\\
        \delta_L\operatorname{sgn}(d_j),&j\in G_1,\\
        0,&j\notin G.
        \end{cases}
\]
Use the single total interval-coordinate tilt
\[
        \Theta=\mathcal T\alpha^-+\mathcal T\alpha^++\lambda.
\]
Generalized H\"older inequality gives
\[
\begin{aligned}
K_g(\Theta)
&\le\frac13K_g(3\mathcal T\alpha^-)
 +\frac13K_g(3\mathcal T\alpha^+)
 +\frac13K_g(3\lambda),
\end{aligned}
\tag{8.20}\label{eq:endpoint-holder-cumulant}
\]
and therefore, by \eqref{eq:endpoint-cumulative-interval-duality},
\[
\begin{aligned}
-\Theta\cdot d+K_g(\Theta)
&\le\frac13\{-3\alpha^-\cdot H+K_g(3\mathcal T\alpha^-)\}\\
&\quad+\frac13\{-3\alpha^+\cdot H+K_g(3\mathcal T\alpha^+)\}\\
&\quad+\frac13\{-3\lambda\cdot d+K_g(3\lambda)\}.
\end{aligned}
\tag{8.21}\label{eq:endpoint-holder-exponent}
\]

The proof of Lemma
\ref{lem:actual-mean-lower-cumulative-chernoff} uses only independence,
nested indicator counts, and the bound that one summand contributes at most
\(r\) to the cumulative linear form.  It therefore applies verbatim to the
present lower-priority counts.  For \(B_R\),
\eqref{eq:endpoint-offender-geometry} gives
\(z_j\le Ce^{-L/\gamma}M_j(g_j)\).  The lower-tail part of Lemma
\ref{lem:actual-mean-lower-cumulative-chernoff}, applied with parameter
\(3\delta\), yields
\[
        -3\alpha^-\cdot H+K_g(3\mathcal T\alpha^-)
        \le-c\sum_{j\in B_R}M_j(g_j)
        \le-c\sum_{j\in B_R}|H_j|.
\tag{8.22}\label{eq:endpoint-right-offender-cost}
\]
For \(B_L\), one has
\(M_j(g_j)\le Ce^{-L/\gamma}z_j\).  The upper-tail part of the same lemma
gives
\[
        -3\alpha^+\cdot H+K_g(3\mathcal T\alpha^+)
        \le-c\sum_{j\in B_L}z_j
        \le-c\sum_{j\in B_L}a_j.
\tag{8.23}\label{eq:endpoint-left-offender-cost}
\]

Under every bounded interval tilt \(\xi\), the categorical support implies
\[
        \Var_\xi(v\cdot W(g))
        \le C_L\sum_{j\in G}a_jv_j^2
\]
for vectors supported on \(G\).  Taylor's formula for the cumulant and a
sufficiently small choice of \(\delta_L\) therefore give
\[
        -3\lambda\cdot d+K_g(3\lambda)
        \le-c_L\sum_{j\in G}
        \min\left\{\frac{d_j^2}{a_j},a_j\right\}.
\tag{8.24}\label{eq:endpoint-local-tilt-cost}
\]

All coordinates of the total tilt \(\Theta\) are bounded.  For
\(j\in G_0\),
\(m_j(g)\ge w_j-\varepsilon a_j\ge(c_\eta/2)a_j\).
Lemma \ref{lem:endpoint-interval-edge-blocks} supplies forward-edge blocks of
order \(a_j\), which may be chosen disjoint.  After projection onto
\(G_0\), a later atom whose coordinate is not retained becomes the anchor
atom.  Lemma \ref{lem:projected-forward-edge-concentration} therefore gives
under the single tilted law
\[
        \Pp_\Theta(W_{G_0}=w_{G_0})
        \le C_L\prod_{j\in G_0}a_j^{-1/2}.
\tag{8.25}\label{eq:endpoint-tilted-local-factor}
\]
The exact change-of-measure identity, together with
\eqref{eq:endpoint-holder-exponent}--\eqref{eq:endpoint-tilted-local-factor},
now yields the same bound as \eqref{eq:endpoint-mixed-cumulative-bound}, but
with interval residuals \(d_j\) in place of cumulative residuals on \(G\),
and initially only the local factors for \(G_0\).  Every index in \(G_1\),
\(B_R\), or \(B_L\) carries a fixed exponential loss of order \(a_j\), which
absorbs its missing factor \(a_j^{-1/2}\).

It remains to compare interval and cumulative residuals without discarding the
boundary values.  Apply Lemma \ref{lem:boundary-truncated-chain} with
\(A_j=a_j\), with its set \(B_-\) equal to the present \(B_R\), and with its
set \(B_+\) equal to the present \(B_L\).  Its assumptions hold by
\(|d_j|\le M_La_j\) on \(G\) and \(|H_j|\le a_j\) on \(B_L\).  Reserving a
fixed fraction of the offender costs in
\eqref{eq:endpoint-right-offender-cost}--\eqref{eq:endpoint-left-offender-cost}
therefore gives \eqref{eq:endpoint-mixed-cumulative-bound}.

Finally, on \(G\), the derivative bounds in
\eqref{eq:endpoint-mean-derivative-bounds}, monotonicity, and compactness imply
\[
        |H_j|\ge c_La_j\min\{|g_j-g_j^*|,1\}.
\tag{8.26}\label{eq:endpoint-mean-displacement}
\]
On \(B_R\), \eqref{eq:endpoint-offender-geometry} gives
\(|H_j|\ge ce^{L/\gamma}a_j\), which may be weakened to
\(ce^{L/(2\gamma)}a_j\).  These observations convert
\eqref{eq:endpoint-mixed-cumulative-bound} into
\eqref{eq:endpoint-mixed-scaled-bound}.
\end{proof}

\begin{lemma}[Endpoint Gumbel product localization]
\label{lem:endpoint-gumbel-localization}
Uniformly over \(\mathfrak E_N(r,\eta)\), the following assertions hold.

\begin{enumerate}[label=(\roman*)]
\item For every fixed compact cube \(K\subset\mathbb R^r\) containing all
saddle vectors for large \(N\), there are \(C_K,c_K>0\) such that, throughout
\(K\cap\mathscr C_n\),
\[
        \Pp(Z(g)=z)
        \le C_K\prod_{j=1}^ra_j^{-1/2}
        \exp\left[-c_K\sum_{j=1}^r
        a_j\min\{(g_j-g_j^*)^2,1\}\right].
\tag{8.27}\label{eq:endpoint-compact-localization}
\]
\item With \(\varphi(g)=e^{-g}e^{-e^{-g}}\),
\[
\lim_{L\to\infty}\limsup_{N\to\infty}
\sup_{\mathfrak E_N(r,\eta)}
\frac{
\displaystyle
\int_{\mathscr C_n\cap\{g\notin[-L,L]^r\}}
\prod_{j=1}^r\varphi(g_j)\Pp(Z(g)=z)\dd g
}{\prod_{j=1}^ra_j^{-1}}=0.
\tag{8.28}\label{eq:endpoint-noncompact-localization}
\]
\end{enumerate}
\end{lemma}

\begin{proof}
Choose \(L\) so that \(K\subset[-L,L]^r\).  With no offenders, the proof of
Lemma \ref{lem:endpoint-mixed-localization} gives
\eqref{eq:endpoint-compact-localization}; equivalently, apply
\eqref{eq:endpoint-mixed-cumulative-bound} with \(G=\{1,\ldots,r\}\).

For the tail, partition the integration region according to the finitely many
sets \(B_R,B_L,G\), and use \eqref{eq:endpoint-mixed-scaled-bound}.  Every
nonoffending coordinate contributes a factor \(O(a_j^{-1/2})\) after
integration on its Gaussian scale, in addition to the prefactor
\(a_j^{-1/2}\).  For a left offender, \(e^{-ca_j}\) absorbs the missing local
factor and
\[
        \int_{-\infty}^{-L}\varphi(g)\dd g=e^{-e^L}.
\]
For a right offender, \(e^{-ce^{L/(2\gamma)}a_j}\) absorbs the missing local
factor and
\[
        \int_L^\infty\varphi(g)\dd g
        =1-e^{-e^{-L}}\le e^{-L}.
\]
At least one offender is present.  Summing over the finitely many offender sets
proves \eqref{eq:endpoint-noncompact-localization}.
\end{proof}

\begin{lemma}[Endpoint covariance stability on the saddle scale]
\label{lem:endpoint-covariance-stability}
Let
\[
        H_a=\diag(a_1^{-1/2},\ldots,a_r^{-1/2}),
        \qquad
        \Sigma_*=\Sigma(g^*).
\]
For every fixed \(R<\infty\),
\[
\sup_{\|u\|\le R}
\left\|
D_a^{-1/2}
\bigl[\Sigma(g^*+H_au)-\Sigma_*\bigr]
D_a^{-1/2}
\right\|_{\mathrm{op}}
\longrightarrow0
\tag{8.29}\label{eq:endpoint-covariance-stability}
\]
uniformly over \(\mathfrak E_N(r,\eta)\).  Consequently, uniformly on the
same sets,
\[
        \frac{\det\Sigma(g^*+H_au)}{\det\Sigma_*}\longrightarrow1
\tag{8.30}\label{eq:endpoint-determinant-stability}
\]
and
\[
\left\|
D_a^{1/2}
\bigl[\Sigma(g^*+H_au)^{-1}-\Sigma_*^{-1}\bigr]
D_a^{1/2}
\right\|_{\mathrm{op}}
\longrightarrow0.
\tag{8.31}\label{eq:endpoint-inverse-stability}
\]
\end{lemma}

\begin{proof}
For \(j\le k\), put
\[
        p_{b,j}(g_j)=
        \exp\left(-e^{-g_j}
        \frac{\lambda_{n,b}}{\lambda_{n,a_j}}\right).
\]
Since the threshold events are nested in the chamber,
\[
        \Sigma_{jk}(g)
        =\sum_{b\notin B}p_{b,j}(g_j)(1-p_{b,k}(g_k)).
\tag{8.32}\label{eq:endpoint-covariance-entry}
\]
Lemma \ref{lem:endpoint-rv-sums}, with \(m=1\), gives uniformly on a fixed
saddle neighborhood
\[
        |\partial_{g_j}\Sigma_{jk}(g)|\le Ca_j,
        \qquad
        |\partial_{g_k}\Sigma_{jk}(g)|\le Ca_k.
\tag{8.33}\label{eq:endpoint-covariance-derivatives}
\]
The same estimate, with the evident interpretation, holds on the diagonal.
Hence, for \(\|u\|\le R\),
\[
        |\Sigma_{jk}(g^*+H_au)-\Sigma_{jk}(g^*)|
        \le C_R(\sqrt{a_j}+\sqrt{a_k}).
\]
After division by \(\sqrt{a_ja_k}\), this is at most
\(C_R(a_j^{-1/2}+a_k^{-1/2})\le2C_Ra_1^{-1/2}\).  Since \(r\) is fixed,
\eqref{eq:endpoint-covariance-stability} follows.  By
\eqref{eq:endpoint-covariance-comparison}, the normalized matrices
\(D_a^{-1/2}\Sigma(g)D_a^{-1/2}\) remain in a compact subset of the positive
definite matrices.  Continuity of determinant and inversion on this set gives
\eqref{eq:endpoint-determinant-stability} and
\eqref{eq:endpoint-inverse-stability}.
\end{proof}

\begin{proposition}[Separated endpoint local estimate]
\label{prop:endpoint-rv-joint}
Assume Definition \ref{def:endpoint-rv}.  Fix \(r\ge1\) and \(\eta>0\).
Then
\[
\lim_{N\to\infty}
\sup_{(n,a)\in\mathfrak E_N(r,\eta)}
\left|
\frac{
\Pp(I_{n,a_1}=\cdots=I_{n,a_r}=1)
}{c_\gamma^r/(a_1\cdots a_r)}-1
\right|=0.
\tag{8.34}\label{eq:endpoint-separated-joint}
\]
\end{proposition}

\begin{proof}
On the joint fixed-point event, the specified priorities are necessarily
ordered as
\(y_1<\cdots<y_r\).  Conditional on their centered Gumbels \(g\) in this
chamber, exactly \(j-1\) specified priorities lie below \(y_j\), so the event
is equivalent to \(Z_j(g)=a_j-j\), \(1\le j\le r\).  Therefore
\[
\begin{aligned}
&\Pp(I_{n,a_1}=\cdots=I_{n,a_r}=1)\\
&\quad=\int_{\mathscr C_n}
\prod_{j=1}^r\varphi(g_j)\,
\Pp(Z(g)=z)\dd g.
\end{aligned}
\tag{8.35}\label{eq:endpoint-exact-integral}
\]
Lemma \ref{lem:endpoint-gumbel-saddle} puts a fixed product neighborhood of
\(g^*\) inside the chamber.  Lemma
\ref{lem:endpoint-gumbel-localization} removes the noncompact region on the
scale \(\prod_j a_j^{-1}\), and its compact bound makes the part outside any
fixed saddle neighborhood exponentially negligible on that scale.

Set
\[
        D_n=\diag(M_1'(g_1^*),\ldots,M_r'(g_r^*)),
        \qquad g=g^*+H_au.
\]
For bounded \(u\), Taylor's theorem and
\eqref{eq:endpoint-mean-derivative-bounds} give
\[
        M(g)-z=D_nH_au+r_n(u),
        \qquad
        \|D_a^{-1/2}r_n(u)\|\longrightarrow0
\tag{8.36}\label{eq:endpoint-mean-stability}
\]
uniformly on fixed \(u\)-balls.  Indeed, the \(j\)-th remainder is
\(O_R(1)\), hence is \(o(\sqrt{a_j})\).  The specified density satisfies
\[
        \prod_{j=1}^r\varphi(g_j^*+u_j/\sqrt{a_j})
        =\prod_{j=1}^r\varphi(g_j^*)(1+o(1))
\tag{8.37}\label{eq:endpoint-density-stability}
\]
uniformly on fixed balls.

Apply Lemma \ref{lem:endpoint-gumbel-llt}, then use Lemma
\ref{lem:endpoint-covariance-stability} to replace \(\Sigma(g)\) by
\(\Sigma_*\).  On every fixed \(u\)-ball, the integrand in
\eqref{eq:endpoint-exact-integral}, including the Jacobian
\(|\det H_a|\), is
\[
\begin{aligned}
&(1+o(1))\prod_{j=1}^r\varphi(g_j^*)
(2\pi)^{-r/2}(\det\Sigma_*)^{-1/2}|\det H_a|\\
&\quad\times
\exp\left[-\frac12u^TH_aD_n\Sigma_*^{-1}D_nH_au\right]\dd u.
\end{aligned}
\tag{8.38}\label{eq:endpoint-saddle-integrand}
\]
The matrices in the exponent are uniformly positive definite and uniformly
bounded.  The compact localization estimate
\eqref{eq:endpoint-compact-localization}, after the change of variables,
provides the integrable majorant
\[
        C\left(\prod_{j=1}^ra_j^{-1}\right)
        \exp\left[-c\sum_{j=1}^r\min\{u_j^2,a_j\}\right].
\]
Thus the fixed-ball limit may be extended uniformly to \(\mathbb R^r\).
The Gaussian determinant cancels exactly:
\[
\begin{aligned}
&(2\pi)^{-r/2}(\det\Sigma_*)^{-1/2}|\det H_a|
\int_{\mathbb R^r}
 e^{-\frac12u^TH_aD_n\Sigma_*^{-1}D_nH_au}\dd u\\
&\qquad=|\det D_n|^{-1}.
\end{aligned}
\tag{8.39}\label{eq:endpoint-determinant-cancellation}
\]
Consequently,
\[
\Pp(I_{n,a_1}=\cdots=I_{n,a_r}=1)
=(1+o(1))\prod_{j=1}^r
\frac{\varphi(g_j^*)}{M_j'(g_j^*)}.
\tag{8.40}\label{eq:endpoint-saddle-product}
\]
By \eqref{eq:endpoint-saddle-asymptotics},
\[
        \varphi(g_j^*)\longrightarrow
        K_\gamma e^{-K_\gamma},
        \qquad
        M_j'(g_j^*)\sim a_j/\gamma.
\]
Each factor in \eqref{eq:endpoint-saddle-product} is therefore asymptotic to
\(c_\gamma/a_j\), uniformly over the admissible family.  This proves
\eqref{eq:endpoint-separated-joint}.
\end{proof}


\begin{theorem}[Endpoint fixed-point universality]
\label{thm:endpoint-fixed-universality}
Assume the terminal rate array is uniformly endpoint regularly varying with
index \(\gamma>0\) in the sense of Definition \ref{def:endpoint-rv}.  Let
\[
        \Xi_n=\sum_{a=1}^n I_{n,a}\delta_{\log(a/b_n)},
\]
where \(b_n\to\infty\) and \(b_n=o(n)\).  For every bounded interval
\(J\subset\mathbb R\), the restriction \(\Xi_n|_J\), regarded as a random
locally finite counting measure with the vague topology, satisfies
\[
        \Xi_n|_J\Rightarrow \PPP(c_\gamma\dd x)|_J.
\tag{8.41}\label{eq:endpoint-ppp-limit}
\]
Equivalently, counts in finitely many disjoint bounded intervals converge
jointly to independent Poisson random variables with means \(c_\gamma\) times
their lengths.  Moreover, for every fixed \(r\ge1\),
\[
        \Ee(F_n)_r\sim(c_\gamma\log n)^r.
\tag{8.42}\label{eq:endpoint-factorial-moments}
\]
In particular,
\[
        \Ee C_1(\sigma_n)\sim c_\gamma\log n,
        \qquad
        \frac{C_1(\sigma_n)}{c_\gamma\log n}
        \longrightarrow1
        \quad\text{in probability}.
\tag{8.43}\label{eq:endpoint-fixed-lln}
\]
\end{theorem}

\begin{proof}
We first prove the local point-process assertion.  For every bounded interval
\(I=(u,v]\), the elementary logarithmic Riemann sum gives
\[
        \sum_{\log(a/b_n)\in I}\frac1a\longrightarrow v-u=|I|.
\tag{8.44}\label{eq:endpoint-harmonic-window}
\]
All coordinates in this sum tend to infinity and are \(o(n)\).  Fix disjoint
bounded intervals \(I_1,\ldots,I_m\) and nonnegative integers
\(k_1,\ldots,k_m\), and put \(k=\sum_\ell k_\ell\).  Expand
\[
        \Ee\prod_{\ell=1}^m(\Xi_n(I_\ell))_{k_\ell}
\]
over ordered distinct terminal coordinates, colored according to their
intervals.  Restrict first to tuples whose logarithms are pairwise more than
\(\eta\) apart.  After sorting the coordinates, Proposition
\ref{prop:endpoint-rv-joint} applies with ratio-separation parameter
\(e^\eta-1\), uniformly over the entire restricted sum.  Its contribution is
therefore
\[
        c_\gamma^k(1+o(1))
        \sum_{\eta\text{-separated}}^*
        \frac1{a_1\cdots a_k}.
\tag{8.45}\label{eq:endpoint-separated-window-sum}
\]

The logarithmic diagonal neighborhood has negligible harmonic weight as
\(\eta\downarrow0\).  Indeed, uniformly for \(a\) in a fixed multiplicative
window,
\[
        \sum_{b:\,|\log b-\log a|\le\eta}\frac1b
        \le C\eta+o(1),
\tag{8.46}\label{eq:endpoint-log-diagonal-bound}
\]
and hence a union bound over pairs gives an \(O_{k,I_1,\ldots,I_m}(\eta)+o(1)\)
bound for the weighted sum of excluded tuples.  Lemma
\ref{lem:endpoint-rv-product-bound} gives the same bound, up to a fixed
constant, for their probability contribution.  Using
\eqref{eq:endpoint-harmonic-window} in
\eqref{eq:endpoint-separated-window-sum} and then letting
\(\eta\downarrow0\), we obtain
\[
        \Ee\prod_{\ell=1}^m(\Xi_n(I_\ell))_{k_\ell}
        \longrightarrow
        \prod_{\ell=1}^m(c_\gamma|I_\ell|)^{k_\ell}.
\tag{8.47}\label{eq:endpoint-mixed-window-factorials}
\]
These are the mixed factorial moments of independent Poisson variables, so the
interval-count vectors converge to the asserted independent Poisson vector.

We now upgrade count convergence to convergence of point measures.  By the
one-point case of Lemma \ref{lem:endpoint-rv-product-bound},
\[
        \sup_n\Ee\Xi_n(J)
        \le C\sup_n\sum_{\log(a/b_n)\in J}\frac1a<\infty
\tag{8.48}\label{eq:endpoint-ppp-tightness}
\]
for every bounded \(J\), after ignoring finitely many initial \(n\).  Let
\(f\ge0\) be continuous with compact support in \(J\).  Choose a finite
partition of \(J\) into half-open intervals and a step function \(f_\pi\),
constant on the partition cells, with
\(\|f-f_\pi\|_\infty\to0\) as the mesh tends to zero.  Since
\(|e^{-x}-e^{-y}|\le|x-y|\) for \(x,y\ge0\),
\[
\left|
\Ee e^{-\langle\Xi_n,f\rangle}
-\Ee e^{-\langle\Xi_n,f_\pi\rangle}
\right|
\le\|f-f_\pi\|_\infty\Ee\Xi_n(J),
\tag{8.49}\label{eq:endpoint-laplace-step-error}
\]
uniformly in \(n\).  For fixed \(\pi\), the interval-count convergence just
proved gives
\[
\Ee e^{-\langle\Xi_n,f_\pi\rangle}
\longrightarrow
\exp\left[-c_\gamma\int_J(1-e^{-f_\pi(x)})\dd x\right].
\]
Letting the mesh tend to zero in view of
\eqref{eq:endpoint-laplace-step-error} yields
\[
\Ee e^{-\langle\Xi_n,f\rangle}
\longrightarrow
\exp\left[-c_\gamma\int_J(1-e^{-f(x)})\dd x\right],
\tag{8.50}\label{eq:endpoint-laplace-functional}
\]
the Laplace functional of \(\PPP(c_\gamma\dd x)|_J\).  Laplace functionals
on nonnegative compactly supported continuous functions determine the law of a
locally finite counting measure and are convergence determining for the vague
topology.  This proves \eqref{eq:endpoint-ppp-limit}; the uniform first-moment
bound \eqref{eq:endpoint-ppp-tightness} also gives tightness directly.

It remains to prove the full factorial moments.  Put
\[
        L_n=\lceil\log n\rceil,
        \qquad
        U_n=\lfloor n/\log n\rfloor,
        \qquad
        H_n=\sum_{L_n\le a\le U_n}\frac1a.
\]
Then \(H_n\sim\log n\).  Fix once and for all a positive separation
\(\eta_0\).  The harmonic weight of ordered \(r\)-tuples in
\([L_n,U_n]\) containing a pair with
\(|\log a_i-\log a_j|\le\eta_0\) is \(O_r(H_n^{r-1})\).  To see this, fix
one coordinate \(a_i\), use
\[
        \sum_{b:\,|\log b-\log a_i|\le\eta_0}\frac1b
        \le C_{\eta_0},
\]
and sum the remaining \(r-2\) coordinates freely.  Repeated coordinates are
included in this bound.  Therefore the separated weighted sum is
\[
        H_n^r+O_r(H_n^{r-1}).
\tag{8.51}\label{eq:endpoint-global-separated-weight}
\]
On the separated tuples, Proposition \ref{prop:endpoint-rv-joint} applies
uniformly and gives the contribution
\[
        c_\gamma^rH_n^r+o(H_n^r).
\]
The nonseparated probability contribution is \(O_r(H_n^{r-1})\) by Lemma
\ref{lem:endpoint-rv-product-bound} and the preceding harmonic estimate.
Thus, for
\(G_n=\sum_{L_n\le a\le U_n}I_{n,a}\),
\[
        \Ee(G_n)_r\sim(c_\gamma\log n)^r.
\tag{8.52}\label{eq:endpoint-truncated-factorials}
\]

Finally, the harmonic mass outside \([L_n,U_n]\) is
\[
        H_n^{\mathrm{out}}
        :=\sum_{a<L_n}\frac1a+
          \sum_{a>U_n}\frac1a
        =O(\log\log n).
\]
By the product envelope, the total contribution to \(\Ee(F_n)_r\) from
ordered tuples containing at least one omitted coordinate is at most
\[
        C_rH_n^{\mathrm{out}}
        \left(\sum_{a=1}^n\frac1a\right)^{r-1}
        =O_r((\log n)^{r-1}\log\log n)
        =o((\log n)^r).
\]
Combining this with \eqref{eq:endpoint-truncated-factorials} proves
\eqref{eq:endpoint-factorial-moments}.

Taking \(r=1,2\) gives
\(\Ee F_n\sim c_\gamma\log n\) and
\(\Var(F_n)=o((\log n)^2)\).  Chebyshev's inequality yields
\eqref{eq:endpoint-fixed-lln}.
\end{proof}


\subsection{Further endpoint limits}

The same regular-variation calculation also gives two boundary limits that are
useful for interpreting the endpoint process.  The first is a mesoscopic
one-coordinate displacement law; the second is a microscopic infinite random
permutation.

\begin{theorem}[Endpoint displacement law]\label{thm:endpoint-displacement}
Assume Definition \ref{def:endpoint-rv}.  Let
$a_{n,1},\ldots,a_{n,m}$ be distinct terminal coordinates, where $m$ is fixed,
$a_{n,j}\to\infty$, and $a_{n,j}=o(n)$ for every $j$.  Then
\[
        \left(\frac{\tau_n(a_{n,1})}{a_{n,1}},\ldots,
        \frac{\tau_n(a_{n,m})}{a_{n,m}}\right)
        \Rightarrow
        \left(A_\gamma E_1^{-1/\gamma},\ldots,
        A_\gamma E_m^{-1/\gamma}\right),
\]
where $E_1,\ldots,E_m$ are independent $\Exp(1)$ variables and
$A_\gamma=\Gamma(1+1/\gamma)$.  Equivalently,
\[
        \log\frac{\tau_n(a_{n,j})}{a_{n,j}}
        \Rightarrow X_\gamma,
\]
where $X_\gamma$ has density
\[
        k_\gamma(x)=\gamma K_\gamma e^{-\gamma x}
        \exp\{-K_\gamma e^{-\gamma x}\},
        \qquad K_\gamma=A_\gamma^\gamma.
\]
\end{theorem}

\begin{proof}
Write
\[
        Y_{n,a}=\log\lambda_{n,a}+G_a,
\]
with independent standard Gumbels.  Conditional on $G_a=g$,
\[
        \tau_n(a)-1
        =\sum_{b\ne a}\one_{\{Y_{n,b}<Y_{n,a}\}}
\]
is a sum of independent Bernoulli variables with mean
\[
        M_{n,a}(g)=\sum_{b\ne a}
        \exp\left(-e^{-g}\frac{\lambda_{n,b}}{\lambda_{n,a}}\right).
\]
Lemma \ref{lem:endpoint-rv-sums} gives, uniformly for $g$ in compact sets,
\[
        \frac{M_{n,a}(g)}a\longrightarrow
        A_\gamma e^{g/\gamma}.
\]
The conditional variance is at most the conditional mean, hence is $O(a)$ on
compact $g$-sets.  Chebyshev's inequality therefore gives, uniformly on such
sets,
\[
        \frac{\tau_n(a)}a-A_\gamma e^{G_a/\gamma}
        \longrightarrow0
        \quad\text{in probability}.
\]
Truncating the Gumbel variables to a fixed compact set and then letting the
truncation grow proves the one-coordinate assertion.  For finitely many
specified coordinates, delete them from the Bernoulli sums.  Lemma
\ref{lem:endpoint-rv-sums} is unchanged after finitely many deletions, and a
union bound over the conditional Chebyshev estimates proves the joint
convergence.  The limiting variables are independent because the specified
Gumbels are independent.

Finally, if $E=e^{-G}\sim\Exp(1)$, then
$A_\gamma e^{G/\gamma}=A_\gamma E^{-1/\gamma}$.  The change of variables
$x=\log A_\gamma+G/\gamma$ gives the displayed density $k_\gamma$.
\end{proof}

\subsection{Mesoscopic endpoint tangent measure}
\label{subsec:endpoint-tangent}

The microscopic limit in Theorem \ref{thm:microscopic-endpoint-permutation} is
an infinite random permutation.  On every mesoscopic scale containing
$b_n\to\infty$ terminal coordinates, the correct object is instead a dense
normalized graph measure.  The unnormalized graph cannot have a locally finite
point-process limit, because a compact source window contains order $b_n$
points.

Define
\[
        \mathcal K_\gamma(x,y)
        =\frac1y k_\gamma\!\left(\log\frac yx\right)
        =\gamma K_\gamma\frac{x^\gamma}{y^{\gamma+1}}
        \exp\left\{-K_\gamma\left(\frac xy\right)^\gamma\right\},
        \qquad x,y>0.
\]

\begin{theorem}[Mesoscopic endpoint tangent measure]
\label{thm:endpoint-tangent-measure}
Assume Definition \ref{def:endpoint-rv}.  Let $b_n\to\infty$ with
$b_n=o(n)$, and put
\[
        \mathcal M_n^{(b_n)}
        =\frac1{b_n}\sum_{a=1}^n
        \delta_{(a/b_n,\tau_n(a)/b_n)}.
\]
Each \(\mathcal M_n^{(b_n)}\) is a Radon measure on compact subsets of
\((0,\infty)^2\), and the density \(\mathcal K_\gamma\) is locally
integrable there.  Then, vaguely in probability on $(0,\infty)^2$,
\[
        \mathcal M_n^{(b_n)}
        \longrightarrow
        \mathcal K_\gamma(x,y)\dd x\dd y.
\]
More precisely, for every $\varphi\in C_c((0,\infty)^2)$,
\[
\begin{aligned}
        \Ee\left|
        \langle\mathcal M_n^{(b_n)},\varphi\rangle
        -\int_0^\infty\int_0^\infty
        \varphi(x,y)\mathcal K_\gamma(x,y)\dd y\dd x
        \right|\longrightarrow0.
\end{aligned}
\]
The limiting kernel is bistochastic with respect to Lebesgue measure.
\end{theorem}

\begin{lemma}[Uniform mesoscopic Laplace sums]
\label{lem:mesoscopic-laplace-sums}
Fix $0<x_-<x_+<\infty$ and $0<u_-<u_+<\infty$.  Uniformly for
$x_-b_n\le a\le x_+b_n$ and $u\in[u_-,u_+]$,
\[
        \frac1{b_n}\sum_{c=1}^n
        \exp\left\{-u\frac{\lambda_{n,c}}{\lambda_{n,a}}\right\}
        =\frac a{b_n}A_\gamma u^{-1/\gamma}+o(1).
\]
For every fixed $u_->0$, the left-hand side is uniformly bounded for all
$u\ge u_-$.
\end{lemma}

\begin{proof}
Put $X_n=n/b_n\to\infty$ and choose
\[
        R_n=(\log(e+X_n))^{4/\gamma},
        \qquad B_n=\lceil2x_+b_nR_n\rceil.
\]
Then $B_n=o(n)$.  Uniform endpoint regular variation applies to all
$a\in[x_-b_n,x_+b_n]$ and $c\le B_n$.  On every fixed multiplicative interval
$\delta a\le c\le Ra$, the uniform convergence theorem for slowly varying
functions gives
\[
        \lambda_{n,c}/\lambda_{n,a}=(c/a)^\gamma(1+o(1))
\]
uniformly.  Hence
\[
\begin{aligned}
        \frac1{b_n}\sum_{\delta a\le c\le Ra}
        e^{-u\lambda_{n,c}/\lambda_{n,a}}
        =\frac a{b_n}\int_\delta^R e^{-uv^\gamma}\dd v+o(1).
\end{aligned}
\]
The range $c<\delta a$ contributes at most $x_+\delta$.  Potter's bound with
exponent $\gamma/2$ gives, for $Ra<c\le B_n$,
\[
        \lambda_{n,c}/\lambda_{n,a}
        \ge c(c/a)^{\gamma/2},
\]
so the normalized contribution is bounded by
\[
        C\frac a{b_n}\int_R^\infty e^{-cu_-v^{\gamma/2}}\dd v.
\]
For $c>B_n$, monotonicity and Potter's bound give
\[
        \lambda_{n,c}/\lambda_{n,a}\ge cR_n^{\gamma/2},
\]
and therefore
\[
        \frac1{b_n}\sum_{c>B_n}
        e^{-u\lambda_{n,c}/\lambda_{n,a}}
        \le X_n\exp\{-cu_-R_n^{\gamma/2}\}=o(1).
\]
Let first $n\to\infty$, then $\delta\downarrow0$ and $R\uparrow\infty$.
Since
\[
        \int_0^\infty e^{-uv^\gamma}\dd v
        =A_\gamma u^{-1/\gamma},
\]
the first assertion follows.  Monotonicity in $u$ gives the final bound.
\end{proof}

Let $E_{n,a}=\lambda_{n,a}T_{n,a}$.  These variables are independent and
standard exponential.  Conditional on $E_{n,a}=u$,
\[
        \tau_n(a)=1+\sum_{c\ne a}\xi_{n,a,c}(u),
        \qquad
        \Pp(\xi_{n,a,c}(u)=1)
        =e^{-u\lambda_{n,c}/\lambda_{n,a}},
\]
with independent summands.  Put
\[
        r_{n,a}(u)=1+\sum_{c\ne a}
        e^{-u\lambda_{n,c}/\lambda_{n,a}}.
\]

\begin{lemma}[Conditional rank self-averaging]
\label{lem:conditional-rank-self-averaging}
Fix $0<x_-<x_+<\infty$.  For every $\varepsilon,\delta>0$,
\[
\begin{aligned}
        \sup_{x_-b_n\le a\le x_+b_n}
        \Pp\bigl(|\tau_n(a)-r_{n,a}(E_{n,a})|>\delta b_n\bigr)
        \le\varepsilon+\frac{C_\varepsilon}{\delta^2b_n}
\end{aligned}
\]
for all large $n$.  Moreover, uniformly for
$a/b_n\in[x_-,x_+]$ and $u$ in a compact subset of $(0,\infty)$,
\[
        \frac{r_{n,a}(u)}{b_n}
        \longrightarrow
        \frac a{b_n}A_\gamma u^{-1/\gamma}.
\]
\end{lemma}

\begin{proof}
Conditionally on $E_{n,a}=u$,
\[
        \Var(\tau_n(a)\mid E_{n,a}=u)\le r_{n,a}(u).
\]
On $\{E_{n,a}\ge\varepsilon\}$, Lemma
\ref{lem:mesoscopic-laplace-sums} and monotonicity in $u$ give
$r_{n,a}(E_{n,a})\le C_\varepsilon b_n$.  Conditional Chebyshev therefore
bounds the deviation probability on this event by
$C_\varepsilon/(\delta^2b_n)$.  The complementary event has probability at
most $\varepsilon$.  The second assertion follows immediately from Lemma
\ref{lem:mesoscopic-laplace-sums}; omitting the $c=a$ summand and adding the
leading one changes the sum by $O(1)$.
\end{proof}

\begin{proof}[Proof of Theorem \ref{thm:endpoint-tangent-measure}]
Fix $\varphi\in C_c((0,\infty)^2)$, and choose
$0<x_-<x_+<\infty$ containing the first-coordinate projection of its support.
Set
\[
        Z_n(\varphi)=\frac1{b_n}\sum_a
        \varphi(a/b_n,\tau_n(a)/b_n)
\]
and
\[
        \widehat Z_n(\varphi)=\frac1{b_n}\sum_a
        \varphi(a/b_n,r_{n,a}(E_{n,a})/b_n).
\]
Only $O(b_n)$ summands are nonzero.  Uniform continuity of $\varphi$ and Lemma
\ref{lem:conditional-rank-self-averaging} give
\[
        \Ee|Z_n(\varphi)-\widehat Z_n(\varphi)|\longrightarrow0.
\]
Indeed, split according as the two second coordinates differ by at most
$\delta$, use the modulus of continuity on that event, and then let
$n\to\infty$, $\delta\downarrow0$, and $\varepsilon\downarrow0$.

The summands in $\widehat Z_n(\varphi)$ are independent functions of the
independent variables $E_{n,a}$.  Hence
\[
        \Var(\widehat Z_n(\varphi))\le C_\varphi/b_n\longrightarrow0.
\]
It remains to compute the expectations.  If $U\sim\Exp(1)$, Lemma
\ref{lem:conditional-rank-self-averaging}, uniformly on compact $u$-sets,
gives
\[
\begin{aligned}
        \Ee\widehat Z_n(\varphi)
        &\longrightarrow
        \int_0^\infty\int_0^\infty
        \varphi(x,xA_\gamma u^{-1/\gamma})e^{-u}\dd u\dd x.
\end{aligned}
\]
To justify the passage, truncate $u$ to $[\varepsilon,M]$, use uniform
convergence and an ordinary Riemann sum, and then use the exponential mass of
$(0,\varepsilon)\cup(M,\infty)$.

For fixed $x$, make the change of variables
$y=xA_\gamma u^{-1/\gamma}$.  Then
$u=K_\gamma(x/y)^\gamma$ and
\[
        e^{-u}\left|\frac{\dd u}{\dd y}\right|
        =\mathcal K_\gamma(x,y).
\]
This proves the $L^1$ convergence of test-function integrals.  A countable
convergence-determining subset of $C_c((0,\infty)^2)$ then gives vague
convergence in probability.

Finally, the first marginal identity follows from the same change of variables:
$\mathcal K_\gamma(x,\cdot)$ is the density of
$xA_\gamma U^{-1/\gamma}$.  For the second marginal, put
$v=K_\gamma(x/y)^\gamma$ to obtain
\[
\begin{aligned}
        \int_0^\infty\mathcal K_\gamma(x,y)\dd x
        &=K_\gamma^{-1/\gamma}
        \int_0^\infty v^{1/\gamma}e^{-v}\dd v=1.
\end{aligned}
\]
Thus the kernel is bistochastic.
\end{proof}

\begin{corollary}[Rectangle law]
\label{cor:endpoint-tangent-rectangles}
For bounded intervals $I,J\Subset(0,\infty)$,
\[
        \frac1{b_n}\#\left\{a:
        a/b_n\in I,\ \tau_n(a)/b_n\in J\right\}
        \xrightarrow{\Pp}
        \int_I\int_J\mathcal K_\gamma(x,y)\dd y\dd x.
\]
\end{corollary}

\begin{proof}
The boundary of $I\times J$ has zero mass under the continuous limiting
density, so the assertion follows from vague convergence and approximation by
compactly supported continuous functions.
\end{proof}

\begin{proposition}[Multiplicative semigroup and diagonal traces]
\label{prop:endpoint-kernel-powers}
For every $r\ge1$,
\[
        \mathcal K_\gamma^{\,r}(x,y)
        =\frac1y k_\gamma^{*r}\!\left(\log\frac yx\right).
\]
In particular,
\[
        \mathcal K_\gamma^{\,r}(x,x)
        =\frac1x k_\gamma^{*r}(0).
\]
\end{proposition}

\begin{proof}
For $r=2$, substitute $z=xe^w$ in the composition integral:
\[
\begin{aligned}
        \int_0^\infty\mathcal K_\gamma(x,z)
        \mathcal K_\gamma(z,y)\dd z
        &=\frac1y\int_{\mathbb R}k_\gamma(w)
        k_\gamma\!\left(\log\frac yx-w\right)\dd w.
\end{aligned}
\]
This is $y^{-1}k_\gamma^{*2}(\log(y/x))$.  Iterate.
\end{proof}


\subsection{Microscopic endpoint permutation}\label{subsec:microscopic-endpoint}

\begin{theorem}[Microscopic endpoint permutation]
\label{thm:microscopic-endpoint-permutation}
Assume Definition \ref{def:endpoint-rv}.  Extend the finite permutation
\(\tau_n\in S_n\) to a permutation \(\bar\tau_n\) of \(\mathbb N\) by
\[
        \bar\tau_n(a)=
        \begin{cases}
        \tau_n(a),&a\le n,\\
        a,&a>n.
        \end{cases}
\]
Then, in the product topology on \(\mathbb N^{\mathbb N}\),
\[
        \bar\tau_n\Rightarrow\tau^{(\gamma)},
\tag{8.53}\label{eq:microscopic-endpoint-convergence}
\]
where, for independent \(E_1,E_2,\ldots\sim\Exp(1)\),
\[
        \tau^{(\gamma)}(a)
        =1+\#\left\{b\ne a:
        \frac{E_b}{b^\gamma}>\frac{E_a}{a^\gamma}\right\}.
\]
Almost surely, \(\tau^{(\gamma)}\) is a permutation of \(\mathbb N\).
\end{theorem}

\begin{proof}
Normalize by the first terminal rate and put
\[
        r_{n,b}=\frac{\lambda_{n,b}}{\lambda_{n,1}},
        \qquad
        Q_{n,b}=\lambda_{n,1}T_{n,b}=\frac{E_b}{r_{n,b}},
\]
using independent \(E_b\sim\Exp(1)\).  For every fixed \(B\), uniform
endpoint regular variation and the uniform convergence theorem for slowly
varying functions give
\[
        (r_{n,1},\ldots,r_{n,B})
        \longrightarrow(1,2^\gamma,\ldots,B^\gamma).
\]
Hence, under this coupling,
\[
        (Q_{n,1},\ldots,Q_{n,B})
        \longrightarrow
        (E_1,E_2/2^\gamma,\ldots,E_B/B^\gamma)
\]
almost surely.

We remove the tail uniformly.  Fix \(m\) and \(\varepsilon>0\).  For fixed
\(a\le m\), the ratios \(r_{n,a}\) remain bounded, so
\[
        \Pp(Q_{n,a}<\varepsilon)\le C_m\varepsilon.
\]
Choose \(W_n=\lfloor n/\log n\rfloor\).  Potter's bound, applied with
reference coordinate \(1\), gives
\[
        r_{n,b}\ge cb^{\gamma/2},
        \qquad B<b\le W_n,
\]
for all sufficiently large \(n\).  By monotonicity, the same lower bound at
\(b=W_n\) applies to every \(b>W_n\).  Consequently,
\[
\begin{aligned}
        \sum_{b>B}\Pp(Q_{n,b}>\varepsilon)
        &=\sum_{b>B}e^{-\varepsilon r_{n,b}}\\
        &\le\sum_{b>B}e^{-c\varepsilon b^{\gamma/2}}
        +n e^{-c\varepsilon W_n^{\gamma/2}}.
\end{aligned}
\tag{8.54}\label{eq:microscopic-tail-truncation}
\]
The second term tends to zero, and the first tends to zero as \(B\to\infty\),
uniformly for all sufficiently large \(n\).  Thus, outside an event of
probability at most \(C_m\varepsilon+o_B(1)\), every comparison determining
the ranks of the first \(m\) coordinates is decided by the first \(B\) clocks.
The same argument applies to the limiting clocks \(E_b/b^\gamma\).  Finite
collections of all these clocks have no ties.  Taking first \(n\to\infty\),
then \(B\to\infty\), and finally \(\varepsilon\downarrow0\), proves
finite-dimensional convergence of
\((\bar\tau_n(1),\ldots,\bar\tau_n(m))\) for every fixed \(m\), which is
\eqref{eq:microscopic-endpoint-convergence}.

For every \(t>0\),
\[
        \sum_{b=1}^\infty
        \Pp(E_b/b^\gamma>t)
        =\sum_{b=1}^\infty e^{-b^\gamma t}<\infty.
\]
Borel--Cantelli implies that only finitely many limiting clocks exceed any
positive threshold.  Moreover,
\[
        \Pp\left(\sup_{b\ge1}\frac{E_b}{b^\gamma}>t\right)
        \le\sum_{b\ge1}e^{-b^\gamma t}\longrightarrow0
        \qquad(t\to\infty),
\]
so the clock set has a finite maximum.  The clocks are almost surely distinct,
and zero is their only accumulation point.  Their decreasing enumeration
therefore contains every label exactly once and assigns every positive integer
rank exactly once.  Thus \(\tau^{(\gamma)}\) is almost surely a permutation of
\(\mathbb N\).
\end{proof}


\begin{remark}[The Gumbel origin of the constant]\label{rem:gumbel-endpoint-constant}
The preceding proof makes the constant transparent.  Conditional on
\(Y_{n,a}=\log\lambda_{n,a}+g\), the expected number of lower terminal priorities
is
\[
        \sum_b\exp\left(-e^{-g}\frac{\lambda_{n,b}}{\lambda_{n,a}}\right)
        \sim aA_\gamma e^{g/\gamma}.
\]
The fixed-point saddle is therefore
\(A_\gamma e^{g/\gamma}=1\), so
\(e^{-g}=K_\gamma=A_\gamma^\gamma\).  The density of the specified centered
Gumbel at the saddle is \(K_\gamma e^{-K_\gamma}\), while the derivative of the
mean with respect to \(g\) is \(a/\gamma\).  Their quotient is exactly
\(c_\gamma/a\), with
\[
        c_\gamma=\gamma K_\gamma e^{-K_\gamma}.
\]
For \(\gamma=1\), \(A_\gamma=K_\gamma=1\), and the constant reduces to the
Sukhatme value \(e^{-1}\).
\end{remark}

\section{Global endpoint Poisson approximation}\label{sec:global-endpoint-poisson}

The local point-process theorem and the fixed-order factorial moments do not, by
themselves, imply a Poisson approximation for the total fixed-point count, whose
mean diverges like $c_\gamma\log n$.  We now prove the global approximation in
total variation.  The argument uses an approximate Palm coupling: to force a
specified terminal coordinate $a$ to be fixed, we interchange its clock with the
clock currently occupying terminal rank $a$.  This changes the total number of
fixed points only through the intended point and, exceptionally, a two-cycle.

Retain the notation of Section \ref{sec:endpoint-fixed-universality}.  Thus
$T_{n,a}\sim\Exp(\lambda_{n,a})$ are independent,
\[
        \tau_n(a)=1+\sum_{b\ne a}\one_{\{T_{n,b}>T_{n,a}\}},
        \qquad
        I_{n,a}=\one_{\{\tau_n(a)=a\}},
        \qquad
        F_n=\sum_{a=1}^n I_{n,a}.
\]
Put
\[
        p_{n,a}=\Pp(I_{n,a}=1),
        \qquad
        m_n=\Ee F_n=\sum_{a=1}^n p_{n,a}.
\]

\begin{theorem}[Global endpoint Poisson approximation]\label{thm:global-endpoint-poisson}
Assume Definition \ref{def:endpoint-rv}.  Then
\[
        d_{\TV}\bigl(\mathcal L(F_n),\Pois(m_n)\bigr)\longrightarrow0.
\]
More precisely,
\[
        d_{\TV}\bigl(\mathcal L(F_n),\Pois(m_n)\bigr)
        =O\!\left(\frac{\log\log n}{\sqrt{\log n}}\right).
\]
Since $m_n\sim c_\gamma\log n$, it follows that
\[
        \frac{F_n-m_n}{\sqrt{m_n}}\Rightarrow N(0,1).
\]
For the Sukhatme rates, this gives
\[
        d_{\TV}\bigl(\mathcal L(F_n),\Pois(\Ee F_n)\bigr)\to0,
        \qquad
        \Ee F_n\sim e^{-1}\log n.
\]
\end{theorem}

\subsection{An approximate-Palm Stein bound}

We first isolate the form of the Stein--Chen argument that will be used.  The
standard solution $f=f_B$ of the Poisson Stein equation satisfies
\[
        \|\Delta f\|_\infty\le\frac{1-e^{-\lambda}}{\lambda},
        \qquad
        \|f\|_\infty\le\min\{1,\lambda^{-1/2}\};
\]
see Barbour, Holst, and Janson \cite[Lemma I.1.1]{BHJ1992}.

\begin{lemma}[Stein bound with approximate Palm laws]\label{lem:approx-palm-stein}
Let $X_1,\ldots,X_N$ be indicators,
\[
        W=\sum_{i=1}^N X_i,
        \qquad p_i=\Ee X_i,
        \qquad \lambda=\Ee W>0.
\]
For every $i$ with $p_i>0$, let $Y^{(i)}$ be an integer-valued random variable
constructed on the same probability space as $W$, and put
\[
        \varepsilon_i=
        d_{\TV}\bigl(\mathcal L(Y^{(i)}),\mathcal L(W\mid X_i=1)\bigr).
\]
Indices with $p_i=0$ are omitted.  Then
\[
\begin{aligned}
 d_{\TV}(\mathcal L(W),\Pois(\lambda))
 &\le \frac{1-e^{-\lambda}}{\lambda}
        \sum_{i=1}^N p_i\Ee|W+1-Y^{(i)}|\\
 &\quad+2\min\{1,\lambda^{-1/2}\}
        \sum_{i=1}^N p_i\varepsilon_i.
\end{aligned}
\]
\end{lemma}

\begin{proof}
Fix $B\subset\mathbb Z_{\ge0}$, and let $f=f_B$ solve
\[
        \lambda f(k+1)-kf(k)
        =\one_{\{k\in B\}}-\Pp(\Pois(\lambda)\in B).
\]
Using
\[
        \Ee[Wf(W)]=\sum_i p_i\Ee[f(W)\mid X_i=1],
\]
we obtain
\[
\begin{aligned}
&\Pp(W\in B)-\Pp(\Pois(\lambda)\in B)\\
&\quad=\sum_i p_i\Bigl(
        \Ee\{f(W+1)-f(Y^{(i)})\}
        +\Ee f(Y^{(i)})-\Ee[f(W)\mid X_i=1]
        \Bigr).
\end{aligned}
\]
The first difference is bounded by
$\|\Delta f\|_\infty\Ee|W+1-Y^{(i)}|$, and the second by
$2\|f\|_\infty\varepsilon_i$.  Apply the displayed Stein factors and take the
supremum over $B$.
\end{proof}

\subsection{The terminal rank transposition}

Fix $n$ temporarily and suppress it from the notation.  The clock-tie set has
product-exponential measure zero.  On its complement, let $b_a(t)$ be the
unique label occupying terminal rank $a$, so that $\tau_t(b_a(t))=a$, and
define $\mathcal T_at$ by interchanging $t_a$ and $t_{b_a(t)}$.  Define
$\mathcal T_a$ arbitrarily, say as the identity, on the tie set.  Then label
$a$ has terminal rank $a$ in $\mathcal T_at$.

Let $\nu_a=\mathcal L(\mathcal T_aT)$, and let
$\mu_a^\star=\mathcal L(T\mid I_a=1)$ be the exact Palm law.

\begin{lemma}[Exact swap likelihood ratio]\label{lem:swap-likelihood}
For a configuration $t$ satisfying $I_a(t)=1$, put
\[
        R_a(t)=\sum_{b=1}^n
        \exp\{-(\lambda_a-\lambda_b)(t_b-t_a)\}.
\]
If
\[
        f(t)=\prod_{j=1}^n\lambda_j e^{-\lambda_jt_j}
\]
is the product-exponential density, then
\[
        \frac{\dd\nu_a}{\dd t}
        =\one_{\{I_a=1\}}f(t)R_a(t),
        \qquad
        \frac{\dd\mu_a^\star}{\dd t}
        =\frac{\one_{\{I_a=1\}}}{p_a}f(t).
\]
Consequently,
\[
        \Ee[R_a\mid I_a=1]=p_a^{-1}
\]
and
\[
        d_{\TV}(\nu_a,\mu_a^\star)
        =\frac12\Ee\bigl[|p_aR_a(T)-1|\mid I_a=1\bigr].
\]
\end{lemma}

\begin{proof}
For $1\le b\le n$, let
\[
        \mathcal R_{a,b}=\{t:b_a(t)=b\}
\]
inside the no-tie region, and let $S_{a,b}$ interchange coordinates $a$ and
$b$.  The sets $\mathcal R_{a,b}$ form a measurable partition.  For every
$t$ with $I_a(t)=1$, the unique preimage of $t$ lying in
$\mathcal R_{a,b}$ is $S_{a,b}t$: in that preimage label $b$ occupies rank
$a$, and the defining transposition returns $t$.  Coordinate transpositions
have Jacobian one, and
\[
        \frac{f(S_{a,b}t)}{f(t)}
        =\exp\{-(\lambda_a-\lambda_b)(t_b-t_a)\}.
\]
Changing variables separately on the rank chambers and summing over $b$ gives
the density of $\nu_a$.  The Palm density is immediate.  Integrating the
first density gives $p_a\Ee[R_a\mid I_a=1]=1$, and integrating the absolute
difference of the two densities gives the total-variation identity.
\end{proof}

\subsection{Uniform endpoint estimates}

Put $\beta=\gamma/2$.  Fix a constant $K_0>4/\gamma$, and define
\[
        h_n=K_0\log\log n,
        \qquad
        U_n=\lfloor ne^{-2h_n}\rfloor,
        \qquad
        V_n=\lfloor ne^{-h_n}\rfloor.
\]
Thus $U_n=n(\log n)^{-2K_0+o(1)}$ and
$V_n=n(\log n)^{-K_0+o(1)}$.  For $a\le U_n$, write
\[
        x_{a,b}=\frac{\lambda_{n,b}}{\lambda_{n,a}}.
\]
We use the uniform convergence theorem and Potter bounds for slowly varying
functions in the forms recorded in Bingham, Goldie, and Teugels
\cite[Theorems 1.2.1 and 1.5.6]{BGT1987}.

\begin{lemma}[Uniform endpoint package]\label{lem:global-endpoint-package}
There are constants $A_0<\infty$,
$0<\kappa_-<\kappa_+<\infty$, and $0<c<C<\infty$ such that, for all
sufficiently large $n$ and $A_0\le a\le U_n$, the following hold.

\begin{enumerate}[label=(\roman*)]
\item For $a\le b\le V_n$,
\[
        c(b/a)^\beta\le x_{a,b}
        \le C(b/a)^{\gamma+\beta},
\]
whereas for $1\le b\le a$,
\[
        c(b/a)^{\gamma+\beta}\le x_{a,b}
        \le C(b/a)^\beta.
\]

\item For $x\ge1$,
\[
        \#\{b\le V_n:x_{a,b}\le x\}
        \le Ca(1+x^{1/\beta}).
\]

\item Let
\[
        \mu_a(t)=\sum_{b\ne a}e^{-\lambda_{n,b}t}.
\]
There is a unique $t_a^*>0$ satisfying $\mu_a(t_a^*)=a-1$.  If
$\kappa_a^*=\lambda_{n,a}t_a^*$, then
\[
        \kappa_-\le\kappa_a^*\le\kappa_+.
\]

\item Uniformly for
$\kappa\in[\kappa_-/4,2\kappa_+]$ and $j=0,1,2,3,4$,
\[
        \sum_{b\ne a}x_{a,b}^{\,j}e^{-\kappa x_{a,b}}\le Ca.
\]
For $j=0,1,2$, the same sums are bounded below by $ca$.

\item At $t=t_a^*$,
\[
\begin{aligned}
        v_a(t)&=\sum_{b\ne a}e^{-\lambda_{n,b}t}
        (1-e^{-\lambda_{n,b}t})\asymp a,\\
        D_a(t)&=-\mu_a'(t)
        =\sum_{b\ne a}\lambda_{n,b}e^{-\lambda_{n,b}t}
        \asymp a\lambda_{n,a}.
\end{aligned}
\]
The same comparisons hold whenever
$\lambda_{n,a}t\in[\kappa_-/2,2\kappa_+]$.  Throughout this scaled
neighborhood, at least $ca$ Bernoulli parameters
$e^{-\lambda_{n,b}t}$ belong to a fixed compact subinterval of $(0,1)$.
\end{enumerate}
\end{lemma}

\begin{proof}
Uniform endpoint regular variation gives, uniformly for $a,b\le V_n=o(n)$,
\[
        x_{a,b}
        =\left(\frac ba\right)^\gamma
        \frac{L(n/b)}{L(n/a)}(1+o(1)).
\]
Potter's bound with exponent $\beta=\gamma/2$ gives both inequalities in
(i).  If $b\ge a$ and $x_{a,b}\le x$, the lower bound in (i) implies
$b\le Ca x^{1/\beta}$; adding the at most $a$ indices below $a$ proves (ii).

We next bracket the saddle.  Since $V_n/U_n=e^{h_n}\to\infty$, the interval
$[\lceil a/2\rceil,\lfloor2a\rfloor]$ lies below $V_n$ for all large $n$.
On this interval the ratios $x_{a,b}$ are bounded above by a fixed constant
$C_0$.  Choose $\kappa_->0$ so small that
$e^{-\kappa_-C_0}\ge3/4$.  The interval contains
$3a/2+O(1)$ indices, and hence, after increasing $A_0$,
\[
        \sum_{b\ne a}e^{-\kappa_-x_{a,b}}>a-1.
\]
For the opposite inequality, the reverse Potter bound gives, for $\kappa\ge1$,
\[
\begin{aligned}
        \sum_{b<a}e^{-\kappa x_{a,b}}
        &\le C a\int_0^1e^{-c\kappa u^{\gamma+\beta}}\dd u
        \le Ca\kappa^{-1/(\gamma+\beta)},\\
        \sum_{a\le b\le V_n}e^{-\kappa x_{a,b}}
        &\le C a\int_1^\infty e^{-c\kappa u^\beta}\dd u
        \le Ca\kappa^{-1/\beta}.
\end{aligned}
\]
For $b>V_n$, monotonicity and (i) give
\[
        x_{a,b}\ge x_{a,V_n}
        \ge c(V_n/U_n)^\beta
        =c e^{\beta h_n}
        =c(\log n)^{K_0\beta}.
\]
Because $K_0\beta>2$,
\[
        \sum_{b>V_n}e^{-\kappa x_{a,b}}
        \le n\exp\{-c\kappa(\log n)^{K_0\beta}\}=o(a)
\]
uniformly for $a\ge A_0$ and fixed $\kappa>0$.  Choose a sufficiently large
fixed $\kappa_+$ so that the first two sums at $\kappa_+$ are together at most
$a/2$, and then take $n$ large enough that the last sum is at most $a/4$.
After increasing $A_0$ once more, the total is less than $a-1$.  The
choices above may be made with fixed margins: for some $c_0>0$,
\[
        \mu_a(\kappa_-/\lambda_{n,a})\ge(1+c_0)a,
        \qquad
        \mu_a(\kappa_+/\lambda_{n,a})\le(1-c_0)a.
\]
Since the sum is continuous and strictly decreasing in $\kappa$, (iii)
follows.

For (iv), the contribution of $b<a$ is at most $a$, since monotonicity gives
$x_{a,b}\le1$.  For $a\le b\le V_n$, dyadic summation using (ii) gives
\[
\begin{aligned}
&\sum_{a\le b\le V_n}x_{a,b}^{\,j}e^{-\kappa x_{a,b}}\\
&\qquad\le Ca\sum_{m\ge0}
        2^{m(j+1/\beta)}e^{-c\kappa2^m}
        \le C_ja
\end{aligned}
\]
uniformly for $\kappa\ge\kappa_-/4$.  On $b>V_n$, the preceding lower bound
on $x_{a,b}$ and the elementary inequality
$x^je^{-\kappa x}\le C_{j,\kappa_-}e^{-(\kappa/2)x}$ for large $x$ give an
$o(a)$ contribution.  This proves the upper bounds.  Restricting to a fixed
subinterval of $[a,2a]$ gives $ca$ terms for which $x_{a,b}$ stays in a fixed
compact subset of $(0,\infty)$, proving the lower bounds for $j=0,1,2$.

Part (v) follows from (iv) after putting $\kappa=\lambda_{n,a}t$.  The lower
variance bound and the final nondegeneracy assertion follow by restricting to
the same fixed subinterval of $[a,2a]$.
\end{proof}

\begin{lemma}[Likelihood-ratio tail sum]\label{lem:global-likelihood-tail}
There are constants $C,c>0$ such that, uniformly under the hypotheses of Lemma
\ref{lem:global-endpoint-package}, for every $M\ge e^2$ and every
$\kappa\in[\kappa_-/2,2\kappa_+]$,
\[
\sum_{\substack{b:\lambda_{n,b}>\lambda_{n,a}}}
(x_{a,b}-1)e^{-\kappa x_{a,b}}
M^{-1/(x_{a,b}-1)}
\le Ca e^{-c\sqrt{\log M}}.
\]
\end{lemma}

\begin{proof}
Put $L=\log M$.  For $x>1$ and
$\kappa\in[\kappa_-/2,2\kappa_+]$,
\[
\begin{aligned}
        \kappa x+\frac{L}{x-1}
        &=\frac{\kappa x}{2}
          +\left(\frac{\kappa x}{2}+\frac{L}{x-1}\right)\\
        &\ge \frac{\kappa x}{2}
          +2\sqrt{\frac{\kappa Lx}{2(x-1)}}
        \ge \frac{\kappa x}{2}+c\sqrt L.
\end{aligned}
\]
Therefore
\[
        (x-1)e^{-\kappa x}M^{-1/(x-1)}
        \le C e^{-c\sqrt{\log M}}x e^{-(\kappa/2)x}.
\]
Summing this inequality and applying Lemma
\ref{lem:global-endpoint-package}(iv) with $j=1$ and
$\kappa/2\in[\kappa_-/4,\kappa_+]$ proves the claim.  This argument includes
the entire rate range and is uniform in $M$.
\end{proof}

\subsection{A quantitative one-dimensional saddle}

For fixed $a$, let
\[
        S_a(t)=\sum_{b\ne a}\xi_b(t),
        \qquad
        \xi_b(t)\sim\operatorname{Bernoulli}
        (e^{-\lambda_{n,b}t})
\]
independently.  Thus $I_{n,a}=1$ is obtained by integrating the event
$S_a(t)=a-1$ against the density of $T_{n,a}$.

\begin{lemma}[Quantitative endpoint saddle]\label{lem:global-one-saddle}
There are constants $\rho,c,C>0$ such that, uniformly for
$A_0\le a\le U_n$,
\[
        p_{n,a}
        =\frac{\lambda_{n,a}e^{-\lambda_{n,a}t_a^*}}
        {D_a(t_a^*)}
        \left(1+O\left(a^{-\rho}
        +e^{-c\sqrt{\log(a+2)}}\right)\right).
\]
Put
\[
        u=\lambda_{n,a}\sqrt a\,(t-t_a^*),
        \qquad
        L_a=(\log(a+2))^{1/4}.
\]
Under the Palm law $\mathcal L(T_{n,a}\mid I_{n,a}=1)$,
\[
        \Pp(|u|>L_a\mid I_{n,a}=1)
        \le Ce^{-c\sqrt{\log(a+2)}}.
\]
Finally, whenever $|u|\le L_a$, uniformly in $b\ne a$ and
$s\in\{0,1\}$,
\[
        \frac{
        \Pp(S_a(t)-\xi_b(t)=a-1-s)}
        {\Pp(S_a(t)=a-1)}
        =1+O(a^{-\rho}).
\]
\end{lemma}

\begin{proof}
Suppress $n$ and write
\[
        q_b(t)=e^{-\lambda_bt},
        \qquad
        \mu(t)=\sum_{b\ne a}q_b(t),
        \qquad
        v(t)=\sum_{b\ne a}q_b(t)(1-q_b(t)).
\]
In the scaled neighborhood
$\lambda_at\in[\kappa_-/2,2\kappa_+]$, Lemma
\ref{lem:global-endpoint-package} gives
\[
        v(t)\asymp a,
        \qquad
        -\mu'(t)\asymp a\lambda_a,
        \qquad
        |\mu^{(j)}(t)|\le Ca\lambda_a^j
        \quad(j=2,3,4),
\]
and at least $ca$ of the $q_b(t)$ lie in a fixed compact subinterval of
$(0,1)$.

We first prove a uniform local expansion.  Let
\[
        \psi_t(\theta)=
        \Ee e^{i\theta(S_a(t)-\mu(t))}.
\]
The nondegenerate Bernoulli block gives, for $|\theta|\le\pi$,
\[
        |\psi_t(\theta)|\le e^{-ca\theta^2}.
\]
Indeed,
$|1-q+qe^{i\theta}|^2=1-4q(1-q)\sin^2(\theta/2)$, and
$\sin^2(\theta/2)\ge c\theta^2$ on $[-\pi,\pi]$.  Uniformly for
$|\theta|\le\theta_0$,
\[
        \log\psi_t(\theta)
        =-\frac12v(t)\theta^2+O(a|\theta|^3).
\]
Fourier inversion, split at $|\theta|=a^{-2/5}$, therefore gives
\begin{equation}
\begin{aligned}
        \Pp(S_a(t)=\ell)
        &=\frac{1}{\sqrt{2\pi v(t)}}
        \exp\left\{-\frac{(\ell-\mu(t))^2}{2v(t)}\right\}
        +O(a^{-1})
\end{aligned}
\label{eq:global-one-local-expansion}
\end{equation}
uniformly for $|\ell-\mu(t)|\le L_a\sqrt a$ in the central window.  On the
inner Fourier interval, the integral of the cubic error is
$O(a\int |\theta|^3e^{-ca\theta^2}\dd\theta)=O(a^{-1})$; the complementary
Fourier integral is exponentially small.  The same estimate holds after
deleting any one summand, with its own mean and variance.

The variance derivative satisfies
\[
        |v'(t)|
        \le \sum_{b\ne a}\lambda_b e^{-\lambda_bt}
        =\lambda_a\sum_{b\ne a}x_{a,b}e^{-\kappa x_{a,b}}
        \le Ca\lambda_a,
        \qquad \kappa=\lambda_at,
\]
by Lemma \ref{lem:global-endpoint-package}(iv).  Hence, if $|u|\le L_a$,
Taylor's theorem gives
\[
        \frac{a-1-\mu(t)}{\sqrt{v(t)}}
        =c_a u+O\left(\frac{1+u^2}{\sqrt a}\right),
\]
where $c_a$ is bounded above and below by positive constants, and
\[
        \frac{v(t)}{v(t_a^*)}=1+O\left(\frac{1+|u|}{\sqrt a}\right).
\]
The Gaussian main term in \eqref{eq:global-one-local-expansion} is at least
$a^{-1/2}e^{-CL_a^2}$.  Hence its absolute $O(a^{-1})$ error is a relative
$O(a^{-\rho})$ error for some fixed $\rho>0$, because
$L_a^2=\sqrt{\log(a+2)}=o(\log a)$.  Deleting one summand changes the mean and
variance by at most one, so the ratio of the two Gaussian main terms differs
from one by
\[
        O\left(\frac{1+|a-1-\mu(t)|}{a}\right)
        =O\left(\frac{L_a}{\sqrt a}\right).
\]
After decreasing $\rho$, this proves the final assertion of the lemma.

We also need a local upper bound.  Put $h=a-1-\mu(t)$ and let
\[
        K_t(\theta)=\log\Ee\exp\{\theta(S_a(t)-\mu(t))\}.
\]
For $|\theta|\le\theta_0$, bounded Bernoulli summands and Lemma
\ref{lem:global-endpoint-package}(iv) give
\[
        K_t''(\theta)=\Var_\theta S_a(t)\le Ca,
        \qquad K_t(0)=K_t'(0)=0,
\]
and therefore $K_t(\theta)\le Ca\theta^2$.  If $|h|\le c_0a$, choose
$\theta=\delta h/a$, with $\delta>0$ sufficiently small.  Exponential tilting
gives
\[
\begin{aligned}
        \Pp(S_a(t)=a-1)
        &=e^{-\theta h+K_t(\theta)}
          \Pp_\theta(S_a(t)=a-1)\\
        &\le Ca^{-1/2}
          \exp\left\{-c\frac{h^2}{a}\right\}.
\end{aligned}
\]
Here the maximal atom under $\Pp_\theta$ is $O(a^{-1/2})$: the tilt is bounded,
so the block of $ca$ Bernoulli parameters in a fixed compact subset of
$(0,1)$ remains nondegenerate, and the same Fourier estimate used above
applies.  If $c_0a<|h|\le Ca$, use instead a fixed tilt
$\theta=\delta\operatorname{sgn}(h)$; choosing $\delta$ sufficiently small
gives exponent at most $-c|h|\le-c'a$, which is bounded above by
$-c''h^2/a$.  Thus, throughout the fixed scaled neighborhood,
\begin{equation}
        \Pp(S_a(t)=a-1)
        \le Ca^{-1/2}\exp\left\{-c\frac{h^2}{a}\right\}.
\label{eq:global-one-gaussian-upper}
\end{equation}
The construction of $\kappa_-$ and $\kappa_+$ in Lemma
\ref{lem:global-endpoint-package} in fact gives fixed positive margins at the
two bracket points.  Monotonicity of $\mu$ therefore implies $|h|\ge ca$
outside the scaled neighborhood, and Bernstein's inequality makes that
integrated contribution exponentially small.

Put
\[
        v_*=v(t_a^*),
        \qquad D_*=D_a(t_a^*),
        \qquad
        w=\frac{D_*}{\sqrt{v_*}}(t-t_a^*).
\]
On $|u|\le L_a$, equivalently $|w|\le C L_a$, Taylor's theorem and
\eqref{eq:global-one-local-expansion} give
\[
\begin{aligned}
&\lambda_ae^{-\lambda_at}
 \Pp(S_a(t)=a-1)\dd t\\
&\quad=\frac{\lambda_ae^{-\lambda_at_a^*}}{D_*}
        \frac{e^{-w^2/2}}{\sqrt{2\pi}}
        \left(1+O(a^{-\rho})\right)\dd w,
\end{aligned}
\]
after decreasing $\rho$ once more.  The error includes the Taylor remainders,
which are polynomially small uniformly on the chosen polylogarithmic window.
By \eqref{eq:global-one-gaussian-upper}, the remainder of the fixed scaled
neighborhood contributes at most
\[
        C\frac{\lambda_ae^{-\lambda_at_a^*}}{D_*}e^{-cL_a^2},
\]
and the region outside that neighborhood is exponentially smaller.  Therefore
\[
\begin{aligned}
        p_{n,a}
        &=\int_0^\infty \lambda_ae^{-\lambda_at}
        \Pp(S_a(t)=a-1)\dd t\\
        &=\frac{\lambda_ae^{-\lambda_at_a^*}}{D_*}
        \left(1+O\left(a^{-\rho}+e^{-cL_a^2}\right)\right).
\end{aligned}
\]
Since $L_a^2=\sqrt{\log(a+2)}$, this is the first assertion.  Dividing the
same central-tail estimate by the saddle asymptotic proves the Palm
concentration bound.
\end{proof}

The preceding lemma and Lemma \ref{lem:endpoint-rv-sums} give, uniformly on
endpoint ranges whose lower endpoint tends to infinity and whose upper endpoint
is $o(n)$,
\[
        p_{n,a}=\frac{c_\gamma+o(1)}{a}.
\]
We shall also use the global bound
\[
        p_{n,a}\le\frac{C}{a},
        \qquad 1\le a\le n,
\]
which follows by exposing the death process and applying Lemma
\ref{lem:endpoint-partial-sum} at the transition with $a$ alive labels.

\begin{lemma}[Logarithmic order of the mean]\label{lem:global-mean-order}
There are constants $0<c<C<\infty$ such that, for all sufficiently large $n$,
\[
        c\log n\le m_n\le C\log n.
\]
\end{lemma}

\begin{proof}
The global point bound gives the upper estimate.  For the lower estimate,
Lemma \ref{lem:global-endpoint-package} and Lemma
\ref{lem:global-one-saddle} imply, after increasing a fixed lower cutoff, that
\[
        p_{n,a}\ge c/a,
        \qquad A_1\le a\le U_n.
\]
Indeed, the saddle prefactor is comparable to $1/a$, and the relative error
tends to zero.  Since
$\log U_n=\log n-O(\log\log n)$, summing over this interval proves the lower
bound.
\end{proof}

\subsection{The swap law is asymptotically Palm}

\begin{proposition}[Quantitative swap--Palm approximation]\label{prop:swap-palm}
There are constants $C,c,\rho>0$ such that, for all sufficiently large $n$ and
$A_0\le a\le U_n$,
\[
        \varepsilon_{n,a}:=
        d_{\TV}(\nu_{n,a},\mu_{n,a}^\star)
        \le C\left(a^{-\rho}+e^{-c\sqrt{\log(a+2)}}\right).
\]
In particular,
\[
        \sum_{a=1}^{U_n}\frac{\varepsilon_{n,a}}a\le C.
\]
\end{proposition}

\begin{proof}
Fix $a$ and suppress $n$.  Put
$t_*=t_a^*$, $\kappa_*=\lambda_at_*$, and $k=a-1$.  Conditional on
$T_a=t$, define
\[
        B_b=\one_{\{T_b>t\}},
        \qquad
        N(t)=\sum_{b\ne a}B_b.
\]
Thus $I_a=1$ is the event $N(t)=k$.  For $b\ne a$, put
\[
        H_b(t,T_b)=
        \exp\{-(\lambda_a-\lambda_b)(T_b-t)\},
        \qquad H_a=1,
\]
so that $R_a=\sum_bH_b$.

A direct integration gives
\[
        \Ee[H_b\mid T_a=t]
        =\frac{\lambda_b}{\lambda_a}
        e^{(\lambda_a-\lambda_b)t}.
\]
Therefore
\[
\begin{aligned}
        M_a(t):=\Ee[R_a\mid T_a=t]
        &=\frac{e^{\lambda_at}}{\lambda_a}
        \sum_{b=1}^n\lambda_be^{-\lambda_bt}\\
        &=1+\frac{e^{\lambda_at}}{\lambda_a}D_a(t).
\end{aligned}
\]
Let
\[
        \eta_a=a^{-\rho_0}+e^{-c_0\sqrt{\log(a+2)}},
\]
where $\rho_0,c_0$ are chosen from Lemma \ref{lem:global-one-saddle}.
That lemma and the identity
\[
        \frac{\lambda_ae^{-\lambda_at_*}}{D_a(t_*)}M_a(t_*)
        =1+\frac{\lambda_ae^{-\lambda_at_*}}{D_a(t_*)}
\]
give
\[
        p_aM_a(t_*)=1+O(\eta_a).
\]
Moreover,
\[
        \frac{\dd}{\dd t}\log M_a(t)
        =\lambda_a-
        \frac{\sum_b\lambda_b^2e^{-\lambda_bt}}
        {\sum_b\lambda_be^{-\lambda_bt}},
\]
so Lemma \ref{lem:global-endpoint-package}(iv) gives
$|\frac{\dd}{\dd t}\log M_a(t)|\le C\lambda_a$ in the saddle
neighborhood.  Consequently, on
\[
        \mathcal C_a(t)=
        \{|\lambda_a\sqrt a(t-t_*)|\le L_a\},
\]
we have
\[
        p_aM_a(t)
        =1+O\left(\eta_a+\frac{L_a}{\sqrt a}\right).
\]

We next compare conditioning on the exact Bernoulli count with the
unconditioned center.  For $s\in\{0,1\}$ and $b\ne a$, let
\[
        A_{b,s}=\Ee[H_b\one_{\{B_b=s\}}\mid T_a=t].
\]
Then
\[
\begin{aligned}
&\Ee[H_b\mid T_a=t,N(t)=k]\\
&\quad=A_{b,1}
        \frac{\Pp(N(t)-B_b=k-1)}{\Pp(N(t)=k)}
        +A_{b,0}
        \frac{\Pp(N(t)-B_b=k)}{\Pp(N(t)=k)}.
\end{aligned}
\]
The deleted-summand assertion of Lemma \ref{lem:global-one-saddle} makes both
ratios $1+O(a^{-\rho_1})$, uniformly in $b$, throughout the central window.
Since all terms are nonnegative and $H_a=1$ is deterministic,
\[
        \Ee[R_a\mid T_a=t,N(t)=k]
        =M_a(t)(1+O(a^{-\rho_1})).
\]

We now control fluctuations of $R_a$.  Let
\[
        B_a=a^{1/8},
        \qquad
        \overline H_b=H_b\wedge B_a,
        \qquad
        \overline R_a=\sum_b\overline H_b.
\]
If $\lambda_b\le\lambda_a$, then
$H_b\le e^{\lambda_at}$, which is bounded in the central window, so such an
index does not contribute to $(H_b-B_a)_+$ for large $a$.  If
$x_b=\lambda_b/\lambda_a>1$ and $\kappa=\lambda_at$, direct integration gives
\[
        \Ee[(H_b-B_a)_+\mid T_a=t]
        =(x_b-1)e^{-\kappa x_b}B_a^{-1/(x_b-1)}.
\]
Lemma \ref{lem:global-likelihood-tail} therefore yields
\[
        \sum_b\Ee[(H_b-B_a)_+\mid T_a=t]
        \le Ca e^{-c\sqrt{\log a}}.
\]
This estimate survives conditioning on $N(t)=k$.  Indeed, for
$Y_b=(H_b-B_a)_+$,
\[
\begin{aligned}
&\Ee[Y_b\mid T_a=t,N(t)=k]\\
&\quad=\sum_{s=0}^1
 \Ee[Y_b\one_{\{B_b=s\}}\mid T_a=t]
 \frac{\Pp(N(t)-B_b=k-s)}{\Pp(N(t)=k)}
 \le C\Ee[Y_b\mid T_a=t],
\end{aligned}
\]
by the same deleted-summand ratio bound.  Summing over $b$ gives
\begin{equation}
        \Ee[R_a-\overline R_a\mid T_a=t,N(t)=k]
        \le Ca e^{-c\sqrt{\log a}}.
\label{eq:global-conditioned-truncation-tail}
\end{equation}

Conditional on $T_a=t$ and $N(t)=k$, the random set
\[
        \mathcal S=\{b\ne a:B_b=1\}
\]
has law
\[
        \Pp(\mathcal S=S\mid |\mathcal S|=k)
        \propto\prod_{b\in S}\frac{q_b}{1-q_b},
        \qquad q_b=e^{-\lambda_bt}.
\]
This is a product Bernoulli measure conditioned on one cardinality.  Product
Bernoulli measures are strongly Rayleigh, and single-cardinality conditioning
preserves this property \cite[Corollary 4.18]{BBL2009}.  Strongly Rayleigh
measures have the stochastic covering property
\cite[Proposition 2.2]{PemantlePeres2014}, so the concentration theorem
\cite[Theorem 3.1]{PemantlePeres2014} applies.

Given $\mathcal S$, the clocks are independent with the corresponding one-sided
truncated exponential laws.  Define
\[
        g_b^+(t)=\Ee[\overline H_b\mid T_a=t,B_b=1],
        \qquad
        g_b^-(t)=\Ee[\overline H_b\mid T_a=t,B_b=0]
\]
for $b\ne a$, and
\[
        G(S)=1+\sum_{b\in S}g_b^+(t)
        +\sum_{\substack{b\ne a\\b\notin S}}g_b^-(t).
\]
Both $g_b^+$ and $g_b^-$ lie in $[0,B_a]$.  Hence, for arbitrary subsets
$S,S'$,
\[
        |G(S)-G(S')|
        \le B_a|S\triangle S'|.
\]
Thus $G/B_a$ is Lipschitz one for Hamming distance.  Pemantle--Peres
concentration gives
\[
        \Pp(|G-\Ee G|>u\mid T_a=t,N(t)=k)
        \le2\exp\left\{-\frac{u^2}{8kB_a^2}\right\},
\]
and therefore
\begin{equation}
        \Ee[|G-\Ee G|\mid T_a=t,N(t)=k]
        \le CB_a\sqrt a.
\label{eq:global-G-concentration}
\end{equation}

Conditional on $\mathcal S$, the truncated clock contributions are independent,
and $0\le\overline H_b\le B_a$.  Consequently,
\[
\begin{aligned}
        \Var(\overline R_a\mid\mathcal S,T_a=t)
        &\le\sum_b\Ee[\overline H_b^2\mid\mathcal S,T_a=t]\\
        &\le B_aG(\mathcal S).
\end{aligned}
\]
Since
$\Ee G=\Ee[\overline R_a\mid T_a=t,N(t)=k]=O(a)$, conditional Cauchy--Schwarz
gives
\begin{equation}
        \Ee[|\overline R_a-G|\mid T_a=t,N(t)=k]
        \le C\sqrt{B_aa}.
\label{eq:global-clock-fluctuation}
\end{equation}
Combining the conditional mean comparison,
\eqref{eq:global-conditioned-truncation-tail},
\eqref{eq:global-G-concentration}, and
\eqref{eq:global-clock-fluctuation}, and using $B_a=a^{1/8}$, yields
\[
\begin{aligned}
&\Ee[|R_a-M_a(t)|\mid T_a=t,N(t)=k]\\
&\qquad\le Ca\left(a^{-\rho_2}+a^{-3/8}
        +e^{-c\sqrt{\log a}}\right).
\end{aligned}
\]
Since $p_a=O(a^{-1})$, the last estimate and the central estimate for
$p_aM_a(t)$ imply, for some $\rho_3>0$,
\[
        \Ee[|p_aR_a-1|\mid T_a=t,I_a=1]
        \le C\left(a^{-\rho_3}+e^{-c\sqrt{\log a}}\right)
\]
throughout the central window.

It remains to remove its complement.  Define the central event on an output
configuration by
\[
        \mathcal C_a=
        \{|\lambda_a\sqrt a\,(t_a-t_*)|\le L_a\}.
\]
Lemma \ref{lem:global-one-saddle} gives
\[
        \Pp(\mathcal C_a^c\mid I_a=1)
        \le Ce^{-cL_a^2}.
\]
Under $\nu_a$, the output coordinate $t_a$ is the $a$-th largest input clock;
write this order statistic as $O_a$.  At
$t_\pm=t_*\pm L_a/(\lambda_a\sqrt a)$, Taylor's theorem and Lemma
\ref{lem:global-endpoint-package}(v) show that the expected number of clocks
above $t_+$ is at most $a-cL_a\sqrt a$, while the expected number above $t_-$
is at least $a+cL_a\sqrt a$.  Bernstein's inequality therefore gives
\[
        \nu_a(\mathcal C_a^c)
        =\Pp(|\lambda_a\sqrt a(O_a-t_*)|>L_a)
        \le Ce^{-cL_a^2}.
\]
The likelihood-ratio identity gives
\[
        p_a\Ee[R_a\one_{\mathcal C_a^c}\mid I_a=1]
        =\nu_a(\mathcal C_a^c).
\]
Using $|p_aR_a-1|\le p_aR_a+1$ on $\mathcal C_a^c$, we conclude that
\[
        \Ee[|p_aR_a-1|\mid I_a=1]
        \le C\left(a^{-\rho_3}
        +e^{-c\sqrt{\log(a+2)}}\right).
\]
Lemma \ref{lem:swap-likelihood} proves the first assertion.  Finally,
\[
        \sum_{a\ge2}a^{-1-\rho_3}<\infty,
        \qquad
        \int_2^\infty\frac{e^{-c\sqrt{\log x}}}{x}\dd x<\infty,
\]
and the finitely many indices below $A_0$ are absorbed into the constant.
\end{proof}

\subsection{The swap changes the count only through a two-cycle}

Let $J_{n,a}$ be the indicator that $a$ belongs to a nontrivial two-cycle of
$\tau_n$, and put $q_{n,a}=\Pp(J_{n,a}=1)$.

\begin{lemma}[Exact effect of the rank swap]\label{lem:swap-count-effect}
For every no-tie clock configuration,
\[
        F_n(\mathcal T_{n,a}T)=F_n+1-I_{n,a}+J_{n,a}.
\]
In particular,
\[
        |F_n+1-F_n(\mathcal T_{n,a}T)|=I_{n,a}+J_{n,a}.
\]
\end{lemma}

\begin{proof}
If $a$ is already fixed, the map does nothing, $I_{n,a}=1$, and
$J_{n,a}=0$.  Suppose $a$ is not fixed, and let $b$ occupy terminal rank $a$.
Interchanging the clocks of $a$ and $b$ exchanges their ranks and leaves every
other rank unchanged.  It fixes $a$.  It fixes $b$ as well precisely when the
old rank of $a$ was $b$, namely when $a$ and $b$ formed a two-cycle.  This gives
the signed identity.  Since $I_{n,a}$ and $J_{n,a}$ are mutually exclusive,
the absolute-value identity follows.
\end{proof}

\begin{lemma}[Two-cycle transition bound]\label{lem:two-cycle-bound}
There is $C<\infty$ such that, for $1\le a<b\le n$,
\[
        \Pp(\tau_n(a)=b,\ \tau_n(b)=a)
        \le C\min\left\{\frac1{ab},
        \frac{\lambda_{n,a}}{b\lambda_{n,b}}\right\}.
\]
\end{lemma}

\begin{proof}
Expose the increasing exponential race as a pure-death process.  When the alive
set is $S$, the next deleted coordinate is $j\in S$ with probability
$\lambda_{n,j}/\sum_{c\in S}\lambda_{n,c}$.  The two-cycle event requires that
$a$ be deleted at the transition from $b$ alive labels to $b-1$, and that $b$
subsequently be deleted at the transition from $a$ alive labels to $a-1$.

At level $b$, any alive set of size $b$ has total rate at least
$\sum_{j=1}^b\lambda_{n,j}\ge c_\Sigma b\lambda_{n,b}$ by Lemma
\ref{lem:endpoint-partial-sum}.  Hence the conditional probability of deleting
$a$ is at most
\[
        C\frac{\lambda_{n,a}}{b\lambda_{n,b}}.
\]
At level $a$, if $b$ is still alive, the total rate is at least
\[
        \lambda_{n,b}+\sum_{j=1}^{a-1}\lambda_{n,j}.
\]
Moreover,
\[
        \sum_{j=1}^{a-1}\lambda_{n,j}
        \ge(c_\Sigma a-1)\lambda_{n,a}
        \ge c'a\lambda_{n,a}
\]
for all sufficiently large $a$; the finitely many smaller values are absorbed
into the constant.  Thus the conditional probability of deleting $b$ at level
$a$ is at most
\[
        C\min\left\{1,
        \frac{\lambda_{n,b}}{a\lambda_{n,a}}\right\}.
\]
Multiplying the two conditional bounds proves the claim.
\end{proof}

\begin{lemma}[Weighted two-cycle summability]\label{lem:two-cycle-sum}
With $U_n,V_n$ as above,
\[
        \sum_{a\le U_n}p_{n,a}q_{n,a}=O(1),
        \qquad
        \sum_{a>U_n}p_{n,a}q_{n,a}=O(\log\log n).
\]
\end{lemma}

\begin{proof}
For an unordered pair $s<\ell$, write
\[
        P_{s,\ell}=
        \Pp(\tau_n(s)=\ell,\tau_n(\ell)=s).
\]
Since $p_{n,a}\le C/a$,
\[
\begin{aligned}
        \sum_{a\le U_n}p_{n,a}q_{n,a}
        &=\sum_{s<\ell}
        \left(p_{n,s}\one_{\{s\le U_n\}}
        +p_{n,\ell}\one_{\{\ell\le U_n\}}\right)P_{s,\ell}\\
        &\le C\sum_{s\le U_n}\frac1s
        \sum_{\ell>s}P_{s,\ell}.
\end{aligned}
\]
For $\ell\le V_n$, Potter's bound gives
\[
        \frac{\lambda_{n,s}}{\lambda_{n,\ell}}
        \le C(s/\ell)^\beta.
\]
Lemma \ref{lem:two-cycle-bound} therefore bounds the summand after the factor
$1/s$ by
\[
        C\min\left\{\frac1{s^2\ell},
        \frac{s^{\beta-1}}{\ell^{1+\beta}}\right\}.
\]
For fixed $s$, split at $\ell=s^{1+1/\beta}$.  Below this point the first term
sums to $O(\log(s+1)/s^2)$; above it the second term sums to $O(s^{-2})$.
These bounds remain valid if the splitting point lies outside $(s,V_n]$.
Summing over $s$ gives a finite constant.

For $\ell>V_n$, monotonicity and Potter's bound give
\[
        \frac{\lambda_{n,s}}{\lambda_{n,\ell}}
        \le\frac{\lambda_{n,s}}{\lambda_{n,V_n}}
        \le C(s/V_n)^\beta.
\]
Hence these pairs contribute at most
\[
\begin{aligned}
        C\sum_{s\le U_n}\frac1s
        \left(\frac{s}{V_n}\right)^\beta
        \sum_{\ell>V_n}\frac1\ell
        &\le C\log(n/V_n)
        \left(\frac{U_n}{V_n}\right)^\beta=o(1).
\end{aligned}
\]
This proves the first estimate.  The second follows from $q_{n,a}\le1$:
\[
        \sum_{a>U_n}p_{n,a}q_{n,a}
        \le C\sum_{a>U_n}\frac1a
        =O(\log(n/U_n))=O(\log\log n).
\]
\end{proof}

\subsection{Proof of the global theorem}

\begin{proof}[Proof of Theorem \ref{thm:global-endpoint-poisson}]
Apply Lemma \ref{lem:approx-palm-stein} with
$W=F_n$, $X_a=I_{n,a}$, and
\[
        Y^{(a)}=F_n(\mathcal T_{n,a}T).
\]
By Lemma \ref{lem:swap-count-effect},
\[
        \Ee|F_n+1-Y^{(a)}|=p_{n,a}+q_{n,a}.
\]
The first Stein term is therefore
\[
        \frac{1-e^{-m_n}}{m_n}
        \sum_{a=1}^n p_{n,a}(p_{n,a}+q_{n,a}).
\]
Since $p_{n,a}\le C/a$,
\[
        \sum_a p_{n,a}^2=O(1).
\]
Lemma \ref{lem:two-cycle-sum} gives
\[
        \sum_a p_{n,a}q_{n,a}=O(\log\log n).
\]
By Lemma \ref{lem:global-mean-order}, the first Stein term is
$O(\log\log n/\log n)$.

For the approximate-Palm term, data processing for total variation gives
\[
 d_{\TV}\bigl(\mathcal L(F_n(\mathcal T_{n,a}T)),
 \mathcal L(F_n\mid I_{n,a}=1)\bigr)
 \le d_{\TV}(\nu_{n,a},\mu_{n,a}^\star)=\varepsilon_{n,a}.
\]
Therefore Proposition \ref{prop:swap-palm} gives
\[
\begin{aligned}
        \sum_{a\le U_n}p_{n,a}\varepsilon_{n,a}
        &\le C\sum_{a\le U_n}\frac1a
        \left(a^{-\rho}+e^{-c\sqrt{\log(a+2)}}\right)\\
        &\le C.
\end{aligned}
\]
On the remaining range, use $\varepsilon_{n,a}\le1$ and $p_{n,a}\le C/a$:
\[
        \sum_{a>U_n}p_{n,a}\varepsilon_{n,a}
        =O(\log\log n).
\]
Lemma \ref{lem:global-mean-order} therefore bounds the second Stein term by
\[
        O\left(m_n^{-1/2}\log\log n\right)
        =O\left(\frac{\log\log n}{\sqrt{\log n}}\right).
\]
Combining the two estimates proves the total-variation bound.

The exact asymptotic $m_n\sim c_\gamma\log n$ is Theorem
\ref{thm:endpoint-fixed-universality}.  If $Z_n\sim\Pois(m_n)$, then
$(Z_n-m_n)/\sqrt{m_n}\Rightarrow N(0,1)$.  Total-variation convergence
transfers this limit from $Z_n$ to $F_n$.  For the Sukhatme profile,
$\gamma=1$ and $c_1=e^{-1}$.
\end{proof}

\begin{remark}[No growing-dimensional dissociation is needed]\label{rem:no-high-order-dissociation}
The proof does not require uniform factorization of joint fixed-point
probabilities for sets of cardinality comparable to $\log n$.  The rank
transposition produces an approximate Palm coupling one coordinate at a time.
The only direct interaction that can change the fixed-point count is a two-cycle,
and its weighted contribution is summable.  Thus the formerly conjectured
high-order dissociation statement is strictly stronger than what is needed for
global Poisson approximation.
\end{remark}


\section{Endpoint short cycles}
\label{sec:endpoint-short-cycles}

Recall that the terminal-coordinate permutation $\tau_n$ is conjugate to
$\sigma_n$ by reversal.  In particular,
\[
        C_r(\tau_n)=C_r(\sigma_n),\qquad r\ge1.
\]
The multiplicative kernel in Proposition
\ref{prop:endpoint-kernel-powers} determines the asymptotic number of every
fixed endpoint cycle length.  Unlike the regular positive-profile regime, each
of these expected cycle counts diverges logarithmically.

For $r\ge1$, put
\[
        \alpha_{\gamma,r}=\frac1r k_\gamma^{*r}(0).
\]
Thus $\alpha_{\gamma,1}=c_\gamma$.

\begin{theorem}[Endpoint short-cycle law]
\label{thm:endpoint-short-cycles}
Assume Definition \ref{def:endpoint-rv}.  For every fixed $R\ge1$ and fixed
nonnegative integers $m_1,\ldots,m_R$,
\[
        \Ee\prod_{r=1}^R(C_r(\sigma_n))_{m_r}
        \sim\prod_{r=1}^R
        (\alpha_{\gamma,r}\log n)^{m_r}.
\]
Consequently, for every fixed $r$,
\[
        \Ee C_r(\sigma_n)
        \sim\frac{k_\gamma^{*r}(0)}r\log n,
\]
and, jointly for every fixed $R$,
\[
        \left(
        \frac{C_1}{\alpha_{\gamma,1}\log n},\ldots,
        \frac{C_R}{\alpha_{\gamma,R}\log n}
        \right)\longrightarrow(1,\ldots,1)
\]
in probability and in every fixed $L^p$, $p<\infty$.
\end{theorem}

The proof has three inputs.  First, we prove a source-to-target endpoint entry
estimate by an endpoint-specific anisotropic local central limit theorem,
mixed noncompact localization under one tilted law, and covariance stability
on the saddle scale.  Second, a death-process envelope localizes every fixed
cycle component in logarithmic coordinates, including components crossing the
mesoscopic cutoff.  Third, a logarithmic trace Riemann sum identifies the
coefficient $k_\gamma^{*r}(0)/r$.

\subsection{General endpoint point constraints}

For fixed $\ell\ge1$, $\eta>0$, and $M<\infty$, define
\[
\begin{aligned}
\mathfrak Q_N(\ell,\eta,M)=\bigl\{&(n,p_1,\ldots,p_\ell,
 q_1,\ldots,q_\ell):
 N\le p_j,q_j\le n/N,\\
& p_1,\ldots,p_\ell\text{ distinct},\quad
 q_1<\cdots<q_\ell,\quad
 q_{j+1}/q_j\ge e^\eta,\\
& e^{-M}\le p_j/q_j\le e^M\quad(1\le j\le\ell)
\bigr\}.
\end{aligned}
\tag{10.1}\label{eq:endpoint-point-admissible}
\]
A uniform assertion below means that the corresponding supremum over
$\mathfrak Q_N(\ell,\eta,M)$ tends to zero as $N\to\infty$.
The coordinates $p_j$ are called \emph{source terminal coordinates}, and the
$q_j$ are the target terminal positions.

Let $P=\{p_1,\ldots,p_\ell\}$ and put
\[
        R_{n,j}=\frac{\lambda_{n,p_j}}{\lambda_{n,q_j}},
        \qquad
        z_j=q_j-j.
\]
For $h=(h_1,\ldots,h_\ell)\in\mathbb R^\ell$, define thresholds
\[
        y_j(h_j)=\log\lambda_{n,q_j}+h_j
\]
and cumulative lower-priority counts
\[
        Z_j(h)=\sum_{b\notin P}
        \one_{\{Y_{n,b}<y_j(h_j)\}},
        \qquad 1\le j\le\ell.
\]
The relevant chamber is
\[
        \mathscr C_{n,q}
        =\{h:y_1(h_1)<\cdots<y_\ell(h_\ell)\}.
\]
Finally, for $R>0$ let
\[
        \varphi_R(h)=R e^{-h}\exp\{-R e^{-h}\}.
\tag{10.2}\label{eq:shifted-gumbel-density}
\]
This is the density of $\log R+G$ when $G$ is standard Gumbel.

\begin{lemma}[Source-to-target endpoint saddle geometry]
\label{lem:endpoint-point-saddle}
Uniformly over $\mathfrak Q_N(\ell,\eta,M)$, the ratios $R_{n,j}$ remain in a
fixed compact subinterval of $(0,\infty)$.  For each $j$, the equation
\[
        M_j(h_j^*)=z_j,
        \qquad
        M_j(h)=\Ee Z_j(h),
\tag{10.3}\label{eq:endpoint-point-saddle-equation}
\]
has a unique solution, and
\[
        h_j^*\longrightarrow g_\gamma:=-\log K_\gamma,
        \qquad
        d_{n,j}:=M_j'(h_j^*)\sim\frac{q_j}{\gamma}.
\tag{10.4}\label{eq:endpoint-point-saddle-asymptotics}
\]
For every fixed $L<\infty$,
\[
        \sup_{|h|\le L}|M_j''(h)|\le C_Lq_j,
        \qquad
        c_Lq_j\le M_j'(h)\le C_Lq_j
\tag{10.5}\label{eq:endpoint-point-mean-derivatives}
\]
whenever $|h-h_j^*|\le1$.  Moreover, there are
$\delta_{\eta,M}>0$, $c_{\eta,M}>0$, and $N_0$ such that, for $N\ge N_0$,
\[
        y_{j+1}(h_{j+1})-y_j(h_j)\ge c_{\eta,M}
\tag{10.6}\label{eq:endpoint-point-chamber-margin}
\]
whenever $|h_j-h_j^*|\le\delta_{\eta,M}$ for every $j$.  Thus this product
neighborhood is contained in $\mathscr C_{n,q}$.
\end{lemma}

\begin{proof}
Uniform regular variation on the fixed multiplicative range
$e^{-M}\le p_j/q_j\le e^M$ gives
\[
        R_{n,j}
        =\left(\frac{p_j}{q_j}\right)^\gamma(1+o(1))
\tag{10.7}\label{eq:endpoint-point-rate-ratio}
\]
uniformly, proving the compactness assertion.  By Lemma
\ref{lem:endpoint-rv-sums}, including its finite-deletion clause, uniformly on
compact $h$-sets,
\[
        M_j(h)=q_jA_\gamma e^{h/\gamma}(1+o(1)),
\tag{10.8}\label{eq:endpoint-point-mean-asymptotic}
\]
\[
        M_j'(h)=\frac{q_jA_\gamma}{\gamma}
        e^{h/\gamma}(1+o(1)).
\tag{10.9}\label{eq:endpoint-point-mean-prime}
\]
The function $M_j$ is continuous and strictly increasing from $0$ to
$n-\ell$, so the saddle exists uniquely.  Since $z_j/q_j\to1$,
\eqref{eq:endpoint-point-mean-asymptotic} gives
$A_\gamma e^{h_j^*/\gamma}\to1$, equivalently
$h_j^*\to-\log K_\gamma$.  Substitution into
\eqref{eq:endpoint-point-mean-prime} gives
$d_{n,j}\sim q_j/\gamma$.

Termwise differentiation expresses $M_j''$ as a linear combination of the
weighted sums in Lemma \ref{lem:endpoint-rv-sums} with powers one and two.
This gives \eqref{eq:endpoint-point-mean-derivatives}.

Finally, the target separation and the Potter lower bound imply that, for some
$c_\eta>0$ and all large $N$,
\[
        \log\frac{\lambda_{n,q_{j+1}}}{\lambda_{n,q_j}}
        \ge2c_\eta.
\tag{10.10}\label{eq:endpoint-point-rate-gap}
\]
Indeed, on a fixed compact range of $q_{j+1}/q_j$ this follows from the uniform
convergence theorem, while larger ratios are handled by
$\lambda_{n,q_{j+1}}/\lambda_{n,q_j}
 \ge c(q_{j+1}/q_j)^{\gamma/2}$.
Since all $h_j^*$ converge to the same constant, a sufficiently small fixed
$\delta_{\eta,M}$ gives \eqref{eq:endpoint-point-chamber-margin}.
\end{proof}

\begin{lemma}[Source-to-target endpoint anisotropic local CLT]
\label{lem:endpoint-point-llt}
Let
\[
        D_q=\diag(q_1,\ldots,q_\ell),
        \qquad
        \Sigma(h)=\Var Z(h),
\]
and let $\mathcal U_n$ be the saddle neighborhood in Lemma
\ref{lem:endpoint-point-saddle}.  Uniformly over
$\mathfrak Q_N(\ell,\eta,M)$ and $h\in\mathcal U_n$,
\[
        cD_q\le\Sigma(h)\le CD_q
\tag{10.11}\label{eq:endpoint-point-covariance-comparison}
\]
in quadratic-form order.  For every fixed $A<\infty$, uniformly for lattice
points $v\in\mathbb Z^\ell$ satisfying
\[
        \|D_q^{-1/2}(v-M(h))\|\le A,
\]
we have
\[
\begin{aligned}
\Pp(Z(h)=v)
&=(2\pi)^{-\ell/2}(\det\Sigma(h))^{-1/2}\\
&\quad\times
\exp\left[-\frac12(v-M(h))^T\Sigma(h)^{-1}(v-M(h))\right]
+o\left(\prod_{j=1}^\ell q_j^{-1/2}\right).
\end{aligned}
\tag{10.12}\label{eq:endpoint-point-llt}
\]
For every fixed $\nu<1/6$, there are $C,c>0$ such that
\[
\Pp(Z(h)=v)
\le C\prod_{j=1}^\ell q_j^{-1/2}
\exp\left[-c\sum_{j=1}^\ell
\frac{(v_j-M_j(h_j))^2}{q_j}\right]
\tag{10.13}\label{eq:endpoint-point-local-upper}
\]
whenever
$|v_j-M_j(h_j)|\le q_j^{1/2+\nu}$ for every $j$.
\end{lemma}

\begin{proof}
For $b\notin P$, the vector
\[
        V_b(h)=
        (\one_{\{Y_{n,b}<y_1(h_1)\}},\ldots,
         \one_{\{Y_{n,b}<y_\ell(h_\ell)\}})
\]
has the cumulative categorical support
\[
        (1,\ldots,1),(0,1,\ldots,1),\ldots,
        (0,\ldots,0,1),0.
\]
By \eqref{eq:endpoint-point-chamber-margin}, the adjacent threshold gaps are
uniformly positive in $\mathcal U_n$.  Choose a small multiplicative interval
$[(1-\delta)q_j,(1+\delta)q_j]$ for each $j$, with $\delta$ so small that the
$\ell$ intervals are disjoint.  Uniform regular variation and
\eqref{eq:endpoint-point-saddle-asymptotics} imply that, for every index $b$ in
the $j$-th block, the two categories immediately below and above the threshold
$y_j(h_j)$ both have probability bounded below by a positive constant.
Removing the fixed set $P$ changes each block size by at most $\ell$.
Consequently, for every $j$, there are at least $cq_j$ independent summands
having two atoms whose difference is the coordinate vector $e_j$, with both
atom probabilities at least $c$.  These bounds persist under every bounded
exponential tilt.

The disjoint blocks give the covariance lower bound in
\eqref{eq:endpoint-point-covariance-comparison}.  The upper bound follows from
\[
        \Var Z_j(h)=O(q_j),
\]
which is Lemma \ref{lem:endpoint-rv-sums} applied at the relevant compact
values of $e^{-h_j}$, followed by Cauchy's inequality.  The same blocks and
Lemma \ref{lem:coord-fourier} give, uniformly under bounded tilts,
\[
\left|\Ee e^{iu\cdot(Z(h)-M(h))}\right|
\le\exp\left[-c\sum_{j=1}^\ell
q_j\operatorname{dist}(u_j,2\pi\mathbb Z)^2\right].
\tag{10.14}\label{eq:endpoint-point-fourier}
\]
On the anisotropic box $|u_j|\le A q_j^{-1/2}$, bounded categorical
increments and the covariance upper bound give
\[
\log\Ee e^{iu\cdot(Z-M)}
=-\frac12u^T\Sigma(h)u
+O\left((\max_j|u_j|)\sum_jq_ju_j^2\right).
\]
After the change of variables $u=D_q^{-1/2}w$, the remainder is
$O_A(q_1^{-1/2})=o(1)$.  Fourier inversion and
\eqref{eq:endpoint-point-fourier} control the complementary region on the
product scale.  This proves \eqref{eq:endpoint-point-llt}.  The tilted local
upper bound \eqref{eq:endpoint-point-local-upper} follows from Lemma
\ref{lem:anisotropic-tilted-upper}: the tilted covariance upper bound follows
from bounded distortion of the categorical probabilities, and the edge blocks
supply the required tilted concentration estimate.
\end{proof}

\begin{lemma}[Source-to-target mixed localization]
\label{lem:endpoint-point-localization}
There is $L_0=L_0(\ell,\eta,M,\gamma)$ such that the following holds for
$L\ge L_0$.  For $h\in\mathscr C_{n,q}$, set
\[
\begin{aligned}
        B_R(h)&=\{j:h_j>L\},\\
        B_L(h)&=\{j:h_j<-L\},\\
        G(h)&=\{1,\ldots,\ell\}\setminus(B_R(h)\cup B_L(h)),
\end{aligned}
\]
and put
\[
        H_j=z_j-M_j(h_j),\qquad H_0=0.
\]
On every region on which these three sets are fixed,
\[
\begin{aligned}
\Pp(Z(h)=z)
&\le C_L\prod_{j=1}^\ell q_j^{-1/2}
\exp\left[-c\sum_{j\in B_R}|H_j|
           -c\sum_{j\in B_L}q_j\right]\\
&\quad\times
\exp\left[-c_L\sum_{j\in G}
\min\left\{\frac{H_j^2}{q_j},q_j\right\}\right].
\end{aligned}
\tag{10.15}\label{eq:endpoint-point-mixed-bound}
\]
Consequently,
\[
\begin{aligned}
\Pp(Z(h)=z)
&\le C_L\prod_{j=1}^\ell q_j^{-1/2}
\exp\left[-ce^{L/(2\gamma)}\sum_{j\in B_R}q_j
           -c\sum_{j\in B_L}q_j\right]\\
&\quad\times
\exp\left[-c_L\sum_{j\in G}
q_j\min\{(h_j-h_j^*)^2,1\}\right].
\end{aligned}
\tag{10.16}\label{eq:endpoint-point-scaled-bound}
\]
Moreover,
\[
\lim_{L\to\infty}\limsup_{N\to\infty}
\sup_{\mathfrak Q_N(\ell,\eta,M)}
\frac{
\displaystyle
\int_{\mathscr C_{n,q}\cap\{h\notin[-L,L]^\ell\}}
\prod_{j=1}^\ell\varphi_{R_{n,j}}(h_j)
\Pp(Z(h)=z)\dd h
}{\prod_{j=1}^\ell q_j^{-1}}=0.
\tag{10.17}\label{eq:endpoint-point-noncompact}
\]
On every fixed compact cube containing the saddle vectors,
\[
\Pp(Z(h)=z)
\le C\prod_{j=1}^\ell q_j^{-1/2}
\exp\left[-c\sum_{j=1}^\ell
q_j\min\{(h_j-h_j^*)^2,1\}\right].
\tag{10.18}\label{eq:endpoint-point-compact-localization}
\]
\end{lemma}

\begin{proof}
Define interval counts
\[
        W_1=Z_1,\qquad W_j=Z_j-Z_{j-1}\quad(2\le j\le\ell),
\]
with targets
\[
        w_1=q_1-1,\qquad
        w_j=q_j-q_{j-1}-1\quad(2\le j\le\ell).
\]
The target separation gives $c_\eta q_j\le w_j\le q_j$.  Let
$m_j=\Ee W_j$ and $d_j=w_j-m_j$; then $d_j=H_j-H_{j-1}$.
For $j\in G$, one has $m_j\le M_j(h_j)=O_L(q_j)$, and hence
$|d_j|=O_L(q_j)$.

Fix the offender sets.  Let
\[
        K_h(\theta)=\log\Ee
        \exp\{\theta\cdot(W(h)-m(h))\}
\]
be the centered interval-coordinate cumulant, and let
\[
        (\mathcal T\alpha)_k=\sum_{j\ge k}\alpha_j.
\]
Then
$(\mathcal T\alpha)\cdot d=\alpha\cdot H$.
Choose $\delta>0$ sufficiently small and set
\[
        \alpha_j^-=-\delta\one_{\{j\in B_R\}},
        \qquad
        \alpha_j^+=\delta\one_{\{j\in B_L\}}.
\]
On $G$, use the local interval tilt
\[
        \lambda_j=
        \begin{cases}
        \delta_Ld_j/q_j,& |d_j|\le\varepsilon q_j,\\
        \delta_L\operatorname{sgn}(d_j),& |d_j|>\varepsilon q_j,
        \end{cases}
\]
where $\varepsilon<c_\eta/4$ and $\delta_L>0$ is sufficiently small; put
$\lambda_j=0$ off $G$.  Use the single total interval-coordinate tilt
\[
        \Theta=\mathcal T\alpha^-+\mathcal T\alpha^++\lambda.
\tag{10.19}\label{eq:endpoint-point-total-tilt}
\]
Generalized H\"older inequality yields
\[
\begin{aligned}
K_h(\Theta)
&\le\frac13K_h(3\mathcal T\alpha^-)
 +\frac13K_h(3\mathcal T\alpha^+)
 +\frac13K_h(3\lambda).
\end{aligned}
\tag{10.20}\label{eq:endpoint-point-holder}
\]

By \eqref{eq:endpoint-point-mean-asymptotic}, for large $L$,
\[
\begin{array}{lll}
 j\in B_R:&M_j(h_j)\ge ce^{L/\gamma}q_j,&H_j<0,\\[1mm]
 j\in B_L:&M_j(h_j)\le Ce^{-L/\gamma}q_j,&0<H_j\le q_j.
\end{array}
\tag{10.21}\label{eq:endpoint-point-offender-geometry}
\]
The proof of Lemma
\ref{lem:actual-mean-lower-cumulative-chernoff} applies to the present nested
lower-priority counts without change.  Applied with parameter $3\delta$, it
gives
\[
        -3\alpha^-\cdot H+K_h(3\mathcal T\alpha^-)
        \le-c\sum_{j\in B_R}|H_j|,
\tag{10.22}\label{eq:endpoint-point-right-cost}
\]
\[
        -3\alpha^+\cdot H+K_h(3\mathcal T\alpha^+)
        \le-c\sum_{j\in B_L}q_j.
\tag{10.23}\label{eq:endpoint-point-left-cost}
\]
Under every bounded interval tilt, the weighted-sum estimate gives
\[
        \Var_\xi(v\cdot W(h))\le C_L\sum_{j\in G}q_jv_j^2
\]
for $v$ supported on $G$.  Taylor's formula for the cumulant and the choice of
$\delta_L$ therefore give
\[
        -3\lambda\cdot d+K_h(3\lambda)
        \le-c_L\sum_{j\in G}
        \min\left\{\frac{d_j^2}{q_j},q_j\right\}.
\tag{10.24}\label{eq:endpoint-point-local-cost}
\]

All coordinates of $\Theta$ are bounded for fixed $L$.  For every locally
regular coordinate $j\in G$ with $|d_j|\le\varepsilon q_j$, the interval mean
is at least $(c_\eta/2)q_j$.  The middle-block argument in Lemma
\ref{lem:endpoint-interval-edge-blocks}, with $q_j$ in place of $a_j$ and with
the fixed deletion set $P$, supplies forward two-atom blocks of order $q_j$.
They can be chosen disjointly, and persist under the total tilt.  After
projection onto the locally regular coordinates, Lemma
\ref{lem:projected-forward-edge-concentration} gives
\[
        \Pp_\Theta(W_{G_0}=w_{G_0})
        \le C_L\prod_{j\in G_0}q_j^{-1/2}.
\tag{10.25}\label{eq:endpoint-point-tilted-factor}
\]
The exact change-of-measure identity, together with
\eqref{eq:endpoint-point-holder}--\eqref{eq:endpoint-point-tilted-factor},
gives the desired exponential estimate first in terms of the interval
residuals $d_j$.  Indices lacking a local factor carry a fixed exponential
loss of order $q_j$, which absorbs the missing factor $q_j^{-1/2}$.

To pass from interval to cumulative residuals without losing boundary terms,
apply Lemma \ref{lem:boundary-truncated-chain} with $A_j=q_j$, with its
$B_-$ equal to the present $B_R$, and its $B_+$ equal to the present $B_L$.
The assumptions follow from $|d_j|=O_L(q_j)$ on $G$ and
$|H_j|\le q_j$ on $B_L$.  Reserving a fixed fraction of the offender costs
proves \eqref{eq:endpoint-point-mixed-bound}.  On $G$, the derivative bounds
and monotonicity imply
\[
        |H_j|\ge c_Lq_j\min\{|h_j-h_j^*|,1\}.
\]
On $B_R$, \eqref{eq:endpoint-point-offender-geometry} gives
$|H_j|\ge ce^{L/\gamma}q_j$.  This proves
\eqref{eq:endpoint-point-scaled-bound}; the compact estimate
\eqref{eq:endpoint-point-compact-localization} is the case with no offenders.

It remains to integrate the tail.  By
\eqref{eq:endpoint-point-rate-ratio}, $R_{n,j}$ lies in a fixed compact
subinterval of $(0,\infty)$.  Uniformly in such $R$,
\[
        \varphi_R(h)\le Ce^{-h}\quad(h\ge0),
        \qquad
        \varphi_R(h)\le Ce^{-h}e^{-ce^{-h}}\quad(h\le0).
\tag{10.26}\label{eq:endpoint-point-density-tails}
\]
Each nonoffending coordinate contributes $O(q_j^{-1/2})$ after integration on
its Gaussian scale, in addition to the prefactor $q_j^{-1/2}$.  For every
offender, the exponential loss in \eqref{eq:endpoint-point-scaled-bound}
absorbs the missing local factor, and the remaining density tail tends to zero
as $L\to\infty$.  Summing over the finitely many offender sets proves
\eqref{eq:endpoint-point-noncompact}.
\end{proof}

\begin{lemma}[Source-to-target covariance stability]
\label{lem:endpoint-point-covariance-stability}
Let
\[
        H_q=\diag(q_1^{-1/2},\ldots,q_\ell^{-1/2}),
        \qquad \Sigma_*=\Sigma(h^*).
\]
For every fixed $A<\infty$,
\[
\sup_{\|u\|\le A}
\left\|
D_q^{-1/2}
\bigl[\Sigma(h^*+H_qu)-\Sigma_*\bigr]
D_q^{-1/2}
\right\|_{\mathrm{op}}\longrightarrow0
\tag{10.27}\label{eq:endpoint-point-covariance-stability}
\]
uniformly over $\mathfrak Q_N(\ell,\eta,M)$.  Consequently,
\[
        \frac{\det\Sigma(h^*+H_qu)}{\det\Sigma_*}\longrightarrow1
\tag{10.28}\label{eq:endpoint-point-determinant-stability}
\]
and
\[
\left\|
D_q^{1/2}
\bigl[\Sigma(h^*+H_qu)^{-1}-\Sigma_*^{-1}\bigr]
D_q^{1/2}
\right\|_{\mathrm{op}}\longrightarrow0
\tag{10.29}\label{eq:endpoint-point-inverse-stability}
\]
uniformly on the same sets.
\end{lemma}

\begin{proof}
For $j\le k$, put
\[
        p_{b,j}(h_j)=
        \exp\left(-e^{-h_j}
        \frac{\lambda_{n,b}}{\lambda_{n,q_j}}\right).
\]
The threshold events are nested in the chamber, so
\[
        \Sigma_{jk}(h)
        =\sum_{b\notin P}p_{b,j}(h_j)(1-p_{b,k}(h_k)).
\tag{10.30}\label{eq:endpoint-point-covariance-entry}
\]
Lemma \ref{lem:endpoint-rv-sums}, with power one, gives on a fixed saddle
neighborhood
\[
        |\partial_{h_j}\Sigma_{jk}(h)|\le Cq_j,
        \qquad
        |\partial_{h_k}\Sigma_{jk}(h)|\le Cq_k.
\]
Hence, for $\|u\|\le A$,
\[
        |\Sigma_{jk}(h^*+H_qu)-\Sigma_{jk}(h^*)|
        \le C_A(\sqrt{q_j}+\sqrt{q_k}).
\]
After division by $\sqrt{q_jq_k}$ this is at most
$C_A(q_j^{-1/2}+q_k^{-1/2})\le2C_Aq_1^{-1/2}$.  Since $\ell$ is fixed,
\eqref{eq:endpoint-point-covariance-stability} follows.  The normalized
covariance matrices remain in a compact subset of the positive-definite
matrices by \eqref{eq:endpoint-point-covariance-comparison}; continuity of
determinant and inversion gives
\eqref{eq:endpoint-point-determinant-stability}--
\eqref{eq:endpoint-point-inverse-stability}.
\end{proof}

\begin{lemma}[Endpoint integrated point-constraint saddle]
\label{lem:endpoint-general-integrated-saddle}
Uniformly over $\mathfrak Q_N(\ell,\eta,M)$,
\[
\begin{aligned}
&\int_{\mathscr C_{n,q}}
\prod_{j=1}^\ell\varphi_{R_{n,j}}(h_j)
\Pp(Z(h)=z)\dd h\\
&\qquad=(1+o(1))
\prod_{j=1}^\ell
\frac{\varphi_{R_{n,j}}(h_j^*)}{d_{n,j}},
\end{aligned}
\tag{10.31}\label{eq:endpoint-point-integrated-saddle}
\]
where $d_{n,j}=M_j'(h_j^*)$.
\end{lemma}

\begin{proof}
By Lemma \ref{lem:endpoint-point-localization}, it suffices first to integrate
over a fixed compact cube and then over a shrinking saddle neighborhood; the
complement is $o(\prod_jq_j^{-1})$.  Put
\[
        h=h^*+H_qu,
        \qquad
        D_n=\diag(d_{n,1},\ldots,d_{n,\ell}).
\]
For bounded $u$, Taylor's theorem and
\eqref{eq:endpoint-point-mean-derivatives} give
\[
        M(h)-z=D_nH_qu+o(D_q^{1/2})
\tag{10.32}\label{eq:endpoint-point-mean-taylor}
\]
uniformly.  The density factors satisfy
\[
        \prod_j\varphi_{R_{n,j}}(h_j)
        =\prod_j\varphi_{R_{n,j}}(h_j^*)(1+o(1))
\]
on bounded $u$-balls.  Lemma \ref{lem:endpoint-point-llt} and Lemma
\ref{lem:endpoint-point-covariance-stability} therefore give, uniformly on
such balls, the Gaussian integrand with covariance $\Sigma_*$ and mean shift
$D_nH_qu$.

Since $\dd h=|\det H_q|\dd u$, the Gaussian integral contributes
\[
\begin{aligned}
&(2\pi)^{-\ell/2}(\det\Sigma_*)^{-1/2}|\det H_q|\\
&\quad\times
\int_{\mathbb R^\ell}
\exp\left[-\frac12u^TH_q^TD_n^T
\Sigma_*^{-1}D_nH_qu\right]\dd u
=|\det D_n|^{-1}.
\end{aligned}
\tag{10.33}\label{eq:endpoint-point-determinant-cancellation}
\]
Indeed,
\[
\det(H_q^TD_n^T\Sigma_*^{-1}D_nH_q)
=|\det H_q|^2|\det D_n|^2(\det\Sigma_*)^{-1}.
\]
The compact localization bound
\eqref{eq:endpoint-point-compact-localization} is an integrable Gaussian
majorant in $u$, and
\eqref{eq:endpoint-point-noncompact} handles the outer region.  Dominated
convergence proves \eqref{eq:endpoint-point-integrated-saddle}.
\end{proof}

\begin{proposition}[Separated endpoint entry estimate]
\label{prop:endpoint-general-entry}
Uniformly over $\mathfrak Q_N(\ell,\eta,M)$,
\[
        \Pp(\tau_n(p_j)=q_j,\ 1\le j\le\ell)
        =(1+o(1))\prod_{j=1}^\ell
        \frac1{q_j}k_\gamma\!\left(\log\frac{q_j}{p_j}\right).
\tag{10.34}\label{eq:endpoint-general-entry}
\]
\end{proposition}

\begin{proof}
Condition on the specified centered Gumbels $G_{p_j}$ and put
\[
        h_j=\log R_{n,j}+G_{p_j}.
\]
Then the specified priority is
$Y_{n,p_j}=\log\lambda_{n,q_j}+h_j$, and $h_j$ has density
$\varphi_{R_{n,j}}$.  Since $q_1<\cdots<q_\ell$, the point constraints force
$Y_{n,p_1}<\cdots<Y_{n,p_\ell}$.  In this chamber exactly $j-1$ specified
priorities lie below the $j$-th threshold, so the event is equivalent to
$Z_j(h)=q_j-j$ for all $j$.  Thus the probability is exactly the integral in
\eqref{eq:endpoint-point-integrated-saddle}.

By Lemma \ref{lem:endpoint-point-saddle},
$h_j^*\to-\log K_\gamma$ and $d_{n,j}\sim q_j/\gamma$.  Therefore
\[
\begin{aligned}
\frac{\varphi_{R_{n,j}}(h_j^*)}{d_{n,j}}
&=\frac{\gamma K_\gamma}{q_j}
R_{n,j}e^{-K_\gamma R_{n,j}}(1+o(1))\\
&=\frac1{q_j}k_\gamma\!\left(\log\frac{q_j}{p_j}\right)(1+o(1)),
\end{aligned}
\]
where the second equality uses
\eqref{eq:endpoint-point-rate-ratio}.  Multiply over $j$.
\end{proof}

\subsection{A global cycle envelope}

The next lemma makes precise the monotonicity used when forced deletions are
interspersed with ordinary Luce deletions.

\begin{lemma}[Scheduled-deletion survival]
\label{lem:scheduled-deletion-survival}
Let $S$ be a finite set of labels with positive rates, and let $p\in S$.  Start
from $S$ and run a schedule in which each step is either an ordinary Luce
deletion from the current set or a forced removal of a chosen competitor
$c\ne p$.  The forced choices may depend on the previously exposed history.
For $m\ge0$, the probability that $p$ is still present after $m$ ordinary
deletions is at most
\[
        H_m(S,p):=
        \Pp\left(\#\{b\in S\setminus\{p\}:T_b<T_p\}\ge m\right)
\tag{10.35}\label{eq:scheduled-survival-function}
\]
for an unmodified exponential race on $S$.
\end{lemma}

\begin{proof}
The function $H_m(S,p)$ is increasing under addition of competitors, by
coupling the clocks.  It also satisfies the exact recursion
\[
        H_m(S,p)=
        \sum_{b\in S\setminus\{p\}}
        \frac{\lambda_b}{\sum_{c\in S}\lambda_c}
        H_{m-1}(S\setminus\{b\},p),
        \qquad m\ge1,
\tag{10.36}\label{eq:scheduled-survival-recursion}
\]
with $H_0=1$: to survive $m$ ordinary deletions, the first ordinary deletion
must be a competitor $b$, after which memorylessness leaves the same problem
with $m-1$ deletions.

Induct on the remaining schedule.  If the next step is a forced removal of
$c$, the induction hypothesis and monotonicity give a bound by
$H_m(S\setminus\{c\},p)\le H_m(S,p)$.  If the next step is ordinary, condition
on the deleted label and use \eqref{eq:scheduled-survival-recursion}.  This
also covers adaptive forced choices after conditioning on the current history.
\end{proof}

\begin{lemma}[Forced-deletion bound]
\label{lem:forced-deletion}
Consider compatible point constraints
$\tau_n(p_i)=q_i$, $1\le i\le L$, with distinct sources and distinct targets.
Then
\[
        \Pp(\tau_n(p_i)=q_i,\ 1\le i\le L)
        \le C^L\prod_{i=1}^L
        \frac{\lambda_{n,p_i}}{q_i\lambda_{n,q_i}}.
\tag{10.37}\label{eq:forced-deletion-basic}
\]
If $(p_*,q_*)$ is one distinguished constraint, the right-hand side may be
multiplied by
\[
        \Pp(\tau_n(p_*)\le q_*+L).
\tag{10.38}\label{eq:forced-deletion-survival-factor}
\]
\end{lemma}

\begin{proof}
Expose the increasing exponential race as a pure-death chain.  When $q$ labels
are alive, every alive set has total rate at least
\[
        \sum_{b\le q}\lambda_{n,b}
        \ge c_\Sigma q\lambda_{n,q}
\]
by Lemma \ref{lem:endpoint-partial-sum}.  Hence the conditional probability of
deleting source $p_i$ at level $q_i$ is at most
$C\lambda_{n,p_i}/(q_i\lambda_{n,q_i})$.  Sum over compatible death histories
and pull out these bounds at the forced levels.  The remaining transition
weights have total mass at most one, proving
\eqref{eq:forced-deletion-basic}.

For the refinement, retain the condition that $p_*$ survives until level
$q_*$.  Before that level at most $L-1$ deletions are forced.  Thus the
residual schedule contains at least $n-q_*-L$ ordinary deletions before
$p_*$ may be removed.  Lemma \ref{lem:scheduled-deletion-survival} bounds the
residual path mass by
\[
\begin{aligned}
H_{n-q_*-L}([n],p_*)
&=\Pp\left(\#\{b\ne p_*:T_b<T_{p_*}\}
\ge n-q_*-L\right)\\
&=\Pp(\tau_n(p_*)\le q_*+L),
\end{aligned}
\]
which proves \eqref{eq:forced-deletion-survival-factor}.
\end{proof}

\begin{lemma}[Endpoint survival tail]
\label{lem:endpoint-survival-tail}
Put
\[
        U_n=\left\lfloor\frac{n}{(\log n)^2}\right\rfloor,
        \qquad V_n=2U_n.
\]
There are constants $C,c,\delta>0$ such that, uniformly for
$1\le q<p\le n$, $q\le U_n$, and every fixed $L$,
\[
        \Pp(\tau_n(p)\le q+L)
        \le C\exp\left\{-c
        \left(\frac{p\wedge V_n}{q+L}\right)^\delta\right\}.
\tag{10.39}\label{eq:endpoint-survival-tail}
\]
\end{lemma}

\begin{proof}
Put $m=q+L$, $\widetilde p=p\wedge V_n$, and choose
$\beta\in(0,\gamma)$ and $\delta=\beta/(1+\beta)$.  If
$d=\widetilde p/m$ is bounded, enlarge $C$.  Otherwise let
$s=\lfloor m d^\delta\rfloor$.  Then $Cm\le s\le\widetilde p/2$, so in
particular the block $\{1,\ldots,s\}$ does not contain $p$.  Take
$t=1/\lambda_{n,s}$.  For every $b\le s$,
\[
        \Pp(T_{n,b}>t)=e^{-\lambda_{n,b}/\lambda_{n,s}}\ge e^{-1}.
\]
If $T_{n,p}<t$ and at least $m$ of these first $s$ clocks exceed $t$, then at
least $m$ clocks exceed $T_{n,p}$, so $\tau_n(p)\ge m+1$.  Therefore
\[
        \Pp(\tau_n(p)\le m)
        \le e^{-\lambda_{n,p}/\lambda_{n,s}}+e^{-cs}
\]
by a binomial Chernoff bound.  Potter's lower bound and monotonicity give
\[
        \lambda_{n,p}/\lambda_{n,s}
        \ge c(\widetilde p/s)^\beta.
\]
By the choice of $s$, both $s$ and $(\widetilde p/s)^\beta$ are bounded below
by a constant times $d^\delta$.  This proves
\eqref{eq:endpoint-survival-tail}.
\end{proof}

Let $\Gamma$ be a prescribed collection of pairwise vertex-disjoint directed
cycles in terminal coordinates, and let $E_\Gamma$ be the event that all those
cycles occur.  For a cycle component $\gamma$ define its capped logarithmic
oscillation by
\[
        \omega_{U_n}(\gamma)
        =\operatorname{osc}_{a\in\gamma}\log(a\wedge U_n).
\tag{10.40}\label{eq:capped-log-oscillation}
\]

\begin{lemma}[Capped path drop]
\label{lem:capped-path-drop}
Let $\gamma$ be a nontrivial directed cycle containing at least one vertex not
exceeding $U_n$.  There is a downward edge $p\mapsto q$ of $\gamma$ such that
$q\le U_n$ and
\[
        \frac{p\wedge V_n}{q}
        \ge\exp\{\omega_{U_n}(\gamma)/|\gamma|\}.
\tag{10.41}\label{eq:capped-path-drop}
\]
\end{lemma}

\begin{proof}
If $\omega_{U_n}(\gamma)=0$, then the component must contain a vertex above
$U_n$ and a vertex equal to $U_n$; any directed entrance from the upper part
to a vertex not exceeding $U_n$ gives a downward edge satisfying
\eqref{eq:capped-path-drop}.  Assume henceforth that the oscillation is
positive.  Choose vertices whose capped logarithms attain the maximum and
minimum, and follow the directed cycle from the former to the latter.  Put
$z_i=v_i\wedge U_n$ along this path.  The product of the ratios
$z_i/z_{i+1}$ over the downward capped steps is at least the net ratio between
the endpoints, namely $e^{\omega_{U_n}(\gamma)}$.  There are at most
$|\gamma|$ steps, so one downward capped step has ratio at least the right-hand
side of \eqref{eq:capped-path-drop}.  For that step,
$z_{i+1}=q\le U_n$ and
$p\wedge V_n\ge p\wedge U_n=z_i$, which proves the claim.
\end{proof}

\begin{corollary}[Cycle envelope]
\label{cor:endpoint-cycle-envelope}
If $L$ is the total number of vertices of $\Gamma$, then
\[
        \Pp(E_\Gamma)
        \le\frac{C_L}{\prod_{a\in V(\Gamma)}a}.
\tag{10.42}\label{eq:cycle-envelope-basic}
\]
If $\gamma$ is a nontrivial component meeting $[1,U_n]$, then
\[
        \Pp(E_\Gamma)
        \le\frac{C_L}{\prod_{a\in V(\Gamma)}a}
        \Psi_L(\omega_{U_n}(\gamma)),
\tag{10.43}\label{eq:cycle-envelope-component}
\]
where
\[
        \Psi_L(x)=C_L\exp\{-c_Le^{c_Lx}\}.
\tag{10.44}\label{eq:cycle-envelope-psi}
\]
Consequently, if
\[
        \Omega_{U_n}(\Gamma)
        =\max_{\substack{\gamma\text{ nontrivial}\
                         \gamma\cap[1,U_n]\ne\varnothing}}
        \omega_{U_n}(\gamma),
\]
then, whenever the maximum is over a nonempty set,
\[
        \Pp(E_\Gamma)
        \le\frac{C_L}{\prod_{a\in V(\Gamma)}a}
        \Psi_L(\Omega_{U_n}(\Gamma)).
\tag{10.45}\label{eq:cycle-envelope-maximum}
\]
\end{corollary}

\begin{proof}
Apply Lemma \ref{lem:forced-deletion}.  Around every directed cycle the rate
ratios telescope:
\[
        \prod_{a\mapsto b}
        \frac{\lambda_{n,a}}{b\lambda_{n,b}}
        =\frac1{\prod_{a\in\gamma}a}.
\]
This proves \eqref{eq:cycle-envelope-basic}.  For a component meeting
$[1,U_n]$, choose the edge supplied by Lemma \ref{lem:capped-path-drop} as the
distinguished constraint in Lemma \ref{lem:forced-deletion}.  Lemma
\ref{lem:endpoint-survival-tail} gives
\eqref{eq:cycle-envelope-component}.  Choosing a component attaining the
maximum gives \eqref{eq:cycle-envelope-maximum}.
\end{proof}

\subsection{The logarithmic trace}

\begin{lemma}[Continuum multiplicative trace]
\label{lem:continuum-endpoint-trace}
For every fixed $r\ge1$,
\[
        I_r(T):=\int_{[0,T]^r}
        \prod_{j=1}^r k_\gamma(x_{j+1}-x_j)
        \dd x_1\cdots\dd x_r
        =T k_\gamma^{*r}(0)+O_r(1),
\tag{10.46}\label{eq:continuum-endpoint-trace}
\]
where $x_{r+1}=x_1$.
\end{lemma}

\begin{proof}
For $r=1$ the identity is exact.  For $r\ge2$, put
$z_j=x_{j+1}-x_j$, $1\le j<r$, and
$z_r=-\sum_{j<r}z_j$.  If $s_0=0$ and
$s_j=z_1+\cdots+z_j$, let
$\mathcal R(z)=\max_js_j-\min_js_j$.  The admissible translations $x_1$ have
length $[T-\mathcal R(z)]_+$, so
\[
        I_r(T)=\int_{\mathbb R^{r-1}}[T-\mathcal R(z)]_+
        \prod_{j=1}^r k_\gamma(z_j)\dd z_1\cdots\dd z_{r-1}.
\]
The product integral without the bracket is $k_\gamma^{*r}(0)$.  The density
$k_\gamma$ has an exponential right tail, a superexponential left tail, and a
finite first absolute moment.  Since
$\mathcal R(z)\le2\sum_{j<r}|z_j|$, the error integral is finite uniformly in
$T$.
\end{proof}

We shall use the elementary harmonic-coordinate bounds
\[
        \sum_{A\le a\le B}\frac1a
        \le1+\log(B/A),
\tag{10.47}\label{eq:harmonic-interval-bound}
\]
and, for $J_a=[\log a,\log(a+1))$,
\[
        |J_a|=\frac1a+O(a^{-2}).
\tag{10.48}\label{eq:log-cell-length}
\]

\begin{lemma}[Discrete endpoint trace]
\label{lem:discrete-endpoint-trace}
Let $L_n\to\infty$, $U_n=o(n)$, and
$T_n=\log(U_n/L_n)\to\infty$.  Then
\[
\begin{aligned}
&\sum_{\substack{L_n\le a_1,\ldots,a_r\le U_n\\
        a_1,\ldots,a_r\text{ distinct}}}
        \prod_{j=1}^r\frac1{a_{j+1}}
        k_\gamma\!\left(\log\frac{a_{j+1}}{a_j}\right)\\
&\qquad=k_\gamma^{*r}(0)T_n+o(T_n),
\end{aligned}
\tag{10.49}\label{eq:discrete-endpoint-trace}
\]
where $a_{r+1}=a_1$.
\end{lemma}

\begin{proof}
Because $\prod_ja_{j+1}^{-1}=\prod_ja_j^{-1}$, the sum is a logarithmic
Riemann sum.  First restrict to
$\operatorname{osc}(\log a_1,\ldots,\log a_r)\le M$.  On this region the
cyclic kernel product and its first derivatives are bounded by a constant
depending on $r,M$.  Using the cells $J_a$ in
\eqref{eq:log-cell-length}, cell integration and
\eqref{eq:harmonic-interval-bound} show that the difference from the continuum
integral over $[\log L_n,\log U_n]^r$ is
\[
        O_{r,M}(T_n/L_n)=o(T_n).
\tag{10.50}\label{eq:trace-riemann-error}
\]
The same bound controls the replacement of $1/a$ by $|J_a|$.

We next record a uniform discrete tail bound.  The local supremum of
$k_\gamma$ over an interval of length one is bounded by an integrable density
$\overline k_\gamma$ having the same exponential right tail, a
superexponential left tail, and a finite first moment.  Enlarging every
logarithmic cell by one and using
\eqref{eq:harmonic-interval-bound}, the discrete cyclic sum with oscillation
larger than $M$ is bounded by a constant times the corresponding continuum
cyclic integral with $k_\gamma$ replaced by $\overline k_\gamma$ and
oscillation larger than $M-r$.  The increment representation in Lemma
\ref{lem:continuum-endpoint-trace} therefore gives
\[
\lim_{M\to\infty}\sup_n\frac1{T_n}
\sum_{\substack{L_n\le a_1,\ldots,a_r\le U_n\\
 \operatorname{osc}(\log a)>M}}
\prod_{j=1}^r\frac1{a_{j+1}}
 k_\gamma\!\left(\log\frac{a_{j+1}}{a_j}\right)=0.
\tag{10.51}\label{eq:trace-discrete-tail}
\]

Finally, an equality $a_i=a_j$ replaces one harmonic factor by a square:
$\sum_{a\ge L_n}a^{-2}=O(L_n^{-1})$.  The remaining localized relative sums
are uniformly bounded by the same convolution majorant, so all equality
diagonals contribute $O_r(L_n^{-1})$.  Combine
\eqref{eq:trace-riemann-error}--\eqref{eq:trace-discrete-tail} with Lemma
\ref{lem:continuum-endpoint-trace}, and then let $M\to\infty$.
\end{proof}

\subsection{Factorial-moment summation}

The next elementary observation is used repeatedly.  If a component has
$r$ vertices, one may choose one logarithmic translation coordinate and
$r-1$ relative coordinates.  By \eqref{eq:harmonic-interval-bound}, if its
logarithmic oscillation is at most $x$, the total harmonic mass of the relative
coordinates is at most $C_r(1+x)^{r-1}$.  For an ordered collection of $m$
components and $L$ vertices, the corresponding bound is
\[
        C_L(1+x)^{L-m}(1+T_n)^m.
\tag{10.52}\label{eq:component-harmonic-volume}
\]
If one component is required to meet $[1,L_n)$, one of the $T_n$ translation
factors in \eqref{eq:component-harmonic-volume} is replaced by
$O(1+\log L_n+x)$.  If a component crosses $U_n$ and its capped oscillation is
at most $x$, its lower anchor belongs to $[U_ne^{-x},U_n]$, so its translation
factor is $O(1+x)$ rather than $T_n$; each vertex above $U_n$ contributes the
harmonic factor
\[
        H_n:=\sum_{a>U_n}\frac1a=O(\log\log n).
\tag{10.53}\label{eq:upper-harmonic-mass}
\]
These bounds follow by summing the nonanchor vertices successively over
multiplicative intervals and will be used without further comment.

\begin{lemma}[Negligible endpoint cycle configurations]
\label{lem:negligible-endpoint-cycle-configurations}
Fix an ordered collection of $m$ cycle components with total number of vertices
$L$.  Put
\[
        L_n=\lceil\log n\rceil,
        \qquad
        U_n=\left\lfloor\frac{n}{(\log n)^2}\right\rfloor,
        \qquad
        T_n=\log(U_n/L_n).
\]
Then:
\begin{enumerate}[label=(\roman*)]
\item among collections with all vertices in $[L_n,U_n]$, the total probability
of those for which some nontrivial component has logarithmic oscillation greater
than $M$ is at most $\varepsilon_M T_n^m$, where
$\varepsilon_M\downarrow0$;
\item if all component oscillations are at most $M$, the total probability of
collections containing two vertices at logarithmic distance at most $\eta$ is
at most
\[
        C_{L,M}\eta T_n^m+O_{L,M}(T_n^{m-1});
\tag{10.54}\label{eq:close-cycle-configurations}
\]
\item collections having at least one vertex outside $[L_n,U_n]$ contribute
$o(T_n^m)$.
\end{enumerate}
\end{lemma}

\begin{proof}
For (i), let $\Omega$ be the maximum logarithmic oscillation among the
nontrivial components.  Since all vertices are at most $U_n$, this is the same
as $\Omega_{U_n}(\Gamma)$.  Corollary
\ref{cor:endpoint-cycle-envelope} and
\eqref{eq:component-harmonic-volume} give, after decomposing
$\Omega\in[k,k+1)$,
\[
\begin{aligned}
&\sum_{\substack{\Gamma:\ V(\Gamma)\subset[L_n,U_n]\\
                  \Omega(\Gamma)>M}}
\Pp(E_\Gamma)\\
&\qquad\le C_LT_n^m
\sum_{k\ge M}(2+k)^{L-m}\Psi_L(k)
=:\varepsilon_MT_n^m.
\end{aligned}
\]
The superexponential decay in \eqref{eq:cycle-envelope-psi} gives
$\varepsilon_M\downarrow0$.

For (ii), parameterize each component by one translation coordinate and its
bounded relative coordinates.  If the close pair lies in one component, it
cuts out a band of width $O(\eta)$ in a bounded relative-coordinate domain.
The logarithmic lattice error is $O(L_n^{-1})$ by
\eqref{eq:harmonic-interval-bound}.  Summing the remaining coordinates gives
$O_{L,M}((\eta+L_n^{-1})T_n^m)$.  If the close pair lies in two different
components, it links two translation coordinates within a bounded interval
whose length depends only on $M$, and therefore removes one free translation
factor.  Its contribution is $O_{L,M}(T_n^{m-1})$.  This proves
\eqref{eq:close-cycle-configurations}.

For (iii), first consider collections with no vertex above $U_n$ but with a
vertex below $L_n$.  Use the maximum-oscillation envelope as in (i).  By the
observation preceding the lemma, the component meeting the lower range has
translation mass $O(1+\log L_n+\Omega)$ rather than $T_n$.  Hence the total is
bounded by
\[
        C_LT_n^{m-1}
        \sum_{k\ge0}(1+\log L_n+k)(2+k)^{L-m}\Psi_L(k)
        =O_L((1+\log L_n)T_n^{m-1})
        =o(T_n^m).
\tag{10.55}\label{eq:lower-boundary-cycle-sum}
\]
The case in which all components are fixed points follows directly from
$\sum_{a<L_n}a^{-1}=O(\log\log n)$ and the same conclusion.

Now suppose that some vertex exceeds $U_n$.  A component lying wholly above
$U_n$ has harmonic mass at most $H_n^r=o(T_n)$ for its fixed length $r$.
For the remaining components, if at least one nontrivial component meets
$[1,U_n]$, choose the one of maximal capped oscillation and use
\eqref{eq:cycle-envelope-maximum}; the summation in (i), together with
\eqref{eq:lower-boundary-cycle-sum}, is then $O_L(T_n^{m-1})$.  If no such
component exists, every remaining component is either a fixed point or wholly
above $U_n$, and its harmonic mass is $O(T_n)$.  Thus any collection
containing a wholly upper component contributes $o(T_n^m)$.

It remains to treat crossing components.  Let $h\ge1$ be the total number of
vertices above $U_n$.  For every component meeting $[1,U_n]$, use the capped
oscillation \eqref{eq:capped-log-oscillation}, and let $\Omega$ be their
maximum.  Corollary \ref{cor:endpoint-cycle-envelope} supplies the factor
$\Psi_L(\Omega)$.  A crossing component has its lower anchor in
$[U_ne^{-\Omega},U_n]$, so it has no free $T_n$ translation factor; its upper
vertices contribute $H_n^h$.  The harmonic-volume bound therefore gives
\[
\begin{aligned}
&\sum_{\Gamma:\ \Gamma\text{ has a crossing component}}
\Pp(E_\Gamma)\\
&\qquad\le C_LH_n^hT_n^{m-1}
\sum_{k\ge0}(2+k)^{L-m+1}\Psi_L(k)
=O_L(H_n^hT_n^{m-1})=o(T_n^m),
\end{aligned}
\tag{10.56}\label{eq:crossing-boundary-cycle-sum}
\]
because $h$ and $L$ are fixed and $(\log\log n)^h=o(\log n)$.  This completes
(iii).
\end{proof}

\begin{proof}[Proof of Theorem \ref{thm:endpoint-short-cycles}]
Let
\[
        m=\sum_{r=1}^Rm_r,
        \qquad
        L=\sum_{r=1}^Rrm_r.
\]
If $m=0$, the assertion is trivial.  Expanding the mixed factorial moment gives
a sum over ordered collections of distinct directed cycles.  Collections
sharing a vertex have probability zero, so only vertex-disjoint collections
contribute.  The falling factorial orders cycles of the same length; hence no
$m_r!$ divisor appears.

Fix $M<\infty$ and $\eta>0$.  Call a collection good if all vertices lie in
$[L_n,U_n]$, every component has logarithmic oscillation at most $M$, and all
$L$ vertices are pairwise separated by more than $\eta$ on logarithmic scale.
For a good collection, every source-to-target ratio lies in
$[e^{-M},e^M]$, and the sorted targets are separated by ratios at least
$e^\eta$.  Proposition \ref{prop:endpoint-general-entry} therefore gives,
uniformly,
\[
        \Pp(E_\Gamma)
        =(1+o(1))\prod_{a\mapsto b\in\Gamma}
        \frac1b k_\gamma\!\left(\log\frac ba\right).
\tag{10.57}\label{eq:good-cycle-entry-product}
\]
The total leading kernel mass of such collections is $O_L(T_n^m)$ by Lemma
\ref{lem:discrete-endpoint-trace} and the nonnegativity of the kernel, so the
accumulated uniform relative error in
\eqref{eq:good-cycle-entry-product} is $o(T_n^m)$.

For one directed cycle of length $r$, Lemma
\ref{lem:discrete-endpoint-trace} gives
\[
        k_\gamma^{*r}(0)T_n+o(T_n)
\]
when a starting vertex is distinguished.  Division by $r$ removes cyclic
rotations.  For an ordered collection of components, the unrestricted leading
kernel sum factors over components.  Equality restrictions and cross-component
close pairs remove $O_L(T_n^{m-1})$, while same-component $\eta$-bands remove
$O_{L,M}(\eta T_n^m)$.  This follows from the harmonic-coordinate proof of
Lemma \ref{lem:negligible-endpoint-cycle-configurations}; on the bounded
oscillation region the leading kernel is bounded, so the same volume estimates
apply verbatim.  Hence the good
contribution is
\[
        \prod_{r=1}^R
        \left(\frac{k_\gamma^{*r}(0)}rT_n\right)^{m_r}
        +o(T_n^m)+O_{L,M}(\eta T_n^m).
\tag{10.58}\label{eq:good-cycle-factorial-main}
\]
The same lemma controls all nongood collections.  Let first $n\to\infty$, then
$\eta\downarrow0$, and finally $M\to\infty$.  Since
\[
        T_n=\log(U_n/L_n)
        =\log n-3\log\log n+O(1)\sim\log n,
\]
\eqref{eq:good-cycle-factorial-main} proves the mixed factorial-moment
asymptotic.

For the $L^p$ conclusion, put
\[
        X_{n,r}=\frac{C_r}{\alpha_{\gamma,r}\log n}.
\]
The first two factorial moments give $X_{n,r}\to1$ in $L^2$, hence in
probability.  Given any fixed $p<\infty$, choose an integer $s>p$.  The
factorial-moment theorem and the identities expressing ordinary powers as
finite linear combinations of falling factorials give
\[
        \sup_n\Ee X_{n,r}^s<\infty.
\]
Thus $|X_{n,r}-1|^p$ is uniformly integrable, and convergence in probability
gives $X_{n,r}\to1$ in $L^p$.  Since $R$ is fixed, the same argument applies
to the vector in any norm on $\mathbb R^R$.
\end{proof}

\begin{remark}[Transpositions]
For $r=2$,
\[
\begin{aligned}
        k_\gamma^{*2}(0)
        &=\int_{\mathbb R}k_\gamma(x)k_\gamma(-x)\dd x\\
        &=\gamma K_\gamma^2
        \int_{\mathbb R}e^{-2K_\gamma\cosh t}\dd t
        =2\gamma K_\gamma^2\mathrm K_0(2K_\gamma),
\end{aligned}
\]
where $t=\gamma x$ and
$\mathrm K_0(z)=\int_0^\infty e^{-z\cosh t}\dd t$.  Thus
\[
        \Ee C_2\sim
        \gamma K_\gamma^2\mathrm K_0(2K_\gamma)\log n.
\]
In the Sukhatme case this coefficient is $\mathrm K_0(2)$.
\end{remark}

\section{Conclusions and theorem map}\label{sec:conclusions}

The results establish a coherent hierarchy of exact, microscopic, mesoscopic,
and macroscopic laws for Luce permutations.

At the two ends of the ordering, Theorems \ref{thm:top-tv} and
\ref{thm:dinfty} give exact top-window comparison identities, while Theorem
\ref{thm:leakage} identifies the total mass and the defect of the limiting
bottom-window law.  The vanishing-block construction in Proposition
\ref{prop:implantation} shows that these microscopic statistics cannot in
general be inferred from a permuton limit alone.

For uniformly positive continuous profiles, Theorem
\ref{thm:regular-cycles} and Corollary \ref{cor:regular-cycles-tv} give joint
Poisson convergence, including total variation on each fixed cycle-count
vector.  Theorem \ref{thm:long-trace} proves that, for some
$q=q(f)\in(0,1)$,
\[
        \lambda_r=\frac1r+O_f(q^r/r).
\]
The diagonal-band, sparse-entry, thin-graph, and deterministic-overlap results
then extend the same microscopic density to broad deterministic collections of
point constraints.  Under uniformly comparable weights, Theorem
\ref{thm:macroscopic-pd} supplies the complementary macroscopic conclusion:
the ranked cycle lengths converge to $\operatorname{PD}(1)$, and their
size-biased order converges to $\operatorname{GEM}(1)$.

For the Sukhatme and endpoint-regular regimes, Theorems \ref{thm:suk-mean},
\ref{thm:suk-ppp}, \ref{thm:suk-fact}, and
\ref{thm:endpoint-fixed-universality} establish the logarithmic fixed-point
mean, local terminal Poisson point processes, and all fixed factorial moments.
Theorem \ref{thm:global-endpoint-poisson} upgrades the full endpoint fixed-point
count to total-variation approximation by a Poisson law with its exact mean.
Theorems \ref{thm:endpoint-displacement},
\ref{thm:endpoint-tangent-measure}, and
\ref{thm:microscopic-endpoint-permutation} identify, respectively, the
finite-dimensional endpoint displacement law, the deterministic mesoscopic
tangent measure, and the microscopic infinite endpoint permutation.  Finally,
Theorem \ref{thm:endpoint-short-cycles} determines the logarithmic means and all
fixed mixed factorial moments of the endpoint cycle counts of every fixed
length, and hence their joint law of large numbers.

\section{Proof mechanisms}\label{sec:proof-mechanisms}

The proofs use two complementary mechanisms.

\subsection{Local-entry and terminal-condition analysis}

The boundary, regular-profile, and endpoint results begin with an exact event
representation and then identify its incoming mass.

\begin{enumerate}[label=(\roman*)]
\item In the top-window problem, the relevant set is the distinct-word support.
The total-variation distance is exactly the independent-sampling collision
probability.

\item In the bottom-window problem, the candidate infinite law has total mass
\[
        \Pp(N_0\ge k),
\]
and its missing mass is exactly $\Pp(N_0<k)$.  The defect is therefore caused
by too few clocks surviving beyond the terminal endpoint.

\item In the regular positive-profile regime, the kernel $\rho(x,y)$ is both
the Luce permuton density and the microscopic point-entry density for
$\sigma_n(i)=j$.  Separated local estimates are completed by endpoint,
collision, and target-crowding bounds before factorial moments are summed.

\item In the Sukhatme profile, the relevant layer is logarithmic in the
terminal coordinate $a=n+1-i$.  The one-point probability is asymptotic to
$e^{-1}/a$, and logarithmic summation yields both the mean and the local Poisson
point process.

\item For endpoint-regular arrays, after an irrelevant common normalization,
the rates behave as
\[
        C_0(a/n)^\gamma L(n/a).
\]
The shifted-Gumbel saddle produces the local coefficient $c_\gamma/a$, while
the multiplicative transition kernel on logarithmic coordinates has diagonal
return density $x^{-1}k_\gamma^{*r}(0)$.

\item The global endpoint Poisson theorem uses an approximate-Palm rank
transposition.  Swapping the clock of coordinate $a$ with the clock occupying
terminal rank $a$ forces $I_{n,a}=1$; the pushforward law is close to the exact
Palm law, and the fixed-point count changes only through the intended point and
a possible two-cycle.
\end{enumerate}

The shifted-Gumbel coordinates are especially effective because the specified
priority density and the derivative of the threshold mean cancel in the
integrated saddle.  This cancellation explains the same local coefficients in
the regular, Sukhatme, and endpoint-regular regimes.

\subsection{Clock-slab mixing and regeneration}

The macroscopic theorem is not a local saddle calculation.  Short clock-time
slabs produce conditional permanental matchings with hereditary edge flatness.
A common minorization yields a uniform $L^2$ contraction, a quantitative prefix
law gives relative balance in each target category, and a two-scale smoothing
argument upgrades global balance to the dependently selected closure slab.
Deletion of an interior macroscopic cycle preserves regular histories, so the
successive residual fractions become asymptotically independent uniforms.  This
regeneration mechanism yields the $\operatorname{GEM}(1)$ and
$\operatorname{PD}(1)$ limits.

\section{Open problems}\label{sec:conjectures}

The results above settle every fixed cycle length in the regular and endpoint
regimes and the macroscopic limit under uniformly comparable weights.  The main
remaining questions concern the scales between these two extremes.

\begin{conjecture}[Mesoscopic cycle process]\label{conj:mesoscopic-cycles}
Under the hypotheses of Theorem \ref{thm:macroscopic-pd}, if
$r_n\to\infty$ and $r_n=o(n)$, then
\[
        \Ee C_{r_n}(\sigma_n)\sim\frac1{r_n}.
\]
More strongly, let $a_n\to\infty$ with $a_n=o(n)$, and let
$[u_j,v_j]\Subset(0,\infty)$, $1\le j\le m$, be pairwise disjoint intervals.
If $L_{n,1},L_{n,2},\ldots$ are the cycle lengths, then the vector
\[
        \left(
        \#\{i:u_ja_n\le L_{n,i}\le v_ja_n\}
        \right)_{j=1}^m
\]
should converge to independent Poisson variables with respective means
$\log(v_j/u_j)$.  Equivalently, the submacroscopic cycle lengths should converge
locally, after scaling by $a_n$, to the scale-invariant Poisson process with
intensity $\dd x/x$.
\end{conjecture}

The scale-invariant process is the corresponding benchmark for uniform
permutations and the Ewens sampling formula; see Arratia, Barbour, and Tavar\'e
\cite{ABT1992}.

\begin{conjecture}[Total number of cycles]\label{conj:total-cycles}
Assume the exact sampled-profile model
\[
        \theta_{n,i}=f(i/n),\qquad
        f\in C[0,1],\qquad 0<m\le f\le M<\infty,
\]
and put
\[
        \mathscr N_n^{\mathrm{cyc}}
        =\sum_{r=1}^n C_r(\sigma_n).
\]
Then
\[
        \Ee\mathscr N_n^{\mathrm{cyc}}
        =\log n+\kappa_f+o(1),
        \qquad
        \frac{\mathscr N_n^{\mathrm{cyc}}
        -\Ee\mathscr N_n^{\mathrm{cyc}}}{\sqrt{\log n}}
        \Rightarrow N(0,1),
\]
where
\[
        \kappa_f
        =\gamma_{\!E}+(\lambda_1-1)
        +\sum_{r=2}^\infty
        \left(\lambda_r-\frac1r\right)
\]
and $\gamma_E$ is the Euler--Mascheroni constant.
\end{conjecture}

The trace correction is absolutely convergent by Theorem
\ref{thm:long-trace}.  The missing input is control uniform in growing cycle
length.  Whether the same constant is stable under every uniformly vanishing
triangular-array perturbation allowed in Section \ref{sec:regular} is a
separate universality question.

\begin{conjecture}[Endpoint-singular macroscopic cycles]
\label{conj:endpoint-pd}
Under Definition \ref{def:endpoint-rv}, the ranked macroscopic cycle lengths
converge to $\operatorname{PD}(1)$.
\end{conjecture}

Theorem \ref{thm:macroscopic-pd} does not apply because the ratio of the largest
to the smallest weight diverges.  Theorem
\ref{thm:endpoint-short-cycles} nevertheless shows that, for every fixed $R$,
\[
        \sum_{r=1}^R rC_r(\sigma_n)
        =O_{\Pp}(\log n)=o_{\Pp}(n).
\]
A proof of Conjecture \ref{conj:endpoint-pd} must additionally control the total
mass in cycle lengths growing with $n$, particularly throughout the mesoscopic
range.

The bulk--endpoint independence and longest-increasing-subsequence conjectures
are recorded in Appendix \ref{app:further-conjectures}.  Proposition
\ref{prop:implantation} shows that the specific sparse-cycle and
longest-increasing-subsequence conclusions formulated here cannot follow from
permuton convergence alone without additional microscopic control.

\appendix
\section{Additional conjectures}\label{app:further-conjectures}

The conjectures below combine the bulk and endpoint mechanisms or concern a
statistic not controlled by the cycle arguments.  They are not used elsewhere
in the paper.

\begin{conjecture}[Bulk--endpoint independence]\label{conj:bulk-endpoint}
Suppose that there are normalizing constants $s_n>0$ and a continuous
function $f:[0,1)\to(0,\infty)$ such that, for every compact
$K\Subset[0,1)$,
\[
        \sup_{i/n\in K}|s_n\theta_{n,i}-f(i/n)|\longrightarrow0,
\]
that
\[
        f(1-u)\sim Cu^\gamma L(1/u),\qquad u\downarrow0,
\]
and that the normalized terminal arrays $s_n\lambda_{n,a}$ satisfy Definition
\ref{def:endpoint-rv} with the same index $\gamma$.  Let $\rho$ be the
associated Luce permuton density.  For a compact set $K\Subset(0,1)$, a bounded
interval $J\Subset\Rr$, and $b_n\to\infty$, $b_n=o(n)$, define
\[
        \Xi_n^{\mathrm{bulk}}
        =\sum_{i=1}^n
        \one_{\{\sigma_n(i)=i\}}\delta_{i/n},
        \qquad
        \Xi_n^{\mathrm{term}}
        =\sum_{a=1}^n I_{n,a}\delta_{\log(a/b_n)}.
\]
Then, in the product vague topology,
\[
        \left(
        \Xi_n^{\mathrm{bulk}}|_K,
        \Xi_n^{\mathrm{term}}|_J
        \right)
        \Rightarrow
        \left(\Xi^{\mathrm{bulk}}|_K,
        \Xi^{\mathrm{term}}|_J\right),
\]
where the two limiting point processes are independent Poisson point processes
with respective intensities
\[
        \rho(x,x)\dd x
        \qquad\text{and}\qquad
        c_\gamma\dd x.
\]
\end{conjecture}

This prediction combines the separated bulk saddle with the endpoint
regular-variation saddle.  A proof should require a mixed-scale integrated
local limit together with uniform domination across the region between the two
scales.

\begin{conjecture}[Longest increasing subsequences under regularity]
\label{conj:regular-lis}
Under the uniformly positive continuous-profile hypotheses of Section
\ref{sec:regular},
\[
        \frac{\operatorname{LIS}(\sigma_n)}{\sqrt n}
        \xrightarrow{\Pp}\mathcal L(\rho),
\]
where the inhomogeneous Hammersley functional is
\[
        \mathcal L(\rho)
        =2\sup_{(x,y)\in\mathcal A}
        \int_0^1
        \sqrt{\rho(x(t),y(t))x'(t)y'(t)}\dd t.
\]
Here $\mathcal A$ is the class of absolutely continuous, coordinatewise
nondecreasing paths $(x(t),y(t))$ from $(0,0)$ to $(1,1)$.
\end{conjecture}

This is the natural nonuniform-permutation analogue of the inhomogeneous
Hammersley variational laws; see Deuschel and Zeitouni
\cite{DeuschelZeitouni1999} and Sj\"ostrand \cite{Sjostrand2023}.  Additional
microscopic regularity is indispensable.  Indeed, in Proposition
\ref{prop:implantation} choose the planted permutation to be increasing and
choose
\[
        \sqrt n\ll m_n\ll n.
\]
The limiting permuton is unchanged, while the planted block forces
$\operatorname{LIS}(\sigma_n)\ge m_n$ with probability tending to one, thereby
changing the $\sqrt n$ scale.

\begin{thebibliography}{99}

\bibitem{ABT1992}
R. Arratia, A. D. Barbour, and S. Tavar\'e.
\newblock Poisson process approximations for the Ewens sampling formula.
\newblock \emph{The Annals of Applied Probability} 2 (1992), no. 3, 519--535.
\newblock \href{https://doi.org/10.1214/aoap/1177005647}
{doi:10.1214/aoap/1177005647}.

\bibitem{BHJ1992}
A. D. Barbour, L. Holst, and S. Janson.
\newblock \emph{Poisson Approximation}.
\newblock Oxford Studies in Probability, vol. 2.
\newblock Clarendon Press, Oxford, 1992.

\bibitem{BHR1999}
T. P. Bidigare, P. J. Hanlon, and D. N. Rockmore.
\newblock A combinatorial description of the spectrum for the Tsetlin library
and its generalization to hyperplane arrangements.
\newblock \emph{Duke Mathematical Journal} 99 (1999), no. 1, 135--174.
\newblock \href{https://doi.org/10.1215/S0012-7094-99-09906-4}
{doi:10.1215/S0012-7094-99-09906-4}.

\bibitem{BGT1987}
N. H. Bingham, C. M. Goldie, and J. L. Teugels.
\newblock \emph{Regular Variation}.
\newblock Encyclopedia of Mathematics and its Applications, vol. 27.
\newblock Cambridge University Press, Cambridge, 1987.

\bibitem{BBL2009}
J. Borcea, P. Br\"and\'en, and T. M. Liggett.
\newblock Negative dependence and the geometry of polynomials.
\newblock \emph{Journal of the American Mathematical Society} 22 (2009),
no. 2, 521--567.
\newblock \href{https://doi.org/10.1090/S0894-0347-08-00618-8}
{doi:10.1090/S0894-0347-08-00618-8}.

\bibitem{BCD2025}
J. Borga, S. Chatterjee, and P. Diaconis.
\newblock Permuton and local limits for the Luce model.
\newblock Preprint, arXiv:2509.07729v2 [math.PR], 2025.
\newblock \url{https://arxiv.org/abs/2509.07729}.

\bibitem{BD1998}
K. S. Brown and P. Diaconis.
\newblock Random walks and hyperplane arrangements.
\newblock \emph{The Annals of Probability} 26 (1998), no. 4, 1813--1854.
\newblock \href{https://doi.org/10.1214/aop/1022855884}
{doi:10.1214/aop/1022855884}.

\bibitem{CDK2024}
S. Chatterjee, P. Diaconis, and G. B. Kim.
\newblock Enumerative theory for the Tsetlin library.
\newblock \emph{Journal of Algebra} 655 (2024), 139--162.
\newblock \href{https://doi.org/10.1016/j.jalgebra.2023.08.009}
{doi:10.1016/j.jalgebra.2023.08.009}.

\bibitem{DeuschelZeitouni1999}
J.-D. Deuschel and O. Zeitouni.
\newblock On increasing subsequences of i.i.d. samples.
\newblock \emph{Combinatorics, Probability and Computing} 8 (1999), no. 3,
247--263.
\newblock \href{https://doi.org/10.1017/S0963548399003776}
{doi:10.1017/S0963548399003776}.

\bibitem{HKMRS2013}
C. Hoppen, Y. Kohayakawa, C. G. Moreira, B. R\'ath, and R. M. Sampaio.
\newblock Limits of permutation sequences.
\newblock \emph{Journal of Combinatorial Theory, Series B} 103 (2013), no. 1,
93--113.
\newblock \href{https://doi.org/10.1016/j.jctb.2012.09.003}
{doi:10.1016/j.jctb.2012.09.003}.

\bibitem{Luce1959}
R. D. Luce.
\newblock \emph{Individual Choice Behavior: A Theoretical Analysis}.
\newblock John Wiley \& Sons, New York, 1959.

\bibitem{Mukherjee2016}
S. Mukherjee.
\newblock Fixed points and cycle structure of random permutations.
\newblock \emph{Electronic Journal of Probability} 21 (2016), Paper No. 40,
18 pp.
\newblock \href{https://doi.org/10.1214/16-EJP4622}
{doi:10.1214/16-EJP4622}.

\bibitem{PemantlePeres2014}
R. Pemantle and Y. Peres.
\newblock Concentration of Lipschitz functionals of determinantal and other
strong Rayleigh measures.
\newblock \emph{Combinatorics, Probability and Computing} 23 (2014), no. 1,
140--160.
\newblock \href{https://doi.org/10.1017/S0963548313000345}
{doi:10.1017/S0963548313000345}.

\bibitem{Petrov1975}
V. V. Petrov.
\newblock \emph{Sums of Independent Random Variables}.
\newblock Ergebnisse der Mathematik und ihrer Grenzgebiete, Band 82.
\newblock Springer-Verlag, Berlin--Heidelberg, 1975.

\bibitem{Pitman2006}
J. Pitman.
\newblock \emph{Combinatorial Stochastic Processes}.
\newblock Lecture Notes in Mathematics, vol. 1875.
\newblock Springer-Verlag, Berlin, 2006.

\bibitem{Renyi1953}
A. R\'enyi.
\newblock On the theory of order statistics.
\newblock \emph{Acta Mathematica Academiae Scientiarum Hungaricae} 4 (1953),
191--231.
\newblock \href{https://doi.org/10.1007/BF02127580}
{doi:10.1007/BF02127580}.

\bibitem{Simon2005}
B. Simon.
\newblock \emph{Trace Ideals and Their Applications}, second edition.
\newblock Mathematical Surveys and Monographs, vol. 120.
\newblock American Mathematical Society, Providence, RI, 2005.

\bibitem{Sjostrand2023}
J. Sj\"ostrand.
\newblock Monotone subsequences in locally uniform random permutations.
\newblock \emph{The Annals of Probability} 51 (2023), no. 4, 1502--1547.
\newblock \href{https://doi.org/10.1214/23-AOP1624}
{doi:10.1214/23-AOP1624}.

\bibitem{Sukhatme1937}
P. V. Sukhatme.
\newblock Tests of significance for samples of the $\chi^2$-population with
two degrees of freedom.
\newblock \emph{Annals of Eugenics} 8 (1937), no. 1, 52--56.
\newblock \href{https://doi.org/10.1111/j.1469-1809.1937.tb02159.x}
{doi:10.1111/j.1469-1809.1937.tb02159.x}.

\end{thebibliography}

\end{document}
